\documentclass[12pt,fontset=fandol]{ctexart}

\makeatletter
\def\input@path{{../}{../../}{../../../}{../../../../}}
\makeatother

\usepackage{researchnotes}
\usepackage{mathrsfs}
\usepackage{pgfplots}
\pgfplotsset{compat=1.18}
\newtheorem{exercise}{习题}[section]
\hypersetup{pdftitle={概率论与数理统计习题集},pdfauthor={公开题源中文编订},bookmarksnumbered=true}
\title{概率论与数理统计习题集\\\large 本科程度·中文选编与详解}
\author{公开题源中文编订}
\date{2026年10月9日}
\begin{document}
\maketitle
\thispagestyle{empty}
\clearpage
\section*{编订说明}
本册面向修读过微积分与基础线性代数的本科生。概率部分从事件与计数出发，逐步进入随机变量、联合分布、数学期望和极限定理；数理统计部分以抽样分布为基础，讨论点估计、区间估计、假设检验及回归分析。每道题目后紧接对应详解，按专题连续编排；题号和书末索引便于检索，适合逐题学习、独立作答后核对关键步骤。

题源包括中文公开习题解答项目和麻省理工学院开放课程 MIT OCW 18.05（2022 年春季）的习题、课堂练习与考试解答。英文题目均已译成中文，并统一记号、补足必要假设；没有直接使用的图形依赖题不予收录。为补齐国内本科数理统计的知识结构，另设明确标注的“编订补充”题。补充题不冒充网络原题，原资料作者也不对本册新增内容背书。

源文中的计算、支撑集或条件存在问题时，以可核查的推导作校订。涉及近似的地方区分正态近似、渐近分布与精确抽样分布。数值结果按所列公式重新计算并作适度舍入。题面中的试验、测量和分类数据用于教学模型，结论只在对应抽样假设下成立。

题源、许可和主要勘误见书末。概率部分保留原题的数学内容，重排并校订文字和公式；统计部分按教学依赖重新编排，因此本册题号与原资料题号不同。
\subsection*{记号约定}
$\operatorname{Bin}(n,p)$、$\operatorname{Pois}(\lambda)$、$U(a,b)$、$N(\mu,\sigma^2)$ 分别表示二项、泊松、均匀、正态分布。$\operatorname{Exp}(\lambda)$ 采用率参数，均值为 $1/\lambda$；$\operatorname{Gamma}(a,b)$ 采用形状 $a$ 和率 $b$，均值为 $a/b$。$\Phi$ 表示标准正态分布函数。除特别说明，$z_q,t_{\nu;q},\chi^2_{\nu;q},F_{a,b;q}$ 一律是\emph{下侧累积概率为 $q$} 的分位数。

样本用大写 $X_i$ 表示，已观测数值用小写 $x_i$ 表示。$\overline X$ 为样本均值，$S^2=\sum_i(X_i-\overline X)^2/(n-1)$ 为无偏样本方差。独立同分布不是由“随机抽取”自动推出的万能条件，使用时均按具体题设理解。
\clearpage
\tableofcontents
\clearpage
\part{习题与详解}


\clearpage\section{事件、计数与几何概率}

\begin{exercise}[事件分解与德摩根律]\label{q:Y1.1}
验证事件的运算公式：
\[
A\cup B=A+\overline{A}B,\quad A=AB+A\overline{B}
\]
\[\begin{aligned}\overline {\bigcup\limits_{1 \leqslant j \leqslant n} {{A_j}} }  &= \bigcap\limits_{1 \leqslant j \leqslant n} {{{\overline A }_j}},\quad  \overline {\bigcap\limits_{1 \leqslant j \leqslant n} {{A_j}} }  \\&= \bigcup\limits_{1 \leqslant j \leqslant n} {{{\overline A }_j}} \end{aligned}\]
\[\begin{aligned}\overline {\bigcup\limits_{j \geqslant 1} {{A_j}} }  &= \bigcap\limits_{j \geqslant 1} {{{\overline A }_j}},\quad  \overline {\bigcap\limits_{j \geqslant 1} {{A_j}} }  \\&= \bigcup\limits_{j \geqslant 1} {{{\overline A }_j}} \end{aligned}\]\par\smallskip {\small 来源：闫一博公开解答，原题 1.1。}
\end{exercise}

\begin{solve}\phantomsection\label{a:Y1.1}这里用乘积表示交、上横线表示补，只有互不相容事件的并才使用加号。对任意样本点 $\omega$，$\omega\in A\cup B$ 等价于“$\omega\in A$，或者 $\omega\notin A$ 且 $\omega\in B$”，所以 $A\cup B=A+\overline AB$。按是否属于 $B$ 划分 $A$，得到 $A=AB+A\overline B$。对于任意有限或可数指标集 $J$，\[\omega\notin\bigcup_{j\in J}A_j\iff\forall j\in J,\ \omega\notin A_j,\]\[\omega\notin\bigcap_{j\in J}A_j\iff\exists j\in J,\ \omega\notin A_j.\]分别把右侧翻译为补集的交与并，即得两条德摩根律。此论证同时适用于题中的有限与可数版本。\end{solve}

\begin{exercise}[用集合语言表示事件]\label{q:Y1.2}
用集合的运算表达以下事件

（a）$A$发生、$B$和$C$不发生；（b）$A$不发生、$B$和$C$发生；（c）$A,B,C$都不发生；（d）仅有$A,B,C$之一发生；（e）至少$A,B,C$之一发生.\par\smallskip {\small 来源：闫一博公开解答，原题 1.2。}
\end{exercise}

\begin{solve}\phantomsection\label{a:Y1.2}
（a）$A\overline{B}\:{\overline{C}}$；（b）$\overline{A}BC$；（c）$\overline{A}\:\overline{B}\:\overline{C}$；（d）$A\overline{B}\:\overline{C}+\overline{A}B\overline{C}+\overline{A}\:\overline{B}C$；（e）$A\cup B\cup C$.
\end{solve}

\begin{exercise}[有放回抽球与样本空间]\label{q:Y1.3}
袋中有大小和质量相同的红、黄、蓝、白球各一个，从中有放回地任取三次，每次一个，并记录下颜色.

（a）写出试验的样本空间;

（b）求取到的三个球颜色互不相同的概率.\par\smallskip {\small 来源：闫一博公开解答，原题 1.3。}
\end{exercise}

\begin{solve}\phantomsection\label{a:Y1.3}
（a）
\[\Omega  = \{ \omega :\omega  = ({a_1},{a_2},{a_3}),{a_i} \in \{\text{红、黄、蓝、白} \} ,i = 1,2,3\} \]

（b）记取到三个球颜色互不相同为事件$A$，则
\[P(A) = \frac{{C_4^3}\cdot 3!}{{{4^3}}} = \frac{{24}}{{64}} = \frac{3}{8}\]
\end{solve}

\begin{exercise}[将事件列化为互斥事件列]\label{q:Y1.4}
设$\{A_j:j=1,2,\cdots\}$是事件列，求互不相容的$\{B_j:j=1,2,\cdots\}$使得对于$n\geqslant 1$，$\bigcup\limits_{j = 1}^n {{A_j}}  = \bigcup\limits_{j = 1}^n {{B_j}} $和$\bigcup\limits_{j = 1}^\infty  {{A_j}}  = \bigcup\limits_{j = 1}^\infty  {{B_j}} $成立.\par\smallskip {\small 来源：闫一博公开解答，原题 1.4。}
\end{exercise}

\begin{solve}\phantomsection\label{a:Y1.4}取 $B_1=A_1$，对 $j\ge2$ 定义 $B_j=A_j\setminus\bigcup_{i<j}A_i$。若 $i<j$，$B_j$ 不含 $A_i$ 中的任何点，而 $B_i\subset A_i$，所以 $B_i\cap B_j=\varnothing$。若某点属于前 $n$ 个 $A_i$ 的并，取它第一次进入的最小指标 $j$，它便属于 $B_j$；反过来 $B_j\subset A_j$，故两个有限并相等。对无限并同样使用最早进入指标，或对有限并取递增极限，得到可数并相等。\end{solve}

\begin{exercise}[两阶段抽样的对称性]\label{q:Y1.5}
$100$件产品中有$3$件次品.

（a）从中任取$2$件，求至少得一件次品的概率；

（b）从中任取$10$件，再从这$10$件中任取$2$件，求至少得一件次品的概率.\par\smallskip {\small 来源：闫一博公开解答，原题 1.5。}
\end{exercise}

\begin{solve}\phantomsection\label{a:Y1.5}
（a）记任取$2$件且无次品为事件$A$，则
\[P(A) = \frac{{C_{97}^2}}{{C_{100}^2}} \]
记任取$2$件至少有一件次品为事件$B$，则
\[
P(B)=1-P(A)=\frac{147}{2475}
\]
（b）可视作从$100$件产品中取出$10$件不拿出，再从$10$件产品中抽$2$件，故概率与（a）相同.
\end{solve}

\begin{exercise}[无放回试钥匙]\label{q:Y1.7}
钥匙串上的$5$把钥匙中只有$1$把可以开房门，现在无放回地试开房门.计算

（a）第三次打开房门的概率；

（b）三次内打开房门的概率；

（c）如果$5$把钥匙中有$2$把可以打开房门，求三次内打开房门的概率.\par\smallskip {\small 来源：闫一博公开解答，原题 1.7。}
\end{exercise}

\begin{solve}\phantomsection\label{a:Y1.7}
记第$i$次打开房门为事件$A_i$，则

（a）记第三次打开房门为事件$B$
\[P(B)=P({\overline A _1}{\overline A _2}{A_3}) = \frac{4}{5} \times \frac{3}{4} \times \frac{1}{3} = \frac{1}{5}\]

（b）记三次内打开房门为事件$C$
\[\begin{aligned}P(C) &= 1 - P({\overline A _1}{\overline A _2}{\overline A _3}) \\&= 1 - \frac{4}{5} \times \frac{3}{4} \times \frac{2}{3} \\&= \frac{3}{5}\end{aligned}\]

（c）若有两把钥匙可开门，记三次内打开房门为事件$D$，则
\[\begin{aligned}P(D) &= 1 - P({\overline A _1}{\overline A _2}{\overline A _3}) \\&= 1 - \frac{3}{5} \times \frac{2}{4} \times \frac{1}{3} \\&= \frac{9}{{10}}\end{aligned}\]
\end{solve}

\begin{exercise}[生日相同的概率]\label{q:Y1.8}
设每个人的生日随机落在$365$天中的任一天，求$n$个人的生日互不相同的概率和至少有两个人生日相同的概率. 假定各人的生日相互独立；分别讨论 $0\le n\le365$ 与 $n>365$。\par\smallskip {\small 来源：闫一博公开解答，原题 1.8。}
\end{exercise}

\begin{solve}\phantomsection\label{a:Y1.8}
记$n$个人生日互不相同为事件$A$，至少有两个人生日相同为事件$B$，则
\[P(A) = \frac{{C_{365}^n}\cdot n!}{{{{365}^n}}},P(B) = 1 - P(A)\]
 当 $n>365$ 时，互不相同的概率为 $0$，至少两人同生日的概率为 $1$。\end{solve}

\begin{exercise}[恰好得到一副手套]\label{q:Y1.9}
有 $6$ 副可区分的手套，每副恰有匹配的左、右手套各一只。从共 $12$ 只手套中等可能地选取 $4$ 只，求恰好含一副匹配手套的概率。\par\smallskip {\small 来源：闫一博公开解答，原题 1.9。}
\end{exercise}

\begin{solve}\phantomsection\label{a:Y1.9}先选择唯一完整的一副，共 $6$ 种；剩下两只必须分别来自另外五副中的不同两副，有 $\binom52$ 种选副方法，每副选左或右各有两种。因此\[P=\frac{6\binom52\,2^2}{\binom{12}4}=\frac{240}{495}=\frac{16}{33}.\]若剩下两只来自同一副，就会形成第二副匹配手套，必须排除。\end{solve}

\begin{exercise}[两次抽样的交集]\label{q:Y1.10}
从标有$1$至$n$的$n$个球中任取$m$个，记下号码后放回.再从这$n$个球中任取$k$个，记下号码.求两组号码中恰有$c$个号码相同的概率,\par\smallskip {\small 来源：闫一博公开解答，原题 1.10。}
\end{exercise}

\begin{solve}\phantomsection\label{a:Y1.10}
记两组号码中恰有$c$个号码相同为事件$A$，则
\[P(A) = \frac{{C_m^cC_{n - m}^{k - c}}}{{C_n^k}}\]
 组合数公式仅在 $\max(0,m+k-n)\le c\le\min(m,k)$ 的整数范围内为正，范围外概率为零。\end{solve}

\begin{exercise}[单位正方形中的圆域]\label{q:Y1.11}
在区间$[0,1]$中随机地取两点，求它们的平方和小于$1$的概率. 假定两点独立且各服从 $U(0,1)$。\par\smallskip {\small 来源：闫一博公开解答，原题 1.11。}
\end{exercise}

\begin{solve}\phantomsection\label{a:Y1.11}
在区间$[0,1]$随机地取两点相互独立，因此可视作在$[0,1]\times [0,1]$中取$(x,y)$，满足$x^2+y^2<1$，也即以原点为圆心$1$为半径的圆，记任取两点平方和小于$1$为事件$A$，则
\[
P(A)=\frac{\pi}{4}
\]
\end{solve}

\begin{exercise}[到达时差与等待]\label{q:Y1.12}
二人相约在$5:00-6:00$点之间见面，如果两人在$5:00-6:00$之间独立地随机到达，求一人至少等另一人半小时的概率.\par\smallskip {\small 来源：闫一博公开解答，原题 1.12。}
\end{exercise}

\begin{solve}\phantomsection\label{a:Y1.12}
记第一人到达的时间为$x$，第二人到达的时间为$y$，即求$P(|x-y|>0.5h)$，记一人至少等另一人半小时为事件$A$，则
\[P(A) = \frac{{2 \times 0.5 \times (0.5 \times 0.5)}}{1\times 1} = \frac{1}{4}\]
\end{solve}

\begin{exercise}[硬币落格与边界区域]\label{q:Y1.13}
一枚直径为 $1$ 的硬币落在边长为 $a>0$ 的无限方格纸上。假定圆心在一个基本方格内均匀分布。求硬币不与任何格线相交的概率，并求使该概率小于 $0.01$ 的 $a$ 的范围。\par\smallskip {\small 来源：闫一博公开解答，原题 1.13。}
\end{exercise}

\begin{solve}\phantomsection\label{a:Y1.13}圆心到每条格线的距离必须大于半径 $1/2$。当 $a\le1$ 时，允许区域面积为零；当 $a>1$ 时，允许区域是边长 $a-1$ 的正方形。因此所求概率为\[p(a)=\begin{cases}0,&0<a\le1,\\(1-1/a)^2,&a>1.\end{cases}\]对 $a>1$，不等式等价于 $1-1/a<0.1$，即 $a<10/9$。合并两种情况得到 $0<a<10/9$。边界接触的概率为零，不影响结果。\end{solve}

\begin{exercise}[两次服务是否相遇]\label{q:Y1.14}
在同一个小时内有两辆汽车独立到达同一个加油站加油，车$A$加油需要$5$min，车$B$需要$8$min.如果每辆车在这一个小时内等可能的到达，计算这两辆车在加油站不能相遇的概率.\par\smallskip {\small 来源：闫一博公开解答，原题 1.14。}
\end{exercise}

\begin{solve}\phantomsection\label{a:Y1.14}设两车到达时刻为 $X,Y$，以一小时起点为零、分钟为单位。它们在 $[0,60]^2$ 上均匀分布。不相遇等价于 $Y>X+5$ 或 $X>Y+8$，即两块互不相交的直角三角形。面积分别为 $55^2/2$ 和 $52^2/2$，故概率为\[\frac{55^2+52^2}{2\cdot60^2}=\frac{5729}{7200}.\]等号边界概率为零；先到哪一车决定应加哪一个服务时长。\end{solve}

\begin{exercise}[立方体内的球域]\label{q:Y1.15}
在$[0,1]$中任取三点$X,Y,Z$，求$X^2+Y^2+Z^2\leqslant 1$的概率和$X^2+Y^2+Z^2=1$的概率. 假定三点相互独立且各服从 $U(0,1)$。\par\smallskip {\small 来源：闫一博公开解答，原题 1.15。}
\end{exercise}

\begin{solve}\phantomsection\label{a:Y1.15}
  在区间$[0,1]$随机地取三点相互独立，因此可视作在$[0,1]\times [0,1]\times [0,1]$中取$(X,Y,Z)$，满足$X^2+Y^2+Z^2\leqslant 1$和$X^2+Y^2+Z^2=1$，也即以原点为球心$1$为半径的球体和球面，则
\[
P(X^2+Y^2+Z^2\leqslant 1)=\frac{4\pi}{3}\times\frac{1}{8}=\frac{\pi}{6}
\]
根据几何概型的特点，区域边界的测度为零，故
\[
P(X^2+Y^2+Z^2=1)=0
\]  
\end{solve}

\begin{exercise}[概率的次可加性]\label{q:Y1.17}
验证概率的次可加性:
\[P(\bigcup\limits_{j = 1}^\infty  {{A_j}) \leqslant } \sum\limits_{j = 1}^\infty  {P({A_j})} \]\par\smallskip {\small 来源：闫一博公开解答，原题 1.17。}
\end{exercise}

\begin{solve}\phantomsection\label{a:Y1.17}
取$A_0=\varnothing$，${B_j} = {A_j} - \bigcup\limits_{i = 0}^{j - 1} {{A_i}} $，$B_j$互不相容，由于$B_j\subset A_j$，故$ \bigcup\limits_{j = 1}^\infty  {{B_j}}  \subset \bigcup\limits_{j = 1}^\infty  {{A_j}} $.下面验证$ \bigcup\limits_{j = 1}^\infty  {{A_j}}  \subset \bigcup\limits_{j = 1}^\infty  {{B_j}} $，任取$x\in \bigcup\limits_{j = 1}^\infty  {{A_j}}$，存在$i\geqslant 1$（取首个包含$x$的$A_j$对应的下标为$i$），使得$x\in A_i$，但$x\notin A_0,A_1,\cdots,A_{i-1}$（若$i=1$，则$x\notin A_0$），故$x\in B_i$，于是有$ \bigcup\limits_{j = 1}^\infty  {{A_j}}  \subset \bigcup\limits_{j = 1}^\infty  {{B_j}} $，这表明
\[\bigcup\limits_{j = 1}^\infty  {{A_j}}  = \bigcup\limits_{j = 1}^\infty  {{B_j}}\] 根据可列可加性和单调性
\[\begin{aligned}
P(\bigcup\limits_{j = 1}^\infty  {{A_j}} ) &= P(\bigcup\limits_{j = 1}^\infty  {{B_j})}  \\&= \sum\limits_{j = 1}^\infty  {P({B_j})}  \leqslant \sum\limits_{j = 1}^\infty  {P({A_j})} 
\end{aligned}\]
\end{solve}

\begin{exercise}[可数个概率一事件]\label{q:Y1.18}
如果事件$A_1,A_2,\cdots$发生的概率都是$1$，证明：
\[P(\bigcap\limits_{j = 1}^\infty  {{A_j}} ) = 1\quad P(\bigcap\limits_{j = 1}^n {{A_j}} ) = 1\]\par\smallskip {\small 来源：闫一博公开解答，原题 1.18。}
\end{exercise}

\begin{solve}\phantomsection\label{a:Y1.18}
应用De Morgan定律
\[\bigcap\limits_{j = 1}^\infty  {{A_j}}  = \overline {\bigcup\limits_{j = 1}^\infty  {\overline {{A_j}} } } \]
故
\[\begin{aligned}P(\bigcap\limits_{j = 1}^\infty  {{A_j}} ) &= P(\overline {\bigcup\limits_{j = 1}^\infty  {\overline {{A_j}} } } ) \\&= 1 - P(\bigcup\limits_{j = 1}^\infty  {\overline {{A_j}} } )\end{aligned}\]
根据概率的次可加性
\[\begin{aligned}0 \leqslant P(\bigcup\limits_{j = 1}^\infty  {\overline {{A_j}} } ) \leqslant \sum\limits_{j = 1}^\infty  {P(\overline {{A_j}} )}  &= \sum\limits_{j = 1}^\infty  {(1 - P({A_j}))}  \\&= \sum\limits_{j = 1}^\infty  0\\&=0 \end{aligned}\]
综上可得
\[P(\bigcap\limits_{j = 1}^\infty  {{A_j}} ) = 1 - P(\bigcup\limits_{j = 1}^\infty  {\overline {{A_j}} } ) = 1\]
同理可证，$P(\bigcap\limits_{j = 1}^n {{A_j}} ) = 1$
\end{solve}

\begin{exercise}[四个互斥区域的分割]\label{q:Y1.19}
对于事件$A,B$，验证
\[
P(AB)+P(\overline{A}B)+P(A\overline{B})+P(\overline{A}\:\overline{B})=1
\]\par\smallskip {\small 来源：闫一博公开解答，原题 1.19。}
\end{exercise}

\begin{solve}\phantomsection\label{a:Y1.19}四个事件按是否属于 $A$、是否属于 $B$ 划分整个样本空间。它们两两互斥，且并为 $\Omega$。因此由有限可加性，\[\begin{aligned}&P(AB)+P(\overline AB)+P(A\overline B)+P(\overline A\,\overline B)\\&\qquad=P(\Omega)=1.\end{aligned}\]这不需要 $A,B$ 独立。\end{solve}

\begin{exercise}[递增事件列的极限]\label{q:Y1.20}
设$A_1\subset A_2\subset \cdots$，$P(A_j)=(1-1/j)^j$.计算$P(\bigcup\limits_{j = 1}^\infty  {{A_j}} )$.\par\smallskip {\small 来源：闫一博公开解答，原题 1.20。}
\end{exercise}

\begin{solve}\phantomsection\label{a:Y1.20}
$\{A_j\}$是单调增序列，故
\[\begin{aligned}
P(\bigcup\limits_{j = 1}^\infty  {{A_j}} ) &= \mathop {\lim }\limits_{n \to \infty } P({A_n}) \\&= \mathop {\lim }\limits_{n \to \infty } {(1 - \frac{1}{j})^j} \\&= \frac{1}{e}
\end{aligned}\]
\end{solve}

\clearpage\section{条件概率、独立性与贝叶斯公式}

\begin{exercise}[至少一次求职成功]\label{q:Y2.1}
甲毕业时在单位A求职成功的概率为$0.4$，在单位B求职成功的概率为$0.6$，在这两家都求职成功的概率为$0.24$.如果甲在这两个单位都求职，计算他至少被一家录用的概率.\par\smallskip {\small 来源：闫一博公开解答，原题 2.1。}
\end{exercise}

\begin{solve}\phantomsection\label{a:Y2.1}
记甲被单位$A$录取为事件$A$，被单位$B$录取为事件$B$，根据加法公式，
\[P(A \cup B) = P(A) + P(B) - P(AB) = 0.4 + 0.6 - 0.24 = 0.76\]
\end{solve}

\begin{exercise}[自习室占用问题]\label{q:Y2.2}
$9$个学生随机走进$3$个自习室.

（a）求至少有一个自习室无学生的概率；

（b）求每个自习室至少有一个学生的概率. 假定每位学生独立且等概率选择每个自习室。\par\smallskip {\small 来源：闫一博公开解答，原题 2.2。}
\end{exercise}

\begin{solve}\phantomsection\label{a:Y2.2}
每个学生走进自习室相互独立

（a）记第$i$个自习室无学生为事件$A_i$，根据容斥公式
\[P({A_1} \cup {A_2} \cup {A_3}) = \sum\limits_{i = 1}^3 {P({A_i})}  - \sum\limits_{1 \leqslant i < j \leqslant 3}^3 {P({A_i}{A_j})}  + P(\bigcap\limits_{i = 1}^3 {{A_i}} )\]
其中
\[\begin{aligned}P({A_i}) &= \frac{{{2^9}}}{{{3^9}}},P({A_i}{A_j}) \\&= \frac{{{1^9}}}{{{3^9}}},P(\bigcap\limits_{i = 1}^3 {{A_i}} ) \\&= 0\end{aligned}\]
代入可得
\[\begin{aligned}P({A_1} \cup {A_2} \cup {A_3}) &= 3 \cdot \frac{{{2^9}}}{{{3^9}}} - 3 \cdot \frac{{{1^9}}}{{{3^9}}} + 0 \\&= \frac{{511}}{{6561}}\end{aligned}\]

（b）记每个自习室至少有一个学生为事件$B$，则
\[P(B) = 1 - P({A_1} \cup {A_2} \cup {A_3}) = \frac{{6050}}{{6561}}\]
\end{solve}

\begin{exercise}[无人拿到自己帽子]\label{q:Y2.3}
$6$个人将各自的帽子混在一起后每人任取一项，求没有人取得自己帽子的概率.\par\smallskip {\small 来源：闫一博公开解答，原题 2.3。}
\end{exercise}

\begin{solve}\phantomsection\label{a:Y2.3}
记第$i$个人没有取到自己的帽子为事件$A_i$，则$\overline{A}_i$表示第$i$个人取到了自己的帽子，根据De Morgan定律
\[P(\bigcap\limits_{i = 1}^6 {{A_i}} ) = 1 - P(\bigcup\limits_{i = 1}^6 {{{\overline A }_i}} )\]
根据容斥公式
\[\begin{aligned}P(\bigcup\limits_{i = 1}^6 {{{\overline A }_i}} ) &= \sum\limits_{n = 1}^6 {{{( - 1)}^{n - 1}}C_6^n\frac{{(6 - n)!}}{{6!}}}  \\&= \sum\limits_{n = 1}^6 {{{( - 1)}^{n - 1}}\frac{1}{{n!}}} \end{aligned}\]
故
\[\begin{aligned}P(\bigcap\limits_{i = 1}^6 {{A_i}} ) &= 1 - \sum\limits_{n = 1}^6 {{{( - 1)}^{n - 1}}\frac{1}{{n!}}}  \\&= \sum\limits_{n = 0}^6 {{{( - 1)}^n}\frac{1}{{n!}}} \end{aligned}\]
\end{solve}

\begin{exercise}[有放回点名的覆盖概率]\label{q:Y2.4}
开学时老师声明本学期要点名$m$个人，每次都是有放回的随机点名.当全班有$n$个学生，计算

（a）你在本学期不被点到的概率；

（b）当$n\leqslant m$，至少有一个学生在本学期不被点到的概率.\par\smallskip {\small 来源：闫一博公开解答，原题 2.4。}
\end{exercise}

\begin{solve}\phantomsection\label{a:Y2.4}
（a）记一名同学一学期不被点到为事件$A$，由于每次点名独立，故
\[P(A) = {\left( {1 - \frac{1}{n}} \right)^m}\]

（b）记第$i$个学生未被点到为事件$A_i$，则至少有一个学生未被点到的概率为$P(\bigcup\limits_{i = 1}^n {{A_i}})$，由于\[P({A_i}) = \frac{{{{(n - 1)}^m}}}{{{n^m}}},P({A_i}{A_j}) = \frac{{{{(n - 2)}^m}}}{{{n^m}}}, \cdots \]根据容斥公式
\[P(\bigcup\limits_{i = 1}^n {{A_i}} ) = \sum\limits_{k = 1}^n {{{( - 1)}^{k - 1}}} C_n^k\frac{{{{(n - k)}^m}}}{{{n^m}}}\]
\end{solve}

\begin{exercise}[信封匹配与容斥]\label{q:Y2.5}
将$n$个不同的信笺随机放入$n$个写好地址的信封，求至少有一封匹配的概率.\par\smallskip {\small 来源：闫一博公开解答，原题 2.5。}
\end{exercise}

\begin{solve}\phantomsection\label{a:Y2.5}记 $A_i$ 为第 $i$ 封信匹配的事件。任意指定 $k$ 封信都匹配时，其余信件可任意排列，故概率为 $(n-k)!/n!$。容斥公式给出\[\begin{aligned}P\left(\bigcup_{i=1}^nA_i\right)&=\sum_{k=1}^n(-1)^{k-1}\binom nk\frac{(n-k)!}{n!}\\&=\sum_{k=1}^n\frac{(-1)^{k-1}}{k!}.\end{aligned}\]特别地，当 $n\to\infty$ 时，它趋于 $1-e^{-1}$。使用容斥的原因是多个匹配事件会同时发生，简单相加会重复计数。\end{solve}

\begin{exercise}[独立任务的共同成功]\label{q:Y2.6}
$6$个人独立破译同一个密码.当第$j$个人能破译密码的概率为$p_j$，计算密码被破译的概率.\par\smallskip {\small 来源：闫一博公开解答，原题 2.6。}
\end{exercise}

\begin{solve}\phantomsection\label{a:Y2.6}
记第$j$个人破译成功为事件$A_j$，根据De Morgan定律
\[P(\bigcup\limits_{j = 1}^6 {{A_j}} ) = 1 - P(\bigcap\limits_{j = 1}^6 {{{\overline A }_j}} )\]
由于$6$个人未成功破译密码相互独立，故
\[P(\bigcup\limits_{j = 1}^6 {{A_j}} ) = 1 - \prod\limits_{j = 1}^6 {(1 - {p_j})} \]
\end{solve}

\begin{exercise}[发现信号的试验次数]\label{q:Y2.7}
在旷野迷路后，甲每隔一小时发出一个求救信号，如果每个信号被搜救队发现的概率是$0.1$，要以$0.95$的概率保证搜救队能发现信号，甲至少要发送多少个信号？ 假设各次信号是否被发现相互独立。\par\smallskip {\small 来源：闫一博公开解答，原题 2.7。}
\end{exercise}

\begin{solve}\phantomsection\label{a:Y2.7}
记甲发出的第$j$个信号被接受为事件$A_j$，根据De Morgan定律
\[P(\bigcup\limits_{j = 1}^n {{A_j}} ) = 1 - P(\bigcap\limits_{j = 1}^n {{{\bar A}_j}} )\]
由于每个信号未被接受相互独立，故
\[P(\bigcup\limits_{j = 1}^n {{A_j}} ) = 1 - \prod\limits_{j = 1}^n {0.9}  = 1 - {0.9^n}\]
由题意知
\[1 - {0.9^n} \geqslant 0.95 \Rightarrow {0.9^n} < 0.05 \Rightarrow n > {\log _{0.9}}(0.05) \Rightarrow n \geqslant 29\]
\end{solve}

\begin{exercise}[保证成功概率的人员数]\label{q:Y2.9}
$12$小时内，每个破译员能够成功破译刚截获的密码的概率（称为破译能力）不小于$0.05$，现在上级派来$5$位破译能力达到$0.2$的专家参加工作，需要再配备多少破译员才能使得密码被破译的概率达到$0.95$. 假设各人的破译结果相互独立。\par\smallskip {\small 来源：闫一博公开解答，原题 2.9。}
\end{exercise}

\begin{solve}\phantomsection\label{a:Y2.9}
记$12$小时内密码未被破译为事件$A$，记第$i$名专家未破译密码为事件$A_i$，第$j$名破译员未破译密码为事件$B_j$，现求破译员人数使得$P(A)<0.05$，即求$n$使得
\[P(A) = P({A_1} \cdots {A_5}{B_1} \cdots {B_n}) < 0.05\]
由于每人未破译密码相互独立，故
\[P(A) = {(0.8)^5}{(0.95)^n}\]
于是有
\[{(0.8)^5}{(0.95)^n} < 0.05 \Rightarrow n \geqslant 37\]
\end{solve}

\begin{exercise}[德梅雷的骰子问题]\label{q:Y2.10}
投掷两个骰子$24$次，至少掷出一对$6$的概率小于二分之一，证明这一结论 假定骰子公平且所有投掷独立。\par\smallskip {\small 来源：闫一博公开解答，原题 2.10。}
\end{exercise}

\begin{proof}\phantomsection\label{a:Y2.10}
记投掷$24$次两个骰子没有出现一对$6$为事件$A$，则
\[P(A) = {\left( {\frac{{35}}{{36}}} \right)^{24}} > \frac{1}{2}\]
这表明掷$24$次两个骰子至少掷出一对$6$的概率小于二分之一.
\end{proof}

\begin{exercise}[条件成绩概率]\label{q:Y2.11}
无论哪门课先考，甲的微积分和线性代数得优的概率分别是$a$和$b$，微积分先考并得优的条件下，线性代数得优的概率是$c$.当微积分先考时，计算

（a）甲的这两门都得优的概率；（b）微积分没得优的条件下，线性代数得优的概率；（c）线性代数得优的概率；（d）至少一门得优的概率；（e）至少一门不得优的概率. 设 $0<a<1$，且给定概率满足联合分布的相容条件。\par\smallskip {\small 来源：闫一博公开解答，原题 2.11。}
\end{exercise}

\begin{solve}\phantomsection\label{a:Y2.11}
记甲微积分得优为事件$A$，线性代数得优为事件$B$

（a）
\[P(AB) = P(A)P(B|A) = ac\]

（b）\[P(B|\bar A) = \frac{{P(\bar AB)}}{{P(\bar A)}} = \frac{{P(B) - P(AB)}}{{P(\bar A)}} = \frac{{b - ac}}{{1 - a}}\]

（c）
\[P(B) = P(A)P(B|A) + P(\bar A)P(B|\bar A) = b\]

（d）
\[P(A \cup B) = P(A) + P(B) - P(AB) = a + b - ac\]

（e）
\[P(\overline A  \cup \overline B ) = 1 - P(AB) = 1 - ac\]
\end{solve}

\begin{exercise}[连续受损的生存概率]\label{q:Y2.12}
一架相机第一次落地摔坏的概率是$0.5$.若第一次没摔坏，第二次落地摔坏的概率是$0.7$.若第二次没摔坏，第三次落地摔坏的概率为$0.9$.求该相机落地三次没有摔坏的概率.\par\smallskip {\small 来源：闫一博公开解答，原题 2.12。}
\end{exercise}

\begin{solve}\phantomsection\label{a:Y2.12}令 $A_i$ 表示第 $i$ 次落地后仍未损坏。题目给出的是在之前完好的条件下的生存概率，故\[\begin{aligned}P(A_1A_2A_3)&=P(A_1)P(A_2\mid A_1)P(A_3\mid A_1A_2)\\&=0.5\times0.3\times0.1=0.015.\end{aligned}\]这里使用条件概率乘法公式，不需要把三次摔落后是否损坏当作独立事件。\end{solve}

\begin{exercise}[已知一枚骰子的点数]\label{q:Y2.13}
投掷$2$个均匀的骰子，已知至少有一个骰子的点数是$3$，

（a）计算这两个骰子的点数之和为$7$的概率；（b）你应当猜这两个骰子的点数之和是$7$还是$6$.\par\smallskip {\small 来源：闫一博公开解答，原题 2.13。}
\end{exercise}

\begin{solve}\phantomsection\label{a:Y2.13}
记第一个骰子掷出点数$i$为事件$A_i$，第二个骰子掷出点数$i$为事件$B_i$，则$A_3\cup B_3$表示至少有一个骰子的点数是$3$

（a）记两个骰子点数之和为$7$为事件$A$，则
\[P(A|{A_3} \cup {B_3}) = \frac{{P({A_4}{B_3} \cup {A_3}{B_4})}}{{P({A_3} \cup {B_3})}} = \frac{2}{{11}}\]

（b）记两个骰子点数之和为$6$为事件$B$，则
\[P(B|{A_3} \cup {B_3}) = \frac{{P({A_3}{B_3})}}{{P({A_3} \cup {B_3})}} = \frac{1}{{11}}\]
应当猜$7$.
\end{solve}

\begin{exercise}[两孩家庭中的条件事件]\label{q:Y2.15}
在一个抽象家庭模型中，两名年龄不同孩子的性别相互独立，男孩概率为 $0.51$，女孩概率为 $0.49$。已知至少有一个女孩，求两人都是女孩的概率。\par\smallskip {\small 来源：闫一博公开解答，原题 2.15。}
\end{exercise}

\begin{solve}\phantomsection\label{a:Y2.15}
记第一个孩子是女孩为事件$A$，第二个孩子是女孩为事件$B$，则
\[P(AB|A \cup B) = \frac{{P(AB)}}{{P(A \cup B)}} = \frac{{{{0.49}^2}}}{{1 - {{0.51}^2}}} = 0.3245\]
\end{solve}

\begin{exercise}[通知渠道的全概率]\label{q:Y2.16}
老板有一个不很负责的秘书.当要秘书通知张经理$5$小时后见面时，秘书马上办理，但是只用某种方式通知一次.设秘书用传真通知的概率是$0.3$，用短信通知的概率是$0.2$，用电子邮件通知的概率为$0.5$，而张经理在$5$小时内能收到传真的概率是$0.8$,能看到短信的概率是$0.9$，能看到电子邮件的概率是$0.4$.

（a）计算张经理收到通知的概率;

（b）如果收到通知的张经理也有$5$\%的概率不能来见老板，计算老板不能按时见到张经理的概率；\par\smallskip {\small 来源：闫一博公开解答，原题 2.16。}
\end{exercise}

\begin{solve}\phantomsection\label{a:Y2.16}
（a）记张经理收到通知为事件$A$，根据全概率公式
\[P(A) = 0.3 \times 0.8 + 0.2 \times 0.9 + 0.5 \times 0.4 = 0.62\]

（b）记张经理按时来见老板为事件$B$，根据全概率公式
\[\begin{aligned}P(B)&=P(A)P(B\mid A)+P(\overline A)P(B\mid\overline A)\\&=0.62\times0.95+0.38\times0=0.589\end{aligned}\]
因此张经理不能按时来见老板的概率为$P(\overline{B})=1-P(B)=0.411$.
\end{solve}

\begin{exercise}[隐私调查的随机化回答]\label{q:Y2.17}
在调查家庭暴力所占家庭的比例$p$时，为得到实际的$p$同时又不侵犯个人隐私，调查人员将袋中放入比例是$p_0$的红球和比例是$q_0=1-p_0$的白球.被调查者在袋中任取一球窥视后放回，并承诺取得红球就讲真话，取到白球就讲假话.被调查者只需在匿名调查表中选“是”或“否”，然后将表放入投票箱.没人能知道被调查者是否讲真话和回答的是什么.

（a）如果声称有家庭暴力的家庭比例是$p_1$，求$p$；

（b）如果袋中装有$30$个红球，$50$个白球，调查了$320$个家庭，其中有$195$个家庭回答了是，估算$p$； 假设取球与家庭状态独立，且 $p_0\ne1/2$。\par\smallskip {\small 来源：闫一博公开解答，原题 2.17。}
\end{exercise}

\begin{solve}\phantomsection\label{a:Y2.17}
（a）记调查家庭回答“是”为事件$A$，抽到红球为事件$B$，则$P(B)=p_0$,$P(\overline{B})=q_0$,$P(A|B)=p$,$P(A|\overline{B})=1-p$，根据全概率公式
\[{p_1} = P(A) = P(B)P(A|B) + P(\overline B )P(A|\overline B ) = {p_0}p + {q_0}(1 - p)\]
于是有
\[p = \frac{{{p_1} - {q_0}}}{{{p_0} + {q_0}}}\]
（b）将数据带入
\[\begin{aligned}p &= \frac{{{p_1} - {q_0}}}{{{p_0} - {q_0}}} \\&= \frac{{\frac{{195}}{{320}} - \frac{5}{8}}}{{\frac{3}{8} - \frac{5}{8}}} \\&= \frac{5}{{80}} \\&= 6.25\% \end{aligned}\]
\end{solve}

\begin{exercise}[比分差达到二的比赛]\label{q:Y2.19}
甲乙独立地比赛，每局甲胜的概率为 $p\in(0,1)$，无平局。每两局构成一轮；若同一人在这一轮连胜两局则此人获胜，否则重新开始。求甲最终获胜的概率。\par\smallskip {\small 来源：闫一博公开解答，原题 2.19。}
\end{exercise}

\begin{solve}\phantomsection\label{a:Y2.19}记甲最终获胜的概率为 $w$。第一轮甲连胜的概率为 $p^2$，乙连胜的概率为 $(1-p)^2$，双方各胜一局的概率为 $2p(1-p)$。在最后一种情况下，后续试验重新开始，故\[w=p^2+2p(1-p)w.\]每轮结束比赛的概率 $p^2+(1-p)^2>0$，从而比赛以概率 $1$ 结束。解得\[w=\frac{p^2}{p^2+(1-p)^2}.\]这里的赛制是固定的两局一轮；若改成任意相邻两局连胜即结束，递推状态和答案都需要改变。\end{solve}

\begin{exercise}[加强抽样的波利亚球袋]\label{q:Y2.20}
袋中起初有 $b$ 个黑球和 $r$ 个红球，$b,r$ 为正整数。每次等可能取一球后放回，再增加同色球一个。记 $B_i$ 为第 $i$ 次取黑球。求 $P(B_1B_2B_3)$、$P(B_2)$、$P(B_3)$，并用归纳法求 $P(B_n)$。\par\smallskip {\small 来源：闫一博公开解答，原题 2.20。}
\end{exercise}

\begin{solve}\phantomsection\label{a:Y2.20}连续三次取黑的概率按条件概率相乘：\[P(B_1B_2B_3)=\frac{b(b+1)(b+2)}{(b+r)(b+r+1)(b+r+2)}.\]对任意初始正整数 $b,r$，将同时归纳证明 $P(B_n)=b/(b+r)$。$n=1$ 成立。若所有初始构成都满足第 $n-1$ 步结论，条件于第一步后，袋子分别变为 $(b+1,r)$ 或 $(b,r+1)$，于是\[\begin{aligned}P(B_n)&=\frac{b}{b+r}\frac{b+1}{b+r+1}+\frac{r}{b+r}\frac{b}{b+r+1}\\&=\frac{b}{b+r}.\end{aligned}\]因此 $P(B_2),P(B_3)$ 也都等于最初黑球比例。各次取球的边缘概率相同，但取球结果并不独立。\end{solve}

\begin{exercise}[已知正点反推交通方式]\label{q:Y2.21}
甲乘汽车、火车的概率分别为$0.6,0.4$.汽车和火车正点到达的概率分别是$0.8$和$0.9$.现在甲已经正点到达，求甲乘的是火车的概率.\par\smallskip {\small 来源：闫一博公开解答，原题 2.21。}
\end{exercise}

\begin{solve}\phantomsection\label{a:Y2.21}
记甲乘火车为事件$A$，甲正点到达为事件$B$，根据全概率公式
\[P(B) = P(A)P(B|A) + P(\overline A )P(B|\overline A ) = 0.4 \times 0.9 + 0.6 \times 0.8 = 0.84\]
再根据贝叶斯公式，有
\[P(A|B) = \frac{{P(A)P(B|A)}}{{P(B)}} = \frac{{0.4 \times 0.9}}{{0.84}} = \frac{3}{7}\]
\end{solve}

\begin{exercise}[观察合格品后的状态]\label{q:Y2.22}
一台机床工作状态良好时，产品的合格率为$99\%$，机床发生故障时的产品合格率为$50\%$.设每次新开机器时机床处于良好状态的概率为$95\%$.如果新开机器后生产的第一件产品是合格品，判断机器处于良好状态的概率.\par\smallskip {\small 来源：闫一博公开解答，原题 2.22。}
\end{exercise}

\begin{solve}\phantomsection\label{a:Y2.22}
记机床工作状态良好为事件$A$，新开机器后生产的产品是合格品为事件$B$，根据全概率公式和贝叶斯公式
\[P(A|B) = \frac{{P(A)P(B|A)}}{{P(B)}} = \frac{{0.95 \times 0.99}}{{0.95 \times 0.99 + 0.05 \times 0.5}} = 0.9741\]
\end{solve}

\begin{exercise}[低先验概率下的连续检测]\label{q:Y2.23}
一种新方法检测出某种疾病的概率是百分之百，但是把没病的人判定有病的概率达到百分之五.设群体中该病的发病率是万分之一.

（a）甲在身体普查中被判断患病，甲的确患病的概率是多少？

（b）甲再次复查又被判断有病时，甲的确患病的概率是多少？ 本题为抽象模型，假定两次检测结果在给定是否患病的条件下相互独立。\par\smallskip {\small 来源：闫一博公开解答，原题 2.23。}
\end{exercise}

\begin{solve}\phantomsection\label{a:Y2.23}记真实患病事件为 $A$，每次阳性为 $B_i$。以下只按题设的抽象条件概率计算。一次阳性后，\[P(A\mid B_1)=\frac{0.0001}{0.0001+0.9999(0.05)}\approx0.00199621.\]两次检测在给定真实状态后独立，所以真阳性的联合条件概率为 $1$，两次假阳性的联合条件概率为 $0.05^2$。因此\[P(A\mid B_1B_2)=\frac{0.0001}{0.0001+0.9999(0.05)^2}\approx0.03846524.\]直接使用原始先验与联合似然，可避免把第一次后验过早舍入后再传入第二次计算。条件独立与无条件独立不是同一假设。\end{solve}

\begin{exercise}[目击证据与先验比例]\label{q:Y2.25}
某城市夏利牌出租车占$85\%$，富康牌出租车占$15\%$.这两种出租车都是红色，富康出租车略大一些，每辆车肇事的概率相同.在一次出租车肇事逃逸案件中，有人指证富康车肇事.为了确定是否富康车肇事，在肇事地点和相似的能见度下警方对证人辨别出租车的能力进行了测验，发现证人正确识别富康车的概率是$90\%$，正确认识夏利车的概率是$80\%$.

（a）如果证人没有撒谎，求肇事者是富康车的概率；

（b）如果两个情况相同的证人独立地作出相同的指证，求肇事者是富康车的概率. 多名证人的识别结果假设在给定肇事车型后条件独立。\par\smallskip {\small 来源：闫一博公开解答，原题 2.25。}
\end{exercise}

\begin{solve}\phantomsection\label{a:Y2.25}记车型为富康的事件为 $A$，证人报告富康为 $B$。根据题设，先验为 $0.15$，两种真实车型下报告的概率为 $0.9$ 与 $0.2$。于是\[P(A\mid B)=\frac{0.15(0.9)}{0.15(0.9)+0.85(0.2)}=\frac{27}{61}\approx0.442623.\]给定车型后，两名证人报告条件独立，故两次相同指证后的后验为\[\frac{0.15(0.9)^2}{0.15(0.9)^2+0.85(0.2)^2}=\frac{243}{311}\approx0.781350.\]这里的先验来自题设中两车型相同的肇事概率假设；若各车型发生事件的概率不同，应先修正事件总体中的车型比例。\end{solve}

\begin{exercise}[无穷多次成功的条件]\label{q:Y2.26}
甲不断地向一个目标射击，第$j$次击中目标的概率为$j^{-\alpha}$，其中$\alpha$是正的常数.如果每次能否击中是相互独立的，当$\alpha$取何值时，甲能够击中目标无穷多次.\par\smallskip {\small 来源：闫一博公开解答，原题 2.26。}
\end{exercise}

\begin{solve}\phantomsection\label{a:Y2.26}
甲每次能否击中目标相互独立，根据博雷尔—坎泰利引理，记甲第$j$次击中目标为事件$A_j$，若
\[\sum\limits_{j = 1}^\infty  {P({A_j})}  = \sum\limits_{j = 1}^\infty  {{j^{ - \alpha }}}  = \infty \]
则$P(A_n\: \text{i.o.})=1$，故$\alpha\leqslant 1$时甲能够击中目标无穷多次.
\end{solve}

\begin{exercise}[最终不再失败的条件]\label{q:Y2.27}
为了得到某种新材料，试验在条件不断改进的情况下进行.如果第$j$次试验成功的概率是$1-j^{-\alpha}$，其中$\alpha$是正常数.计算某次试验后的所有试验不再失败的概率$P$. 假设各次试验是否成功相互独立。\par\smallskip {\small 来源：闫一博公开解答，原题 2.27。}
\end{exercise}

\begin{solve}\phantomsection\label{a:Y2.27}
记第$j$次试验失败为事件$A_j$，则$P(A_j)=j^{-\alpha}$，根据博雷尔—坎泰利引理
\[\left\{ \begin{gathered}
  P({A_n}\,{\text{i}}.{\text{o}}{\text{.) = 0}} \hfill \quad \alpha>1 \\
  P({A_n}\,{\text{i}}.{\text{o}}{\text{.) = 1}} \hfill \quad \alpha\leqslant 1\\ 
\end{gathered}  \right.\]
这表明$\alpha>1$时，某次试验后所有试验不再失败的概率为$1$，$\alpha\leqslant 1$时，某次试验后所有试验不再失败的概率为$0$.
\end{solve}

\begin{exercise}[加红球后的黑球出现次数]\label{q:Y2.28}
设想口袋中只有一个黑球和一个红球.现在在袋中每次任取一个球后放回，且第$j$次取球后再加入红球$j$个.如果一直做下去，能够取到黑球多少次？\par\smallskip {\small 来源：闫一博公开解答，原题 2.28。}
\end{exercise}

\begin{solve}\phantomsection\label{a:Y2.28}
记第$j$次取球摸到黑球为事件$A_j$，根据题意
\[P({A_j}) = \frac{1}{{1 + (1 + 1 + 2 +  \cdots  + (j - 1))}} = \frac{1}{{2 + \frac{{j(j - 1)}}{2}}}\]
根据博雷尔—坎泰利引理，由于
\[\sum\limits_{j = 1}^\infty  {\frac{2}{{4 + j(j - 1)}}}  < \infty \]
故$P({A_n}\,{\text{i}}.{\text{o}}{\text{.) = 0}}$，能够取到黑球有限次.
\end{solve}

\begin{exercise}[黑红同时增长的抽球]\label{q:Y2.29}
袋中最初有一个黑球和一个红球。每次随机取一球后放回，并在第 $j$ 次抽取后确定地增加 $j$ 个红球和一个黑球。各次使用独立的随机抽取机制。证明黑球几乎必然被抽到无穷多次。\par\smallskip {\small 来源：闫一博公开解答，原题 2.29。}
\end{exercise}

\begin{solve}\phantomsection\label{a:Y2.29}第 $j$ 次抽取前，黑球有 $j$ 个，红球有 $1+j(j-1)/2$ 个，所以\[p_j=P(B_j)=\frac{2j}{j^2+j+2}.\]由于每次增加的球数与抽到的颜色无关，袋中构成是确定的；独立随机抽取机制使事件 $B_j$ 相互独立。又 $p_j\sim2/j$，故 $\sum_jp_j=\infty$。由第二博雷尔—坎泰利引理，$P(B_j\text{ 无穷多次发生})=1$。若增球数量依赖抽取结果，就不能未经检查直接使用这个独立性论证。\end{solve}

\begin{exercise}[全概率公式的证明]\label{q:Y2.30}
证明全概率公式：如果事件$A_1,A_2,\cdots$互不相容，$B\subset \bigcup\limits_{j = 1}^\infty  {{A_j}} $，有
\[P(B) = \sum\limits_{j = 1}^\infty  {P({A_j})P(B|{A_j})} \] 对出现在条件概率中的事件，假设 $P(A_j)>0$。\par\smallskip {\small 来源：闫一博公开解答，原题 2.30。}
\end{exercise}

\begin{proof}\phantomsection\label{a:Y2.30}
根据题意有
\[\begin{aligned}B \subset \bigcup\limits_{j = 1}^\infty  {{A_j}}  \Rightarrow B &= (\bigcup\limits_{j = 1}^\infty  {{A_j}} )B \Rightarrow B \\&= \bigcup\limits_{j = 1}^\infty  {({A_j}} B)\end{aligned}\]
由于$A_1,A_2,\cdots$互不相容，故$A_jB$也互不相容，因此有
\[P(B) = P(\bigcup\limits_{j = 1}^\infty  {({A_j}} B)) = \sum\limits_{j = 1}^\infty  {P({A_j}B)} \]
再根据乘法公式
\[P(B) = \sum\limits_{j = 1}^\infty  {P({A_j}B)}  = \sum\limits_{j = 1}^\infty  {P({A_j})P(B|{A_j})} \]
\end{proof}

\begin{exercise}[贝叶斯公式的证明]\label{q:Y2.31}
证明贝叶斯公式：如果事件$A_1,A_2,\cdots$互不相容，$B\subset \bigcup\limits_{j = 1}^\infty  {{A_j}} $，则$P(B)>0$时，有
\[P({A_j}|B) = \frac{{P({A_j})P(B|{A_j})}}{{\sum\limits_{i = 1}^\infty  {P({A_i})P(B|{A_i})} }},j \geqslant 1\] 对出现在条件概率中的事件，假设 $P(A_j)>0$。\par\smallskip {\small 来源：闫一博公开解答，原题 2.31。}
\end{exercise}

\begin{proof}\phantomsection\label{a:Y2.31}
根据条件概率公式
  \[P({A_j}|B) = \frac{{P({A_j}B)}}{{P(B)}} = \frac{{P({A_j})P(B|{A_j})}}{{P(B)}}\]
由于$A_1,A_2,\cdots$互不相容，故$A_jB$也互不相容，再根据全概率公式
\[\begin{aligned}P({A_j}|B) &= \frac{{P({A_j})P(B|{A_j})}}{{P(B)}} \\&= \frac{{P({A_j})P(B|{A_j})}}{{\sum\limits_{i = 1}^\infty  {P({A_i})P(B|{A_i})} }}\end{aligned}\]

\end{proof}

\begin{exercise}[串联与并联可靠性]\label{q:Y2.32}
设元件$A$和$B$在一年内烧断的概率分别是$p_1$和$p_2$.设它们是否烧断是相互独立的.在一年内，

（a）当$A,B$并联时，求$A,B$构成的系统断电的概率；

（b）当$A,B$串联时，求$A,B$构成的系统断电的概率.\par\smallskip {\small 来源：闫一博公开解答，原题 2.32。}
\end{exercise}

\begin{solve}\phantomsection\label{a:Y2.32}
记元件$A$在一年中烧坏为事件$A$，元件$B$在一年中烧坏为事件$B$

（a）并联时，两元件均烧坏则系统断电，两元件烧断是相互独立的
\[P(AB) = P(A)P(B) = {p_1}{p_2}\]
（b）串联时，两元件至少有一个烧坏则系统断电，根据加法公式
\[P(A \cup B) = P(A) + P(B) - P(AB) = {p_1} + {p_2} - {p_1}{p_2}\]
\end{solve}

\begin{exercise}[比赛局数与胜率]\label{q:Y2.33}
如果你的水平高于对手，为保证比赛的胜利，你期望三局两胜的比赛规则还是五局三胜的比赛规则？ 设各局结果独立，每局获胜概率均为 $p\in(1/2,1)$，无平局。\par\smallskip {\small 来源：闫一博公开解答，原题 2.33。}
\end{exercise}

\begin{solve}\phantomsection\label{a:Y2.33}令 $p>1/2$ 为甲每局获胜概率，$q=1-p$。三局两胜的概率为 $P_3=p^2+2p^2q$；五局三胜的概率为 $P_5=p^3+3p^3q+6p^3q^2$。按决胜局前已经取得的胜负数计数，以上各项互不重叠。化简得\[P_5-P_3=3p^2q^2(2p-1)>0.\]所以在各局独立且胜率固定的模型下，实力较强的一方更偏好五局三胜；若每局胜率随体力或其他状态变化，这个比较不能直接沿用。\end{solve}

\begin{exercise}[集体迟到的提前量]\label{q:Y2.34}
设开会时每个人是否迟到是相互独立的.如果每个人迟到至少$t$min的概率是$(t+1)^{-2}/100$，在组织一个有$1000$人参加的大会时，要以$0.99$的概率保证无人迟到，应当通知提前多长时间入场.\par\smallskip {\small 来源：闫一博公开解答，原题 2.34。}
\end{exercise}

\begin{solve}\phantomsection\label{a:Y2.34}
由于每个人是否迟到相互独立，记无人迟到$t$min以上为事件$A_t$，则
\[P({A_t}) = {(1 - \frac{1}{{100{{(t + 1)}^2}}})^{1000}} > 0.99\]
计算可得$t\geqslant 31$min.
\end{solve}

\clearpage\section{随机变量及其分布}

\begin{exercise}[分布函数的无穷远极限]\label{q:Y3.1}
用概率的连续性证明分布函数的性质：
\[\begin{aligned}\mathop {\lim }\limits_{x \to \infty } F(x) &= F(\infty ) \\&= 1,\mathop {\lim }\limits_{x \to  - \infty } F(x) \\&= F( - \infty ) \\&= 0\end{aligned}\]\par\smallskip {\small 来源：闫一博公开解答，原题 3.1。}
\end{exercise}

\begin{solve}\phantomsection\label{a:Y3.1}
\[\begin{aligned}1&=P( - \infty  < X < \infty ) \\&= \sum\limits_{n =  - \infty }^\infty  {P(n < X \leqslant n + 1)}  \\&= \sum\limits_{n =  - \infty }^\infty  {F(n + 1) - F(n)} \end{aligned}\]
于是有
\[\mathop {\lim }\limits_{n \to \infty } F(n) - \mathop {\lim }\limits_{m \to  - \infty } F(m) = 1\]
由于分布函数是单调不减的有界函数，故极限存在（上式有意义），且
\[\begin{aligned}\mathop {\lim }\limits_{n \to   \infty } F\left( n \right) &= \mathop {\lim }\limits_{x \to   \infty } F(x),\mathop {\lim }\limits_{m \to  - \infty } F\left( m \right) \\&= \mathop {\lim }\limits_{x \to  - \infty } F(x)\end{aligned}\]
由于$0\leqslant F(x)\leqslant 1$，由$F(x)$的单调性可知
\[\begin{aligned}\mathop {\lim }\limits_{x \to \infty } F(x) &= F(\infty ) \\&= 1,\mathop {\lim }\limits_{x \to  - \infty } F(x) \\&= F( - \infty ) \\&= 0\end{aligned}\]
\end{solve}

\begin{exercise}[指数型密度的归一化]\label{q:Y3.2}
设$g(x)=\lambda e^{-3x},x\geqslant 0$是概率密度，求常数$\lambda$.\par\smallskip {\small 来源：闫一博公开解答，原题 3.2。}
\end{exercise}

\begin{solve}\phantomsection\label{a:Y3.2}
根据归一化条件
\[1 = \int_{ - \infty }^\infty  {g(x)dx}  = \int_0^\infty  {\lambda {e^{ - 3x}}dx}  = \frac{\lambda }{3}\]
故$\lambda=3$.
\end{solve}

\begin{exercise}[二次指数可积的条件]\label{q:Y3.3}
如果常数$a,b,c$使得$g(x)=\exp(ax^2+bx+c),x\in (-\infty,\infty)$是概率密度，则必有$a<0$.\par\smallskip {\small 来源：闫一博公开解答，原题 3.3。}
\end{exercise}

\begin{proof}\phantomsection\label{a:Y3.3}
反证假设$a\geqslant 0$，则有$\exp(ax^2)\geqslant 1$，于是有
\[\begin{aligned}\int_{ - \infty }^\infty  {g(x)dx}  &= \int_{ - \infty }^\infty  {\exp (a{x^2} + bx + c)dx}  \geqslant \int_{ - \infty }^\infty  {\exp (bx + c)dx}  \\&= \left. {\frac{1}{b}\exp (bx + c)} \right|_{ - \infty }^\infty \end{aligned}\]
由于$a,b,c$均为常数，故无论上式中$b$取何值，均不满足归一化条件，反正假设不真，$a<0$.
\end{proof}

\begin{exercise}[最大值与最小值恒等式]\label{q:Y3.4}
设$X,Y$是随机变量，验证

（a）$\max(X,Y)=(|X-Y|+X+Y)/2$；

（b）$\min(X,Y)=(X+Y-|X-Y|)/2$；

（c）$\max(X,Y)+\min(X,Y)=X+Y$.\par\smallskip {\small 来源：闫一博公开解答，原题 3.4。}
\end{exercise}

\begin{solve}\phantomsection\label{a:Y3.4}对每个具体样本点，若 $X\ge Y$，则 $|X-Y|=X-Y$，代入两式分别得到 $X$ 与 $Y$；若 $X<Y$，则 $|X-Y|=Y-X$，分别得到 $Y$ 与 $X$。因此\[\max(X,Y)=\frac{|X-Y|+X+Y}{2},\]\[\min(X,Y)=\frac{X+Y-|X-Y|}{2}.\]将这两个恒等式相加，绝对值项抵消，得到 $\max(X,Y)+\min(X,Y)=X+Y$。结论逐点成立，不需要任何分布、矩或独立性假设。\end{solve}

\begin{exercise}[分段密度的归一化]\label{q:Y3.5}
设$X$有概率密度$f(x) = \left\{ \begin{gathered}
  (x - 1)/a \hfill \quad x\in[1,2]\\
  (x - 3)/a \hfill \quad x\in[3,5]\\
  0 \hfill \quad \text{其他}\\ 
\end{gathered}  \right.$，计算$a$.\par\smallskip {\small 来源：闫一博公开解答，原题 3.5。}
\end{exercise}

\begin{solve}\phantomsection\label{a:Y3.5}
由于$f(x)$是$X$的概率密度，根据归一化条件
\[1 = \int_{ - \infty }^\infty  {f(x)dx}  = \int_1^2 {\frac{{(x - 1)}}{a}dx}  + \int_3^5 {\frac{{(x - 3)}}{a}dx} \]
计算可得
\[1 = \frac{1}{{2a}} - \frac{5}{{2a}} + \frac{9}{{2a}} = \frac{5}{{2a}} \Rightarrow a = \frac{5}{2}\]
\end{solve}

\begin{exercise}[二项计数的尾概率]\label{q:Y3.6}
设某手机一天收到了$8$个短信，每个短信是公事短信的概率是$0.2$，是私事短信的概率是$0.8$.用$X$表示这天收到的公事短信数，计算$X$的概率分布和$P(X\geqslant 3)$. 假设各短信类别相互独立。\par\smallskip {\small 来源：闫一博公开解答，原题 3.6。}
\end{exercise}

\begin{solve}\phantomsection\label{a:Y3.6}设 $X$ 为八条短信中相应类别的条数。独立性及每条概率为 $0.2$ 给出 $X\sim\operatorname{Bin}(8,0.2)$。因此\[P(X\ge3)=1-\sum_{k=0}^{2}\binom8k0.2^k0.8^{8-k}=0.20308224.\]取补事件时必须包含 $X=0$，不能从 $k=1$ 开始相减。\end{solve}

\begin{exercise}[两组二项计数的比较]\label{q:Y3.7}
甲、乙击中目标的概率分别是$0.6,0.7$，如果各射击$3$次，

（a）计算击中次数相同的概率；

（b）计算甲击中的次数多的概率. 假设全部射击结果相互独立。\par\smallskip {\small 来源：闫一博公开解答，原题 3.7。}
\end{exercise}

\begin{solve}\phantomsection\label{a:Y3.7}
记甲击中目标的次数为随机变量$X$，乙击中目标的次数为随机变量$Y$，$X\sim B(3,0,6),Y\sim B(3,0,7)$

（a）甲乙击中目标相互独立，故
\[P(X = Y) = \sum\limits_{k = 0}^3 {P(X = k)P(Y = k)} \]
计算可得
\[P(X = Y) = 0.32076\]

（b）甲乙击中目标相互独立，故
\[P(X > Y) = \sum\limits_{0 \leqslant j < i \leqslant 3} {P(X = i)P(Y = j)} \]
计算可得
\[P(X > Y) = 0.243\]
\end{solve}

\begin{exercise}[首次与第多次成功]\label{q:Y3.8}
编辑正在校对一本书稿.假设该书稿的每页有一个打印错误的概率是$0.05$，没有打印错误的概率是$0.95$.计算编辑

（a）在第$k$页发现第一个打印错误的概率

（b）在第$k$页发现第$r$个打印错误的概率

（c）发现第$r$个错误前，恰有$k$页无打印错误的概率

（d）如果这本书一共有$n$页，计算全书打印错误数的概率分布 假设各页错误情况相互独立，且每页至多一个错误。\par\smallskip {\small 来源：闫一博公开解答，原题 3.8。}
\end{exercise}

\begin{solve}\phantomsection\label{a:Y3.8}
（a）记第一个打印错误所在页数为随机变量$X$，则$X$服从参数为$0.05$的几何分布，故
\[P(X=k)=0.95^{k-1} 0.05\]
（b）记第$r$打印错误所在页数为随机变量$Y$，则$Y$服从参数为$(r,0.05)$的Pascal分布，故
\[P(Y = k) = C_{k - 1}^{r - 1}{0.95^{k - r}}{0.05^r}\]
（c）记发现第$r$个错误前，无打印错误的页数为随机变量$Z$，则$Z$服从参数为$(r,0.05)$的负二项分布，故 
\[P(Z = k) = C_{k + r - 1}^{r - 1}{0.95^k}{0.05^r}\]
（d）记$n$页书的打印错误数为随机变量$\xi$，则$\xi\sim B(n,0.05)$
\end{solve}

\begin{exercise}[泊松分裂与独立性]\label{q:Y3.9}
甲每天收到的电子邮件数服从泊松分布$P(\lambda)$，且每封电子邮件被过滤掉的概率是$0.2$.

（a）当有$n$封电子邮件发给甲，计算他见到其中$k$封的概率$h_k$；

（b）计算甲每天见到的电子邮件数的分布；

（c）已知甲看到了自己的$k$封电子邮件，计算他有$m$封被过滤掉的概率；

（d）甲每天见到的邮件数和被滤掉的邮件数是否独立. 给定邮件总数后，假设各封邮件独立地被过滤。\par\smallskip {\small 来源：闫一博公开解答，原题 3.9。}
\end{exercise}

\begin{solve}\phantomsection\label{a:Y3.9}记总邮件数为 $N$，看到的数量为 $X$，过滤数量为 $Z=N-X$。给定 $N=n$，有 $X\sim\operatorname{Bin}(n,0.8)$，所以 $h_k=\binom nk0.8^k0.2^{n-k}$。对非负整数 $k,m$，\[\begin{aligned}P(X=k,Z=m)&=P(N=k+m)P(X=k\mid N=k+m)\\&=e^{-\lambda}\frac{\lambda^{k+m}}{k!m!}0.8^k0.2^m\\&=\left[e^{-0.8\lambda}\frac{(0.8\lambda)^k}{k!}\right]\left[e^{-0.2\lambda}\frac{(0.2\lambda)^m}{m!}\right].\end{aligned}\]联合质量函数分解为两个泊松质量函数，故 $X\sim\operatorname{Pois}(0.8\lambda)$、$Z\sim\operatorname{Pois}(0.2\lambda)$ 且独立。因此已知 $X=k$ 后，$Z$ 仍为率 $0.2\lambda$ 的泊松分布。\end{solve}

\begin{exercise}[周期信号的均匀等待]\label{q:Y3.13}
求救者每间隔$2$min发出一次瞬时呼叫，随机到达的救援者在收到呼叫信号的范围内至少停留多长时间才能以$0.95$的概率收到呼叫.\par\smallskip {\small 来源：闫一博公开解答，原题 3.13。}
\end{exercise}

\begin{solve}\phantomsection\label{a:Y3.13}
记救援者在$t$时刻到达可收到呼叫信号的范围，求救者发出求救信号为随机变量$X$，则$X\sim U(t,t+2)$，故
\[0.95 \leqslant P(t < X \leqslant t + \Delta t) = \int_t^{t + \Delta t} {\frac{1}{2}dt}  = \frac{1}{2}\Delta t\]
这表明至少停留$1.9$min才能以$0.95$的概率收到呼叫.
\end{solve}

\begin{exercise}[恒定风险率与指数寿命]\label{q:Y3.14}
一个使用了$t$小时的热敏电阻在$\Delta t$内失效的概率是$\lambda \Delta t+o(\Delta t)$，设该热敏电阻的使用寿命是连续性随机变量，求该热敏电阻的寿命分布. 此处的失效概率指已存活到 $t$ 时的条件概率，且初始寿命大于零。\par\smallskip {\small 来源：闫一博公开解答，原题 3.14。}
\end{exercise}

\begin{solve}\phantomsection\label{a:Y3.14}
记热敏电阻的寿命（失效时间）为随机变量$X$，根据题意
\[P(X \leqslant t + \Delta t|X > t) = \lambda \Delta t + o(\Delta t)\]
记$X$的分布函数为$F(t)$，相应的概率密度为$f(t)$，于是有
\[\begin{aligned}\lambda \Delta t + o(\Delta t) &= \frac{{P(t < X \leqslant t + \Delta t)}}{{P(X > t)}} \\&= \frac{{F(t + \Delta t) - F(t)}}{{1 - F(t)}}\end{aligned}\]
两端同除$\Delta t$，取极限可得
\[\begin{aligned}\mathop {\lim }\limits_{\Delta t \to 0} \lambda  + o(1) &= \frac{1}{{1 - F(t)}}\mathop {\lim }\limits_{\Delta t \to 0} \frac{{F(t + \Delta t) - F(t)}}{{\Delta t}} \Rightarrow \lambda  \\&= \frac{{f(t)}}{{1 - F(t)}}\end{aligned}\]
上式右端是一个变量分离方程，积分可得
\[ - \ln (1 - F(t)) = \lambda t + C \Rightarrow F(t) = 1 - C{e^{ - \lambda t}}\]
根据归一化条件可得
\[\int_0^\infty  {f(t)dt}  = 1 \Rightarrow C = 1\]
于是得到
\[F(t) = 1 - {e^{ - \lambda t}}\]
这表明$X\sim \varepsilon(\lambda)$.
\end{solve}

\begin{exercise}[正态尺寸的合格概率]\label{q:Y3.15}
一台机床加工的部件长度服从正态分布$N(8,36\times 10^{-8})$.当部件的长度在$8\pm 0.0015$内为合格品，求一部件为合格品的概率.\par\smallskip {\small 来源：闫一博公开解答，原题 3.15。}
\end{exercise}

\begin{solve}\phantomsection\label{a:Y3.15}
记机床加工的部件长度为随机变量$X$，则$X\sim N(8,36\times 10^{-8})$，于是有
\[\begin{aligned}P(8 - 0.0015 \leqslant X \leqslant 8 + 0.0015) &= P( - 2.5 \leqslant \frac{{X - 8}}{{0.0006}} \leqslant 2.5) \\&= 2\Phi (2.5) - 1 \\&= 0.9876\end{aligned}\]
\end{solve}

\begin{exercise}[泊松变量的平方根]\label{q:Y3.16}
设$X\sim P(\lambda)$，计算$Y=\sqrt{X}$的概率分布.\par\smallskip {\small 来源：闫一博公开解答，原题 3.16。}
\end{exercise}

\begin{solve}\phantomsection\label{a:Y3.16}
$P(Y\geqslant 0)=1$，由离散变量的一一对应变换
\[P(Y = \sqrt k ) = P(X = k) = \frac{{{\lambda ^k}}}{{k!}}{e^{ - \lambda }},k = 0,1,2, \cdots \]
\end{solve}

\begin{exercise}[均匀电流的平方变换]\label{q:Y3.17}
设电流$I$在$8\sim 9$A之间均匀分布.当电流通过$2\Omega$的电阻时，消耗的功率（单位：W）是$W=2I^2$，求$W$的概率密度.\par\smallskip {\small 来源：闫一博公开解答，原题 3.17。}
\end{exercise}

\begin{solve}\phantomsection\label{a:Y3.17}
根据题意$I\sim U(8,9)$，故有概率密度$f_I(i)=1.i\in (8,9)$，$P(W>0)=1$，于是有
\[P(W = w) = P(2{I^2} = w) = P(I = \sqrt {\frac{w}{2}} ) = \frac{1}{4}\sqrt {\frac{2}{w}} dw\]
这表明$W$有概率密度
\[{f_W}(w) = \frac{1}{4}\sqrt {\frac{2}{w}},w\in(128,162) \]
\end{solve}

\begin{exercise}[指数变量的仿射变换]\label{q:Y3.18}
设$X\sim \varepsilon(\lambda)$，$Y=aX+b,a>0$，求$Y$的概率密度.\par\smallskip {\small 来源：闫一博公开解答，原题 3.18。}
\end{exercise}

\begin{solve}\phantomsection\label{a:Y3.18}
根据题意$X\sim \varepsilon(\lambda)$，故有概率密度$f_X(x)=\lambda e^{-\lambda x}$，由于$a>0$，故$P(Y>b)=1$，于是有
\[P(Y = y) = P(aX + b = y) = P(X = \frac{{y - b}}{a}) = \frac{\lambda }{a}{e^{ - \frac{{\lambda (y - b)}}{a}}}dy\]
这表明$Y$有概率密度
\[{f_Y}(y) = \frac{\lambda }{a}{e^{ - \frac{{\lambda (y - b)}}{a}}},y > b\]
\end{solve}

\begin{exercise}[半柯西变量的对数]\label{q:Y3.19}
$X$有概率密度
\[f(x)=\frac{2}{\pi(1+x^2)},x\geqslant 0\]
求$Y=\ln X$的概率密度.\par\smallskip {\small 来源：闫一博公开解答，原题 3.19。}
\end{exercise}

\begin{solve}\phantomsection\label{a:Y3.19}
$P(Y\in \mathbb{R})=1$，于是有 
\[P(Y = y) = P(\ln X = y) = P(X = {e^y}) = \frac{{2{e^y}}}{{\pi (1 + {e^{2y}})}}dy\]
这表明$Y$有概率密度
\[{f_Y}(y) = \frac{{2{e^y}}}{{\pi (1 + {e^{2y}})}},y \in ( - \infty ,\infty )\]
\end{solve}

\begin{exercise}[倒数、绝对值与正切变换]\label{q:Y3.20}
设$X$有分段连续的概率密度$f(x)$.

（a）求$Y=X^{-1}$的概率密度；

（b）求$Y=|X|$的概率密度；

（c）求$Y=\tan(X)$的概率密度.\par\smallskip {\small 来源：闫一博公开解答，原题 3.20。}
\end{exercise}

\begin{solve}\phantomsection\label{a:Y3.20}使用分段单调变换的密度公式。连续分布在零点及正切奇点处的总概率为零，所以可在这些点任意定义变换。\par（a）对 $y\ne0$，逆变换为 $x=1/y$，导数绝对值为 $1/y^2$，故\[f_{1/X}(y)=y^{-2}f(1/y).\]（b）对 $y>0$，两个逆像为 $y,-y$，因此\[f_{|X|}(y)=f(y)+f(-y),\]负半轴密度为零。\par（c）正切函数在各区间 $(-\pi/2+k\pi,\pi/2+k\pi)$ 上单调，逆像为 $x=\arctan y+k\pi$。因此\[f_{\tan X}(y)=\frac1{1+y^2}\sum_{k\in\mathbb Z}f(\arctan y+k\pi).\]各分支贡献非负，可以用非负积分交换定理求和；所得密度公式几乎处处成立。\end{solve}

\begin{exercise}[二项众数与参数估计]\label{q:Y3.21}
设$X$服从二项分布$B(n,p)$.

（a）已知$n=19,p=0.7$，求$p_k=P(X=k)$的最大值点$k$；

（b）已知$n=19,X=9$，求使得$P(X=9)$达到最大的$p$.\par\smallskip {\small 来源：闫一博公开解答，原题 3.21。}
\end{exercise}

\begin{solve}\phantomsection\label{a:Y3.21}二项质量函数相邻比值为\[\frac{p_{k+1}}{p_k}=\frac{n-k}{k+1}\frac{p}{1-p}.\]它大于一等价于 $k+1<(n+1)p$。本题 $(n+1)p=20(0.7)=14$ 是整数，故 $k=13$ 与 $k=14$ 都达到最大概率，不能只保留一个。观察到 $X=9$ 后，以 $p$ 为变量的似然核为 $p^9(1-p)^{10}$，对数导数为 $9/p-10/(1-p)$，驻点为 $p=9/19$，二阶导数为负，故它是极大似然估计。\end{solve}

\begin{exercise}[迟到人数的众数和估计]\label{q:Y3.22}
全班有$95$个学生，每个学生上课迟到的概率为$p$.假设每个学生是否迟到是相互独立的.

（a）当$p=0.05$时，有多少个学生迟到的概率最大；

（b）如果有$7$个人迟到，你认为$p$是多少？\par\smallskip {\small 来源：闫一博公开解答，原题 3.22。}
\end{exercise}

\begin{solve}\phantomsection\label{a:Y3.22}迟到人数 $X\sim\operatorname{Bin}(95,p)$。二项分布的相邻比值表明，当 $(n+1)p$ 不是整数时唯一众数为其整数部分；若是整数则有相邻两个众数。本题 $96(0.05)=4.8$，所以最可能迟到人数为 $4$。观察到七人迟到时，似然核为 $p^7(1-p)^{88}$，对数导数过零点为 $p=7/95\approx0.073684$，即极大似然估计。众数求的是固定参数下最可能的结果，估计求的是固定结果下最合适的参数。\end{solve}

\begin{exercise}[泊松众数与似然]\label{q:Y3.23}
设$X$是服从参数是$\lambda$的泊松分布.

（a）已知$\lambda=23.8$，求$p_k=P(X=k)$的最大值点$k$；

（b）已知$X=21$，求使得$P(X=21)$达到最大的$\lambda$.\par\smallskip {\small 来源：闫一博公开解答，原题 3.23。}
\end{exercise}

\begin{solve}\phantomsection\label{a:Y3.23}泊松相邻概率比为 $p_{k+1}/p_k=\lambda/(k+1)$。当 $\lambda=23.8$ 时，唯一众数为 $23$。观察到 $X=21$ 后，对数似然为 $-\lambda+21\log\lambda-\log(21!)$，导数为 $-1+21/\lambda$，故极大似然估计为 $\widehat\lambda=21$。若泊松参数恰为正整数 $r$，则 $r-1,r$ 同为众数，非整数情形才具有唯一的整数部分众数。\end{solve}

\begin{exercise}[计数结果与参数的区别]\label{q:Y3.24}
假设你每天收到的短信数服从参数是$\lambda$的泊松分布.

（a）已知$\lambda=8.9$时，你今天收到多少个短信的概率最大？

（b）如果你今天收到了$12$个短信，你认为$\lambda$应当是多少？\par\smallskip {\small 来源：闫一博公开解答，原题 3.24。}
\end{exercise}

\begin{solve}\phantomsection\label{a:Y3.24}若 $X\sim\operatorname{Pois}(\lambda)$，相邻概率之比为\[\frac{P(X=k+1)}{P(X=k)}=\frac{\lambda}{k+1}.\]当 $\lambda=8.9$ 时，概率先增加后减少，唯一众数为 $8$。另一方面，已观测到 $X=12$ 时，关于 $\lambda$ 的似然为 $L(\lambda)=e^{-\lambda}\lambda^{12}/12!$，其对数导数为 $-1+12/\lambda$，故极大似然估计为 $\widehat\lambda=12$。前问固定参数比较不同结果，后问固定结果比较不同参数，两个优化变量不同。\end{solve}

\begin{exercise}[骰子样本的最大最小值]\label{q:Y3.25}
将一个骰子投掷$n$次，用$m$表示掷得的最小点数，用$M$表示掷得的最大点数，计算：

（a）$P(m=k),1\leqslant k\leqslant 6$；

（b）$P(M=k),1\leqslant k\leqslant 6$；

（c）$P(m=2,M=5)$.\par\smallskip {\small 来源：闫一博公开解答，原题 3.25。}
\end{exercise}

\begin{solve}\phantomsection\label{a:Y3.25}假定各次骰子公平且独立。最小值至少为 $k$ 的概率为 $((7-k)/6)^n$，所以\[P(m=k)=\left(\frac{7-k}{6}\right)^n-\left(\frac{6-k}{6}\right)^n.\]最大值不超过 $k$ 的概率为 $(k/6)^n$，所以\[P(M=k)=\left(\frac{k}{6}\right)^n-\left(\frac{k-1}{6}\right)^n.\]要求最小值为二、最大值为五，等价于所有点数在 $\{2,3,4,5\}$ 且二与五都出现。容斥得\[P(m=2,M=5)=\left(\frac46\right)^n-2\left(\frac36\right)^n+\left(\frac26\right)^n.\]对 $n=1$ 该式为零，与两个极值不可能不同相符。\end{solve}

\begin{exercise}[圆周投影的反正弦密度]\label{q:Y3.26}
在半径为 $R>0$ 的圆周上按弧长均匀取一点，圆心为原点。求该点横坐标 $X$ 的分布函数及密度。\par\smallskip {\small 来源：闫一博公开解答，原题 3.26。}
\end{exercise}

\begin{solve}\phantomsection\label{a:Y3.26}令 $\Theta\sim U(0,2\pi)$，则 $X=R\cos\Theta$。对 $-R<x<R$，满足 $\cos\Theta\le x/R$ 的角度集合长度为 $2\pi-2\arccos(x/R)$，所以\[F_X(x)=1-\frac{1}{\pi}\arccos\frac{x}{R}.\]当 $x\le-R$ 时分布函数为 $0$，当 $x\ge R$ 时为 $1$。在内部求导得\[f_X(x)=\frac{1}{\pi\sqrt{R^2-x^2}},\quad |x|<R,\]其余处密度为零。端点处密度可任意定义。连续随机变量的单点概率为零，因此不能把 $P(X=x)$ 当作密度。\end{solve}

\begin{exercise}[正弦变换的两条分支]\label{q:Y3.28}
设$X$有概率密度$f(x)=cx/\pi^2,x\in (0,\pi)$.求$Y=\sin X$的概率密度.\par\smallskip {\small 来源：闫一博公开解答，原题 3.28。}
\end{exercise}

\begin{solve}\phantomsection\label{a:Y3.28}归一化条件为 $c\int_0^\pi x\,dx/\pi^2=c/2=1$，所以 $c=2$。当 $0<y<1$ 时，两个逆像为 $x_1=\arcsin y$ 和 $x_2=\pi-\arcsin y$，导数绝对值都为 $1/\sqrt{1-y^2}$。于是\[f_Y(y)=\frac{f_X(x_1)+f_X(x_2)}{\sqrt{1-y^2}}=\frac{2}{\pi\sqrt{1-y^2}}.\]其余处密度为零。两条分支都要计入；只保留 $\arcsin y$ 会漏掉 $X>\pi/2$ 的贡献。\end{solve}

\begin{exercise}[逆分布函数生成正态数]\label{q:Y3.29}
设$X \sim  U(0,1)$，$\Phi^{-1}(x)$是$\Phi(x)$的反函数，求$Y=\Phi^{-1} (X)$的分布函数. 其中 $\Phi$ 是标准正态分布函数。\par\smallskip {\small 来源：闫一博公开解答，原题 3.29。}
\end{exercise}

\begin{solve}\phantomsection\label{a:Y3.29}因 $\Phi$ 严格递增、连续且从实轴映到 $(0,1)$，对任意实数 $y$，\[P(Y\le y)=P(\Phi^{-1}(X)\le y)=P(X\le\Phi(y))=\Phi(y).\]因此 $Y$ 的分布函数正是标准正态分布函数，故 $Y\sim N(0,1)$。这个论证直接使用分布函数，不需要把零概率单点事件当作密度。\end{solve}

\begin{exercise}[分段单调变换]\label{q:Y3.31}
设$X$服从参数为$\lambda$的指数分布，求
\[Y = \left\{ \begin{gathered}
  X \hfill \quad X\geqslant 1\\
  {X^2} \hfill \quad  X<1\\ 
\end{gathered}  \right.\]

的概率密度.\par\smallskip {\small 来源：闫一博公开解答，原题 3.31。}
\end{exercise}

\begin{solve}\phantomsection\label{a:Y3.31}变换在正半轴严格增加且在 $1$ 处连续，逆函数在 $0<y<1$ 为 $\sqrt y$，在 $y>1$ 为 $y$。因此\[f_Y(y)=\begin{cases}\lambda e^{-\lambda\sqrt y}/(2\sqrt y),&0<y<1,\\\lambda e^{-\lambda y},&y>1,\\0,&y<0.\end{cases}\]在 $y=1$ 处没有原子，密度值可任意指定。两段积分分别为 $P(X<1)=1-e^{-\lambda}$ 和 $P(X>1)=e^{-\lambda}$，总和为一。\end{solve}

\clearpage\section{联合分布与变量变换}

\begin{exercise}[联合概率表与边缘分布]\label{q:Y4.3}
设离散型概率分布$p_{ij}=P(X=i,Y=j)$如下表所示，行指标为 $X$ 的取值，列指标为 $Y$ 的取值。

\begin{center}
    
    \begin{tabular}{|c|c|c|c|c|c|} \hline 
 $p_{ij}$&  $1$&$2$&$3$&$4$&$5$\\ \hline 
 $1$&  $0.06$&$0.05$    &$0.04$&$0.01$&$0.02$\\ \hline 
         $2$&    $0.05$&$0.10$&$0.10$&$0.05$&$0.03$\\ \hline 
         $3$& 
      $0.07$&$0.05$& $0.01$&$0.02$&$0.02$\\ \hline
 $4$& $0.05$& $0.02$& $0.01$&$0.01$&$0.03$\\\hline
 $5$& $0.05$& $0.06$& $0.05$& $0.02$&$0.02$\\\hline\end{tabular}
\end{center}

求$X,Y$的边缘分布；\par\smallskip {\small 来源：闫一博公开解答，原题 4.3。}
\end{exercise}

\begin{solve}\phantomsection\label{a:Y4.3}固定一个变量的取值，对另一个变量求和即可得到边缘概率。具体地，$P(X=i)=\sum_{j=1}^5p_{ij}$，所以按行相加得到
\begin{center}
    
    \begin{tabular}{|c|c|c|c|c|c|} \hline 
         $X$&  $1$& $2$&$3$& $4$&$5$\\ \hline 
         $P$& 
     $0.18$& $0.33$& $0.17$& $0.12$&$0.20$\\ \hline\end{tabular}

\end{center}
同理，$P(Y=j)=\sum_{i=1}^5p_{ij}$，按列相加得到\begin{center}
    
    \begin{tabular}{|c|c|c|c|c|c|} \hline 
         $Y$&  $1$& $2$&$3$& $4$&$5$\\ \hline 
         $P$& 
     $0.28$& $0.28$& $0.21$& $0.11$&$0.12$\\ \hline\end{tabular}

\end{center}

两组边缘概率各自之和均为 $1$，可据此检查求和是否遗漏。\end{solve}

\begin{exercise}[同参数二项变量之和]\label{q:Y4.6}
设$X\sim B(n,p),Y\sim B(m,p)$，当$X,Y$独立，计算$X+Y$的概率分布.\par\smallskip {\small 来源：闫一博公开解答，原题 4.6。}
\end{exercise}

\begin{solve}\phantomsection\label{a:Y4.6}概率生成函数在独立和下相乘：\[E s^{X+Y}=(1-p+ps)^n(1-p+ps)^m=(1-p+ps)^{n+m}.\]这正是 $\operatorname{Bin}(n+m,p)$ 的生成函数，所以\[P(X+Y=k)=\binom{n+m}{k}p^k(1-p)^{n+m-k}.\]也可以把 $X,Y$ 分别表示成 $n,m$ 个独立伯努利试验的成功数，合并即为 $n+m$ 次试验。两组成功概率必须相同；参数不同的二项和一般不再是二项分布。\end{solve}

\begin{exercise}[函数依赖的离散向量]\label{q:Y4.8}
如果随机向量$(X,Y)$有如下概率分布：
\[
P(X=i,Y=\frac{1}{i})=c,i=1,2,\cdots,8
\]
确定常数$c$并求$X,Y$的概率分布.\par\smallskip {\small 来源：闫一博公开解答，原题 4.8。}
\end{exercise}

\begin{solve}\phantomsection\label{a:Y4.8}
    \[\sum\limits_{i = 1}^8 {P(X = i,Y = \frac{1}{i})}  = 1 \Rightarrow c = \frac{1}{8}\]
    于是得到$X,Y$的概率分布为
    \[P(X = i) = \frac{1}{8},P(Y = \frac{1}{i}) = \frac{1}{8},i = 1,2, \cdots ,8\]
\end{solve}

\begin{exercise}[最小值的尾概率]\label{q:Y4.9}
设随机变量$X,Y$独立同分布，证明
\[
P(a<\min (X,Y)\leqslant b)=P^2(X>a)-P^2(X>b)
\]\par\smallskip {\small 来源：闫一博公开解答，原题 4.9。}
\end{exercise}

\begin{solve}\phantomsection\label{a:Y4.9}设 $Z=\min(X,Y)$，则 $\{Z>a\}=\{X>a,Y>a\}$。独立同分布给出 $P(Z>a)=P(X>a)^2$。又因 $a<b$，事件 $\{Z>b\}$ 包含于 $\{Z>a\}$，故\[\begin{aligned}P(a<Z\le b)&=P(Z>a)-P(Z>b)\\&=P(X>a)^2-P(X>b)^2.\end{aligned}\]此证明同时适用于连续、离散和混合分布；端点的严格与非严格关系由事件差确定，不需要另加连续性假设。\end{solve}

\begin{exercise}[三角支撑上的边缘密度]\label{q:Y4.15}
设随机向量$(X,Y)$有联合密度
\[f(x,y) = \left\{ \begin{gathered}
  {e^{ - x}} \hfill \quad 0<y<x\\
  0 \hfill \quad \text{其他}\\ 
\end{gathered}  \right.\]
求$X,Y$的边缘密度.\par\smallskip {\small 来源：闫一博公开解答，原题 4.15。}
\end{exercise}

\begin{solve}\phantomsection\label{a:Y4.15}
\[{f_X}(x) = \int_0^x {{e^{ - x}}dy}  = x{e^{ - x}}\]
    \[{f_Y}(y) = \int_y^\infty  {{e^{ - x}}dx}  = {e^{ - y}}\]
 以上分别在 $x>0$、$y>0$ 成立，其余处密度为零。\end{solve}

\begin{exercise}[随机家庭规模与男孩数]\label{q:Y4.16}
设$p=1-q\in (0,1)$，$0<\alpha<q/p$.假设一个家庭有$n$个小孩的概率是
\[{p_n} = \left\{ \begin{gathered}
  \alpha {p^n} \hfill \quad n\geqslant 1\\
  1 - \alpha p/q \hfill \quad n=0\\ 
\end{gathered}  \right.\]
如果男孩和女孩的出生是等可能的，求一个家庭有$k$个男孩的概率.\par\smallskip {\small 来源：闫一博公开解答，原题 4.16。}
\end{exercise}

\begin{solve}\phantomsection\label{a:Y4.16}
用$X$表示一个家庭的男孩个数，$Y$表示一个家庭的小孩数.    
\[P(X = k) = P(Y = n)P(X = k|Y = n)\]
    （i）$k=0$时，根据全概率公式
    \[\begin{aligned}P(X &= 0) \\&= 1 - \alpha \frac{p}{q} + {\sum\limits_{n = 1}^\infty  {\alpha {p^n}\left( {\frac{1}{2}} \right)} ^n} \\&= 1 - \alpha \frac{p}{q} + \alpha \frac{{\frac{p}{2}}}{{1 - \frac{p}{2}}} \\&= 1 - \frac{{\alpha p}}{q} + \frac{{\alpha p}}{{2 - p}}\end{aligned}\]
    （ii）$k\geqslant 1$时，根据全概率公式
    \[\begin{aligned}P(X &= k) \\&= \alpha {\sum\limits_{n = k}^\infty  {C_n^k\left( {\frac{p}{2}} \right)} ^n} \\&= \alpha {\sum\limits_{m = 0}^\infty  {C_{m + k}^k\left( {\frac{p}{2}} \right)} ^{m + k}} \\&= \alpha {\left( {1 - \frac{p}{2}} \right)^{ - (k + 1)}}{\left( {\frac{p}{2}} \right)^k}\end{aligned}\]
\end{solve}

\begin{exercise}[连续变量的离散随机缩放]\label{q:Y4.22}
设随机变量$X,Y$独立，$X$有概率密度$f(x)$，$Y$有离散型概率分布$P(Y=a_i)=p_i>0,i=1,2,\cdots$，若$a_i$均不为零，证明$Z=XY$有概率密度
\[h(z) = \sum\limits_{i = 1}^\infty  {\frac{{{p_i}}}{{|{a_i}|}}f\left( {\frac{z}{{{a_i}}}} \right)} \]\par\smallskip {\small 来源：闫一博公开解答，原题 4.22。}
\end{exercise}

\begin{solve}\phantomsection\label{a:Y4.22}给定 $Y=a_i\ne0$ 时，$Z=a_iX$ 的密度为 $|a_i|^{-1}f(z/a_i)$。对任意可测集合 $D$，全概率公式给出\[P(Z\in D)=\sum_i p_i\int_D\frac1{|a_i|}f(z/a_i)\,dz.\]被积项非负，交换求和与积分合法，故\[h(z)=\sum_i\frac{p_i}{|a_i|}f(z/a_i).\]它的积分等于 $\sum_i p_i=1$，因此是概率密度。绝对值保证负缩放时密度仍非负。\end{solve}

\begin{exercise}[博克斯—穆勒正态生成]\label{q:Y4.25}
设随机变量$U$和$V$独立，都服从$(0,1)$上的均匀分布.

（a）计算$r=\sqrt{-2\ln U}$和$\theta=2\pi V$的概率密度.

（b）证明$X=r\cos(\theta),Y=r\sin(\theta)$独立，都服从标准正态分布.\par\smallskip {\small 来源：闫一博公开解答，原题 4.25。}
\end{exercise}

\begin{solve}\phantomsection\label{a:Y4.25}对 $r>0$，$P(R>r)=P(U<e^{-r^2/2})=e^{-r^2/2}$，所以 $f_R(r)=re^{-r^2/2}$。又 $\Theta=2\pi V\sim U(0,2\pi)$，且它与 $R$ 独立。极坐标变换的雅可比绝对值为 $r$，除原点与角度接缝这些零面积集合外可逐区域使用变换公式，得到\[f_{X,Y}(x,y)=\frac1{2\pi}e^{-(x^2+y^2)/2}.\]在整个平面上，该密度等于两个标准正态密度的乘积，所以 $X,Y$ 独立且均为标准正态。这里没有把 $\arctan(y/x)$ 当成全平面唯一角度；角度应按象限选取。\end{solve}

\begin{exercise}[指数变量的极值与比值]\label{q:Y4.30}
设随机变量$X,Y$独立，$X\sim \mathscr{E}(\lambda),Y\sim \mathscr{E}(\mu)$.分别求$U=\min (X,Y),V=\max(X,Y),W=X/Y$的概率密度.\par\smallskip {\small 来源：闫一博公开解答，原题 4.30。}
\end{exercise}

\begin{solve}\phantomsection\label{a:Y4.30}当 $u>0$，最小值的生存函数为 $e^{-(\lambda+\mu)u}$，所以\[f_U(u)=(\lambda+\mu)e^{-(\lambda+\mu)u}.\]最大值的分布函数为 $(1-e^{-\lambda v})(1-e^{-\mu v})$，求导得\[f_V(v)=\lambda e^{-\lambda v}+\mu e^{-\mu v}-(\lambda+\mu)e^{-(\lambda+\mu)v},\quad v>0.\]对于比值 $W=X/Y$，采用 $x=wy$，逆雅可比为 $y$，故\[\begin{aligned}f_W(w)&=\lambda\mu\int_0^\infty ye^{-(\lambda w+\mu)y}\,dy\\&=\frac{\lambda\mu}{(\lambda w+\mu)^2},\quad w>0.\end{aligned}\]各密度在负半轴为零。若写最小值与最大值的联合密度，其支撑必须是 $0<u<v$，不能写成整个正象限。\end{solve}

\begin{exercise}[平方和与平方比的变换]\label{q:Y4.34}
设 $X,Y$ 独立且均服从率为 $1$ 的指数分布。求 $U=X^2+Y^2,V=X^2/Y^2$ 的联合密度。\par\smallskip {\small 来源：闫一博公开解答，原题 4.34。}
\end{exercise}

\begin{solve}\phantomsection\label{a:Y4.34}正象限内变换一一对应于 $u>0,v>0$，逆变换为\[x=\sqrt{\frac{uv}{1+v}},\qquad y=\sqrt{\frac{u}{1+v}}.\]计算逆变换的雅可比绝对值，得\[\left|\frac{\partial(x,y)}{\partial(u,v)}\right|=\frac{1}{4\sqrt v(1+v)}.\]由于原密度为 $e^{-x-y}$，故\[f_{U,V}(u,v)=\frac{\exp\{-\sqrt{u/(1+v)}(1+\sqrt v)\}}{4\sqrt v(1+v)},\quad u,v>0.\]在其余区域上密度为零。雅可比乘在密度上，而不是乘在单点事件概率上。\end{solve}

\begin{exercise}[线性变换的密度公式]\label{q:Y4.36}
设 $(X,Y)$ 有联合密度 $f$，可逆矩阵 $A=\begin{pmatrix}a&b\\c&d\end{pmatrix}$，逆矩阵的元素记为 $a_1,b_1,c_1,d_1$。令 $(U,V)^\mathsf T=A(X,Y)^\mathsf T$，证明其联合密度为\[g(u,v)=\frac{f(a_1u+b_1v,c_1u+d_1v)}{|\det A|}.\]\par\smallskip {\small 来源：闫一博公开解答，原题 4.36。}
\end{exercise}

\begin{solve}\phantomsection\label{a:Y4.36}可逆性保证每个 $(u,v)$ 对应唯一逆像 $(a_1u+b_1v,c_1u+d_1v)$，逆变换导数矩阵恒为 $A^{-1}$。因此雅可比绝对值为 $|\det A^{-1}|=1/|\det A|$。对任意可测集合 $D$，变量替换公式给出\[\begin{aligned}P((U,V)\in D)&=\iint_{A^{-1}D}f(x,y)\,dx\,dy\\&=\iint_D\frac{f(a_1u+b_1v,c_1u+d_1v)}{|\det A|}\,du\,dv.\end{aligned}\]这证明了密度公式。当 $A$ 不可逆时，像可能集中在直线上，一般不存在二维勒贝格密度，不能沿用此式。\end{solve}

\begin{exercise}[常数与任意变量独立]\label{q:A4.4}
常数 $a$ 是随机变量, 按定义证明 $a$ 与任何随机变量 $Y$ 独立.\par\smallskip {\small 来源：ayhe123 公开讲义，原题 4.4。}
\end{exercise}

\begin{proof}\phantomsection\label{a:A4.4}
    $\because$ 当 $b\geq a$ 时有 $P(a\leq b)=1,P(a\leq b,Y\leq y)=P(Y\leq y)$, 当 $b<a$ 时有 $P(a\leq b)=P(a\leq b,Y\leq y)=0$, $\therefore\forall b,y$ 都有
    \[P(a\leq b,Y\leq y)=P(Y\leq y)P(a\leq b).\]

    由分布函数的乘积分解判据得 $a$ 与 $Y$ 独立.
\end{proof}

\begin{exercise}[联合分布函数的右连续性]\label{q:A4.5}
证明: 对于固定的 $x$, 联合分布函数 $F(x,y)$ 关于 $y$ 单调不减且右连续.\par\smallskip {\small 来源：ayhe123 公开讲义，原题 4.5。}
\end{exercise}

\begin{solve}\phantomsection\label{a:A4.5}固定 $x$。若 $y_1\le y_2$，事件 $\{X\le x,Y\le y_1\}$ 包含于 $\{X\le x,Y\le y_2\}$，由概率的单调性，$F(x,y_1)\le F(x,y_2)$。对任意序列 $y_n\downarrow y$，定义 $A_n=\{X\le x,Y\le y_n\}$，则 $A_n\downarrow\{X\le x,Y\le y\}$。概率的从上连续性给出\[F(x,y_n)=P(A_n)\longrightarrow P\left(\bigcap_nA_n\right)=F(x,y).\]故关于第二个坐标右连续。这个性质不要求随机向量具有密度。\end{solve}

\begin{exercise}[曲边区域上的联合密度]\label{q:A4.11}
设 $a$ 是常数, $(X,Y)$ 有概率密度
    \[f(x,y)=\begin{cases}
        ax^2y, & x^2<y<1, \\
        0, & \text{其他}.
    \end{cases}\]

    求 $X,Y$ 的边缘密度, 说明 $X,Y$ 不独立.\par\smallskip {\small 来源：ayhe123 公开讲义，原题 4.11。}
\end{exercise}

\begin{solve}\phantomsection\label{a:A4.11}
    由
    \[\begin{aligned}\int_{-\infty}^\infty\int_{-\infty}^\infty f(x,y)\mathrm{d}x\mathrm{d}y&=\iint\limits_{x^2<y<1}ax^2y\mathrm{d}x\mathrm{d}y\\&=1\end{aligned}\]

    得 $a=\dfrac{21}{4}$. 有
    \[\begin{aligned}f_X(x)&=\int_{-\infty}^\infty f(x,y)\mathrm{d}y\\&=\int_{x^2}^1\dfrac{21}{4}x^2y\mathrm{d}y\\&=\dfrac{21}{8}x^2(1-x^2),\quad -1<x<1,\end{aligned}\]
    \[\begin{aligned}f_Y(y)&=\int_{-\infty}^\infty f(x,y)\mathrm{d}x\\&=\int_{-\sqrt{y}}^{\sqrt{y}}\dfrac{21}{4}x^2y\mathrm{d}x\\&=\dfrac{7}{2}y^2\sqrt{y},\quad 0<y<1.\end{aligned}\]

    $\because f_X(x)f_Y(y)\neq f(x,y)$, $\therefore X,Y$ 不独立.
\end{solve}

\begin{exercise}[指数寿命的比较]\label{q:A4.14}
设随机变量 $X,Y$ 相互独立, $X\sim\mathcal{E}(\lambda),Y\sim\mathcal{E}(\mu)$, 计算 $P(X>Y)$.\par\smallskip {\small 来源：ayhe123 公开讲义，原题 4.14。}
\end{exercise}

\begin{solve}\phantomsection\label{a:A4.14}条件于 $Y=y>0$，指数变量的生存概率为 $P(X>y)=e^{-\lambda y}$。独立性保证给定 $Y$ 不改变 $X$ 的分布，故\[\begin{aligned}P(X>Y)&=\int_0^\infty e^{-\lambda y}\mu e^{-\mu y}\,dy\\&=\frac\mu{\lambda+\mu}.\end{aligned}\]可用对称性检查：当两率相同时概率为 $1/2$；当 $\lambda$ 很大、$X$ 往往很短时，该概率趋于零。\end{solve}

\begin{exercise}[不独立变量的平方独立]\label{q:A4.17}
设 $(X,Y)$ 在 $(-1,1)^2$ 上的联合密度为 $(1+xy)/4$，其余处为零。证明 $X,Y$ 不独立，但 $X^2,Y^2$ 独立。\par\smallskip {\small 来源：ayhe123 公开讲义，原题 4.17。}
\end{exercise}

\begin{solve}\phantomsection\label{a:A4.17}对任一变量积分，得到两个边缘密度均为区间 $(-1,1)$ 上的 $1/2$。联合密度不几乎处处等于边缘密度之积 $1/4$，所以不独立。另一方面，当 $0<u,v<1$ 时，\[\begin{aligned}P(X^2\le u,Y^2\le v)&=\int_{-\sqrt u}^{\sqrt u}\int_{-\sqrt v}^{\sqrt v}\frac{1+xy}{4}\,dy\,dx\\&=\sqrt{uv}.\end{aligned}\]含 $xy$ 的积分因对称性为零；而边缘分布函数分别为 $\sqrt u,\sqrt v$。在区间外按分布函数的 $0$、$1$ 延拓仍成立，因此平方变量独立。\end{solve}

\begin{exercise}[曲线下均匀点的投影]\label{q:A4.18}
设 $D$ 是非负连续函数 $g(x)$ 与 $x$ 轴所夹的区域, $D$ 的面积 $m(D)\in(0,\infty)$. 设 $(X,Y)$ 在 $D$ 上均匀分布, 求 $X$ 的概率密度.\par\smallskip {\small 来源：ayhe123 公开讲义，原题 4.18。}
\end{exercise}

\begin{solve}\phantomsection\label{a:A4.18}
    设 $x_0\in\mathbb{R}$. 有
    \begin{align*}
        P(X\leq x_0) & =\iint_D\dfrac{1}{m(D)}I[X\leq x_0]\mathrm{d}x\mathrm{d}y \\
        & =\int_{-\infty}^{x_0}\mathrm{d}x\int_0^{g(x)}\dfrac{1}{m(D)}\mathrm{d}y \\
        & =\int_{-\infty}^{x_0}\dfrac{g(x)}{m(D)}\mathrm{d}x.
    \end{align*}

    $\therefore X$ 有概率密度
    \[f_X(x)=\dfrac{g(x)}{m(D)},\]

    其中 $m(D)=\int_{-\infty}^\infty g(x)\mathrm{d}x$.
\end{solve}

\begin{exercise}[离散与连续变量之差]\label{q:A4.19}
设随机变量 $X$ 服从二项分布 $\mathcal{B}(n,p)$, $Y$ 服从指数分布 $\mathcal{E}(\lambda)$. 当 $X,Y$ 独立时, 求 $Z=Y-X$ 的分布函数和概率密度.\par\smallskip {\small 来源：ayhe123 公开讲义，原题 4.19。}
\end{exercise}

\begin{solve}\phantomsection\label{a:A4.19}记 $b_k=\binom nkp^k(1-p)^{n-k}$。给定 $X=k$ 后，$Z=Y-k$ 是平移后的指数变量，所以\[F_Z(z)=\sum_{k=0}^n b_k(1-e^{-\lambda(z+k)})\mathbf1_{\{z\ge-k\}},\]\[f_Z(z)=\sum_{k=0}^n b_k\lambda e^{-\lambda(z+k)}\mathbf1_{\{z>-k\}}.\]第二式几乎处处成立；各端点处密度值可任意指定。它是有限个平移指数密度的混合，不能在求导后丢掉对 $k$ 的求和。\end{solve}

\begin{exercise}[求和变换后的边缘积分]\label{q:A4.23}
设随机向量 $(X,Y)$ 有联合密度
    \[f(x,y)=\begin{cases}
        \dfrac{1}{2}(x+y)e^{-(x+y)}, & x,y>0 \\
        0 & \text{其他}.
    \end{cases}\]

    求 $Z=X+Y$ 的概率密度.\par\smallskip {\small 来源：ayhe123 公开讲义，原题 4.23。}
\end{exercise}

\begin{solve}\phantomsection\label{a:A4.23}令 $Z=X+Y,W=X$，逆变换为 $x=w,y=z-w$，雅可比绝对值为一。原支撑 $x,y>0$ 化成 $z>0,0<w<z$，因此\[f_{Z,W}(z,w)=\frac12ze^{-z}\mathbf1_{\{z>0,0<w<z\}}.\]求边缘必须对 $w$ 积分：\[f_Z(z)=\int_0^z\frac12ze^{-z}\,dw=\frac12z^2e^{-z},\quad z>0.\]这正是形状为三、率为一的伽马密度。原解漏掉边缘积分中的区间长度 $z$，其结果不能归一化，本题据此校正。\end{solve}

\begin{exercise}[独立正态的线性组合]\label{q:A4.27}
设 $X_1,\ldots,X_n$ 相互独立，$X_j\sim N(\mu_j,\sigma_j^2)$。求 $Y=\sum_{j=1}^n a_jX_j+b$ 的分布。\par\smallskip {\small 来源：ayhe123 公开讲义，原题 4.27。}
\end{exercise}

\begin{solve}\phantomsection\label{a:A4.27}正态分布的特征函数为 $\exp(i\mu t-\sigma^2t^2/2)$。独立性使联合乘积的期望分解，故\[\varphi_Y(t)=\exp\left\{it\left(b+\sum_ja_j\mu_j\right)-\frac{t^2}{2}\sum_ja_j^2\sigma_j^2\right\}.\]因此 $Y$ 服从均值 $b+\sum_ja_j\mu_j$、方差 $\sum_ja_j^2\sigma_j^2$ 的正态分布。若该方差为零，则分布退化在其均值处。仅知道每个边缘分布是正态而不假定独立或联合正态，不足以推出线性组合正态。\end{solve}

\begin{exercise}[二维正态的条件分布]\label{q:A4.28}
设 $(X_1,X_2)$ 联合正态，均值为 $\mu_1,\mu_2$，标准差 $\sigma_1,\sigma_2>0$，相关系数 $|\rho|<1$。记 $z_i=(x_i-\mu_i)/\sigma_i$、$Q=z_1^2-2\rho z_1z_2+z_2^2$，联合密度为\[f(x_1,x_2)=\frac{e^{-Q/[2(1-\rho^2)]}}{2\pi\sigma_1\sigma_2\sqrt{1-\rho^2}}.\]求 $X_2\mid X_1=x_1$ 的条件分布。\par\smallskip {\small 来源：ayhe123 公开讲义，原题 4.28。}
\end{exercise}

\begin{solve}\phantomsection\label{a:A4.28}把标准化二次型配方：\[z_1^2-2\rho z_1z_2+z_2^2=(1-\rho^2)z_1^2+(z_2-\rho z_1)^2.\]将联合密度除以 $X_1$ 的 $N(\mu_1,\sigma_1^2)$ 边缘密度，第一项指数和相应归一化因子抵消，剩下关于 $x_2$ 的正态密度。因此\[X_2\mid X_1=x_1\sim N(m(x_1),(1-\rho^2)\sigma_2^2),\]\[m(x_1)=\mu_2+\rho\frac{\sigma_2}{\sigma_1}(x_1-\mu_1).\]条件期望对 $x_1$ 线性，条件方差不依赖 $x_1$。连续条件分布通过密度比定义，不是以正概率事件 $X_1=x_1$ 作初等条件概率除法。\end{solve}

\begin{exercise}[尾分布幂次与竞争概率]\label{q:A4.32}
随机变量 $X,Y$ 相互独立, $X$ 有分布函数 $F_X(x)$ 和概率密度 $f_X(x)=F_X'(x)$. 如果 $P(Y>y)=(P(X>y))^\beta$, 其中 $\beta$ 是正常数, 求 $P(X\geq Y)$.\par\smallskip {\small 来源：ayhe123 公开讲义，原题 4.32。}
\end{exercise}

\begin{solve}\phantomsection\label{a:A4.32}由题设 $F_Y(y)=1-[1-F_X(y)]^\beta$。独立性允许条件于 $X$，所以\[P(X\ge Y)=E[F_Y(X)]=1-E[1-F_X(X)]^\beta.\]由于 $X$ 的分布连续，概率积分变换给出 $F_X(X)\sim U(0,1)$，即使 $F_X$ 不处处严格增加也成立。因此\[P(X\ge Y)=\int_0^1[1-(1-u)^\beta]du=\frac{\beta}{1+\beta}.\]该方法避免声称“有密度便处处可导”；绝对连续分布函数一般只在几乎处处具有导数。\end{solve}

\clearpage\section{数学期望、方差与协方差}

\begin{exercise}[广告信息之间的等待]\label{q:A5.5}
一部手机收到的短信中有 $p=2\%$ 是广告, 你期望相邻的两次广告短信中有多少个不是广告短信?\par\smallskip {\small 来源：ayhe123 公开讲义，原题 5.5。}
\end{exercise}

\begin{solve}\phantomsection\label{a:A5.5}
    设 $X$ 为相邻的两次广告短信中不是广告短信的数目, 则 $P(X=k)=p^k(1-p)$. 有
    \begin{align*}
        EX & =\sum\limits_{k=0}^\infty kP(X=k) \\
        & =\sum\limits_{k=1}^\infty kp^k(1-p) \\
        & =p(1-p)\sum\limits_{k=1}^\infty kp^{k-1} \\
        & =p(1-p)\left(\sum\limits_{k=1}^\infty\dfrac{\partial}{\partial x}x^k\right)\Bigg|_{x=p} \\
        & =p(1-p)\left(\dfrac{\partial}{\partial x}\sum\limits_{k=0}^\infty x^k\right)\Bigg|_{x=p} \\
        & =p(1-p)\left(\dfrac{\partial}{\partial x}\dfrac{1}{1-x}\right)\Bigg|_{x=p} \\
        & =p(1-p)\cdot\dfrac{1}{(1-p)^2}=\dfrac{p}{1-p}.
    \end{align*}

    把 $p=2\%$ 代入得 $EX=49$.
\end{solve}

\begin{exercise}[标准正态半径的期望]\label{q:A5.8}
设 $X,Y$ 相互独立且均服从标准正态分布。求 $R=\sqrt{X^2+Y^2}$ 的密度及数学期望。\par\smallskip {\small 来源：ayhe123 公开讲义，原题 5.8。}
\end{exercise}

\begin{solve}\phantomsection\label{a:A5.8}二维联合密度为 $(2\pi)^{-1}\exp[-(x^2+y^2)/2]$。用极坐标 $x=r\cos\theta,y=r\sin\theta$，雅可比为 $r$，积分角度得\[f_R(r)=r e^{-r^2/2},\qquad r>0.\]这称为瑞利分布。由于 $R$ 非负，期望积分的范围必须是正半轴：\[ER=\int_0^\infty r^2e^{-r^2/2}\,dr=\sqrt{\frac\pi2}.\]最后一步可分部积分，将它化为 $\int_0^\infty e^{-r^2/2}\,dr$。此外 $ER^2=EX^2+EY^2=2$，所以 $\operatorname{Var}(R)=2-\pi/2$。\end{solve}

\begin{exercise}[交换对称与比例期望]\label{q:A5.9}
设 $X_1,X_2,\cdots,X_n$ 独立同分布, 且 $P(X_i>0)=1$. 对于 $k\leq n$, 计算
    \[E\dfrac{X_1+X_2+\cdots+X_k}{X_1+X_2+\cdots+X_n}.\]\par\smallskip {\small 来源：ayhe123 公开讲义，原题 5.9。}
\end{exercise}

\begin{solve}\phantomsection\label{a:A5.9}
    设
    \[Y_i=\dfrac{X_i}{X_1+X_2+\cdots+X_n},\quad (n=1,2,\cdots,n).\]
    
    $\because X_1,X_2,\cdots,X_n$ 独立同分布, $\therefore Y_1,Y_2,\cdots,Y_n$ 同分布. $\because P(X_i>0)=1,\therefore P(Y_i>0)=1$. 有
    \[1=E\dfrac{X_1+X_2+\cdots+X_n}{X_1+X_2+\cdots+X_n}=EY_1+EY_2+\cdots+EY_n=nEY_1,\]
    \[E\dfrac{X_1+X_2+\cdots+X_k}{X_1+X_2+\cdots+X_n}=EY_1+EY_2+\cdots+EY_k=kEY_1,\]

    $\therefore$
    \[E\dfrac{X_1+X_2+\cdots+X_k}{X_1+X_2+\cdots+X_n}=\dfrac{k}{n}.\]
\end{solve}

\begin{exercise}[离散极值的尾和公式]\label{q:A5.12}
设 $X_1,\ldots,X_n$ 为相互独立、同分布的正整数值变量，$p_k=P(X_1=k)$。记 $u_k=P(X_1<k)$、$v_k=1-u_k$。证明\[E\min_iX_i=\sum_{k=1}^\infty v_k^n,\qquad E\max_iX_i=\sum_{k=1}^\infty(1-u_k^n).\]\par\smallskip {\small 来源：ayhe123 公开讲义，原题 5.12。}
\end{exercise}

\begin{solve}\phantomsection\label{a:A5.12}对任意非负整数值变量 $Z$，恒等式 $Z=\sum_{k\ge1}\mathbf1_{\{Z\ge k\}}$ 与非负积分交换定理给出 $EZ=\sum_{k\ge1}P(Z\ge k)$，允许两边都为无穷。最小值不小于 $k$ 当且仅当所有 $X_i\ge k$，由独立性其概率为 $v_k^n$。最大值小于 $k$ 当且仅当所有 $X_i<k$，概率为 $u_k^n$，所以最大值不小于 $k$ 的概率为 $1-u_k^n$。分别代入尾和公式即得结论。\end{solve}

\begin{exercise}[曲边密度的归一化勘误]\label{q:A5.14}
设 $(X,Y)$ 的联合密度为\[f(x,y)=\frac{3}{2x^3y^2},\qquad x>1,\quad 1<xy<x^2,\]其余处为零。验证归一化并求 $EY$。（此处按原解的积分式校订原题密度的幂次。）\par\smallskip {\small 来源：ayhe123 公开讲义，原题 5.14。}
\end{exercise}

\begin{solve}\phantomsection\label{a:A5.14}固定 $x>1$，积分区域为 $1/x<y<x$，所以\[\int_1^\infty\!\int_{1/x}^{x}\frac{3}{2x^3y^2}\,dy\,dx=\frac32\int_1^\infty(x^{-2}-x^{-4})\,dx=1.\]由非负函数的积分交换定理，\[\begin{aligned}EY&=\frac32\int_1^\infty x^{-3}\!\int_{1/x}^x y^{-1}\,dy\,dx\\&=3\int_1^\infty x^{-3}\log x\,dx=\frac34.\end{aligned}\]末式分部积分即可得到。原资料题面中的 $3/(2x^2y^3)$ 在该区域上不能归一化，故不能直接作为概率密度。\end{solve}

\begin{exercise}[全部指数寿命结束的等待]\label{q:A5.17}
设 $5$ 台计算机独立工作, 每台计算机感染病毒前的时间(分别记作 $X_1,\cdots,X_5$)都服从参数为 $\lambda$ 的指数分布 $\mathcal{E}(\lambda)$. 对 $5$ 台计算机都感染病毒前的时间的数学期望是多少?\par\smallskip {\small 来源：ayhe123 公开讲义，原题 5.17。}
\end{exercise}

\begin{solve}\phantomsection\label{a:A5.17}
    记 $Y=\max(X_1,\cdots,X_5)$, 则
    \begin{align*}
        P(Y>y) & =1-P(Y\leq y) \\
        & =1-P(X_1\leq y,X_2\leq y,\cdots,X_5\leq y) \\
        & =1-P(X_1\leq y)P(X_2\leq y)\cdots P(X_5\leq y) \\
        & =1-(1-e^{-\lambda y})^5 \\
        & =1-\left(1+\sum\limits_{i=1}^5C_5^i(-e^{-\lambda y})^i\right) \\
        & =\sum\limits_{i=1}^5C_5^i(-1)^{i-1}e^{-\lambda iy}.
    \end{align*}

    由相应的分布或期望公式可得
    \begin{align*}
        EY & =\int_0^\infty P(Y>y)\mathrm{d}y \\
        & =\int_0^\infty\sum\limits_{i=1}^5C_5^i(-1)^{i-1}e^{-\lambda iy}\mathrm{d}y \\
        & =\sum\limits_{i=1}^5C_5^i(-1)^{i-1}\int_0^\infty e^{-\lambda iy}\mathrm{d}y \\
        & =\sum\limits_{i=1}^5C_5^i\dfrac{1}{\lambda i}(-1)^{i-1}=\dfrac{137}{60\lambda}.
    \end{align*}
\end{solve}

\begin{exercise}[样本均值与方差的无偏性]\label{q:A5.20}
设 $n\geq 2$, $X_1,X_2,\cdots,X_n$ 独立同分布,
    \[\begin{aligned}\hat{\mu}&=\dfrac{1}{n}\sum\limits_{j=1}^nX_j,\quad\hat{\sigma}^2\\&=\dfrac{1}{n-1}\sum\limits_{j=1}^n(X_j-\hat{\mu})^2,\end{aligned}\]

    验证 $E\hat{\mu}=EX_1,E\hat{\sigma}^2=\operatorname{Var}(X_1)$. 假设 $E(X_1^2)<\infty$。\par\smallskip {\small 来源：ayhe123 公开讲义，原题 5.20。}
\end{exercise}

\begin{solve}\phantomsection\label{a:A5.20}记共同均值为 $\mu$、方差为 $\sigma^2$。期望的线性性给出 $E\widehat\mu=\mu$。平方和恒等式为\[\sum_{j=1}^n(X_j-\widehat\mu)^2=\sum_{j=1}^n(X_j-\mu)^2-n(\widehat\mu-\mu)^2.\]因独立性使 $\operatorname{Var}(\widehat\mu)=\sigma^2/n$，两边取期望得到 $(n-1)\sigma^2$。除以 $n-1$ 即有 $E\widehat\sigma^2=\sigma^2$。有限二阶矩是这里取方差与平方和期望的前提。\end{solve}

\begin{exercise}[最优逆方差加权]\label{q:A5.22}
设独立随机变量 $X_1,\ldots,X_n$ 都有均值 $\mu$，方差分别为 $\sigma_i^2>0$。在 $\sum_iw_i=1$ 的约束下，求线性无偏估计量 $\sum_iw_iX_i$ 的最小方差及最优权重。\par\smallskip {\small 来源：ayhe123 公开讲义，原题 5.22。}
\end{exercise}

\begin{solve}\phantomsection\label{a:A5.22}独立性给出方差 $\sum_iw_i^2\sigma_i^2$。由柯西不等式，\[1=\left(\sum_i(w_i\sigma_i)\sigma_i^{-1}\right)^2\le\left(\sum_iw_i^2\sigma_i^2\right)\left(\sum_i\sigma_i^{-2}\right).\]等号成立当且仅当 $w_i\sigma_i=c/\sigma_i$。结合权重和约束得到\[w_i=\frac{\sigma_i^{-2}}{\sum_j\sigma_j^{-2}},\qquad\min\operatorname{Var}=\frac{1}{\sum_j\sigma_j^{-2}}.\]方差较小的观测获得较大权重；等方差时退化为算术平均。\end{solve}

\begin{exercise}[非负变量的尾积分]\label{q:A5.25}
设非负随机变量 $X$ 有离散分布 $p_j=P(X=x_j),j\geq 1$. 证明:
    
    (1) $P(X>x)=\sum\limits_{j=1}^\infty p_jI[x<x_j]$;

    (2) $EX=\int_0^\infty P(X>x)\mathrm{d}x$.\par\smallskip {\small 来源：ayhe123 公开讲义，原题 5.25。}
\end{exercise}

\begin{proof}\phantomsection\label{a:A5.25}
    (1) 有 $\forall j,P(X>x|X=x_j)=I[x<x_j]$. 由全概率公式得
    \[P(X>x)=\sum\limits_{j=1}^\infty P(X>x|X=x_j)P(X=x_j)=\sum\limits_{j=1}^\infty p_jI[x<x_j].\]

    (2) 由 (1) 得
    \[\int_0^\infty P(X>x)\mathrm{d}x=\int_0^\infty\sum\limits_{j=1}^\infty p_jI[x<x_j]\mathrm{d}x.\]

    $\because\sum\limits_{j=1}^\infty p_jI[x<x_j]\leq\sum\limits_{j=1}^\infty p_j=1$, 由 Weierstrass 判别法得 $\sum\limits_{j=1}^\infty p_jI[x<x_j]$ 在 $(0,\infty)$ 上一致收敛, $\therefore$ 求和号和积分号可以交换次序. 有
    \begin{align*}
        \int_0^\infty\sum\limits_{j=1}^\infty p_jI[x<x_j]\mathrm{d}x & =\sum\limits_{j=1}^\infty p_j\int_0^\infty I[x<x_j]\mathrm{d}x \\
        & =\sum\limits_{j=1}^\infty p_j\int_0^{x_j}\mathrm{d}x \\
        & =\sum\limits_{j=1}^\infty p_jx_j=EX.\qedhere
    \end{align*}
\end{proof}

\begin{exercise}[和与差的相关系数]\label{q:A5.26}
如果 $\operatorname{Cov}(X,Y)=0$, 验证 $\rho(X+Y,X-Y)=\dfrac{\operatorname{Var}(X)-\operatorname{Var}(Y)}{\operatorname{Var}(X)+\operatorname{Var}(Y)}$. 假设二阶矩有限且 $\operatorname{Var}(X)+\operatorname{Var}(Y)>0$。\par\smallskip {\small 来源：ayhe123 公开讲义，原题 5.26。}
\end{exercise}

\begin{solve}\phantomsection\label{a:A5.26}由协方差的双线性与对称性，\[\operatorname{Cov}(X+Y,X-Y)=\operatorname{Var}(X)-\operatorname{Var}(Y).\]又因为 $\operatorname{Cov}(X,Y)=0$，\[\operatorname{Var}(X+Y)=\operatorname{Var}(X-Y)=\operatorname{Var}(X)+\operatorname{Var}(Y)>0.\]将协方差除以两个标准差之积即可得到题设公式。尤其当原变量方差相同时，和与差不相关；只有在额外联合正态等条件下，才能进一步断言它们独立。\end{solve}

\begin{exercise}[正方形密度的协方差]\label{q:A5.28}
设 $(X,Y)$ 有概率密度
    \[f(x,y)=\begin{cases}
        x+y & x\in(0,1),y\in(0,1), \\
        0 & \text{其他}.
    \end{cases}\]

    计算 $\operatorname{Cov}(X,Y)$.\par\smallskip {\small 来源：ayhe123 公开讲义，原题 5.28。}
\end{exercise}

\begin{solve}\phantomsection\label{a:A5.28}
    有
    \[f_X(x)=\int_{-\infty}^\infty f(x,y)\mathrm{d}y=\int_0^1(x+y)\mathrm{d}y=x+\dfrac{1}{2},\quad x\in(0,1),\]
    \[EX=\int_{-\infty}^\infty xf_X(x)\mathrm{d}x=\dfrac{7}{12}.\]

    对称地, 有 $EY=\dfrac{7}{12}$.

    \[\begin{aligned}E(XY)&=\iint_{\mathbb{R}^2}xyf(x,y)\mathrm{d}x\mathrm{d}y\\&=\int_0^1\mathrm{d}x\int_0^1xy(x+y)\mathrm{d}y\\&=\dfrac{1}{3}.\end{aligned}\]

    $\therefore$
    \[\operatorname{Cov}(X,Y)=E(XY)-EXEY=-\dfrac{1}{144}.\]
\end{solve}

\begin{exercise}[正态尺寸差与装配]\label{q:A5.30}
设活塞的直径 $X$ 的平均值为 $20.00\text{cm}$, 标准差为 $0.02$; 气缸的直径 $Y$ 的平均值为 $20.10\text{cm}$, 标准差为 $0.02$. 设 $X$ 与 $Y$ 独立且都服从 正态分布, 计算活塞能装入气缸的概率.\par\smallskip {\small 来源：ayhe123 公开讲义，原题 5.30。}
\end{exercise}

\begin{solve}\phantomsection\label{a:A5.30}活塞能装入等价于 $X<Y$。独立正态变量之差仍为正态，\[Y-X\sim N(0.10,0.02^2+0.02^2).\]故所求概率为\[P(Y-X>0)=\Phi\left(\frac{0.10}{0.02\sqrt2}\right)\approx0.9997965.\]直接对差变量标准化可避免二维积分。方差相加依赖两个尺寸独立；若存在同批次相关性，应在差的方差中减去两倍协方差。\end{solve}

\begin{exercise}[不相关但不独立的反例]\label{q:A5.32}
设随机变量 $X$ 具有偶对称密度，且 $0<EX^2<\infty$。证明 $X$ 与 $|X|$ 不相关但不独立。\par\smallskip {\small 来源：ayhe123 公开讲义，原题 5.32。}
\end{exercise}

\begin{solve}\phantomsection\label{a:A5.32}由柯西不等式，$E|X|<\infty$；由密度对称性，$EX=0$，且奇函数 $x|x|$ 的积分为零。因此\[\operatorname{Cov}(X,|X|)=E[X|X|]-EX\,E|X|=0.\]若 $X$ 与 $|X|$ 独立，则对任意 $a$，事件 $B_a=\{|X|\le a\}$ 既由 $X$ 决定又由 $|X|$ 决定，故 $P(B_a)=P(B_a)^2$。于是每个这样的概率只能为 $0$ 或 $1$，迫使 $|X|$ 几乎处处为常数。这与 $X$ 具有密度相矛盾。因此二者不独立。\end{solve}

\begin{exercise}[中心正态的高阶矩]\label{q:A5.33}
设 $X\sim N(0,\sigma^2)$, $n\in\mathbb{N}_+$, 验证
    \[EX^n=\begin{cases}
        \sigma^n(n-1)!!, & n=2m, \\
        0 & \text{其他}.
    \end{cases}\]\par\smallskip {\small 来源：ayhe123 公开讲义，原题 5.33。}
\end{exercise}

\begin{solve}\phantomsection\label{a:A5.33}正态分布的所有绝对矩有限，故可用对称性处理奇次幂，得到 $EX^{2m+1}=0$。设正态密度为 $f$，则 $f^{\prime}(x)=-xf(x)/\sigma^2$。对偶次矩分部积分，边界项因高斯衰减为零：\[\begin{aligned}EX^{2m}&=-\sigma^2\int_{-\infty}^{\infty}x^{2m-1}f^{\prime}(x)dx\\&=(2m-1)\sigma^2EX^{2m-2}.\end{aligned}\]从 $EX^0=1$ 递推得 $EX^{2m}=\sigma^{2m}(2m-1)!!$，其中双阶乘为所有正奇数 $1\cdot3\cdots(2m-1)$ 的乘积。\end{solve}

\begin{exercise}[圆周与圆内均匀模型]\label{q:A5.35}
在半径为 $R>0$ 的圆周上按弧长均匀取一点 $(X,Y)$。求 $EX,EY,\operatorname{Var}(X),\operatorname{Var}(Y)$ 及 $\operatorname{Cov}(X,Y)$，并判断独立性。\par\smallskip {\small 来源：ayhe123 公开讲义，原题 5.35。}
\end{exercise}

\begin{solve}\phantomsection\label{a:A5.35}取 $\Theta\sim U(0,2\pi)$，令 $X=R\cos\Theta,Y=R\sin\Theta$。三角函数在整周期上的积分给出\[EX=EY=0,\quad EX^2=EY^2=\frac{R^2}{2},\quad EXY=0.\]故两个方差均为 $R^2/2$，协方差为零。但它们不独立：事件 $|X|>\sqrt3R/2$ 与 $|Y|>R/2$ 都有正概率，而其交集由 $X^2+Y^2=R^2$ 可知为空。圆内按面积均匀取点是另一种模型，此时两个方差均为 $R^2/4$。原资料解答混用了这两种均匀分布，本题已校订。\end{solve}

\clearpage\section{条件期望与生成函数}

\begin{exercise}[随机组数下的条件二项]\label{q:A6.1}
设 $N_1,N_2,\cdots$ 独立同分布, $N_1\sim\mathcal{B}(8,0.6)$, $M\sim\mathcal{P}(10)$, $M$ 与 $\{N_j\}$ 独立. 对于 $T=\sum\limits_{j=1}^MN_j,k\geq1$, 给出 $T|\{M=k\}$ 的概率分布.\par\smallskip {\small 来源：ayhe123 公开讲义，原题 6.1。}
\end{exercise}

\begin{solve}\phantomsection\label{a:A6.1}
    $N_j$ 表示成功概率为 $0.6$ 的 $8$ 次独立重复试验中成功的次数. $\because N_1,N_2,\cdots$ 独立, $\therefore T|\{M=k\}=N_1+\cdots+N_k$ 表示成功概率为 $0.6$ 的 $8k$ 次独立重复试验中成功的次数, 服从二项分布 $\mathcal{B}(8k,0.6)$.
\end{solve}

\begin{exercise}[几何个指数等待之和]\label{q:A6.3}
设 $N$ 服从参数为 $p$ 的几何分布, $X_1,X_2,\cdots$ 相互独立, 都服从参数为 $\lambda$ 的指数分布, 当 $N$ 和 $\{X_j\}$ 独立时, 计算 $W=\sum\limits_{j=1}^NX_j$ 的数学期望. 这里几何分布取值 $1,2,\ldots$，质量函数为 $p(1-p)^{k-1}$。\par\smallskip {\small 来源：ayhe123 公开讲义，原题 6.3。}
\end{exercise}

\begin{solve}\phantomsection\label{a:A6.3}
    记 $q=1-p$。由独立指数变量之和的伽马分布公式, $W|\{N=n\}$ 有概率密度
    \[f_W(w;n)=\dfrac{\lambda^n}{(n-1)!}e^{-\lambda w}w^{n-1},\quad w>0.\]

    $\therefore$
    \[E(W|N=n)=\int_0^\infty wf_W(w;n)\mathrm{d}w=\int_0^\infty\dfrac{\lambda^n}{(n-1)!}e^{-\lambda w}w^n\mathrm{d}w.\]

    令 $t=\lambda w$, 得
    \begin{align*}
        \int_0^\infty\dfrac{\lambda^n}{(n-1)!}e^{-\lambda w}w^n\mathrm{d}w & =\dfrac{1}{\lambda(n-1)!}\int_0^\infty e^{-t}t^n\mathrm{d}t \\
        & =\dfrac{1}{\lambda(n-1)!}\Gamma(n+1)=\dfrac{n}{\lambda}.
    \end{align*}

    由相应的分布或期望公式可得
    \begin{align*}
        EW & =\sum\limits_{n=1}^\infty E(W|N=n)P(N=n) \\
        & =\sum\limits_{n=1}^\infty\dfrac{n}{\lambda}q^{n-1}p \\
        & =\dfrac{p}{\lambda}\left(\dfrac{\mathrm{d}}{\mathrm{d}x}\sum\limits_{n=0}^\infty x^n\right)\Bigg|_{x=q} \\
        & =\dfrac{p}{\lambda}\left(\dfrac{1}{(1-x)^2}\right)\Bigg|_{x=q}=\dfrac{1}{\lambda p}.
    \end{align*}
\end{solve}

\begin{exercise}[复合泊松等待的均值方差]\label{q:A6.4}
某出租车在一天内遇到的红灯数 $N$ 满足参数为 $\lambda$ 的 泊松 分布. 如果在每个红灯处的等待时间相互独立, 都在 $(0,2)$ 中均匀分布, 计算他一天内用于等候红灯的时间 $X$ 的数学期望和方差. 假设红灯数与各次等待时间独立。\par\smallskip {\small 来源：ayhe123 公开讲义，原题 6.4。}
\end{exercise}

\begin{solve}\phantomsection\label{a:A6.4}
    设 $X_1,X_2,\cdots$ 是在第 $i$ 个红灯处的等待时间, 则 $X|\{N=n\}=X_1+X_2+\cdots+X_n$, $X_i$ 独立同分布, $X_1\sim\mathcal{U}(0,2)$, $N$ 与 $\{X_i\}$ 独立. 有
    \[E(X|N=n)=EX_1+EX_2+\cdots+EX_n=nEX_1=n,\]
    \begin{align*}
        E(X^2|N=n) & =E(X_1+X_2+\cdots+X_n)^2 \\
        & =E\left(X_1^2+X_2^2+\cdots+X_n^2+\sum\limits_{i=1}^n\sum\limits_{\substack{j=1\\j\neq i}}^nX_iX_j\right) \\
        & =nEX_1^2+\sum\limits_{i=1}^n\sum\limits_{\substack{j=1\\j\neq i}}^nEX_iEX_j \\
        & =\dfrac{4n}{3}+n(n-1).
    \end{align*}

    由相应的分布或期望公式可得
    \[EX=\sum\limits_{n=0}^\infty P(N=n)E(X|N=n)=\lambda,\]
    \begin{align*}
        EX^2 & =\sum\limits_{n=0}^\infty P(N=n)E(X^2|N=n) \\
        & =\sum\limits_{n=0}^\infty\left(\dfrac{4n}{3}+n(n-1)\right)\dfrac{\lambda^n}{n!}e^{-\lambda}  \\
        & =\dfrac{4}{3}\sum\limits_{n=1}^\infty\dfrac{\lambda^n}{(n-1)!}e^{-\lambda}+\sum\limits_{n=2}^\infty\dfrac{\lambda^n}{(n-2)!}e^{-\lambda}  \\
        & =\dfrac{4}{3}\lambda+\lambda^2.
    \end{align*}

    $\therefore$
    \[\operatorname{Var}(X)=EX^2-(EX)^2=\dfrac{4}{3}\lambda.\]
\end{solve}

\begin{exercise}[条件指数与条件伽马]\label{q:A6.9}
设 $(X,Y)$ 有联合密度 $f(x,y)=ye^{-y(x+1)},x>0,y>0$. 计算 $X|\{Y=y\},Y|\{X=x\}$ 的概率密度 $f_{X|Y}(x|y),f_{Y|X}(y|x)$.\par\smallskip {\small 来源：ayhe123 公开讲义，原题 6.9。}
\end{exercise}

\begin{solve}\phantomsection\label{a:A6.9}
    $f_{X|Y}(x|y),f_{Y|X}(y|x)$ 可以由 $f(x,y)$ 除以一个归一化因子得到. 有
    \[f_{X|Y}(x|y)=\dfrac{f(x,y)}{\int_0^\infty f(x,y)\mathrm{d}x}=\dfrac{ye^{-y(x+1)}}{e^{-y}}=ye^{-yx},\]
    \[f_{Y|X}(y|x)=\dfrac{f(x,y)}{\int_0^\infty f(x,y)\mathrm{d}y}=(x+1)^2ye^{-y(x+1)}.\]
\end{solve}

\begin{exercise}[威布尔寿命的剩余等待]\label{q:A6.14}
假设某设备的使用寿命 $Y$ 服从 威布尔（Weibull）分布, 有概率密度
    \[f(y)=aby^{b-1}e^{-ay^b},\quad y\geq0.\]

    对 $b=1/2$, 计算平均寿命 $E(Y-t|Y>t)$. 其中 $a>0$，所求量是已经存活到 $t$ 后的平均剩余寿命。\par\smallskip {\small 来源：ayhe123 公开讲义，原题 6.14。}
\end{exercise}

\begin{solve}\phantomsection\label{a:A6.14}
    由相应的分布或期望公式可得
    \begin{align*}
        E(Y|Y>t) & =\dfrac{E(YI[Y>t])}{P(Y>t)} \\
        & =\dfrac{\int_t^\infty yf(y)\mathrm{d}y}{\int_t^\infty f(y)\mathrm{d}y} \\
        & =\dfrac{ab\int_t^\infty y^{1/2}e^{-ay^{1/2}}\mathrm{d}y}{ab\int_t^\infty y^{-1/2}e^{-ay^{1/2}}\mathrm{d}y} \\
        & =\dfrac{\int_t^\infty y^{1/2}e^{-ay^{1/2}}\mathrm{d}y}{\int_t^\infty y^{-1/2}e^{-ay^{1/2}}\mathrm{d}y}.
    \end{align*}

    令 $u=ay^{1/2}$, 则 $y=(u/a)^2,\mathrm{d}y=\dfrac{2}{a^2}u\mathrm{d}u$. 有
    \begin{align*}
        \dfrac{\int_t^\infty y^{1/2}e^{-ay^{1/2}}\mathrm{d}y}{\int_t^\infty y^{-1/2}e^{-ay^{1/2}}\mathrm{d}y} & =\dfrac{\int_{a\sqrt{t}}^\infty(u/a)e^{-u}(2/a^2)u\mathrm{d}u}{\int_{a\sqrt{t}}^\infty(u/a)^{-1}e^{-u}(2/a^2)u\mathrm{d}u} \\
        & =\dfrac{\int_{a\sqrt{t}}^\infty u^2e^{-u}\mathrm{d}u}{a^2\int_{a\sqrt{t}}^\infty e^{-u}\mathrm{d}u} \\
        & =\dfrac{e^{-a\sqrt{t}}(2+2a\sqrt{t}+a^2t)}{a^2e^{-a\sqrt{t}}} \\
        & =\dfrac{2+2a\sqrt{t}+a^2t}{a^2}.
    \end{align*}

    $\therefore$
    \[E(Y-t|Y>t)=E(Y|Y>t)-E(t|Y>t)=\dfrac{2+2a\sqrt{t}+a^2t}{a^2}-t.\]
\end{solve}

\begin{exercise}[泊松观察后的伽马后验]\label{q:A6.15}
设 $\alpha,\beta>0$。随机强度 $Y$ 服从 $\operatorname{Gamma}(\alpha,\beta)$ 分布，其中 $\beta$ 为率参数。给定 $Y=y$ 后，计数 $N$ 的条件分布为 $\operatorname{Pois}(y)$。观察到 $N=n$ 后，求 $Y$ 的后验分布和后验均值。\par\smallskip {\small 来源：ayhe123 公开讲义，原题 6.15。}
\end{exercise}

\begin{solve}\phantomsection\label{a:A6.15}先验密度与似然分别为\[\pi(y)=\frac{\beta^\alpha}{\Gamma(\alpha)}y^{\alpha-1}e^{-\beta y},\quad P(N=n\mid y)=\frac{y^ne^{-y}}{n!}.\]由贝叶斯公式，其乘积中依赖 $y$ 的部分为 $y^{\alpha+n-1}e^{-(\beta+1)y}$。将其归一化得到\[Y\mid N=n\sim\operatorname{Gamma}(\alpha+n,\beta+1).\]所以 $E(Y\mid N=n)=(\alpha+n)/(\beta+1)$。须始终区分率参数与尺度参数；这里的伽马均值是形状除以率。\end{solve}

\begin{exercise}[七个接口的第三次失效]\label{q:A6.17}
一台计算机有 $7$ 个 USB 接口, 至少有 $5$ 个接口正常时, 就认为计算机的 USB 接口能用. 假设开始时每个 USB 接口都正常, 其寿命 $X_1,X_2,\cdots,X_7$ 相互独立, 都服从参数为 $\lambda$ 的指数分布. 计算这台计算机的 USB 能用的时间 $Y$ 的概率分布.\par\smallskip {\small 来源：ayhe123 公开讲义，原题 6.17。}
\end{exercise}

\begin{solve}\phantomsection\label{a:A6.17}系统在第三个接口失效时停止工作，所以 $Y=X_{(3)}$ 是七个寿命的第三小值。对 $y\ge0$，记 $q=1-e^{-\lambda y}$，则在时刻 $y$ 已失效数量服从 $\operatorname{Bin}(7,q)$。故\[P(Y>y)=\sum_{k=0}^{2}\binom7k(1-e^{-\lambda y})^k e^{-(7-k)\lambda y}.\]利用顺序统计量密度或对其求导，得到\[f_Y(y)=105\lambda(1-e^{-\lambda y})^2e^{-5\lambda y},\quad y>0.\]指数中的系数是 $5$；原解写为 $2$ 会导致积分不等于一。此处还可由独立指数间隔求 $EY=(1/7+1/6+1/5)/\lambda$。\end{solve}

\begin{exercise}[指数顺序统计量的间隔]\label{q:A6.20}
设 $X_1,X_2,\cdots,X_n$ 独立同分布, 都服从参数为 $\lambda$ 的指数分布, 证明随机变量
    \[Y_1=X_{(1)},\quad Y_k=X_{(k)}-X_{(k-1)},\quad k=2,3,\cdots,n\]
    相互独立, 且 $Y_i\sim\mathcal{E}((n+1-i)\lambda)$.\par\smallskip {\small 来源：ayhe123 公开讲义，原题 6.20。}
\end{exercise}

\begin{solve}\phantomsection\label{a:A6.20}原样本连续，所以并列值的概率为零。顺序统计量在 $0<x_{(1)}<\cdots<x_{(n)}$ 上的联合密度为\[n!\lambda^n\exp\left(-\lambda\sum_{i=1}^nx_{(i)}\right).\]令 $x_{(i)}=\sum_{j=1}^iy_j$，变换矩阵为下三角且对角元均为一，故雅可比为一；新支撑为所有 $y_j>0$。又\[\sum_{i=1}^nx_{(i)}=\sum_{j=1}^n(n-j+1)y_j.\]因此联合密度为\[\prod_{j=1}^n\left[(n-j+1)\lambda e^{-(n-j+1)\lambda y_j}\right].\]每个因子是一个完整指数密度，支撑为直积，故各间隔独立，率分别为 $n\lambda,(n-1)\lambda,\ldots,\lambda$。\end{solve}

\begin{exercise}[仿射变换的概率母函数]\label{q:A6.21}
设取非负整数值的随机变量 $X$ 有母函数 $g(s)$, 对非负整数 $a,b$, 求 $Y=aX+b$ 的母函数 $h(s)$.\par\smallskip {\small 来源：ayhe123 公开讲义，原题 6.21。}
\end{exercise}

\begin{solve}\phantomsection\label{a:A6.21}
    $Y$ 的取值为 $\{ja+b|j\in\mathbb{R}\}$. 有
    \begin{align*}
        g(s) & =\sum\limits_{j=0}^\infty s^jP(X=j) \\
        & =\sum\limits_{j=0}^\infty(s^{1/a})^{ja}P(Y=ja+b) \\
        & =(s^{1/a})^{-b}\sum\limits_{j=0}^\infty(s^{1/a})^{ja+b}P(Y=ja+b) \\
        & =s^{-b/a}h(s^{1/a}).
    \end{align*}

    $\therefore$
    \[h(s^{1/a})=g(s)s^{b/a}\Rightarrow h(s)=g(s^a)s^b.\]
\end{solve}

\begin{exercise}[两组硬币成功数之差]\label{q:A6.24}
两人各抛均匀的硬币 $n$ 次, 计算甲的正面次数恰好大于乙的正面次数 $k$ 次的概率.\par\smallskip {\small 来源：ayhe123 公开讲义，原题 6.24。}
\end{exercise}

\begin{solve}\phantomsection\label{a:A6.24}设甲乙的正面次数为 $H_1,H_2$，彼此独立且均为 $\operatorname{Bin}(n,1/2)$。由于 $n-H_2$ 仍服从相同二项分布且与 $H_1$ 独立，\[H_1-H_2+n=H_1+(n-H_2)\sim\operatorname{Bin}(2n,1/2).\]因此对整数 $-n\le k\le n$，\[P(H_1-H_2=k)=\frac{\binom{2n}{n+k}}{2^{2n}},\]范围外概率为零。将负幂生成函数转化为这个非负二项和，可以避免移位时错写幂次或求和下标。\end{solve}

\begin{exercise}[均匀和的特征函数极限]\label{q:A6.29}
设 $X_1,X_2,\cdots$ 独立同分布, 都在 $(-1,1)$ 中均匀分布. 定义
    \[Y_n=\sqrt{\dfrac{3}{n}}\sum\limits_{j=1}^nX_j.\]

    (1) 计算 $X_1$ 的特征函数 $\phi_X(t)$;

    (2) 计算 $Y_n$ 的特征函数 $\phi_{Y_n}(t)$, 证明 $Y_n\xrightarrow{d}N(0,1)$.\par\smallskip {\small 来源：ayhe123 公开讲义，原题 6.29。}
\end{exercise}

\begin{solve}\phantomsection\label{a:A6.29}直接积分可得 $\varphi_X(t)=\sin t/t$，$t\ne0$，并在 $t=0$ 连续定义为 $1$。独立性及缩放性质给出\[\varphi_{Y_n}(t)=\left[\frac{\sin(t\sqrt{3/n})}{t\sqrt{3/n}}\right]^n.\]对每个固定 $t$，由 $\sin u/u=1-u^2/6+O(u^4)$，括号内为 $1-t^2/(2n)+O(n^{-2})$。所以\[\varphi_{Y_n}(t)\to e^{-t^2/2}.\]极限是标准正态特征函数且在零点连续，由莱维连续性定理，$Y_n\xrightarrow dN(0,1)$。余项展开是对 $n\to\infty$ 而 $t$ 固定进行的。\end{solve}

\begin{exercise}[指数和的特征函数极限]\label{q:A6.30}
设 $X_1,X_2,\cdots$ 独立同分布, $X_1\sim\mathcal{E}(\lambda)$. 定义
    \[S_n=\sum\limits_{j=1}^nX_j,\quad Y_n=\dfrac{S_n-n/\lambda}{\sqrt{n/\lambda^2}}.\]

    (1) 计算 $X_1$ 的特征函数 $\phi_X(t)$;

    (2) 计算 $Y_n$ 的特征函数 $\phi_{Y_n}(t)$;

    (3) 证明 $Y_n\xrightarrow{d}N(0,1)$.\par\smallskip {\small 来源：ayhe123 公开讲义，原题 6.30。}
\end{exercise}

\begin{solve}\phantomsection\label{a:A6.30}指数变量的特征函数为\[\varphi_X(t)=\int_0^\infty\lambda e^{-(\lambda-it)x}dx=\frac{\lambda}{\lambda-it}.\]独立性及中心化缩放给出\[\varphi_{Y_n}(t)=\left[\frac{e^{-it/\sqrt n}}{1-it/\sqrt n}\right]^n.\]对固定 $t$，在原点附近使用复对数的连续分支，\[\log\varphi_{Y_n}(t)=n[-it/\sqrt n-\log(1-it/\sqrt n)]\to-\frac{t^2}2.\]线性项相消，二次项给出上述极限，余项为 $O(n^{-1/2})$。莱维连续性定理得到 $Y_n\xrightarrow dN(0,1)$。\end{solve}

\begin{exercise}[二项分布的泊松极限]\label{q:A6.34}
设 $X_n\sim\operatorname{Bin}(n,p_n)$，且 $np_n\to\lambda\in[0,\infty)$。证明 $X_n\xrightarrow{d}\operatorname{Pois}(\lambda)$。\par\smallskip {\small 来源：ayhe123 公开讲义，原题 6.34。}
\end{exercise}

\begin{solve}\phantomsection\label{a:A6.34}对固定实数 $t$，特征函数为\[\varphi_n(t)=[1+p_n(e^{it}-1)]^n.\]由 $np_n\to\lambda$ 可知 $p_n\to0$、$np_n^2\to0$。在 $1$ 的邻域展开对数，得到\[\log\varphi_n(t)=np_n(e^{it}-1)+O(np_n^2)\to\lambda(e^{it}-1).\]因此特征函数逐点趋于 $\exp\{\lambda(e^{it}-1)\}$，后者是泊松分布的特征函数且在原点连续。由莱维连续性定理得到结论。条件 $np_n\to\lambda$ 不等同于每个 $n$ 都满足 $p_n=\lambda/n$。\end{solve}

\begin{exercise}[给定总和后的部分和]\label{q:A6.35}
设 $X_1,\ldots,X_n$ 独立同分布且 $E|X_1|<\infty$，$S_k=\sum_{j=1}^kX_j$。证明 $E(S_k\mid S_n)=kS_n/n$，$1\le k\le n$。\par\smallskip {\small 来源：ayhe123 公开讲义，原题 6.35。}
\end{exercise}

\begin{solve}\phantomsection\label{a:A6.35}联合分布在交换坐标后不变，而 $S_n$ 也不变，因此各个 $E(X_j\mid S_n)$ 几乎处处相等。条件期望的线性性给出\[\sum_{j=1}^nE(X_j\mid S_n)=E(S_n\mid S_n)=S_n.\]于是每一项等于 $S_n/n$，再对前 $k$ 项求和即得结论。可积性保证这些条件期望存在；一般不需要条件密度存在。\end{solve}

\clearpage\section{大数定律与中心极限定理}

\begin{exercise}[截尾期望的概率上界]\label{q:A7.2}
证明: 设非负随机变量 $X$ 有概率密度 $f(x)$, 则 $\forall$ 正数 $M$, 有
    \[P(X\geq M)\leq\dfrac{1}{M}\int_M^\infty xf(x)\mathrm{d}x.\]\par\smallskip {\small 来源：ayhe123 公开讲义，原题 7.2。}
\end{exercise}

\begin{proof}\phantomsection\label{a:A7.2}
    $\because X$ 有概率密度 $f(x)$, $\therefore$
    \begin{align*}
        MP(X\geq M) & =M\int_M^\infty f(x)\mathrm{d}x \\
        & =\int_M^\infty Mf(x)\mathrm{d}x \\
        & \leq\int_M^\infty xf(x)\mathrm{d}x.\qedhere
    \end{align*}
\end{proof}

\begin{exercise}[柯西不等式与尾概率]\label{q:A7.3}
证明: 设非负随机变量 $X$ 有分布函数 $F(x)$, 则 $\forall$ 正数 $M$, 有
    \[P(X>M)\leq\dfrac{\sqrt{(1-F(M))EX^2}}{M}.\] 假设有关二阶矩有限，且 $M>0$。\par\smallskip {\small 来源：ayhe123 公开讲义，原题 7.3。}
\end{exercise}

\begin{solve}\phantomsection\label{a:A7.3}记 $A=\{X>M\}$，则 $P(A)=1-F(M)$。由非负性与柯西不等式，\[MP(A)\le E[X\mathbf1_A]\le\sqrt{EX^2\,P(A)}.\]除以 $M>0$ 即得\[P(X>M)\le\frac{\sqrt{(1-F(M))EX^2}}M.\]若 $P(A)=0$，结论仍成立。这里不需要密度存在，也不需要把弱不等式误写成对所有情形都严格的不等式。\end{solve}

\begin{exercise}[最小值的概率界]\label{q:A7.4}
设 $X_1,\ldots,X_n$ 独立同分布且非负，$EX_1^2<\infty$。对 $M>0$，证明\[P(\min_iX_i\ge M)\le\left(\frac{EX_1^2\,P(X_1\ge M)}{M^2}\right)^{n/2}.\]\par\smallskip {\small 来源：ayhe123 公开讲义，原题 7.4。}
\end{exercise}

\begin{solve}\phantomsection\label{a:A7.4}令 $A=\{X_1\ge M\}$。因 $X_1\mathbf1_A\ge M\mathbf1_A$，有 $MP(A)\le E(X_1\mathbf1_A)$。再用柯西不等式得到\[MP(A)\le\sqrt{EX_1^2\,P(A)}.\]将这一不等式两侧除以 $M$ 并取 $n$ 次幂，同时使用独立性 $P(\min_iX_i\ge M)=P(A)^n$，即得结论。当 $P(A)=0$ 时两侧均为零，无须进行除以 $P(A)$ 的操作。\end{solve}

\begin{exercise}[伽马矩与切比雪夫界]\label{q:A7.5}
设 $X$ 有概率密度 $f(x)=\dfrac{x^m}{m!}e^{-x},\ x\geq0$, 计算 $EX,\operatorname{Var}(X)$, 证明
    \[P(0<X<2(m+1))\geq\dfrac{m}{m+1}.\]\par\smallskip {\small 来源：ayhe123 公开讲义，原题 7.5。}
\end{exercise}

\begin{solve}\phantomsection\label{a:A7.5}该密度为形状 $m+1$、率一的伽马分布，其中 $m$ 为非负整数。由伽马积分，\[EX=m+1,\quad EX^2=(m+1)(m+2),\]\[\operatorname{Var}(X)=m+1.\]切比雪夫不等式给出\[P(|X-(m+1)|\ge m+1)\le\frac1{m+1}.\]取补事件，得到 $P(0<X<2(m+1))\ge m/(m+1)$。这是下界，最后一步不能写成概率恰好等于 $m/(m+1)$。\end{solve}

\begin{exercise}[三角函数样本均值之比]\label{q:A7.10}
设 $X_1,X_2,\cdots$ 独立同分布, 都在 $(0,\pi/2)$ 中均匀分布,
    \[Y_n=\dfrac{\sin X_1+\sin X_2+\cdots+\sin X_n}{\cos X_1+\cos X_2+\cdots+\cos X_n},\quad n\geq1,\]
    
    当 $n\to\infty$ 时, 计算 $Y_n$ 的 a.s. 极限.\par\smallskip {\small 来源：ayhe123 公开讲义，原题 7.10。}
\end{exercise}

\begin{solve}\phantomsection\label{a:A7.10}由均匀密度积分，\[E\sin X_1=\frac2\pi\int_0^{\pi/2}\sin x\,dx=\frac2\pi,\]\[E\cos X_1=\frac2\pi\int_0^{\pi/2}\cos x\,dx=\frac2\pi.\]两个函数均有界，所以强大数定律分别适用；在两个概率一事件的交集上，两组样本平均都收敛到 $2/\pi$。分母极限非零，除法在该点连续，故比值 $Y_n$ 几乎必然收敛到 $1$。\end{solve}

\begin{exercise}[比例调查的样本量]\label{q:A7.14}
利用独立伯努利抽样估计未知比例 $p$。希望样本比例与 $p$ 的误差不超过 $0.01$，置信水平约为 $99\%$。用中心极限定理求不依赖 $p$ 的保守样本量，并说明该结果是否是有限样本的严格保证。\par\smallskip {\small 来源：ayhe123 公开讲义，原题 7.14。}
\end{exercise}

\begin{solve}\phantomsection\label{a:A7.14}样本比例 $\widehat p$ 的方差为 $p(1-p)/n\le1/(4n)$。正态近似要求\[z_{0.995}\frac{1}{2\sqrt n}\le0.01.\]采用 $z_{0.995}=2.575829\ldots$，得到 $n\ge16588$；若将分位数上调为 $2.58$，得到 $n\ge16641$。这是正态近似下的规划值，不是所有 $p$ 的有限样本精确覆盖保证。若仅用切比雪夫不等式，\[P(|\widehat p-p|\ge0.01)\le\frac{1}{4n(0.01)^2},\]取 $n\ge250000$ 可给出至少 $99\%$ 的严格但较粗的保证。\end{solve}

\begin{exercise}[两组样本均值差的近似]\label{q:A7.15}
两组样本相互独立；第一组为来自均值 $\mu_1$、有限正方差 $\sigma_1^2$ 的独立同分布样本，容量 $n$；第二组相应为 $\mu_2,\sigma_2^2,m$。当 $n,m$ 都较大时，求两个样本均值之差的近似分布。\par\smallskip {\small 来源：ayhe123 公开讲义，原题 7.15。}
\end{exercise}

\begin{solve}\phantomsection\label{a:A7.15}由中心极限定理，各组标准化样本均值渐近正态；两组独立保证两项独立。于是\[\overline X-\overline Y\ \dot\sim\ N\left(\mu_1-\mu_2,\frac{\sigma_1^2}{n}+\frac{\sigma_2^2}{m}\right).\]更严格地，把差除以其真实标准差；独立性和两个分量的正态极限使该标准化差趋于 $N(0,1)$。若两总体本来就是正态分布，上式为任意 $n,m$ 下的精确分布。样本容量不同，第二项分母应为 $m$。\end{solve}

\begin{exercise}[泊松概率和的二分极限]\label{q:A7.20}
证明\[\begin{aligned}\lim_{n\to\infty}e^{-n}\sum_{k=0}^{n}\frac{n^k}{k!}&=\frac12,\qquad\lim_{n\to\infty}e^{-n}\sum_{k=0}^{n-1}\frac{n^k}{k!}\\&=\frac12.\end{aligned}\]\par\smallskip {\small 来源：ayhe123 公开讲义，原题 7.20。}
\end{exercise}

\begin{solve}\phantomsection\label{a:A7.20}令 $Z_1,Z_2,\ldots$ 独立同分布，且 $Z_j\sim\operatorname{Pois}(1)$，则 $S_n=\sum_{j=1}^nZ_j\sim\operatorname{Pois}(n)$。由于 $EZ_j=\operatorname{Var}(Z_j)=1$，中心极限定理给出\[\frac{S_n-n}{\sqrt n}\xrightarrow{d}N(0,1).\]第一个和等于 $P((S_n-n)/\sqrt n\le0)$，故极限为 $\Phi(0)=1/2$。第二个和的标准化阈值为 $-1/\sqrt n\to0$。对任意 $\varepsilon>0$，当 $n$ 足够大时可将该概率夹在阈值 $-\varepsilon$ 与 $0$ 的概率之间，再令 $\varepsilon\downarrow0$，仍得到 $1/2$。\end{solve}

\begin{exercise}[缩放收敛与紧性]\label{q:A7.21}
如果 $\sqrt{n}\xi_n\xrightarrow{d}N(0,1)$, 证明 $\xi_n\xrightarrow{p}0$.\par\smallskip {\small 来源：ayhe123 公开讲义，原题 7.21。}
\end{exercise}

\begin{solve}\phantomsection\label{a:A7.21}令 $Z_n=\sqrt n\xi_n$。依分布收敛于一个有限实值分布意味着分布族渐近紧：对任意 $\delta>0$，可选 $M$ 使 $\limsup_nP(|Z_n|>M)<\delta$，因为标准正态两侧尾概率可任意小，且 $\pm M$ 为连续点。固定 $\varepsilon>0$，当 $n$ 足够大时 $\varepsilon\sqrt n>M$，故\[P(|\xi_n|>\varepsilon)=P(|Z_n|>\varepsilon\sqrt n)\le P(|Z_n|>M).\]取上极限后再令 $\delta\downarrow0$，得到依概率收敛到零。不能直接把随 $n$ 改变的阈值代入只对固定点成立的分布函数收敛式。\end{solve}

\begin{exercise}[不同参数二项和的极限]\label{q:A7.22}
设 $X_j$ 相互独立且 $X_j\sim\operatorname{Bin}(j,p_j)$，其中 $0.1\le p_j\le0.9$。令 $S_n=\sum_{j=1}^nX_j$。证明 $(S_n-ES_n)/\sqrt{\operatorname{Var}(S_n)}$ 依分布收敛于标准正态分布。\par\smallskip {\small 来源：ayhe123 公开讲义，原题 7.22。}
\end{exercise}

\begin{solve}\phantomsection\label{a:A7.22}可在等分布的模型中将每个 $X_j$ 写成 $j$ 个参数为 $p_j$ 的独立伯努利变量之和，并令不同组也相互独立。记 $B_n^2=\operatorname{Var}(S_n)$，则\[B_n^2=\sum_{j=1}^n jp_j(1-p_j)\ge0.09\frac{n(n+1)}2\longrightarrow\infty.\]每个中心化伯努利变量的绝对值不超过 $1$。对任意 $\varepsilon>0$，最终有 $\varepsilon B_n>1$，所以林德伯格条件中的所有截尾二阶矩均为零。林德伯格中心极限定理遂给出所需结论。使用这一形式的定理时，须把二项变量拆成有界的独立小项，不能只因原变量独立便直接断言结论。\end{solve}

\begin{exercise}[小扰动不改变分布极限]\label{q:A7.25}
如果 $\xi_n\xrightarrow{d}N(0,1),\eta_n\xrightarrow{p}0$, 证明 $\xi_n+\eta_n\xrightarrow{d}N(0,1)$.\par\smallskip {\small 来源：ayhe123 公开讲义，原题 7.25。}
\end{exercise}

\begin{solve}\phantomsection\label{a:A7.25}固定 $x\in\mathbb R$ 和 $\varepsilon>0$。由事件包含关系，\[P(\xi_n\le x-\varepsilon)-P(|\eta_n|>\varepsilon)\le P(\xi_n+\eta_n\le x),\]\[P(\xi_n+\eta_n\le x)\le P(\xi_n\le x+\varepsilon)+P(|\eta_n|>\varepsilon).\]令 $n\to\infty$，外侧概率分别趋于 $\Phi(x-\varepsilon)$ 与 $\Phi(x+\varepsilon)$；再由 $\Phi$ 连续令 $\varepsilon\downarrow0$，得到中间分布函数趋于 $\Phi(x)$。因此依分布收敛成立，整个证明不需要 $\xi_n$ 与 $\eta_n$ 独立。\end{solve}

\clearpage\section{样本与抽样分布}

\begin{exercise}[统计量与未知参数]\label{q:M-F-C4}
地铁延误时间 $X_1,\ldots,X_5$ 独立同分布，均服从未知率参数 $\lambda>0$ 的指数分布。判断以下量是否为统计量，并说明理由：$1/\lambda$、样本均值 $\overline X$、$\overline X-1/\lambda$、样本标准差 $S$。\par\smallskip {\small 来源：MIT 期末试卷，概念题 4。}
\end{exercise}

\begin{solve}\phantomsection\label{a:M-F-C4}统计量是样本的可测函数，其定义不能包含未知参数。因此 $\overline X$ 和 $S$ 是统计量；$1/\lambda$ 是总体均值，包含未知参数；$\overline X-1/\lambda$ 虽然也是随机量，但不属于本模型的统计量。注意“分布依赖参数”和“计算公式包含未知参数”不同：$\overline X$ 的抽样分布依赖 $\lambda$，仍然是统计量。观测前用大写表示随机统计量，观测后用 $\overline x,s$ 表示已经计算出的数值。\end{solve}

\begin{exercise}[样本均值的均值与方差]\label{q:B8.1}
设 $X_1,\ldots,X_n$ 独立同分布，$EX_i=\mu$，$\operatorname{Var}(X_i)=\sigma^2<\infty$。求 $E\overline X$、$\operatorname{Var}(\overline X)$ 及 $\operatorname{Cov}(X_i,\overline X)$，并说明若样本相关，哪些公式仍成立。\par\smallskip {\small 来源：编订补充。}
\end{exercise}

\begin{solve}\phantomsection\label{a:B8.1}由期望的线性性，不需要独立性即可得到 $E\overline X=\mu$。独立性使交叉协方差为零，所以\[\operatorname{Var}(\overline X)=\frac{1}{n^2}\sum_{j=1}^n\sigma^2=\frac{\sigma^2}{n},\qquad\operatorname{Cov}(X_i,\overline X)=\frac{\sigma^2}{n}.\]若各观测仍有相同均值、方差但存在相关性，则\[\operatorname{Var}(\overline X)=\frac{\sigma^2}{n}+\frac{2}{n^2}\sum_{i<j}\operatorname{Cov}(X_i,X_j).\]正相关可能显著降低平均的降噪效果，故不能把独立样本的标准误机械用于相关数据。\end{solve}

\begin{exercise}[样本方差的无偏性]\label{q:B8.2}
在独立同分布且方差有限的条件下，令 $S^2=(n-1)^{-1}\sum_{i=1}^n(X_i-\overline X)^2$，$n\ge2$。证明 $ES^2=\sigma^2$，并求使用分母 $n$ 的样本二阶中心矩的偏差。\par\smallskip {\small 来源：编订补充。}
\end{exercise}

\begin{solve}\phantomsection\label{a:B8.2}先作恒等变形：\[\sum_i(X_i-\overline X)^2=\sum_i(X_i-\mu)^2-n(\overline X-\mu)^2.\]右侧两项的期望分别为 $n\sigma^2$ 和 $n\operatorname{Var}(\overline X)=\sigma^2$，所以平方和的期望为 $(n-1)\sigma^2$，除以 $n-1$ 即得无偏性。若分母为 $n$，估计量的期望为 $(n-1)\sigma^2/n$，偏差为 $-\sigma^2/n$。自由度损失来自用同一组数据估计了均值；上述无偏性不要求总体正态。\end{solve}

\begin{exercise}[正态样本均值和残差的独立性]\label{q:B8.3}
设 $X_i$ 独立同分布于 $N(\mu,\sigma^2)$，$\sigma>0$。证明 $\overline X$ 与整个残差向量 $(X_1-\overline X,\ldots,X_n-\overline X)$ 独立，并推出 $\overline X$ 与 $S^2$ 独立。\par\smallskip {\small 来源：编订补充。}
\end{exercise}

\begin{solve}\phantomsection\label{a:B8.3}样本向量为联合正态向量，任何线性变换后的向量仍联合正态。对每个 $i$，\[\operatorname{Cov}(\overline X,X_i-\overline X)=\sigma^2/n-\sigma^2/n=0.\]因此样本均值与残差向量之间的交叉协方差矩阵为零。联合正态向量中，零交叉协方差意味着两个子向量独立；这个结论也允许残差协方差矩阵奇异。$S^2$ 是残差向量的可测函数，因此也与 $\overline X$ 独立。一般总体下，不相关并不推出独立。\end{solve}

\begin{exercise}[卡方分布及自由度]\label{q:B8.4}
设 $X_i$ 为容量 $n\ge2$ 的正态样本。证明 $(n-1)S^2/\sigma^2\sim\chi^2_{n-1}$，并解释为何不是 $\chi^2_n$。\par\smallskip {\small 来源：编订补充。}
\end{exercise}

\begin{solve}\phantomsection\label{a:B8.4}标准化后记 $Z_i=(X_i-\mu)/\sigma$。取正交矩阵，其第一行为 $n^{-1/2}(1,\ldots,1)$。变换后得到独立标准正态变量 $U_1,\ldots,U_n$，其中 $U_1=\sqrt n(\overline X-\mu)/\sigma$。正交变换保持平方和，故\[\frac{(n-1)S^2}{\sigma^2}=\sum_{i=1}^nZ_i^2-U_1^2=\sum_{j=2}^nU_j^2.\]最后一项是 $n-1$ 个独立标准正态变量平方和，按定义服从 $\chi^2_{n-1}$。估计未知均值去掉了沿全一向量的一个自由方向。\end{solve}

\begin{exercise}[学生 t 统计量]\label{q:B8.5}
设 $X_1,\ldots,X_n$ 为来自 $N(\mu,\sigma^2)$ 的样本。推导 $T=\sqrt n(\overline X-\mu)/S$ 的分布，并说明为什么直接将 $S$ 代替 $\sigma$ 后仍使用标准正态分布不是精确做法。\par\smallskip {\small 来源：编订补充。}
\end{exercise}

\begin{solve}\phantomsection\label{a:B8.5}令 $Z=\sqrt n(\overline X-\mu)/\sigma$，$V=(n-1)S^2/\sigma^2$。正态抽样理论给出 $Z\sim N(0,1)$、$V\sim\chi^2_{n-1}$，且两者独立。于是\[T=\frac{Z}{\sqrt{V/(n-1)}}\sim t_{n-1}.\]分母含有估计方差的不确定性，使 $t$ 分布尾部较厚。$n$ 很大时它趋近标准正态；有限 $n$ 时若使用正态分位数，均值区间会偏窄。非正态总体下，这一有限样本精确结论通常不成立。\end{solve}

\begin{exercise}[方差比与 F 分布]\label{q:B8.6}
两组相互独立的正态样本容量分别为 $n,m\ge2$，方差为 $\sigma_1^2,\sigma_2^2$。求 $(S_1^2/\sigma_1^2)/(S_2^2/\sigma_2^2)$ 的分布，并推导 $F$ 分位数的倒数关系。\par\smallskip {\small 来源：编订补充。}
\end{exercise}

\begin{solve}\phantomsection\label{a:B8.6}两个独立的标准化平方和分别服从 $\chi^2_{n-1}$ 和 $\chi^2_{m-1}$。因此\[\frac{S_1^2/\sigma_1^2}{S_2^2/\sigma_2^2}\sim F_{n-1,m-1}.\]若 $W\sim F_{a,b}$，则按定义交换分子分母可得 $1/W\sim F_{b,a}$。记 $F_{a,b;q}$ 为下侧累积概率 $q$ 的分位数，则连续性给出\[F_{a,b;q}=\frac{1}{F_{b,a;1-q}}.\]使用表格时必须核对分子和分母自由度顺序，也要核对采用的是下侧还是上侧分位数。\end{solve}

\begin{exercise}[顺序统计量的分布]\label{q:B8.7}
设 $X_i$ 独立同分布于连续分布函数 $F$，密度为 $f$。推导最大值 $X_{(n)}$、最小值 $X_{(1)}$ 的分布函数，以及第 $k$ 个顺序统计量的密度。\par\smallskip {\small 来源：编订补充。}
\end{exercise}

\begin{solve}\phantomsection\label{a:B8.7}由独立性，所有观测不超过 $x$ 的概率为 $F(x)^n$，所以 $F_{X_{(n)}}(x)=F(x)^n$；类似地 $P(X_{(1)}>x)=[1-F(x)]^n$。第 $k$ 个顺序统计量落在 $x$ 附近，要求恰有 $k-1$ 个观测小于 $x$，一个落在微小区间，其余大于 $x$。求极限或对二项尾和求导，得到\[f_{(k)}(x)=\frac{n!}{(k-1)!(n-k)!}F(x)^{k-1}[1-F(x)]^{n-k}f(x).\]连续性使相等的并列观测概率为零。离散总体中不能直接使用这个密度公式。\end{solve}

\begin{exercise}[均匀样本最大值]\label{q:B8.8}
设 $X_i$ 独立同分布于 $U(0,\theta)$，$\theta>0$。令 $M=X_{(n)}$。求 $EM,\operatorname{Var}(M)$，并构造 $\theta$ 的无偏估计量。\par\smallskip {\small 来源：编订补充。}
\end{exercise}

\begin{solve}\phantomsection\label{a:B8.8}$M$ 的密度为 $nx^{n-1}/\theta^n$，$0<x<\theta$。对 $r>0$ 积分得到\[E(M^r)=\frac{n}{n+r}\theta^r.\]故 $EM=n\theta/(n+1)$，而\[\operatorname{Var}(M)=\left[\frac{n}{n+2}-\frac{n^2}{(n+1)^2}\right]\theta^2=\frac{n\theta^2}{(n+1)^2(n+2)}.\]将最大值乘以 $(n+1)/n$ 得到无偏估计量 $\widetilde\theta=(n+1)M/n$。注意它可能大于实际观测最大值，这不违反参数作为总体上界的含义。\end{solve}

\begin{exercise}[指数样本的最小值与和]\label{q:B8.9}
设 $X_i$ 独立同分布于率为 $\lambda$ 的指数分布。求最小值 $M$ 和总和 $T$ 的分布，并比较 $nM$ 与 $\overline X$ 的均值和方差。\par\smallskip {\small 来源：编订补充。}
\end{exercise}

\begin{solve}\phantomsection\label{a:B8.9}对 $t\ge0$，$P(M>t)=\prod_i e^{-\lambda t}=e^{-n\lambda t}$，所以 $M\sim\operatorname{Exp}(n\lambda)$。拉普拉斯变换相乘给出 $T\sim\operatorname{Gamma}(n,\lambda)$，这里采用率参数。因此\[E(nM)=E\overline X=1/\lambda,\quad\operatorname{Var}(nM)=1/\lambda^2,\quad\operatorname{Var}(\overline X)=1/(n\lambda^2).\]两者都无偏估计总体均值，但 $nM$ 的方差不随样本容量缩小，$\overline X$ 则具有一致性。这说明无偏性本身不能保证估计质量。\end{solve}

\begin{exercise}[中心极限定理的正确标准化]\label{q:M-F-C13}
设独立同分布样本的均值为 $\mu$，方差为 $0<\sigma^2<\infty$。判断：大样本时每个观测近似正态；样本均值近似正态；$(\overline X-\mu)/\sigma$ 近似标准正态；$\sqrt n(\overline X-\mu)/\sigma$ 近似标准正态。\par\smallskip {\small 来源：MIT 期末试卷，概念题 13（补足有限正方差条件）。}
\end{exercise}

\begin{solve}\phantomsection\label{a:M-F-C13}依次为错、对、错、对。中心极限定理讨论的是加和后的标准化量，并不改变每个原始观测的分布。其精确渐近表述为\[\frac{\sqrt n(\overline X-\mu)}{\sigma}\xrightarrow{d}N(0,1).\]因此 $\overline X$ 可近似视为 $N(\mu,\sigma^2/n)$，而没有乘 $\sqrt n$ 的量方差是 $1/n$。这里的“近似”没有自动给出某个固定样本容量的误差保证，重尾或强偏态会使收敛较慢。\end{solve}

\clearpage\section{矩估计、偏差与一致性}

\begin{exercise}[伯努利矩估计]\label{q:B9.1}
独立样本来自 $\operatorname{Bern}(p)$。用矩法估计 $p$，求估计量的偏差、方差和均方误差，并证明一致性。\par\smallskip {\small 来源：编订补充。}
\end{exercise}

\begin{solve}\phantomsection\label{a:B9.1}总体一阶矩为 $EX=p$，以样本一阶矩替代得到 $\widehat p=\overline X$。于是\[E\widehat p=p,\qquad\operatorname{MSE}(\widehat p)=\operatorname{Var}(\widehat p)=\frac{p(1-p)}n.\]切比雪夫不等式给出对任意 $\varepsilon>0$，\[P(|\widehat p-p|>\varepsilon)\le\frac{p(1-p)}{n\varepsilon^2}\longrightarrow0.\]因此估计量依概率一致。矩估计是一个构造方法，无偏性和一致性则是需要分别验证的性质。\end{solve}

\begin{exercise}[均匀分布的矩估计与极值估计]\label{q:B9.2}
样本来自 $U(0,\theta)$。求矩估计 $\widehat\theta_M$，并比较它与 $\widetilde\theta=(n+1)X_{(n)}/n$ 的均方误差。\par\smallskip {\small 来源：编订补充。}
\end{exercise}

\begin{solve}\phantomsection\label{a:B9.2}由 $EX=\theta/2$ 得 $\widehat\theta_M=2\overline X$。它无偏且方差为 $4\theta^2/(12n)=\theta^2/(3n)$。修正最大值估计也无偏，其方差为\[\operatorname{Var}(\widetilde\theta)=\frac{\theta^2}{n(n+2)}.\]二者均方误差之比为 $(n+2)/3$，所以 $n>1$ 时修正最大值更好，$n=1$ 时相同。这里充分利用参数限制支撑集的信息，比只匹配均值获得更快的误差下降。\end{solve}

\begin{exercise}[伽马分布的两个矩参数]\label{q:B9.3}
设样本来自 $\operatorname{Gamma}(\alpha,\beta)$，$\alpha,\beta>0$，$\beta$ 为率。设 $m_1=\overline X$，$v_n=n^{-1}\sum_i(X_i-\overline X)^2$。由前两个矩构造 $\alpha,\beta$ 的估计。\par\smallskip {\small 来源：编订补充。}
\end{exercise}

\begin{solve}\phantomsection\label{a:B9.3}总体均值、方差分别为 $\alpha/\beta$ 和 $\alpha/\beta^2$。将它们等同于 $m_1,v_n$，解方程得\[\widehat\alpha=\frac{m_1^2}{v_n},\qquad\widehat\beta=\frac{m_1}{v_n}.\]有限样本需 $v_n>0$；当 $n\ge2$ 且总体连续时，样本全相等的概率为零。由大数定律，样本一阶、二阶原点矩收敛，进而 $m_1\to\alpha/\beta$、$v_n\to\alpha/\beta^2$；连续映射定理给出两个估计的一致性。比值形式一般并非无偏。\end{solve}

\begin{exercise}[正态分布的矩估计]\label{q:B9.4}
样本来自 $N(\mu,\sigma^2)$，两个参数均未知。用一阶和二阶原点矩求矩估计，并比较其方差估计与无偏样本方差。\par\smallskip {\small 来源：编订补充。}
\end{exercise}

\begin{solve}\phantomsection\label{a:B9.4}理论矩为 $EX=\mu$、$EX^2=\mu^2+\sigma^2$。配对样本矩可得\[\widehat\mu=\overline X,\qquad\widehat{\sigma^2}=\frac1n\sum_iX_i^2-\overline X^2=\frac1n\sum_i(X_i-\overline X)^2.\]后者等于 $(n-1)S^2/n$，期望为 $(n-1)\sigma^2/n$，有向下偏差。矩法并不保证无偏。若改用 $S^2$，则得到无偏方差估计，但它不再是严格按照这两个原点矩方程得到的估计。\end{solve}

\begin{exercise}[偏差与均方误差分解]\label{q:B9.5}
对二阶矩有限的估计量 $T$，证明 $E_\theta(T-\theta)^2=\operatorname{Var}_\theta(T)+[E_\theta T-\theta]^2$。若 $X_i\sim N(\mu,\sigma^2)$ 且 $\sigma^2$ 已知，比较 $\overline X$ 与 $a\overline X$ 的均方误差。\par\smallskip {\small 来源：编订补充。}
\end{exercise}

\begin{solve}\phantomsection\label{a:B9.5}在 $T-\theta=(T-ET)+(ET-\theta)$ 中平方后取期望，交叉项因 $E(T-ET)=0$ 消失，得到分解。对 $a\overline X$，偏差为 $(a-1)\mu$，方差为 $a^2\sigma^2/n$，故\[\operatorname{MSE}_\mu(a\overline X)=a^2\sigma^2/n+(a-1)^2\mu^2.\]若 $0<a<1$ 且 $\mu$ 靠近零，缩小均值可能降低均方误差；但对较大的 $|\mu|$，偏差平方会占主导。不能只比较方差便断言一个有偏估计统一优于无偏估计。\end{solve}

\begin{exercise}[样本方差的最小均方误差缩放]\label{q:B9.6}
正态样本容量 $n\ge2$。在形如 $aS^2$ 的估计量中，求估计 $\sigma^2$ 的最小均方误差系数 $a$。\par\smallskip {\small 来源：编订补充。}
\end{exercise}

\begin{solve}\phantomsection\label{a:B9.6}因为 $(n-1)S^2/\sigma^2\sim\chi^2_{n-1}$，且卡方分布方差为自由度的两倍，故 $\operatorname{Var}(S^2)=2\sigma^4/(n-1)$。于是\[\operatorname{MSE}(aS^2)=\sigma^4\left[\frac{2a^2}{n-1}+(a-1)^2\right].\]对 $a$ 求导并令其为零，得到 $a=(n-1)/(n+1)$；二阶导数为正。最小均方误差为 $2\sigma^4/(n+1)$。这个最优性只针对所指定的常数倍估计量类，且依赖正态假设。\end{solve}

\begin{exercise}[一致性与渐近正态性]\label{q:B9.7}
设 $X_i$ 独立同分布于 $\operatorname{Pois}(\lambda)$，$\lambda>0$。证明 $\overline X$ 一致估计 $\lambda$，并求 $\log\overline X$ 的渐近分布；解释全零样本的处理。\par\smallskip {\small 来源：编订补充。}
\end{exercise}

\begin{solve}\phantomsection\label{a:B9.7}有限均值保证大数定律适用，故 $\overline X\xrightarrow{P}\lambda$。又由中心极限定理，$\sqrt n(\overline X-\lambda)\xrightarrow{d}N(0,\lambda)$。函数 $g(x)=\log x$ 在 $\lambda>0$ 处可微，导数为 $1/\lambda$，德尔塔方法给出\[\sqrt n(\log\overline X-\log\lambda)\xrightarrow{d}N(0,1/\lambda).\]全零样本概率为 $e^{-n\lambda}$，趋于零；可在该事件上将估计另定义为任意有限数，从而得到处处有定义的渐近结论。不能对有限样本直接计算 $\log0$。\end{solve}

\begin{exercise}[样本分位数的直观比较]\label{q:B9.8}
对均匀总体 $U(0,\theta)$，比较由 $2\overline X$ 与 $2X_1$ 估计 $\theta$。两者是否无偏、是否一致？由此说明“大样本可用更多数据”与“估计公式实际利用更多数据”的区别。\par\smallskip {\small 来源：编订补充。}
\end{exercise}

\begin{solve}\phantomsection\label{a:B9.8}两者期望均为 $\theta$，故都无偏；方差分别为 $\theta^2/(3n)$ 与 $\theta^2/3$。前者依概率趋于 $\theta$。后者的分布不随 $n$ 改变；例如对 $0<\varepsilon<\theta$，\[P(|2X_1-\theta|>\varepsilon)=1-\varepsilon/\theta>0,\]所以不一致。拥有越来越多观测，并不使一个始终只用第一个观测的估计量自动变好。此例还说明无偏与一致是不同性质。\end{solve}

\clearpage\section{极大似然估计}

\begin{exercise}[离散参数空间中的硬币似然]\label{q:M-C10-1}
独立抛硬币 $80$ 次，观察到 $49$ 次正面、$31$ 次反面。先假定正面概率只可能是 $1/3,1/2,2/3$，求极大似然估计；再允许 $p\in[0,1]$，重新求估计。\par\smallskip {\small 来源：MIT 第 10 讲，题 1。}
\end{exercise}

\begin{solve}\phantomsection\label{a:M-C10-1}似然为 $L(p)=\binom{80}{49}p^{49}(1-p)^{31}$。在三个候选点的值约为 $2.08\times10^{-7},0.0118,0.0545$，故离散参数空间中的极大似然估计为 $2/3$。连续参数空间内，\[\ell^{\prime}(p)=49/p-31/(1-p)=0\]给出 $\widehat p=49/80=0.6125$；二阶导数严格为负，且两个边界似然为零，所以这是全局最大值。最大值取决于允许的参数集合，而不只是由数据决定。\end{solve}

\begin{exercise}[指数寿命的极大似然估计]\label{q:M-C10-2}
五个独立寿命观测为 $2,3,1,3,4$，假定来自率参数未知的指数分布。写出似然、对数似然和极大似然估计，并给出容量为 $n$ 的一般公式。\par\smallskip {\small 来源：MIT 第 10 讲，题 2。}
\end{exercise}

\begin{solve}\phantomsection\label{a:M-C10-2}五个密度相乘得到 $L(\lambda)=\lambda^5e^{-13\lambda}$，所以 $\ell(\lambda)=5\log\lambda-13\lambda$。导数为 $5/\lambda-13$，二阶导数为 $-5/\lambda^2<0$，因而 $\widehat\lambda=5/13$。一般正观测下，\[\widehat\lambda=\frac{n}{\sum_i x_i}=\frac1{\overline x}.\]这里估计的是单位时间事件率；平均寿命的极大似然估计为 $\overline x$。率和平均寿命是互为倒数的两个参数，须避免混用。\end{solve}

\begin{exercise}[已知形状的伽马率估计]\label{q:M-E2-II1}
总体密度为 $f(x;\beta)=\beta^3x^2e^{-\beta x}/2$，$x>0$，$\beta>0$。五个独立观测为 $1,1,2,3,5$。求 $\beta$ 的极大似然估计，并推广到已知形状 $\alpha$ 的伽马模型。\par\smallskip {\small 来源：MIT 第二次考试，第二部分题 1。}
\end{exercise}

\begin{solve}\phantomsection\label{a:M-E2-II1}不含 $\beta$ 的因子记为 $C>0$，则 $L(\beta)=C\beta^{15}e^{-12\beta}$。对数似然为 $15\log\beta-12\beta+\log C$，故驻点满足 $15/\beta-12=0$，得到 $\widehat\beta=5/4$。严格负的二阶导数保证最大值。一般地，$\operatorname{Gamma}(\alpha,\beta)$ 样本的率参数估计为\[\widehat\beta=\frac{n\alpha}{\sum_iX_i}.\]这也等于已知形状时的矩估计；若形状未知，两个参数的似然方程不再具有同样的简单闭式。\end{solve}

\begin{exercise}[正态模型的联合极大似然]\label{q:B10.1}
对正态样本，同时估计未知的 $\mu,\sigma^2$。推导联合极大似然估计，并说明样本所有值相同的特殊情形。\par\smallskip {\small 来源：编订补充。}
\end{exercise}

\begin{solve}\phantomsection\label{a:B10.1}令 $v=\sigma^2>0$，对数似然为\[\ell(\mu,v)=-\frac n2\log(2\pi v)-\frac{1}{2v}\sum_i(x_i-\mu)^2.\]固定 $v$，平方和在 $\mu=\overline x$ 时最小。代回后对 $v$ 求导，得到\[\widehat\mu=\overline x,\qquad\widehat v=\frac1n\sum_i(x_i-\overline x)^2.\]当平方和为正时，该驻点是全局最大值。若所有观测相同，则沿 $\mu$ 等于共同值、$v\downarrow0$，似然无界，在 $v>0$ 的参数空间中没有有限极大值；连续正态样本容量至少二时此事件概率为零。\end{solve}

\begin{exercise}[均匀端点与支撑集]\label{q:B10.2}
设样本来自闭区间密度表示 $f(x;\theta)=\theta^{-1}\mathbf1_{[0,\theta]}(x)$。求 $\theta$ 的极大似然估计，并解释只对 $-n\log\theta$ 求导为什么不够。\par\smallskip {\small 来源：编订补充。}
\end{exercise}

\begin{solve}\phantomsection\label{a:B10.2}联合似然为\[L(\theta)=\theta^{-n}\mathbf1_{\{\theta\ge M\}},\qquad M=\max_i x_i.\]在可行域 $\theta\ge M$ 内似然随 $\theta$ 递减，所以最大值在边界 $\widehat\theta=M$。若忽略指示函数，就会错误地把参数推向零，即使这样会使观测落在支撑集之外。密度在端点的取值不影响概率分布，但可影响极大值是否达到；这里明确采用包含端点的密度版本。\end{solve}

\begin{exercise}[几何分布的估计与边界]\label{q:B10.3}
设 $P(X=k)=p(1-p)^{k-1}$，$k=1,2,\ldots$，$0<p\le1$。独立样本容量为 $n$，求 $p$ 的极大似然估计。若改成记录首次成功前的失败次数，公式如何改变？\par\smallskip {\small 来源：编订补充。}
\end{exercise}

\begin{solve}\phantomsection\label{a:B10.3}第一种记法的似然为 $p^n(1-p)^{T-n}$，$T=\sum_i x_i\ge n$。若 $T>n$，对数导数 $n/p-(T-n)/(1-p)=0$，得到 $\widehat p=n/T=1/\overline x$。若 $T=n$，所有观测都是 $1$，最大值在 $p=1$。若记失败次数 $Y=X-1$，则 $\widehat p=1/(1+\overline Y)$。几何分布存在两种通行支撑约定，估计公式必须随之改变。\end{solve}

\begin{exercise}[带删失的指数寿命]\label{q:B10.4}
独立寿命 $X_i\sim\operatorname{Exp}(\lambda)$，每个个体最多观察到固定时间 $c>0$。记录 $Y_i=\min(X_i,c)$ 和失效指示 $\delta_i=\mathbf1_{\{X_i\le c\}}$。推导 $\lambda$ 的极大似然估计。\par\smallskip {\small 来源：编订补充。}
\end{exercise}

\begin{solve}\phantomsection\label{a:B10.4}已失效个体的贡献为 $\lambda e^{-\lambda y_i}$，未失效个体贡献为生存概率 $e^{-\lambda c}$。合并得\[L(\lambda)=\lambda^d\exp\{-\lambda\sum_i y_i\},\qquad d=\sum_i\delta_i.\]当 $d>0$ 时，对数似然求导得到 $\widehat\lambda=d/\sum_i y_i$，即失效次数除以累计观察时间。当 $d=0$ 时，似然随 $\lambda$ 递减，在严格正参数空间中只有 $\lambda\downarrow0$ 的上确界；若扩展模型允许零率，可记边界估计为零。不能把所有未失效寿命误记成恰好在 $c$ 失效。\end{solve}

\begin{exercise}[极大似然的不变性]\label{q:B10.5}
若 $\widehat\theta$ 是参数 $\theta$ 的极大似然估计，且 $g$ 为一一映射，证明 $g(\widehat\theta)$ 是 $g(\theta)$ 的极大似然估计。以指数模型的中位数为例应用该结论。\par\smallskip {\small 来源：编订补充。}
\end{exercise}

\begin{solve}\phantomsection\label{a:B10.5}记新参数 $\eta=g(\theta)$，新似然只是重新参数化：$\widetilde L(\eta)=L(g^{-1}(\eta))$。一一对应保持候选分布不变，因此最大值在 $\eta=g(\widehat\theta)$ 达到。指数分布中位数 $m$ 满足 $1-e^{-\lambda m}=1/2$，故 $m=\log2/\lambda$，极大似然估计为\[\widehat m=\frac{\log2}{\widehat\lambda}=\overline X\log2.\]对非一一映射也有适当的轮廓似然版本，但不能把“估计量变换”与“变换后仍无偏”混为一谈。\end{solve}

\begin{exercise}[指数率估计的有限样本偏差]\label{q:B10.6}
指数率为 $\lambda$ 的样本总和 $T\sim\operatorname{Gamma}(n,\lambda)$。求 $E(n/T)$，并构造 $\lambda$ 的无偏估计。说明所需样本容量条件。\par\smallskip {\small 来源：编订补充。}
\end{exercise}

\begin{solve}\phantomsection\label{a:B10.6}将伽马密度积分可得，对 $n>1$，\[E(T^{-1})=\frac{\lambda^n}{\Gamma(n)}\int_0^\infty t^{n-2}e^{-\lambda t}\,dt=\frac{\lambda}{n-1}.\]故极大似然估计 $n/T$ 的期望为 $n\lambda/(n-1)$，向上有偏；无偏修正为 $(n-1)/T$。当 $n=1$ 时倒数期望发散，不能沿用此结论。若 $n>2$，还可得 $E(T^{-2})=\lambda^2/[(n-1)(n-2)]$，从而无偏修正的方差为 $\lambda^2/(n-2)$。\end{solve}

\clearpage\section{充分统计量与无偏估计}

\begin{exercise}[伯努利样本的充分统计量]\label{q:B11.1}
设 $X_i\sim\operatorname{Bern}(p)$ 独立。证明总成功次数 $T=\sum_iX_i$ 是 $p$ 的充分统计量，并给出给定 $T=t$ 时样本向量的条件分布。\par\smallskip {\small 来源：编订补充。}
\end{exercise}

\begin{solve}\phantomsection\label{a:B11.1}联合概率为 $p^t(1-p)^{n-t}$，对所有含恰好 $t$ 个 $1$ 的向量相同。因为 $P(T=t)=\binom nt p^t(1-p)^{n-t}$，条件概率为\[P((X_1,\ldots,X_n)=x\mid T=t)=\binom nt^{-1}\]（当 $\sum x_i=t$），其他向量为零。它不含 $p$，这正说明给定 $T$ 后样本排列不再提供关于 $p$ 的信息，也可用因子分解定理证明充分性。\end{solve}

\begin{exercise}[泊松样本的充分性与条件二项分布]\label{q:B11.2}
设 $X_i\sim\operatorname{Pois}(\lambda)$ 独立。证明 $T=\sum_iX_i$ 充分，并求 $X_1\mid T=t$ 的分布。\par\smallskip {\small 来源：编订补充。}
\end{exercise}

\begin{solve}\phantomsection\label{a:B11.2}联合质量函数为 $e^{-n\lambda}\lambda^{\sum x_i}/\prod_i x_i!$，按因子分解定理 $T$ 充分。又 $X_1\sim\operatorname{Pois}(\lambda)$ 与 $T-X_1\sim\operatorname{Pois}((n-1)\lambda)$ 独立，故\[P(X_1=k\mid T=t)=\binom tk(1/n)^k(1-1/n)^{t-k}.\]因此条件分布为 $\operatorname{Bin}(t,1/n)$，不含 $\lambda$。这个条件分布还能用于将只使用一个观测的估计改进为使用全样本的信息。\end{solve}

\begin{exercise}[拉奥—布莱克改进]\label{q:B11.3}
在泊松样本模型下，$U=\mathbf1_{\{X_1=0\}}$ 无偏估计 $e^{-\lambda}$。计算 $E(U\mid T)$，其中 $T=\sum_iX_i$，并证明改进后的均方误差不大于原估计。\par\smallskip {\small 来源：编订补充。}
\end{exercise}

\begin{solve}\phantomsection\label{a:B11.3}由条件二项分布，\[\widetilde U=E(U\mid T)=P(X_1=0\mid T)=(1-1/n)^T.\]当 $n=1,T=0$ 时约定 $0^0=1$。迭代期望保证 $E\widetilde U=EU=e^{-\lambda}$。全方差公式给出\[\operatorname{Var}(U)=E\{\operatorname{Var}(U\mid T)\}+\operatorname{Var}(\widetilde U),\]所以方差不会增大。由于二者均无偏，方差比较也就是均方误差比较。这个过程利用条件期望去掉在充分信息给定后仍存在的无关随机波动。\end{solve}

\begin{exercise}[二项族的完全性]\label{q:B11.4}
设 $T\sim\operatorname{Bin}(n,p)$，$p\in(0,1)$。若函数 $g$ 满足对所有 $p$ 都有 $E_pg(T)=0$，证明 $g(t)=0$，$t=0,\ldots,n$。说明该结论对最小方差无偏估计的意义。\par\smallskip {\small 来源：编订补充。}
\end{exercise}

\begin{solve}\phantomsection\label{a:B11.4}假设等价于\[\sum_{t=0}^ng(t)\binom nt p^t(1-p)^{n-t}=0.\]除以正数 $(1-p)^n$，令 $z=p/(1-p)\in(0,\infty)$，得到一个在正半轴恒为零的多项式 $\sum_tg(t)\binom ntz^t$。多项式各系数必须为零，所以 $g(t)=0$。因此 $T$ 完全；加上充分性，莱曼—谢费定理说明任何作为 $T$ 的函数的无偏估计，都是对应参数函数的唯一几乎处处的均匀最小方差无偏估计。\end{solve}

\begin{exercise}[估计平方概率]\label{q:B11.5}
伯努利样本容量 $n\ge2$。构造 $p^2$ 的无偏估计，并与直接代入估计 $\widehat p^2$ 比较偏差。\par\smallskip {\small 来源：编订补充。}
\end{exercise}

\begin{solve}\phantomsection\label{a:B11.5}总成功次数 $T$ 满足阶乘矩 $E[T(T-1)]=n(n-1)p^2$，故\[U=\frac{T(T-1)}{n(n-1)}\]无偏。它是完全充分统计量 $T$ 的函数，因此也是 $p^2$ 的均匀最小方差无偏估计。代入估计的期望为\[E(T/n)^2=p^2+\frac{p(1-p)}n,\]所以向上有偏。无偏估计可以理解为对所有不同观测对 $X_iX_j$ 求平均；使用不同观测避免将 $X_i^2=X_i$ 混入平方概率的估计。\end{solve}

\begin{exercise}[伯努利模型的信息量与下界]\label{q:B11.6}
对 $p\in(0,1)$ 的伯努利样本，计算费希尔信息量，给出无偏估计 $p$ 的克拉美—拉奥下界，并判断样本比例是否达到该下界。\par\smallskip {\small 来源：编订补充。}
\end{exercise}

\begin{solve}\phantomsection\label{a:B11.6}单个观测的对数似然导数为 $U_p=X/p-(1-X)/(1-p)=(X-p)/[p(1-p)]$。其期望为零，方差为 $1/[p(1-p)]$，所以 $n$ 个独立观测的信息量为\[I_n(p)=\frac{n}{p(1-p)}.\]在可交换微分与求和等正则条件下，无偏估计 $p$ 的方差至少为 $1/I_n(p)=p(1-p)/n$。样本比例恰好具有这个方差，因而在此正则模型中达到下界。边界 $p=0,1$ 不属于这里的正则参数区间。\end{solve}

\begin{exercise}[信息下界并非总能达到]\label{q:B11.7}
在指数率模型中计算 $I_n(\lambda)$。将克拉美—拉奥下界与 $n>2$ 时无偏估计 $(n-1)/T$ 的方差比较，解释为何二者不必相等。\par\smallskip {\small 来源：编订补充。}
\end{exercise}

\begin{solve}\phantomsection\label{a:B11.7}单个对数似然为 $\log\lambda-\lambda X$，二阶导数为 $-1/\lambda^2$，故 $I_n(\lambda)=n/\lambda^2$。无偏估计率参数的方差下界为 $\lambda^2/n$。而 $(n-1)/T$ 的方差为 $\lambda^2/(n-2)$，严格大于下界。下界是所有符合条件无偏估计的共同限制，并不保证存在达到它的估计量。两者之比为 $n/(n-2)\to1$，显示渐近效率与有限样本恰好达到下界是不同概念。\end{solve}

\begin{exercise}[非正则模型中的错误下界]\label{q:B11.8}
有人对 $U(0,\theta)$ 模型计算 $\partial\log f/\partial\theta=-1/\theta$，据此声称任何无偏估计 $\theta$ 的方差至少为 $\theta^2/n$。指出错误，并给出反证性比较。\par\smallskip {\small 来源：编订补充。}
\end{exercise}

\begin{solve}\phantomsection\label{a:B11.8}均匀模型的支撑集随 $\theta$ 改变。只在支撑内部求导会遗漏边界贡献，得分的期望也不是零，所以通常的正则克拉美—拉奥推导不适用。修正最大值 $U=(n+1)X_{(n)}/n$ 是无偏的，方差为\[\operatorname{Var}(U)=\frac{\theta^2}{n(n+2)}<\frac{\theta^2}{n}.\]这不违反任何有效定理，而是说明所声称的下界缺少正则假设。应用信息不等式之前，必须检查微分与积分交换以及参数相关支撑问题。\end{solve}

\clearpage\section{贝叶斯估计与预测}

\begin{exercise}[三类别观察的贝叶斯更新]\label{q:M-E2-II2}
独立试验产生 $A,B,C$ 三种结果，概率分别为 $\theta/2,\theta/2,1-\theta$，$0<\theta<1$。先验密度为 $3\theta^2$。观察序列 $A,B,C,A,B$ 后，求后验密度及后验均值。\par\smallskip {\small 来源：MIT 第二次考试，第二部分题 2。}
\end{exercise}

\begin{solve}\phantomsection\label{a:M-E2-II2}该有序序列的似然为 $(\theta/2)^4(1-\theta)$。与先验相乘，关于 $\theta$ 的部分为 $\theta^6(1-\theta)$，故后验属于 $\operatorname{Beta}(7,2)$。因为\[\int_0^1\theta^6(1-\theta)\,d\theta=\frac17-\frac18=\frac1{56},\]后验密度为 $56\theta^6(1-\theta)$，后验均值为 $7/9$。若只记录三个类别的计数，似然多一个与参数无关的多项式系数，后验不变。\end{solve}

\begin{exercise}[先验赔率与后验赔率]\label{q:M-E2-II3}
某对象具有特征 $H$ 的先验赔率为 $1:100$。筛选规则在 $H$ 下必为阳性，在 $H^c$ 下有 $10\%$ 概率为阳性。阳性后，求后验赔率及后验概率。\par\smallskip {\small 来源：MIT 第二次考试，第二部分题 3（抽象分类表述）。}
\end{exercise}

\begin{solve}\phantomsection\label{a:M-E2-II3}赔率 $1:100$ 表示 $P(H)/P(H^c)=1/100$，并不是 $P(H)=1/100$。阳性的似然比为 $1/0.1=10$，所以\[\frac{P(H\mid+)}{P(H^c\mid+)}=\frac1{100}\times10=\frac1{10}.\]因此后验赔率为 $1:10$，后验概率为 $1/11$。高灵敏度并不使一次阳性自动意味着很高的后验概率，先验比例仍参与计算。本题只是给定概率模型下的推断，不把筛选结果解释成确定标签。\end{solve}

\begin{exercise}[贝塔与二项共轭]\label{q:M-F-C14}
成功概率 $\theta$ 的先验为 $\operatorname{Beta}(5,5)$。独立进行 $30$ 次伯努利试验，观察到 $20$ 次成功。求后验分布、后验均值、后验众数，以及下一次成功的预测概率。\par\smallskip {\small 来源：MIT 期末试卷，概念题 14（扩展计算）。}
\end{exercise}

\begin{solve}\phantomsection\label{a:M-F-C14}先验核为 $\theta^4(1-\theta)^4$，似然核为 $\theta^{20}(1-\theta)^{10}$，所以后验为 $\operatorname{Beta}(25,15)$。于是\[E(\theta\mid D)=25/40=5/8,\qquad\operatorname{mode}(\theta\mid D)=24/38=12/19.\]下一次成功的后验预测概率是对未知参数积分：$P(X_{31}=1\mid D)=E(\theta\mid D)=5/8$。后验均值、众数和极大似然估计 $2/3$ 是不同的概念，有限样本下通常不相同。\end{solve}

\begin{exercise}[平方损失下的贝叶斯估计]\label{q:B12.1}
给定数据后，参数 $\Theta$ 的后验二阶矩有限。在平方损失 $(a-\Theta)^2$ 下，证明最优决策为后验均值；并说明绝对损失对应的解。\par\smallskip {\small 来源：编订补充。}
\end{exercise}

\begin{solve}\phantomsection\label{a:B12.1}在后验分布下作平方分解，\[E[(a-\Theta)^2\mid D]=\operatorname{Var}(\Theta\mid D)+[a-E(\Theta\mid D)]^2.\]第一项不随 $a$ 改变，第二项在后验均值处唯一最小。因此贝叶斯估计依赖损失函数。对绝对损失，在连续后验下求导得到 $2F_{\Theta\mid D}(a)-1$，其过零点为后验中位数；含原子的情形使用左右导数，任何后验中位数都是最优解。\end{solve}

\begin{exercise}[伽马与泊松共轭的多次观察]\label{q:B12.2}
独立计数 $X_i\mid\lambda\sim\operatorname{Pois}(t_i\lambda)$，其中暴露时间 $t_i>0$ 已知。先验 $\lambda\sim\operatorname{Gamma}(a,b)$，采用率参数。求后验及其均值，并解释暴露时间的作用。\par\smallskip {\small 来源：编订补充。}
\end{exercise}

\begin{solve}\phantomsection\label{a:B12.2}似然中依赖 $\lambda$ 的部分是 $\lambda^{\sum_i x_i}\exp[-\lambda\sum_i t_i]$。因此\[\lambda\mid D\sim\operatorname{Gamma}\left(a+\sum_i x_i,\ b+\sum_i t_i\right),\]后验均值为 $(a+\sum_i x_i)/(b+\sum_i t_i)$。若令 $T=\sum_i t_i$，可将其写成 $\frac{b}{b+T}(a/b)+\frac{T}{b+T}(\sum_i x_i/T)$，即先验均值与数据率的加权平均。计数增加形状，累计观察时间增加率；把次数与暴露时间都简单记成样本容量会出错。\end{solve}

\begin{exercise}[正态均值的共轭更新]\label{q:B12.3}
设 $X_i\mid\mu\sim N(\mu,\sigma^2)$ 独立，$\sigma^2$ 已知；先验 $\mu\sim N(m_0,\tau_0^2)$。求后验分布及下一次观测的预测分布。\par\smallskip {\small 来源：编订补充。}
\end{exercise}

\begin{solve}\phantomsection\label{a:B12.3}将先验和似然指数中的二次项配方，后验精度等于精度之和：\[v_n=\left(\tau_0^{-2}+n\sigma^{-2}\right)^{-1},\quad m_n=v_n\left(\frac{m_0}{\tau_0^2}+\frac{n\overline x}{\sigma^2}\right).\]故 $\mu\mid D\sim N(m_n,v_n)$。下一次观测可写成 $X_{n+1}=\mu+\varepsilon$，其中给定数据后噪声独立且方差为 $\sigma^2$，所以预测分布为 $N(m_n,\sigma^2+v_n)$。预测方差比参数后验方差更大，因为还包含新观测本身的随机性。\end{solve}

\begin{exercise}[主持人规则改变三门问题]\label{q:M-P7-1}
奖品等概率位于 $A,B,C$ 三门之一，选手先选 $A$。观察到主持人打开 $B$ 且里面无奖品。分别求以下规则下的后验概率：主持人知道奖品位置，总打开未选的空门，有两扇可开时等概率选择；主持人不知道位置，随机打开 $B,C$；主持人以 $0.7$ 概率按第一规则、以 $0.3$ 概率按第二规则，且规则选择独立于奖品位置。\par\smallskip {\small 来源：MIT 习题 7，题 1。}
\end{exercise}

\begin{solve}\phantomsection\label{a:M-P7-1}令 $D$ 为“打开 $B$ 且无奖品”。第一种规则下，$D$ 在奖品分别位于 $A,B,C$ 时的似然为 $(1/2,0,1)$。乘以相等先验并归一化，后验为 $(1/3,0,2/3)$，应换到 $C$。第二种规则的似然为 $(1/2,0,1/2)$，后验为 $(1/2,0,1/2)$，换与不换相同。混合规则下似然为 $(1/2,0,0.7+0.3/2)=(1/2,0,17/20)$，归一化得\[(P(A\mid D),P(B\mid D),P(C\mid D))=(10/27,0,17/27).\]仍应换门。观察表面相同，但观察是如何产生的决定了似然。\end{solve}

\begin{exercise}[未知面数骰子的预测]\label{q:M-P7-2}
从面数为 $4,6,8,12,20$ 的五枚公平骰子中等概率选一枚；每枚各面标为 $1$ 至其面数。选定后独立投掷，连续 $n\ge1$ 次都得到 $7$。求面数的后验分布、下一次仍为 $7$ 的预测概率及其极限。\par\smallskip {\small 来源：MIT 习题 7，题 2（选编）。}
\end{exercise}

\begin{solve}\phantomsection\label{a:M-P7-2}面数为 $4,6$ 时似然为零。对 $d\in\{8,12,20\}$，似然为 $d^{-n}$，所以\[P(D=d\mid7,\ldots,7)=\frac{d^{-n}}{8^{-n}+12^{-n}+20^{-n}}.\]下一次为 $7$ 的概率是条件概率的后验平均，等于\[\frac{8^{-(n+1)}+12^{-(n+1)}+20^{-(n+1)}}{8^{-n}+12^{-n}+20^{-n}}.\]将分子分母同乘 $8^n$ 可见后验集中到八面骰，预测概率趋于 $1/8$。连续看到同一结果不会使下一次结果的概率趋于一，因为给定骰子之后仍有投掷随机性。\end{solve}

\begin{exercise}[显著性与后验概率的区别]\label{q:M-P8-1fg}
独立试验 $250$ 次，正面 $140$ 次。先只考虑 $p=0.5$ 和 $p=0.55$ 两个假设，先验各为 $1/2$，求后者后验概率；再取 $p\sim\operatorname{Beta}(1,1)$，写出后验并求 $P(p>0.5\mid D)$ 的表达式。为什么这些量不是双侧检验的 $p$ 值？\par\smallskip {\small 来源：MIT 习题 8，题 1(f,g)（选编）。}
\end{exercise}

\begin{solve}\phantomsection\label{a:M-P8-1fg}两个点假设的似然比为 $r=0.55^{140}0.45^{110}/0.5^{250}$，所以后验为 $r/(1+r)\approx0.8524$。连续均匀先验给出 $p\mid D\sim\operatorname{Beta}(141,111)$，于是\[P(p>0.5\mid D)=1-F_{\operatorname{Beta}(141,111)}(0.5)\approx0.97.\]这些是在指定先验后对参数的概率陈述。频率学派 $p$ 值则是在零假设成立时，出现至少同样极端数据的概率；它没有把概率分配给零假设本身。选择不同先验或参数空间，后验可以不同。\end{solve}

\clearpage\section{区间估计与样本量}

\begin{exercise}[小样本正态总体的三类区间]\label{q:M-C22-1}
独立正态样本观测为 $4,1,2,3$。先假定总体标准差已知为 $2$，求均值的 $90\%$ 置信区间；再假定标准差未知，求均值、方差和标准差各自的 $90\%$ 置信区间。\par\smallskip {\small 来源：MIT 第 22 讲，题 1(a—d)。}
\end{exercise}

\begin{solve}\phantomsection\label{a:M-C22-1}样本容量 $n=4$，均值 $\overline x=2.5$，离差平方和为 $5$，故 $s^2=5/3$。已知标准差时，均值区间为\[2.5\pm z_{0.95}\frac2{\sqrt4}=(0.8551,4.1449).\]未知标准差时使用 $t_3$ 分位数 $2.3534$，得到\[2.5\pm2.3534\sqrt{5/12}=(0.9809,4.0191).\]方差区间由卡方枢轴得到\[\left(\frac5{\chi^2_{3;0.95}},\frac5{\chi^2_{3;0.05}}\right)\approx(0.6398,14.2107).\]对两个端点取平方根，标准差区间约为 $(0.7999,3.7697)$。方差与标准差区间不可混用；原解中的一个标准差下端点数值已校订。\end{solve}

\begin{exercise}[已知标准差的置信水平比较]\label{q:M-P10-1}
样本容量 $30$，总体为 $N(\theta,9^2)$，样本均值为 $15$。求 $95\%$ 和 $70\%$ 均值置信区间。如果样本容量增至 $90$ 而其他条件不变，区间宽度如何变化？\par\smallskip {\small 来源：MIT 习题 10，题 1。}
\end{exercise}

\begin{solve}\phantomsection\label{a:M-P10-1}标准误为 $9/\sqrt{30}$。使用下侧分位数，两个区间分别为\[15\pm z_{0.975}\frac9{\sqrt{30}}=15\pm3.2205,\]\[15\pm z_{0.85}\frac9{\sqrt{30}}=15\pm1.7030.\]更高覆盖要求更宽区间。样本容量从 $30$ 变为 $90$ 时，宽度按 $1/\sqrt n$ 缩放，成为原来的 $1/\sqrt3$。样本均值只影响区间中心，不影响已知标准差模型中的宽度。\end{solve}

\begin{exercise}[灌装量的正态均值区间]\label{q:M-P10-2}
九个灌装量观测为 $352,351,361,353,352,358,360,358,359$。假设观测独立且总体正态。分别在总体标准差已知为 $3$、总体标准差未知的情况下，求总体均值的 $95\%$ 与 $98\%$ 置信区间。\par\smallskip {\small 来源：MIT 习题 10，题 2（校订舍入）。}
\end{exercise}

\begin{solve}\phantomsection\label{a:M-P10-2}计算得 $\overline x=356$，离差平方和 $124$，所以 $s^2=15.5$。已知标准差时标准误为 $1$，故两个区间为 $356\pm1.9600$、$356\pm2.3263$。未知标准差时标准误为 $\sqrt{15.5}/3$，自由度为 $8$，故\[I_{0.95}=356\pm t_{8;0.975}\frac{\sqrt{15.5}}3,\]\[I_{0.98}=356\pm t_{8;0.99}\frac{\sqrt{15.5}}3.\]分别约为 $356\pm3.0263$ 与 $356\pm3.8011$。这里未知标准差版本较宽，同时受样本标准差大于 $3$ 和 $t$ 分位数大于正态分位数影响。\end{solve}

\begin{exercise}[正态方差的区间]\label{q:M-P10-3}
十二个观测为 $6,6.4,7,5.8,6,5.8$；$5.9,6.7,6.1,6.5,6.3,5.8$。假定总体正态，求方差的 $95\%$ 置信区间，并解释区间为何不关于样本方差对称。\par\smallskip {\small 来源：MIT 习题 10，题 3。}
\end{exercise}

\begin{solve}\phantomsection\label{a:M-P10-3}样本均值约为 $6.1917$，无偏方差 $s^2\approx0.1535606$。因为 $11S^2/\sigma^2\sim\chi^2_{11}$，故\[P\left(\chi^2_{11;0.025}\le\frac{11S^2}{\sigma^2}\le\chi^2_{11;0.975}\right)=0.95.\]对正量求倒数并反向排列，得到方差区间\[\left(\frac{11s^2}{\chi^2_{11;0.975}},\frac{11s^2}{\chi^2_{11;0.025}}\right)\approx(0.07706,0.44268).\]卡方分布偏斜，且反演涉及倒数，所以区间通常不以 $s^2$ 为对称中心。\end{solve}

\begin{exercise}[身高区间与宽度要求]\label{q:M-P11-1}
正态总体的 $20$ 个独立身高观测，样本均值 $69.55$ 英寸，样本方差 $14.26$ 平方英寸。求未知方差时的 $90\%$ 均值区间。若另有可靠信息表明总体标准差为 $3.77$ 英寸，求相应区间，并求使区间总宽度不超过 $1$ 英寸所需样本量。\par\smallskip {\small 来源：MIT 习题 11，题 1。}
\end{exercise}

\begin{solve}\phantomsection\label{a:M-P11-1}未知方差时使用 $t_{19;0.95}=1.7291$，得到\[69.55\pm1.7291\sqrt{14.26/20}\approx(68.09,71.01).\]已知标准差时使用 $z_{0.95}=1.6449$，得到\[69.55\pm1.6449(3.77/\sqrt{20})\approx(68.16,70.94).\]总宽度为 $2z_{0.95}3.77/\sqrt n$，令其不超过 $1$，得到 $n\ge(2z_{0.95}3.77)^2\approx153.81$，所以至少需要 $154$ 个观测。这里“宽度”是两个端点之差，不是误差半径。\end{solve}

\begin{exercise}[顺序统计量构造端点区间]\label{q:M-P11-2}
从 $U(0,a)$ 中独立抽样 $n\ge2$ 次，记第二小的观测为 $Y_{(2)}$。已知 $Y_{(2)}/a\sim\operatorname{Beta}(2,n-1)$。利用该枢轴构造 $a$ 的 $95\%$ 置信区间；当 $n=9,y_{(2)}=1.5$ 时给出数值。\par\smallskip {\small 来源：MIT 习题 11，题 2。}
\end{exercise}

\begin{solve}\phantomsection\label{a:M-P11-2}令 $b_L,b_U$ 分别为 $\operatorname{Beta}(2,n-1)$ 的 $0.025,0.975$ 下侧分位数。则\[P\left(b_L\le\frac{Y_{(2)}}a\le b_U\right)=0.95.\]因所有量为正，反演得\[a\in\left[\frac{Y_{(2)}}{b_U},\frac{Y_{(2)}}{b_L}\right].\]当 $n=9$ 时，$b_L\approx0.02814,b_U\approx0.48250$，故区间约为 $[3.11,53.30]$。区间很宽，说明只使用第二小值时，上端点信息有限。也可以与必然约束 $a\ge X_{(n)}$ 相交而保持覆盖率，因为真实参数以概率一满足该约束。\end{solve}

\begin{exercise}[比例误差与样本量]\label{q:M-C22-3}
利用独立伯努利样本估计未知比例 $p$，要求正态近似的 $95\%$ 误差半径不超过 $0.01$。用 $p(1-p)\le1/4$ 规划样本量；改为 $80\%$ 时如何变化？若样本容量固定为 $900$，给出两种置信水平的保守误差半径。\par\smallskip {\small 来源：MIT 第 22 讲，题 3（采用精确正态分位数）。}
\end{exercise}

\begin{solve}\phantomsection\label{a:M-C22-3}正态近似误差半径上界为 $z_{1-\alpha/2}/(2\sqrt n)$。因此 $95\%$ 水平要求\[n\ge\frac{z_{0.975}^2}{4(0.01)^2},\]向上取整得 $9604$；$80\%$ 水平使用 $z_{0.9}=1.28155$，得到 $4106$。若 $n=900$，两个半径分别约为 $0.03267$、$0.02136$。用 $z\approx2$ 的简便规则会得到 $10000$，这与精确正态分位数的规划值不同。这里仍是近似区间规划，尤其接近边界比例时应另查有限样本覆盖。\end{solve}

\begin{exercise}[指数率的精确区间]\label{q:B13.1}
指数样本总和为 $T$，容量 $n$。利用 $2\lambda T\sim\chi^2_{2n}$ 构造率 $\lambda$ 与平均寿命 $1/\lambda$ 的 $1-\alpha$ 置信区间。\par\smallskip {\small 来源：编订补充。}
\end{exercise}

\begin{solve}\phantomsection\label{a:B13.1}枢轴分布不含未知参数，因此\[P\left(\chi^2_{2n;\alpha/2}\le2\lambda T\le\chi^2_{2n;1-\alpha/2}\right)=1-\alpha.\]在 $T>0$ 时除以 $2T$，率区间为\[\left[\frac{\chi^2_{2n;\alpha/2}}{2T},\frac{\chi^2_{2n;1-\alpha/2}}{2T}\right].\]对端点取倒数并交换顺序，平均寿命区间为\[\left[\frac{2T}{\chi^2_{2n;1-\alpha/2}},\frac{2T}{\chi^2_{2n;\alpha/2}}\right].\]这里是指数抽样假设下的有限样本精确覆盖，不需要中心极限定理。\end{solve}

\begin{exercise}[均匀端点的单侧区间]\label{q:B13.2}
设 $M=X_{(n)}$ 来自 $U(0,\theta)$ 样本。利用 $M/\theta$ 的分布构造形如 $[M,cM]$ 的精确 $1-\alpha$ 置信区间。\par\smallskip {\small 来源：编订补充。}
\end{exercise}

\begin{solve}\phantomsection\label{a:B13.2}对 $0<u<1$，$P(M/\theta\le u)=u^n$。由于 $M\le\theta$ 以概率一成立，区间包含参数只需满足 $\theta\le cM$，即 $M/\theta\ge1/c$。其概率为 $1-c^{-n}$。令它等于 $1-\alpha$，得到 $c=\alpha^{-1/n}$，所以\[[M,M/\alpha^{1/n}].\]此区间不以某个点估计为中心，而是利用总体支撑给出的确定性下界。构造置信区间不必总是“估计值加减标准误”。\end{solve}

\begin{exercise}[两个均值差的合并方差区间]\label{q:B13.3}
两组独立正态样本容量为 $n,m$，均值为 $\mu_1,\mu_2$，共同未知方差为 $\sigma^2$。构造均值差的 $1-\alpha$ 置信区间，并指出共同方差假设的位置。\par\smallskip {\small 来源：编订补充。}
\end{exercise}

\begin{solve}\phantomsection\label{a:B13.3}定义\[S_p^2=\frac{(n-1)S_1^2+(m-1)S_2^2}{n+m-2}.\]由独立卡方变量相加，$(n+m-2)S_p^2/\sigma^2\sim\chi^2_{n+m-2}$；它与均值差独立。因此标准化统计量\[\frac{(\overline X-\overline Y)-(\mu_1-\mu_2)}{S_p\sqrt{1/n+1/m}}\sim t_{n+m-2}.\]反演得到均值差区间为 $(\overline X-\overline Y)\pm t_{n+m-2;1-\alpha/2}S_p\sqrt{1/n+1/m}$。合并平方和共用同一个真实方差，正是共同方差假设的用途。\end{solve}

\begin{exercise}[韦尔奇近似区间]\label{q:B13.4}
两组独立正态样本的方差可能不同。写出均值差的韦尔奇区间及近似自由度。用 $n=10,m=15,s_1^2=4,s_2^2=9,\overline x-\overline y=2$ 给出 $95\%$ 区间的计算式。\par\smallskip {\small 来源：编订补充。}
\end{exercise}

\begin{solve}\phantomsection\label{a:B13.4}差的估计标准误为 $SE=\sqrt{s_1^2/n+s_2^2/m}$，萨特思韦特近似自由度为\[\nu=\frac{(s_1^2/n+s_2^2/m)^2}{(s_1^2/n)^2/(n-1)+(s_2^2/m)^2/(m-1)}.\]本题 $SE=1$，$\nu=1/(0.16/9+0.36/14)\approx22.991$，区间为 $2\pm t_{\nu;0.975}$，约 $(-0.069,4.069)$。它通常是近似区间，不应标成自由度整数的精确 $t$ 结论。相比先做方差检验再决定方法，事先合理选定允许异方差的方法更容易保持清晰的推断规则。\end{solve}

\begin{exercise}[配对差的区间]\label{q:B13.5}
对五名受试者作配对测量，所得差值依次为 $2,1,3,0,4$。假定不同受试者相互独立，差值来自正态总体。求平均差的 $95\%$ 置信区间；解释为何不应把前后数据当作两组独立样本。\par\smallskip {\small 来源：编订补充。}
\end{exercise}

\begin{solve}\phantomsection\label{a:B13.5}差值均值为 $2$，离差平方和为 $10$，所以 $s_D^2=10/4=2.5$。标准误为 $\sqrt{2.5/5}=\sqrt{0.5}$，自由度为 $4$，区间为\[2\pm t_{4;0.975}\sqrt{0.5}\approx2\pm1.9632.\]同一受试者前后观测通常相关。差值方差包含 $\operatorname{Var}(X)+\operatorname{Var}(Y)-2\operatorname{Cov}(X,Y)$；把前后当作独立会遗漏协方差。配对分析所需正态假设针对差值，而不是仅分别针对前后边缘分布。\end{solve}

\begin{exercise}[置信区间与后验区间不能互换]\label{q:M-P10-5}
参数 $\mu$ 只取 $0,1,2,3$，先验概率依次为 $0.94,0.02,0.02,0.02$。条件于 $\mu$，$64$ 个观测独立来自 $N(\mu,25)$，得到 $\overline x=1.6$。构造通常的 $95\%$ 正态均值区间，并计算该已实现区间在指定先验与后验下包含参数的概率。\par\smallskip {\small 来源：MIT 习题 10，题 5（选编）。}
\end{exercise}

\begin{solve}\phantomsection\label{a:M-P10-5}标准误为 $5/8=0.625$，区间为 $1.6\pm1.959964(0.625)$，即约 $(0.3750,2.8250)$。它包含候选值 $1,2$，所以先验概率和为 $0.04$。后验权重为\[w_j=\pi_j\exp\left[-\frac{(1.6-j)^2}{2(0.625)^2}\right],\quad j=0,1,2,3.\]归一化后约为 $(0.537,0.191,0.247,0.025)$，所以区间的后验概率约为 $0.438$，不是 $0.95$。频率覆盖是对固定真实参数反复抽样时的随机区间而言；后验概率是在给定数据和先验以后对参数的不确定性而言。\end{solve}

\begin{exercise}[自助法百分位区间]\label{q:B13.6}
给定来自同一总体的独立样本，容量为 $n$。希望估计总体中位数。说明非参数自助法构造百分位区间的步骤、随机性来源以及一种可能失效的情形。\par\smallskip {\small 来源：编订补充。}
\end{exercise}

\begin{solve}\phantomsection\label{a:B13.6}把经验分布视为总体近似。每次从已有 $n$ 个观测中有放回抽取 $n$ 个值，计算自助样本中位数，重复 $B$ 次得到 $T_1^*,\ldots,T_B^*$。以其经验 $\alpha/2$ 与 $1-\alpha/2$ 分位数作为百分位区间端点。原始抽样随机性产生数据，自助重抽样随机性近似统计量的抽样波动，有限 $B$ 还引入蒙特卡洛误差。若总体在中位数附近有正且连续的密度，通常具有良好渐近基础；若分位数不唯一、存在大原子或样本极小，普通百分位法可能覆盖不佳。增大 $B$ 不能弥补原始样本代表性不足。\end{solve}

\clearpage\section{假设检验与功效}

\begin{exercise}[硬币的精确双侧检验]\label{q:M-P8-1ab}
独立旋转硬币 $250$ 次，观察到 $140$ 次正面。检验 $H_0:p=0.5$ 对 $H_1:p\ne0.5$。计算精确双侧 $p$ 值，并判断在 $0.10$ 与 $0.05$ 水平是否拒绝。比较不作连续性修正的正态近似。\par\smallskip {\small 来源：MIT 习题 8，题 1(a,b)。}
\end{exercise}

\begin{solve}\phantomsection\label{a:M-P8-1ab}在零假设下 $X\sim\operatorname{Bin}(250,0.5)$，分布关于 $125$ 对称。按偏离 $125$ 的绝对距离度量极端程度，\[p_{\rm exact}=P(X\le110)+P(X\ge140)=2P(X\ge140)\approx0.06642.\]因此在 $0.10$ 水平拒绝，在 $0.05$ 水平不拒绝。未修正正态近似使用 $z=(140-125)/\sqrt{62.5}$，给出约 $0.05778$；若右尾阈值改用 $139.5$，通常更接近精确概率。离散分布下不同双侧定义可能不同，本题因对称性定义很自然。\end{solve}

\begin{exercise}[离散拒绝域与功效]\label{q:M-P8-1cd}
继续考虑 $X\sim\operatorname{Bin}(250,p)$。构造水平不超过 $0.01$ 的对称双侧拒绝域；再对水平不超过 $0.05$ 的对称拒绝域，分别计算 $p=0.55$ 和 $p=0.60$ 时的功效。\par\smallskip {\small 来源：MIT 习题 8，题 1(c,d)。}
\end{exercise}

\begin{solve}\phantomsection\label{a:M-P8-1cd}$0.01$ 水平的拒绝域取 $X\le104$ 或 $X\ge146$，每个尾的零假设概率不超过 $0.005$。$0.05$ 水平对应 $X\le109$ 或 $X\ge141$。因此在任意给定 $p$ 下，功效为\[\pi(p)=\sum_{k=0}^{109}\binom{250}{k}p^k(1-p)^{250-k}+\sum_{k=141}^{250}\binom{250}{k}p^k(1-p)^{250-k}.\]计算得 $\pi(0.55)\approx0.35237$，$\pi(0.60)\approx0.88963$。效应远离零假设时，分布落入拒绝域的概率增大。离散检验的实际大小通常小于指定水平，不能总写成严格等号。\end{solve}

\begin{exercise}[第一类与第二类错误]\label{q:M-P8-2}
有 $140$ 个真实阴性对象和 $140$ 个真实阳性对象。规则将真实阴性中的 $9$ 个判为阳性，将真实阳性中的 $15$ 个判为阴性。以“对象真实阴性”为零假设，估计第一类错误率、第二类错误率及功效，并解释为什么不能直接得到阳性结果为假的概率。\par\smallskip {\small 来源：MIT 习题 8，题 2（抽象分类表述）。}
\end{exercise}

\begin{solve}\phantomsection\label{a:M-P8-2}第一类错误是对真实零假设作拒绝，因此估计错误率为 $9/140\approx0.06429$。第二类错误是阳性对象未被识别，其估计概率为 $15/140\approx0.10714$，功效为 $125/140\approx0.89286$。阳性结果为假的概率是逆向条件概率，需结合目标人群中的真实阳性比例。此处两类各抽取 $140$ 个是实验设计，不保证目标人群也各占一半。错误率的条件方向不可交换。\end{solve}

\begin{exercise}[复合零假设与功效规划]\label{q:M-P8-3}
独立仪器读数 $X_i\sim N(\mu,5^2)$，检验 $H_0:\mu\le40$ 对 $H_1:\mu>40$。取 $n=3$，要求所有零假设参数下误拒概率不超过 $0.04$，求样本均值阈值及 $\mu=45$ 时的功效。再求使该功效至少为 $0.9$ 的最小样本量。\par\smallskip {\small 来源：MIT 习题 8，题 3（单位改称测量单位）。}
\end{exercise}

\begin{solve}\phantomsection\label{a:M-P8-3}对固定右尾拒绝域，误拒概率随 $\mu$ 增大，故零假设中的最不利值为 $40$。阈值为\[c=40+z_{0.96}\frac5{\sqrt3}\approx45.0538.\]在 $\mu=45$ 时，功效为 $1-\Phi((c-45)/(5/\sqrt3))\approx0.4926$。一般样本量下，功效为 $\Phi(\sqrt n-z_{0.96})$；要求其至少 $0.9$，即\[\sqrt n\ge z_{0.96}+z_{0.9}\approx3.0322.\]所以最小整数样本量为 $10$。复合零假设要求控制其中所有参数值下的错误率，不能任取某个内部点代替最不利点。\end{solve}

\begin{exercise}[直接计算两类错误]\label{q:M-P8-4}
只观察一次 $X\sim U(0,\theta)$。检验 $H_0:\theta=2$ 对 $H_1:\theta\ne2$，规定当 $X\le0.1$ 或 $X\ge1.9$ 时拒绝。求第一类错误概率，以及 $\theta=2.5$ 时的第二类错误概率；写出一般功效函数。\par\smallskip {\small 来源：MIT 习题 8，题 4。}
\end{exercise}

\begin{solve}\phantomsection\label{a:M-P8-4}零假设下拒绝域在支撑内总长度为 $0.2$，所以第一类错误概率为 $0.2/2=0.1$。当 $\theta=2.5$ 时，不拒绝区间长度为 $1.8$，第二类错误概率为 $1.8/2.5=0.72$。一般功效为\[\pi(\theta)=\begin{cases}1,&0<\theta\le0.1,\\0.1/\theta,&0.1<\theta\le1.9,\\1-1.8/\theta,&\theta>1.9.\end{cases}\]这是一条关于真实参数的函数，而不是单一数字。功效在某些接近零假设的备择处也可能偏低。\end{solve}

\begin{exercise}[单侧与双侧 t 检验]\label{q:M-P9-1}
正态样本容量 $16$，样本均值 $11$，无偏样本方差 $4$。检验总体均值为 $10$，显著性水平 $0.05$。分别考虑双侧备择与预先指定的右侧备择，计算统计量、$p$ 值和结论。\par\smallskip {\small 来源：MIT 习题 9，题 1。}
\end{exercise}

\begin{solve}\phantomsection\label{a:M-P9-1}零假设下\[T=\frac{\overline X-10}{S/\sqrt{16}}\sim t_{15},\qquad t_{\rm obs}=2.\]双侧 $p$ 值为 $2[1-F_{t_{15}}(2)]\approx0.06395$，故不拒绝；右侧 $p$ 值为 $1-F_{t_{15}}(2)\approx0.03197$，故拒绝。两种检验回答不同问题、采用不同拒绝域。不能看到样本均值偏高后才临时把双侧检验改成右侧检验并仍声称维持原先的检验水平。\end{solve}

\begin{exercise}[停止规则改变抽样分布]\label{q:M-P9-2}
观察到 $12$ 次抛掷，其中 $3$ 次正面、$9$ 次反面，最后一次为第三个正面。检验硬币 $p=0.5$ 对 $p<0.5$。分别按“固定抛 $12$ 次”和“直到第三个正面才停止”的计划求单侧 $p$ 值；若采用贝塔先验，比较后验。\par\smallskip {\small 来源：MIT 习题 9，题 2。}
\end{exercise}

\begin{solve}\phantomsection\label{a:M-P9-2}固定次数方案以正面数 $H\sim\operatorname{Bin}(12,0.5)$ 为统计量，$p$ 值为 $P(H\le3)\approx0.07300$。停止方案记录第三个正面前反面数 $N$，其负二项概率为 $P(N=k)=\binom{k+2}{2}2^{-(k+3)}$，故\[P(N\ge9)\approx0.03271.\]两个抽样空间不同，频率检验结果可以不同。对于预先规定且不另含参数依赖因素的上述停止规则，观测数据似然都与 $p^3(1-p)^9$ 成比例；$\operatorname{Beta}(a,b)$ 先验均更新为 $\operatorname{Beta}(a+3,b+9)$。这不意味着任意自适应筛选和未报告数据都可忽略。\end{solve}

\begin{exercise}[正态右尾检验的分位数]\label{q:M-E2-II4}
某标准正态右尾检验的水平为 $0.1$，观测统计量为 $z=2$。求临界值、$p$ 值和结论，并说明分布函数与分位数函数分别承担什么作用。\par\smallskip {\small 来源：MIT 第二次考试，第二部分题 4。}
\end{exercise}

\begin{solve}\phantomsection\label{a:M-E2-II4}临界值满足 $P(Z>c)=0.1$，所以 $c=\Phi^{-1}(0.9)\approx1.28155$。观测 $p$ 值为\[p=P(Z\ge2)=1-\Phi(2)\approx0.02275.\]因 $2>c$ 或等价地 $p<0.1$，拒绝零假设。分位数函数由给定尾概率反求横轴阈值；分布函数由给定横轴值计算累积概率。把 $\Phi(0.9)$ 当作临界值，是将两种函数方向弄反。\end{solve}

\begin{exercise}[检验与置信区间对偶]\label{q:M-F-C10}
一个双侧两样本均值检验得到 $p=0.08$。假设置信区间与检验使用同一模型和同一统计量，判断：在 $0.05$ 水平拒绝均值相等；$90\%$ 均值差区间包含零；$95\%$ 区间包含零；参数以 $95\%$ 概率属于已计算的 $95\%$ 区间。\par\smallskip {\small 来源：MIT 期末试卷，概念题 10。}
\end{exercise}

\begin{solve}\phantomsection\label{a:M-F-C10}依次为错、错、对、不能作这种频率学派概率解释。$p=0.08>0.05$，故在 $0.05$ 水平不拒绝，零位于对应 $95\%$ 区间；但 $p<0.10$，所以在 $0.10$ 水平拒绝，零不在 $90\%$ 区间。区间由接受域反演而来时，这种对偶精确成立。最后一句把随机区间的重复抽样覆盖与给定数据后的参数后验概率混在一起，需要额外先验才能作后者陈述。\end{solve}

\begin{exercise}[奈曼—皮尔逊检验]\label{q:B14.1}
独立样本来自 $N(\mu,\sigma^2)$，$\sigma^2$ 已知。检验简单假设 $H_0:\mu=\mu_0$ 对 $H_1:\mu=\mu_1>\mu_0$。利用似然比推导最强水平 $\alpha$ 检验及其功效。\par\smallskip {\small 来源：编订补充。}
\end{exercise}

\begin{solve}\phantomsection\label{a:B14.1}备择与零假设的对数似然比为\[\log\frac{L(\mu_1)}{L(\mu_0)}=\frac{n(\mu_1-\mu_0)}{\sigma^2}\overline X-\frac{n(\mu_1^2-\mu_0^2)}{2\sigma^2}.\]它随 $\overline X$ 严格增加。奈曼—皮尔逊引理说明应在大均值时拒绝，阈值由零假设大小确定为 $c=\mu_0+z_{1-\alpha}\sigma/\sqrt n$。在 $\mu_1$ 下功效为\[1-\Phi\left(z_{1-\alpha}-\frac{(\mu_1-\mu_0)\sqrt n}{\sigma}\right).\]最强性针对指定的简单备择以及所有同水平检验。似然比排序确定拒绝方向，零假设分布确定阈值。\end{solve}

\begin{exercise}[正态均值的似然比检验]\label{q:B14.2}
正态样本方差已知，检验 $H_0:\mu=\mu_0$ 对 $H_1:\mu\ne\mu_0$。求广义似然比 $\Lambda=L(\mu_0)/\sup_\mu L(\mu)$，并说明它如何给出双侧正态检验。\par\smallskip {\small 来源：编订补充。}
\end{exercise}

\begin{solve}\phantomsection\label{a:B14.2}全参数空间中的最大似然在 $\widehat\mu=\overline X$ 达到。利用平方和分解，\[\sum_i(X_i-\mu_0)^2-\sum_i(X_i-\overline X)^2=n(\overline X-\mu_0)^2.\]所以\[-2\log\Lambda=\frac{n(\overline X-\mu_0)^2}{\sigma^2}=Z^2.\]零假设下它精确服从 $\chi^2_1$，等价拒绝域为 $|Z|>z_{1-\alpha/2}$。本例无需调用渐近的威尔克斯定理，因为统计量的有限样本分布已经精确已知。\end{solve}

\begin{exercise}[多重检验与邦费罗尼控制]\label{q:B14.3}
同时进行 $20$ 个检验，每个检验在真实零假设下误拒概率为 $0.05$。若全部零假设真实且检验相互独立，求至少一次误拒的概率。若不保证独立，怎样把至少一次误拒的概率控制在 $0.05$？\par\smallskip {\small 来源：编订补充；主题参照 MIT 习题 8，题 5(b)。}
\end{exercise}

\begin{solve}\phantomsection\label{a:B14.3}独立时无一次误拒的概率为 $0.95^{20}$，所以至少一次误拒的概率为\[1-0.95^{20}\approx0.64151.\]若不保证独立，可令每个检验的显著性水平不超过 $0.05/20=0.0025$。由并事件上界，至少一次误拒的概率不超过各真实零假设误拒概率之和，因而不超过 $0.05$。这个邦费罗尼控制不需要独立性，并且在只有部分零假设真实时仍成立。代价是单个检验的功效可能降低。\end{solve}

\begin{exercise}[不拒绝不等于等效]\label{q:B14.4}
某均值差估计为 $0.2$，标准误为 $1$，大样本双侧检验未拒绝差为零。能否据此认定两个总体等效？若预先规定等效界为 $[-0.5,0.5]$，给出一种合适的判断方法。\par\smallskip {\small 来源：编订补充。}
\end{exercise}

\begin{solve}\phantomsection\label{a:B14.4}未拒绝可能由差确实小，也可能由数据精度不足造成。此例的 $95\%$ 近似区间为 $0.2\pm1.96$，相当宽，包含许多实际重要的差异。等效检验把“不等效”作为零假设，即差不大于 $-0.5$ 或不小于 $0.5$。在双单侧检验中，各单侧使用 $\alpha=0.05$，等价于检查 $90\%$ 双侧区间是否完全包含在 $(-0.5,0.5)$ 内。本例 $90\%$ 区间为 $0.2\pm1.6449$，并未落入等效界，不能认定等效。界限应在观察结果前由问题背景确定。\end{solve}

\clearpage\section{拟合优度与列联表}

\begin{exercise}[单个正态总体的方差检验]\label{q:M-P9-3}
十二个独立正态观测为 \[\begin{gathered}1.76,-2.28,-0.56,1.46,0.59,1.26,\\-1.94,-0.79,-0.86,-1.41,2.07,1.30.\end{gathered}\]。在 $0.05$ 水平检验 $H_0:\sigma^2=1$ 对 $H_1:\sigma^2>1$。\par\smallskip {\small 来源：MIT 习题 9，题 3。}
\end{exercise}

\begin{solve}\phantomsection\label{a:M-P9-3}样本均值为 $0.05$，离差平方和为 $25.7372$，故 $s^2\approx2.33975$。在零假设下\[Q=\frac{(n-1)S^2}{1}\sim\chi^2_{11}.\]观察值 $q=25.7372$，右尾 $p$ 值约为 $0.0071$，小于 $0.05$，所以拒绝零假设。均值是未知并由样本估计的，因此自由度为 $11$。如果均值已知而用相对已知均值的平方和，统计量和自由度都不同。\end{solve}

\begin{exercise}[两个总体的类别比例比较]\label{q:M-E2-II5}
两组各 $50$ 个独立对象，第一组两类别计数为 $(15,35)$，第二组为 $(5,45)$。利用未作连续性修正的皮尔逊卡方检验，在 $0.05$ 水平检验类别比例是否相同。说明结论的范围。\par\smallskip {\small 来源：MIT 第二次考试，第二部分题 5（教学数据）。}
\end{exercise}

\begin{solve}\phantomsection\label{a:M-E2-II5}合并后的类别比例估计为 $20/100=0.2$ 与 $0.8$，所以每组期望计数均为 $(10,40)$。皮尔逊统计量为\[Q=\frac{(15-10)^2}{10}+\frac{(35-40)^2}{40}+\frac{(5-10)^2}{10}+\frac{(45-40)^2}{40}=6.25.\]自由度为 $(2-1)(2-1)=1$，渐近 $p$ 值约为 $0.01242$，故拒绝比例相同的零假设。结论是模型和抽样条件下的比例差异证据；不能仅凭观察列联表推出一个类别对另一个变量有因果作用。\end{solve}

\begin{exercise}[多项分类的拟合优度]\label{q:B15.1}
公平六面骰子独立掷 $120$ 次，各面计数为 $18,22,19,21,17,23$。用皮尔逊拟合优度检验公平性，写出统计量、自由度和结论。\par\smallskip {\small 来源：编订补充。}
\end{exercise}

\begin{solve}\phantomsection\label{a:B15.1}零假设指定每类概率均为 $1/6$，所以期望计数都是 $20$。统计量为\[Q=\sum_{j=1}^6\frac{(O_j-20)^2}{20}=\frac{4+4+1+1+9+9}{20}=1.4.\]六个计数和受总数约束，未从数据估计额外概率参数，所以自由度为 $6-1=5$。$\chi^2_{5;0.95}\approx11.0705$，观察值远小于临界值，不拒绝公平性。这个结果不证明骰子绝对公平，只说明当前样本未给出足够的偏离证据。\end{solve}

\begin{exercise}[估计参数后的自由度]\label{q:B15.2}
把来自一个泊松分布的独立计数按预先规定的 $k$ 个互斥类别分组，用样本估计未知率参数，再做拟合优度检验。通常使用什么自由度？说明使用该结论需注意的条件。\par\smallskip {\small 来源：编订补充。}
\end{exercise}

\begin{solve}\phantomsection\label{a:B15.2}若零假设中一个自由参数通过适当的同一分组多项似然方法估计，且满足常规正则性、各类别概率为正并有足够期望计数，则皮尔逊统计量渐近自由度为 $k-1-1=k-2$。其中一个自由度来自总概率为一，另一个来自估计参数。若用未分组信息估计参数而又对分组计数作检验，精确渐近分布未必简单地等于该卡方分布，不能无条件扣一。分组应覆盖全部支撑，尾部低期望类别可合并；根据观察偏差反复调整分组会影响原定检验。\end{solve}

\begin{exercise}[列联表独立性检验]\label{q:B15.3}
某 $2\times3$ 表的两行计数为 $(20,30,50)$ 与 $(30,30,40)$。假定样本可按多项模型处理，检验行变量与列变量独立。求期望计数、统计量和自由度。\par\smallskip {\small 来源：编订补充。}
\end{exercise}

\begin{solve}\phantomsection\label{a:B15.3}总数为 $200$，行和均为 $100$，列和为 $50,60,90$。在独立模型下，期望计数为行和乘列和除以总数，两行均为 $(25,30,45)$。故\[Q=2\left(\frac{25}{25}+0+\frac{25}{45}\right)=\frac{28}{9}\approx3.1111.\]自由度为 $(2-1)(3-1)=2$，右尾概率为 $e^{-Q/2}\approx0.2111$。在 $0.05$ 水平不拒绝独立性。独立模型含有行比例与列比例，约束扣除后才得到上述自由度。\end{solve}

\begin{exercise}[小样本的费希尔精确检验]\label{q:B15.4}
一个 $2\times2$ 表的第一行为 $(4,1)$，第二行为 $(1,4)$。在固定行列边际和的条件下，以第一行第一格 $A$ 为统计量，求右侧精确 $p$ 值和按概率不大于观察表概率定义的双侧 $p$ 值。\par\smallskip {\small 来源：编订补充。}
\end{exercise}

\begin{solve}\phantomsection\label{a:B15.4}两个行和与列和均为 $5$，总数为 $10$。在零假设下，条件分布为超几何分布：\[P(A=a)=\frac{\binom5a\binom5{5-a}}{\binom{10}5},\quad a=0,\ldots,5.\]观测 $A=4$，右侧 $p$ 值为 $(25+1)/252=13/126\approx0.10317$。双侧按题设概率排序纳入 $a=0,1,4,5$，概率总和为 $52/252=13/63\approx0.20635$。小期望计数时，不能把大样本卡方近似当作精确分布；双侧精确检验还须说明采用的极端程度定义。\end{solve}

\begin{exercise}[两个正态方差的 F 检验]\label{q:M-P9-5}
两组独立正态样本各有 $20$ 个观测，样本方差分别为 $1.1302$ 与 $1.1647$。检验 $H_0:\sigma_1^2=\sigma_2^2$ 对双侧备择，计算 $F$ 统计量、双侧 $p$ 值并构造方差比的 $95\%$ 区间。\par\smallskip {\small 来源：MIT 习题 9，题 5（给出原数据的汇总量）。}
\end{exercise}

\begin{solve}\phantomsection\label{a:M-P9-5}统计量取固定顺序的比值 $F=S_1^2/S_2^2$，零假设下服从 $F_{19,19}$；观测值约为 $0.97038$。按双倍较小尾概率定义，\[p=2\min\{F_{19,19}^{\rm cdf}(f),1-F_{19,19}^{\rm cdf}(f)\}\approx0.9484.\]不拒绝等方差。方差比 $r=\sigma_1^2/\sigma_2^2$ 的枢轴为 $F/r$，故区间为\[\left[\frac{f}{F_{19,19;0.975}},\frac{f}{F_{19,19;0.025}}\right]\approx[0.3841,2.4516].\]区间较宽，说明“不显著”也可能伴随相当大的不确定性；此检验对偏离正态较敏感。\end{solve}

\clearpage\section{回归与方差分析}

\begin{exercise}[三个点的最小二乘直线]\label{q:M-C26-1a}
给定三个点 $(1,3)$、$(2,1)$、$(4,4)$。写出含截距的直线模型，最小化残差平方和求拟合直线，并验证拟合线经过样本均值点。\par\smallskip {\small 来源：MIT 第 26 讲，题 1(a) 与题 2(a—c)。}
\end{exercise}

\begin{solve}\phantomsection\label{a:M-C26-1a}模型为 $y_i=a x_i+b+e_i$。最小化\[Q(a,b)=(3-a-b)^2+(1-2a-b)^2+(4-4a-b)^2.\]对两个参数求偏导得到正规方程\[21a+7b=21,\qquad7a+3b=8.\]解得 $a=1/2,b=3/2$，所以 $\widehat y=x/2+3/2$。$\overline x=7/3,\overline y=8/3$，代入直线确有 $\widehat y(7/3)=8/3$。设计矩阵两列线性无关，平方和的 Hessian 正定，因此该驻点是唯一全局最小值。\end{solve}

\begin{exercise}[二次多项式仍是线性回归]\label{q:M-C26-1b}
对点 $(1,3)$、$(2,1)$、$(4,4)$ 拟合 $y=ax^2+bx+c$。给出设计矩阵和解，并解释为什么二次、三次曲线也属于线性回归。三个观测拟合三个参数有什么统计限制？\par\smallskip {\small 来源：MIT 第 26 讲，题 1(b,d)（选编）。}
\end{exercise}

\begin{solve}\phantomsection\label{a:M-C26-1b}设计矩阵的三行为 $(1,1,1),(4,2,1),(16,4,1)$，响应向量为 $(3,1,4)^\mathsf T$。三个横坐标不同，矩阵可逆；解得\[a=7/6,\qquad b=-11/2,\qquad c=22/3.\]曲线经过所有三个点，残差平方和为零。所谓线性回归是指对未知参数线性，而不要求对自变量 $x$ 线性；三次模型只需增加 $x^3$ 这一列。但三个观测估计三个参数时残差自由度为零，不能再用残差估计误差方差，也不能据零残差断言泛化预测完美。\end{solve}

\begin{exercise}[对数尺度的指数拟合]\label{q:M-C26-1c}
仍使用三个正响应点 $(1,3)$、$(2,1)$、$(4,4)$，在对数尺度拟合 $\log y=ax+b$。写出最小化目标和参数公式，并解释它与直接最小化原尺度误差 $\sum_i(y_i-e^{ax_i+b})^2$ 是否相同。\par\smallskip {\small 来源：MIT 第 26 讲，题 1(c)（补充模型区别）。}
\end{exercise}

\begin{solve}\phantomsection\label{a:M-C26-1c}令 $z_i=\log y_i$。对数尺度的目标是\[Q(a,b)=\sum_i(z_i-ax_i-b)^2,\]解为 $a=\sum_i(x_i-\overline x)(z_i-\overline z)/\sum_i(x_i-\overline x)^2$，$b=\overline z-a\overline x$。代入本题得 $a\approx0.18122,b\approx0.40547$。两种目标通常给出不同拟合，因为对数尺度侧重相对误差，原尺度侧重绝对误差。若模型为 $\log Y=ax+b+\varepsilon$、$\varepsilon\sim N(0,\sigma^2)$，则 $e^{ax+b}$ 是条件中位数，条件均值还需乘 $e^{\sigma^2/2}$。\end{solve}

\begin{exercise}[正态误差下的最小二乘与似然]\label{q:M-C26-2d}
在固定设计模型 $Y_i=ax_i+b+\varepsilon_i$ 中，误差独立同分布于 $N(0,\sigma^2)$。证明 $a,b$ 的极大似然估计等于最小二乘估计，并求误差方差的极大似然估计。\par\smallskip {\small 来源：MIT 第 26 讲，题 2(d)。}
\end{exercise}

\begin{solve}\phantomsection\label{a:M-C26-2d}对数似然为\[\ell(a,b,\sigma^2)=-\frac n2\log(2\pi\sigma^2)-\frac{1}{2\sigma^2}\sum_i(y_i-ax_i-b)^2.\]固定正方差时，最大化似然等价于最小化残差平方和 $SSE$。代回最小二乘参数，再对方差求导，得到 $\widehat\sigma^2_{\rm ML}=SSE/n$。当设计秩为二且 $n>2$ 时，无偏估计则为 $SSE/(n-2)$。若 $SSE=0$，在严格正方差空间可能不存在有限极大值。只说“独立同方差”不足以推出似然等价，还需要正态误差或另行指定的误差分布。\end{solve}

\begin{exercise}[正规方程与残差正交]\label{q:B16.1}
线性模型写成 $Y=X\beta+\varepsilon$，其中固定设计矩阵 $X$ 为 $n\times p$ 满列秩矩阵。推导最小二乘解，证明残差与设计矩阵各列正交，并说明包含截距时残差和为零。\par\smallskip {\small 来源：编订补充。}
\end{exercise}

\begin{solve}\phantomsection\label{a:B16.1}平方和为 $(Y-X\beta)^\mathsf T(Y-X\beta)$，梯度为 $-2X^\mathsf T(Y-X\beta)$。令其为零得到正规方程 $X^\mathsf TX\widehat\beta=X^\mathsf TY$。满列秩保证 $X^\mathsf TX$ 正定可逆，故\[\widehat\beta=(X^\mathsf TX)^{-1}X^\mathsf TY.\]残差 $e=Y-X\widehat\beta$ 满足 $X^\mathsf Te=0$。若全一向量是 $X$ 的一列，则 $\sum_ie_i=0$，于是样本响应均值等于拟合值均值。没有截距项时，残差和为零和通常的中心化平方和分解都不一定成立。\end{solve}

\begin{exercise}[高斯—马尔可夫最优性]\label{q:B16.2}
在固定满列秩设计模型中，假设 $E\varepsilon=0$、$\operatorname{Cov}(\varepsilon)=\sigma^2I$。证明最小二乘估计在线性无偏估计中协方差最小，不要求误差正态。\par\smallskip {\small 来源：编订补充。}
\end{exercise}

\begin{solve}\phantomsection\label{a:B16.2}任意线性无偏估计可写为 $AY$，其中 $AX=I_p$。记 $A_0=(X^\mathsf TX)^{-1}X^\mathsf T$，令 $D=A-A_0$，则 $DX=0$。因为 $A_0D^\mathsf T=(X^\mathsf TX)^{-1}X^\mathsf TD^\mathsf T=0$，\[\operatorname{Cov}(AY)=\sigma^2A_0A_0^\mathsf T+\sigma^2DD^\mathsf T.\]第二项半正定，所以任何线性组合的方差都不小于最小二乘对应方差。这是在线性无偏估计类中的最优性，不是对所有可能的非线性或有偏估计的统一最优性。证明只用一二阶矩条件。\end{solve}

\begin{exercise}[斜率估计的方差与区间]\label{q:B16.3}
在含截距的一元正态线性回归中，记 $S_{xx}=\sum_i(x_i-\overline x)^2>0$。推导斜率估计的均值、方差及置信区间，并说明扩大自变量跨度的作用。\par\smallskip {\small 来源：编订补充。}
\end{exercise}

\begin{solve}\phantomsection\label{a:B16.3}斜率估计为 $\widehat a=\sum_i(x_i-\overline x)Y_i/S_{xx}$。代入模型，利用中心化和为零，得 $E\widehat a=a$；独立同方差误差给出\[\operatorname{Var}(\widehat a)=\frac{\sigma^2}{S_{xx}}.\]正态假设使它为正态变量，且与 $s^2=SSE/(n-2)$ 独立，因此\[\frac{\widehat a-a}{s/\sqrt{S_{xx}}}\sim t_{n-2}.\]区间为 $\widehat a\pm t_{n-2;1-\alpha/2}s/\sqrt{S_{xx}}$。在模型可靠和误差方差不变时，增大 $S_{xx}$ 可提高斜率精度，但超出已观察范围的外推还面临模型形式风险。\end{solve}

\begin{exercise}[均值响应区间与预测区间]\label{q:B16.4}
在一元正态线性回归中，对固定 $x_0$，分别构造均值响应 $a x_0+b$ 的区间和一个独立新响应 $Y_0$ 的预测区间。解释额外的方差项。\par\smallskip {\small 来源：编订补充。}
\end{exercise}

\begin{solve}\phantomsection\label{a:B16.4}令 $\widehat y_0=\widehat a x_0+\widehat b$，其方差为\[\sigma^2h_0,\qquad h_0=\frac1n+\frac{(x_0-\overline x)^2}{S_{xx}}.\]均值响应区间为 $\widehat y_0\pm t_{n-2;1-\alpha/2}s\sqrt{h_0}$。独立新响应还含新误差 $\varepsilon_0$，所以预测误差方差为 $\sigma^2(1+h_0)$，预测区间为\[\widehat y_0\pm t_{n-2;1-\alpha/2}s\sqrt{1+h_0}.\]额外的 $1$ 对应个体噪声。即使样本无限多，个体预测区间通常也不会收缩为零，而均值响应估计的不确定性可随数据增加而缩小。\end{solve}

\begin{exercise}[单因素方差分析]\label{q:B16.5}
三组独立正态观测分别为 $(2,3,4)$、$(4,5,6)$、$(6,7,8)$，假定共同方差。用单因素方差分析在 $0.05$ 水平检验三个总体均值相等。\par\smallskip {\small 来源：编订补充。}
\end{exercise}

\begin{solve}\phantomsection\label{a:B16.5}三组均值为 $3,5,7$，总均值为 $5$。组间平方和\[SS_B=3[(3-5)^2+0+(7-5)^2]=24,\]自由度为 $2$；各组内平方和都是 $2$，故 $SS_W=6$，自由度为 $9-3=6$。统计量为\[F=\frac{SS_B/2}{SS_W/6}=12.\]零假设下服从 $F_{2,6}$，右尾概率为 $(6/(6+2\times12))^3=0.008$，因此拒绝均值全相等。该结论说明至少有一对均值不同；具体哪些组不同需另作有误差控制的比较。\end{solve}

\begin{exercise}[相关系数与回归解释度]\label{q:B16.6}
对含截距的一元最小二乘回归，证明拟合平方和与残差平方和正交分解，并证明 $R^2=r^2$。说明该等式需要排除哪些退化情形。\par\smallskip {\small 来源：编订补充。}
\end{exercise}

\begin{solve}\phantomsection\label{a:B16.6}记 $y_i-\overline y=(\widehat y_i-\overline y)+e_i$。由残差与常数列、自变量列正交，交叉项和为零，所以 $SST=SSR+SSE$。又 $\widehat a=S_{xy}/S_{xx}$，因此 $SSR=\widehat a^2S_{xx}=S_{xy}^2/S_{xx}$，从而\[R^2=\frac{SSR}{S_{yy}}=\frac{S_{xy}^2}{S_{xx}S_{yy}}=r^2.\]需要 $S_{xx}>0,S_{yy}>0$，否则斜率或相关系数可能未定义。该等式针对含截距的一元模型；多元回归中不能用某一个自变量与响应的简单相关系数平方代替整个模型的 $R^2$。\end{solve}

\clearpage\section{综合训练与解题辨析}

\begin{exercise}[帕累托模型从估计到区间]\label{q:M-F-5}
独立观测来自密度 $f(x;\theta)=\theta x^{-\theta-1}$，$x\ge1,\theta>0$。（1）求极大似然估计；（2）证明 $Y=\log X$ 服从指数分布；（3）构造 $\theta$ 的精确置信区间；（4）说明总体均值何时存在。\par\smallskip {\small 来源：MIT 期末试卷，题 5；区间为编订延伸。}
\end{exercise}

\begin{solve}\phantomsection\label{a:M-F-5}设 $T=\sum_i\log x_i$。似然与 $\theta^n\exp(-\theta T)$ 成比例，故对数导数为 $n/\theta-T$，当 $T>0$ 时极大似然估计为 $n/T$。在 $y\ge0$ 上，\[P(Y>y)=P(X>e^y)=e^{-\theta y},\]所以 $Y\sim\operatorname{Exp}(\theta)$，总和 $T\sim\operatorname{Gamma}(n,\theta)$。因此\[\left[\frac{\chi^2_{2n;\alpha/2}}{2T},\frac{\chi^2_{2n;1-\alpha/2}}{2T}\right]\]是精确 $1-\alpha$ 区间。最后，$EX=\int_1^\infty\theta x^{-\theta}\,dx$ 仅在 $\theta>1$ 时有限，值为 $\theta/(\theta-1)$。率参数始终可估计，并不要求原始变量存在均值；不能对所有 $\theta>0$ 都使用基于 $\overline X$ 的有限均值渐近理论。\end{solve}

\begin{exercise}[瑞利参数与变换统计量]\label{q:M-PF-26}
样本密度为 $f(x;\tau)=\tau x\exp(-\tau x^2/2)$，$x>0,\tau>0$。（1）求 $\tau$ 的极大似然估计；（2）求 $Z=X^2/2$ 的分布；（3）给出无偏率估计与精确置信区间。\par\smallskip {\small 来源：MIT 模拟期末，题 26；后两问为编订延伸。}
\end{exercise}

\begin{solve}\phantomsection\label{a:M-PF-26}对数似然中参数相关项为 $n\log\tau-(\tau/2)\sum_i x_i^2$，故\[\widehat\tau=\frac{2n}{\sum_i x_i^2}.\]对 $z>0$，$P(X^2/2>z)=e^{-\tau z}$，因此 $Z$ 为率 $\tau$ 的指数变量，$T=\sum_iX_i^2/2\sim\operatorname{Gamma}(n,\tau)$。当 $n>1$，无偏率估计为 $(n-1)/T$；又 $2\tau T=\tau\sum_iX_i^2\sim\chi^2_{2n}$，所以区间为\[\left[\frac{\chi^2_{2n;\alpha/2}}{\sum_iX_i^2},\frac{\chi^2_{2n;1-\alpha/2}}{\sum_iX_i^2}\right].\]同一模型可通过变量变换复用指数抽样理论，但必须明确瑞利分布这里采用的参数化。\end{solve}

\begin{exercise}[定位中的连续更新与预测]\label{q:M-F-7}
四扇门依次为 $A,B,C,D$，颜色分别为橙、橙、蓝、蓝。机器人初始等概率位于各门；扫描颜色正确率为 $0.7$。首次读到橙色后，求位置后验；随后确定地向右移一扇门（$D$ 后接 $A$），求新位置先验和下一次读到蓝色的预测概率。各次扫描在给定真实位置后独立。\par\smallskip {\small 来源：MIT 期末试卷，题 7。}
\end{exercise}

\begin{solve}\phantomsection\label{a:M-F-7}初始先验为 $(1/4,1/4,1/4,1/4)$。橙色读数的似然为 $(0.7,0.7,0.3,0.3)$，归一化后得到\[(0.35,0.35,0.15,0.15).\]运动后概率按位置循环右移，成为 $(0.15,0.35,0.35,0.15)$。下一次蓝色读数的似然为 $(0.3,0.3,0.7,0.7)$，全概率公式给出预测概率\[0.15(0.3)+0.35(0.3)+0.35(0.7)+0.15(0.7)=0.5.\]若实际又读到蓝色，未归一化后验为 $(0.045,0.105,0.245,0.105)$，归一化得 $(0.09,0.21,0.49,0.21)$。这是“运动预测—观测更新”的两步，不能把移动前的门标签直接沿用。\end{solve}

\begin{exercise}[六类别概率的综合检验]\label{q:M-F-8}
一个教学模型声称六类概率为 $(0.20,0.20,0.20,0.15,0.15,0.10)$。独立抽取 $1000$ 个对象，观察计数为 $(180,190,185,\allowbreak160,165,120)$。在 $0.05$ 水平检验模型，求统计量、自由度与结论，并指出若对象按包装聚集抽样可能出现的问题。\par\smallskip {\small 来源：MIT 期末试卷，题 8（删除品牌及时效性陈述）。}
\end{exercise}

\begin{solve}\phantomsection\label{a:M-F-8}期望计数为 $(200,200,200,150,150,100)$。皮尔逊统计量为\[Q=2+0.5+1.125+\frac23+1.5+4=9.791666\ldots.\]零假设完全指定概率，故自由度为 $5$。右尾概率约为 $0.08136$，在 $0.05$ 水平不拒绝。若同一包装内的对象颜色相关，计数未必服从简单独立多项模型；单把总数当成独立样本容量会误估波动。题设独立性是采用该参考分布的重要依据，而不是由对象数量大自动保证的条件。\end{solve}

\begin{exercise}[两样本检验的完整流程]\label{q:M-F-9}
甲乙两组各有 $11$ 个独立观测，样本均值分别为 $115,110$，样本标准差为 $8,6$。在正态且共同方差的模型下进行双侧均值检验，显著性水平 $0.05$；再求甲总体均值的 $95\%$ 置信区间。\par\smallskip {\small 来源：MIT 期末试卷，题 9（改称甲乙两总体）。}
\end{exercise}

\begin{solve}\phantomsection\label{a:M-F-9}先说明抽样独立和共同方差是假设，不能由两个样本标准差接近自动证明。合并方差为\[s_p^2=\frac{10(64)+10(36)}{20}=50.\]均值差标准误为 $\sqrt{50(2/11)}=\sqrt{100/11}$，统计量为 $t=5/\sqrt{100/11}\approx1.65831$，零假设下自由度为 $20$。双侧 $p$ 值约 $0.1129$，故不拒绝。甲总体均值区间独立地使用甲样本的 $t_{10}$ 分布：\[115\pm t_{10;0.975}\frac8{\sqrt{11}}\approx(109.63,120.37).\]两组比较的自由度与单组区间的自由度不同，不能照搬。\end{solve}

\begin{exercise}[双侧检验在指定备择下的功效]\label{q:M-F-11}
独立抽取 $100$ 个 $N(\mu,1)$ 观测，以 $0.05$ 水平检验 $H_0:\mu=0$ 对双侧备择。求当 $\mu=0.4$ 时的功效，并说明原资料将接受区间写成拒绝区间会造成什么错误。\par\smallskip {\small 来源：MIT 期末试卷，题 11。}
\end{exercise}

\begin{solve}\phantomsection\label{a:M-F-11}零假设下标准误为 $0.1$，不拒绝区间为 $[-0.1960,0.1960]$（端点已舍入），其补集才是拒绝域。在真实均值为 $0.4$ 时，将端点相对 $0.4$ 标准化，功效为\[\begin{aligned}\pi(0.4)&=P_{0.4}(|\overline X|>0.195996)\\&=\Phi(-5.959964)+1-\Phi(-2.040036)\\&\approx0.97933.\end{aligned}\]若误把中央区间叫作拒绝域，会计算成第二类错误概率，约为 $0.02067$。核对区间方向可用直观检验：当真实均值与零假设相差四个标准误时，功效理应较高。\end{solve}

\begin{exercise}[折线基函数的最小二乘]\label{q:M-F-12}
观测为 $(4,1)$、$(-2,1.5)$、$(0,0.5)$，模型为 $Y_i=b|x_i-2|+\varepsilon_i$。求最小二乘估计；在独立零均值同方差误差下求其方差；描述拟合曲线的形状。\par\smallskip {\small 来源：MIT 期末试卷，题 12（以文字替代作图）。}
\end{exercise}

\begin{solve}\phantomsection\label{a:M-F-12}记 $z_i=|x_i-2|$，本题为 $(2,4,2)$，模型对参数 $b$ 线性且无截距。平方和导数给出\[\widehat b=\frac{\sum_i z_iY_i}{\sum_i z_i^2}=\frac{2(1)+4(1.5)+2(0.5)}{4+16+4}=\frac38.\]独立同方差误差保证 $E\widehat b=b$、$\operatorname{Var}(\widehat b)=\sigma^2/24$。拟合曲线为 $\widehat y=(3/8)|x-2|$，顶点 $(2,0)$，左段斜率 $-3/8$，右段斜率 $3/8$。这里自变量变换后的列向量只有一列，不应额外加入未在题设中的截距。\end{solve}

\begin{exercise}[自助分位数的读法]\label{q:M-F-13}
原样本容量 $36$，均值 $7.0$，标准差 $0.84$。进行 $500$ 次有放回自助重抽样并计算样本均值。排序后第 $25$ 个值为 $6.740$，第 $475$ 个值为 $7.224$。采用经验分布逆函数的分位数约定，构造 $90\%$ 百分位自助区间，并讨论重抽样次数的影响。\par\smallskip {\small 来源：MIT 期末试卷，题 13（保留所需次序统计量）。}
\end{exercise}

\begin{solve}\phantomsection\label{a:M-F-13}经验分布逆函数将概率 $q$ 的分位数定义为排序后的第 $\lceil Bq\rceil$ 个值。对 $B=500$，$0.05$ 与 $0.95$ 分位数分别对应第 $25$ 和第 $475$ 个值，所以区间为 $[6.740,7.224]$。不同软件可能采用插值分位数，因此末位可略有差异。增加 $B$ 主要减少自助分位数的蒙特卡洛误差，并不增加原始样本包含的总体信息。百分位区间不是自动精确的 $90\%$ 区间，其覆盖依赖统计量的正则性、样本代表性和近似质量。\end{solve}

\begin{exercise}[抽球机制与逆向条件概率]\label{q:M-PF-18}
袋中有 $3$ 个橙球、$2$ 个蓝球。第一次若取到橙球，则取出不放回；若取到蓝球，则放回并另加一个橙球。随后取第二个球。求两次均橙的概率；已知第二次为橙，求第一次为蓝的概率。\par\smallskip {\small 来源：MIT 模拟期末，题 18。}
\end{exercise}

\begin{solve}\phantomsection\label{a:M-PF-18}令 $O_i,B_i$ 分别表示第 $i$ 次为橙、蓝。第一种路径概率为\[P(O_1O_2)=\frac35\frac24=\frac3{10}.\]若第一次为蓝，第二次袋中有四橙二蓝，故 $P(B_1O_2)=(2/5)(4/6)=4/15$。于是 $P(O_2)=3/10+4/15=17/30$，贝叶斯公式给出\[P(B_1\mid O_2)=\frac{4/15}{17/30}=\frac8{17}.\]两次取球不独立；把第二次橙球概率固定成最初的 $3/5$ 会忽略袋中构成随路径变化。\end{solve}

\begin{exercise}[停时分布的边界计算]\label{q:M-PF-19}
反复独立投公平六面骰子，直到累计点数严格大于 $6$ 才停止，记投掷次数为 $N$。求 $P(N\le1),P(N\le2),P(N\le7)$，再利用尾和求 $EN$。\par\smallskip {\small 来源：MIT 模拟期末，题 19；均值为编订延伸。}
\end{exercise}

\begin{solve}\phantomsection\label{a:M-PF-19}一次点数最多为 $6$，所以第一项为零。两次总点数大于 $6$ 的有序组合有 $21$ 个，因此第二项为 $21/36=7/12$。七次即使都为一点，总和也为 $7$，所以第三项为一。对 $1\le k\le6$，事件 $N>k$ 等价于 $k$ 个正整数点数之和不超过 $6$；每个自动不超过六。正整数分拆的有序计数为 $\binom6k$，故\[EN=\sum_{k=0}^6P(N>k)=\sum_{k=0}^6\binom6k6^{-k}=(7/6)^6.\]其中 $k=0$ 项为一。这个停时有界，所以尾和和期望交换没有收敛困难。\end{solve}

\begin{exercise}[分段密度与分位数]\label{q:M-PF-20}
成绩 $X$ 在 $[0,1]$ 上具有密度：$0\le x\le1/2$ 时为 $4x$，$1/2<x\le1$ 时为 $4-4x$。求及格线为 $0.55$ 时的通过概率、$87.5\%$ 分位数，并写出整个分布函数。\par\smallskip {\small 来源：MIT 模拟期末，题 20。}
\end{exercise}

\begin{solve}\phantomsection\label{a:M-PF-20}两个三角形面积各为 $1/2$，密度归一化。对 $0\le x\le1/2$，$F(x)=2x^2$；对 $1/2<x\le1$，$F(x)=1-2(1-x)^2$。在区间左、右侧分别延拓为 $0,1$。所以通过概率为\[P(X\ge0.55)=2(1-0.55)^2=0.405.\]分位数位于右半区，解 $1-2(1-q)^2=0.875$ 得 $q=0.75$。另一个代数根 $1.25$ 不在支撑内，必须舍去。求分位数时先判断所在的分段，不能对整段使用同一公式。\end{solve}

\begin{exercise}[从分布函数回到密度与期望]\label{q:M-PF-21}
分布函数在 $x<0$ 时为零，在 $0\le x\le1$ 时为 $x(2-x)$，在 $x>1$ 时为一。求密度、$EX$、$\operatorname{Var}(X)$ 和 $P(X<0.4)$。说明端点是否有原子概率。\par\smallskip {\small 来源：MIT 模拟期末，题 21；二阶矩为编订延伸。}
\end{exercise}

\begin{solve}\phantomsection\label{a:M-PF-21}分布函数在两个端点连续，因此没有端点原子。在内部求导得 $f(x)=2-2x$。直接积分可得\[EX=\int_0^1x(2-2x)\,dx=\frac13,\]\[EX^2=\int_0^1x^2(2-2x)\,dx=\frac16.\]因此方差为 $1/6-1/9=1/18$。因无点质量，$P(X<0.4)=F(0.4)=0.64$。一般地，分布函数某点的跳跃大小才是该点概率，导数密度只捕捉绝对连续部分。\end{solve}

\begin{exercise}[参数自助法与基本区间]\label{q:M-PF-29}
连续独立伯努利试验的成功概率为 $p$，记录直到首次失败的总试验次数 $X$（包括失败），故 $P(X=k)=p^{k-1}(1-p)$。给定容量 $n$ 的样本，估计 $p$，并说明参数自助法构造 $95\%$ 基本区间的步骤。\par\smallskip {\small 来源：MIT 模拟期末，题 29。}
\end{exercise}

\begin{solve}\phantomsection\label{a:M-PF-29}均值为 $1/(1-p)$，矩估计也是极大似然估计：$\widehat p=1-1/\overline X$。以该拟合参数为真，从上述几何分布模拟独立自助样本，容量仍为 $n$；计算每次 $\widehat p^*=1-1/\overline X^*$，重复足够多次。令估计误差 $\Delta^*=\widehat p^*-\widehat p$ 的经验 $0.025,0.975$ 分位数为 $d_L,d_U$，基本区间为\[[\widehat p-d_U,\widehat p-d_L].\]它是对估计误差区间作反演，与直接取 $\widehat p^*$ 分位数的百分位区间不同。可与参数空间 $[0,1]$ 相交；若估计落在边界导致拟合分布退化，普通自助近似需特别谨慎。\end{solve}

\begin{exercise}[极小尾概率与数值稳定性]\label{q:M-PF-30}
独立抽取 $100$ 个 $N(\mu,4)$ 观测，双侧检验 $H_0:\mu=3$，水平 $0.05$。若样本均值为 $5$，求 $p$ 值；再求 $\mu=4$ 时的功效。解释为什么不能把很小的正态尾概率写成严格的零。\par\smallskip {\small 来源：MIT 模拟期末，题 30（校订非零尾概率）。}
\end{exercise}

\begin{solve}\phantomsection\label{a:M-PF-30}标准误为 $2/10=0.2$，拒绝域为 $|\overline X-3|>1.959964(0.2)$。当均值为 $5$，观测 $z=10$，\[p=2[1-\Phi(10)]\approx1.52397\times10^{-23}.\]数值计算宜使用生存函数直接计算右尾，避免浮点数中 $1-\Phi(10)$ 的消减误差。在真实均值为 $4$ 时，接受区间相对该均值标准化的两端为 $-6.959964,-3.040036$，所以\[\pi(4)=1-\Phi(-3.040036)+\Phi(-6.959964)\approx0.99882.\]有限实数处正态尾概率严格为正；数值打印为零只表示精度不足。\end{solve}

\begin{exercise}[正态测量的后验精度规划]\label{q:M-PF-31}
未知量 $\theta$ 的先验为 $N(80,4)$，每次测量误差独立且服从 $N(0,0.5)$。一次读数为 $85$，求后验分布。若继续独立测量，总共至少测量多少次可使后验方差小于 $0.01$？\par\smallskip {\small 来源：MIT 模拟期末，题 31。}
\end{exercise}

\begin{solve}\phantomsection\label{a:M-PF-31}先验精度为 $1/4$，单次观测精度为 $2$。因此后验方差为 $v=(1/4+2)^{-1}=4/9$，均值为\[m=\frac{80/4+85/0.5}{1/4+2}=\frac{760}{9}.\]故后验为 $N(760/9,4/9)$。在 $n$ 次测量后，后验方差为 $v_n=(1/4+2n)^{-1}=4/(1+8n)$。要求 $v_n<0.01$，即 $n>399/8=49.875$，所以最少总测量次数为 $50$。方差只依赖次数和误差方差，后验均值则还依赖实际测量值。\end{solve}

\begin{exercise}[泊松分裂与参数推断]\label{q:B17.1}
一天到达的对象数 $N\sim\operatorname{Pois}(\lambda)$，每个对象独立以概率 $p$ 标为甲类。记两类数量为 $X,Y$。（1）推导联合分布并判断独立性；（2）连续观察 $m$ 天两类总数 $T_X,T_Y$，求 $\lambda,p$ 的极大似然估计；（3）说明零总数时的问题。\par\smallskip {\small 来源：编订补充。}
\end{exercise}

\begin{solve}\phantomsection\label{a:B17.1}给定 $N=n$，$X\sim\operatorname{Bin}(n,p)$，且 $Y=n-X$。对非负整数 $x,y$，\[P(X=x,Y=y)=e^{-\lambda}\frac{\lambda^{x+y}}{x!y!}p^x(1-p)^y.\]它可分解为率 $\lambda p$ 与 $\lambda(1-p)$ 的两个泊松质量函数，故 $X,Y$ 独立。多天数据的参数相关似然为\[e^{-m\lambda}\lambda^{T_X+T_Y}p^{T_X}(1-p)^{T_Y}.\]当总数 $T=T_X+T_Y>0$，极大似然为 $\widehat\lambda=T/m$、$\widehat p=T_X/T$。若 $T=0$，关于 $p$ 的似然恒定，数据不能识别分类比例；率参数的最大值位于零边界（若参数空间允许）。\end{solve}

\begin{exercise}[有限总体抽样与修正因子]\label{q:B17.2}
有限总体有 $N$ 个对象，其中 $M$ 个具有某特征。无放回简单随机抽样 $n$ 个，令 $X$ 为特征数、$\widehat p=X/n$。求分布、均值和方差，比较与有放回伯努利抽样的差别。\par\smallskip {\small 来源：编订补充。}
\end{exercise}

\begin{solve}\phantomsection\label{a:B17.2}每个大小为 $n$ 的子集等可能，所以\[P(X=k)=\frac{\binom Mk\binom{N-M}{n-k}}{\binom Nn}.\]令 $p=M/N$，单个抽样位置的指示变量均值为 $p$；不同位置的联合成功概率为 $M(M-1)/[N(N-1)]$，故协方差为 $-p(1-p)/(N-1)$。于是\[E\widehat p=p,\qquad\operatorname{Var}(\widehat p)=\frac{p(1-p)}n\frac{N-n}{N-1}.\]最后的因子是有限总体修正。抽样比例很小时近似有放回模型；当 $n=N$ 时普查没有抽样误差，方差恰好为零。\end{solve}

\begin{exercise}[随机化回答的估计误差]\label{q:B17.3}
为保护隐私，每名对象独立掷一枚公平硬币：正面时回答真实的是非问题，反面时强制回答“是”。真实“是”的总体比例为 $p$。对 $n$ 人记录“是”的比例 $\widehat q$。构造 $p$ 的无偏估计，求方差并解释参数越界问题。\par\smallskip {\small 来源：编订补充。}
\end{exercise}

\begin{solve}\phantomsection\label{a:B17.3}单人回答“是”的概率为 $q=p/2+1/2$，所以 $p=2q-1$。估计量 $\widehat p=2\widehat q-1$ 无偏，且\[\operatorname{Var}(\widehat p)=4q(1-q)/n=(1-p^2)/n.\]因为有限样本可能出现 $\widehat q<1/2$，该无偏估计可能为负。将它截断到 $[0,1]$ 可满足参数约束并减小针对该区间真值的平方误差，但通常引入偏差。该机制牺牲一部分统计效率以保护个体回答隐私；独立硬币机制与真实状态的关系必须明确。\end{solve}

\begin{exercise}[尾和公式与收集全套]\label{q:B17.4}
独立抽取 $m$ 种等概率卡片，每次一张。记收集到全部 $m$ 种所需次数为 $T$。把过程分解为等待新类型的阶段，求 $ET$ 和 $\operatorname{Var}(T)$，并说明独立性依据。\par\smallskip {\small 来源：编订补充。}
\end{exercise}

\begin{solve}\phantomsection\label{a:B17.4}已经有 $k$ 种类型时，下一次得到新类型的概率为 $(m-k)/m$。记此阶段等待次数为 $W_k$，包括成功一次。独立抽样与到达阶段时的重启性质保证各阶段等待时间相互独立，且 $W_k$ 为支撑 $1,2,\ldots$ 的几何变量。因此\[ET=\sum_{k=0}^{m-1}\frac{m}{m-k}=m\sum_{j=1}^m\frac1j.\]几何方差为 $(1-p)/p^2$，故\[\operatorname{Var}(T)=\sum_{j=1}^m\left(\frac{m^2}{j^2}-\frac mj\right).\]最后几种类型的成功概率很小，导致收集末期等待较久。只用“每种平均等 $m$ 次”相加会忽略状态变化。\end{solve}

\begin{exercise}[随机和的复合泊松结构]\label{q:B17.5}
设 $N\sim\operatorname{Pois}(\lambda)$，$Z_i$ 独立同分布且与 $N$ 独立，$EZ_i=\mu$、$\operatorname{Var}(Z_i)=\sigma^2$。令 $S=\sum_{i=1}^NZ_i$，空和为零。求 $ES,\operatorname{Var}(S)$ 和特征函数，并说明只给前两矩是否足以精确求 $P(S\le c)$。\par\smallskip {\small 来源：编订补充。}
\end{exercise}

\begin{solve}\phantomsection\label{a:B17.5}给定 $N$，条件均值与方差分别为 $N\mu,N\sigma^2$。全期望、全方差给出\[ES=\lambda\mu,\qquad\operatorname{Var}(S)=\lambda(\sigma^2+\mu^2).\]若 $\varphi_Z(t)=Ee^{itZ_1}$，则\[\varphi_S(t)=E[\varphi_Z(t)^N]=\exp\{\lambda[\varphi_Z(t)-1]\}.\]只知道 $\mu,\sigma^2$ 不能确定整个特征函数，也就不能精确确定一般尾概率。可以给出切比雪夫界或在额外大强度条件下作正态近似，但必须标明近似性质。\end{solve}

\begin{exercise}[正态变量平方和与比值]\label{q:B17.6}
设 $Z_1,Z_2,Z_3$ 独立标准正态。令 $U=Z_1^2+Z_2^2$、$V=Z_3^2$。求 $U,V$ 的分布、$U/(2V)$ 的分布，以及 $U/(U+V)$ 的分布。说明后两个比值的不同支撑。\par\smallskip {\small 来源：编订补充。}
\end{exercise}

\begin{solve}\phantomsection\label{a:B17.6}$U\sim\chi^2_2$、$V\sim\chi^2_1$ 且独立，所以 $(U/2)/(V/1)\sim F_{2,1}$，支撑为正半轴。将卡方视为同率 $1/2$ 的伽马变量，变换 $R=U/(U+V),T=U+V$ 的逆式为 $U=RT,V=(1-R)T$，雅可比为 $T$。联合密度可分解为贝塔核和伽马核，得到\[R\sim\operatorname{Beta}(1,1/2),\qquad T\sim\chi^2_3,\]并且独立。$R$ 的支撑是 $(0,1)$。二者虽都由平方和比值形成，但缩放比与份额比的分布不同。\end{solve}

\begin{exercise}[指数竞争与获胜标记]\label{q:B17.7}
设 $X\sim\operatorname{Exp}(\lambda)$、$Y\sim\operatorname{Exp}(\mu)$ 独立。令 $T=\min(X,Y)$、$J=\mathbf1_{\{X<Y\}}$。求 $T,J$ 的分布并证明独立性。\par\smallskip {\small 来源：编订补充。}
\end{exercise}

\begin{solve}\phantomsection\label{a:B17.7}生存函数相乘得到 $P(T>t)=e^{-(\lambda+\mu)t}$，故 $T\sim\operatorname{Exp}(\lambda+\mu)$。对小区间 $dt$，联合密度标记为\[P(T\in dt,J=1)=\lambda e^{-\lambda t}e^{-\mu t}\,dt.\]将其写成 $[\lambda/(\lambda+\mu)]\,[(\lambda+\mu)e^{-(\lambda+\mu)t}]dt$，可见 $P(J=1)=\lambda/(\lambda+\mu)$ 且条件时间分布与无条件相同；$J=0$ 同理。因此 $T$ 与先发生哪一类事件独立。这一性质依赖指数分布的恒定风险率，一般寿命分布不一定成立。\end{solve}

\begin{exercise}[条件期望作为最小平方预测]\label{q:B17.8}
设 $EY^2<\infty$，以 $X$ 的可测函数 $g(X)$ 预测 $Y$。证明 $E(Y\mid X)$ 最小化均方预测误差，并给出误差分解。若 $(X,Y)$ 联合正态，写出最优预测。\par\smallskip {\small 来源：编订补充。}
\end{exercise}

\begin{solve}\phantomsection\label{a:B17.8}记 $m(X)=E(Y\mid X)$。对二阶可积的 $g(X)$，展开平方并使用\[E[(Y-m(X))(m(X)-g(X))]=0,\]其理由是先条件于 $X$，残差的条件期望为零。因而\[E(Y-g(X))^2=E(Y-m(X))^2+E(m(X)-g(X))^2.\]最小值在 $g=m$ 几乎处处达到。联合正态且 $\operatorname{Var}(X)>0$ 时，\[m(X)=EY+\frac{\operatorname{Cov}(X,Y)}{\operatorname{Var}(X)}(X-EX).\]对一般联合分布，这个线性式只是最优线性预测，未必等于全部可测函数中的最优预测。\end{solve}

\begin{exercise}[混合正态的方差分解]\label{q:B17.9}
先以概率 $w$ 选第一组、概率 $1-w$ 选第二组，随后分别从 $N(\mu_1,\sigma_1^2)$ 或 $N(\mu_2,\sigma_2^2)$ 生成 $X$。求均值和方差。若两个组的均值相同，混合分布是否必为正态？\par\smallskip {\small 来源：编订补充。}
\end{exercise}

\begin{solve}\phantomsection\label{a:B17.9}令组标记为 $G$。全期望给出 $EX=w\mu_1+(1-w)\mu_2$。全方差分为组内波动和组间波动：\[\begin{aligned}\operatorname{Var}(X)&=w\sigma_1^2+(1-w)\sigma_2^2\\&\quad+w(1-w)(\mu_1-\mu_2)^2.\end{aligned}\]即使均值相同而方差不同，混合分布一般也非正态。例如共同均值为零时四阶矩为 $3[w\sigma_1^4+(1-w)\sigma_2^4]$，而相同总体方差的正态四阶矩是 $3[w\sigma_1^2+(1-w)\sigma_2^2]^2$；非退化权重及不同方差使前者严格更大。\end{solve}

\begin{exercise}[德尔塔方法估计风险函数]\label{q:B17.10}
独立伯努利样本参数为 $p\in(0,1)$，样本比例为 $\widehat p$。推导对数赔率 $g(p)=\log[p/(1-p)]$ 的代入估计渐近分布，构造近似区间，并解释边界观测问题。\par\smallskip {\small 来源：编订补充。}
\end{exercise}

\begin{solve}\phantomsection\label{a:B17.10}中心极限定理给出 $\sqrt n(\widehat p-p)\xrightarrow dN(0,p(1-p))$。导数为 $g^{\prime}(p)=1/[p(1-p)]$，故\[\sqrt n[g(\widehat p)-g(p)]\xrightarrow dN\left(0,\frac1{p(1-p)}\right).\]代入估计标准误为 $1/\sqrt{n\widehat p(1-\widehat p)}$，可在对数赔率尺度构造正态区间，再用逻辑函数 $e^u/(1+e^u)$ 映回比例区间。当样本全成功或全失败时公式未定义；固定内部真值下这类事件概率趋于零，但有限样本仍需使用适合边界的区间或明确的平滑方案。\end{solve}

\begin{exercise}[威尔逊比例区间的反演]\label{q:B17.11}
对伯努利比例 $p$，从正态得分不等式 $|\widehat p-p|\le z\sqrt{p(1-p)/n}$ 出发推导威尔逊区间。指出它与把 $p$ 直接替成 $\widehat p$ 的沃尔德区间有何区别。\par\smallskip {\small 来源：编订补充。}
\end{exercise}

\begin{solve}\phantomsection\label{a:B17.11}两侧平方并移项得到\[(n+z^2)p^2-(2n\widehat p+z^2)p+n\widehat p^2\le0.\]求二次方程的两个根，区间中心为\[c=\frac{\widehat p+z^2/(2n)}{1+z^2/n},\]半径为\[h=\frac{z}{1+z^2/n}\sqrt{\frac{\widehat p(1-\widehat p)}n+\frac{z^2}{4n^2}}.\]因此区间为 $[c-h,c+h]$。它保留零假设方差中的未知 $p$ 再反演，因此即使 $\widehat p=0$ 或 $1$ 也不会退化成单点；沃尔德区间在这些情形半径为零。该推导仍基于正态近似，并非一般有限样本精确覆盖。\end{solve}

\begin{exercise}[正态模型的充分统计量]\label{q:B17.12}
正态样本的 $\mu,\sigma^2$ 都未知。证明 $(\sum_iX_i,\sum_iX_i^2)$ 是充分统计量，说明它与 $(\overline X,S^2)$ 的等价关系，并讨论只保留均值是否足够。\par\smallskip {\small 来源：编订补充。}
\end{exercise}

\begin{solve}\phantomsection\label{a:B17.12}展开平方和后，联合密度可写为\[(2\pi\sigma^2)^{-n/2}\exp\left[-\frac{\sum_i x_i^2-2\mu\sum_i x_i+n\mu^2}{2\sigma^2}\right].\]样本仅通过两个给定总和进入参数相关部分，因子分解定理给出充分性。容量固定且 $n\ge2$ 时，由\[\sum_iX_i=n\overline X,\quad\sum_iX_i^2=(n-1)S^2+n\overline X^2\]可互相恢复两个表示。只保留均值会丢失离差信息，一般不能同时充分反映均值和方差两个未知参数。充分性说的是模型内参数信息，不意味着保留了所有可视化或模型诊断信息。\end{solve}

\begin{exercise}[最小值模型的非对称估计]\label{q:B17.13}
设 $X_i=\theta+E_i$，其中 $E_i$ 独立服从率 $\lambda$ 已知的指数分布，$\theta\in\mathbb R$。求 $\theta$ 的极大似然估计、无偏修正和一个精确 $1-\alpha$ 置信区间。\par\smallskip {\small 来源：编订补充。}
\end{exercise}

\begin{solve}\phantomsection\label{a:B17.13}似然中参数相关部分为 $e^{n\lambda\theta}\mathbf1_{\{\theta\le M\}}$，$M=\min_iX_i$，所以最大值在 $\widehat\theta=M$ 达到。因为 $M-\theta\sim\operatorname{Exp}(n\lambda)$，有 $EM=\theta+1/(n\lambda)$，无偏修正为 $M-1/(n\lambda)$。又\[P\left(0\le M-\theta\le\frac{-\log\alpha}{n\lambda}\right)=1-\alpha,\]所以区间为 $[M+\log\alpha/(n\lambda),M]$。它以最小值为上界，体现模型只允许观测位于参数右侧；对称正态区间在这里并非自然构造。\end{solve}

\begin{exercise}[贝塔—二项预测的额外变异]\label{q:B17.14}
给定参数 $p$ 后，未来 $m$ 次试验独立伯努利；当前后验为 $p\mid D\sim\operatorname{Beta}(a,b)$。令未来成功数为 $K$。求预测质量函数、均值和方差，并与固定 $p=a/(a+b)$ 的二项分布比较。\par\smallskip {\small 来源：编订补充。}
\end{exercise}

\begin{solve}\phantomsection\label{a:B17.14}对参数积分得到\[P(K=k\mid D)=\binom mk\frac{B(a+k,b+m-k)}{B(a,b)}.\]令 $s=a+b$，则 $Ep=a/s$，$\operatorname{Var}(p)=ab/[s^2(s+1)]$。全期望给出 $EK=ma/s$；全方差给出\[\begin{aligned}\operatorname{Var}(K)&=mE[p(1-p)]+m^2\operatorname{Var}(p)\\&=\frac{mab(s+m)}{s^2(s+1)}.\end{aligned}\]固定均值参数的二项方差是 $mab/s^2$，预测方差多了因子 $(s+m)/(s+1)$。当 $m>1$ 时它大于一，因为共同未知参数使未来试验在边际预测下相关。\end{solve}

\begin{exercise}[正态性、独立性与零相关]\label{q:B17.15}
令 $Z\sim N(0,1)$，独立符号变量 $S$ 等概率取 $-1,1$。定义 $X=Z,Y=SZ$。证明两个边缘都为标准正态且协方差为零，但不独立。说明这为何不违反联合正态中的零相关结论。\par\smallskip {\small 来源：编订补充。}
\end{exercise}

\begin{solve}\phantomsection\label{a:B17.15}正态分布对称，所以随机乘以独立符号不改变边缘分布，$X,Y$ 都为标准正态。又 $EXY=E(SZ^2)=ES\,EZ^2=0$，两个均值均为零，故协方差为零。另一方面 $Y^2=X^2$ 以概率一成立；若二者独立，则 $X^2,Y^2$ 也独立，但同一个非退化随机量不可能与自身独立。比如事件 $\{X^2\le1\}$ 与 $\{Y^2\le1\}$ 实际相同，其概率严格介于零与一。该向量集中在两条直线上，并非联合正态；边缘分别正态不等于联合正态。\end{solve}

\begin{exercise}[支撑集破坏独立性]\label{q:B17.16}
设 $(X,Y)$ 在三角形 $0<x<y<1$ 上密度为 $2$，其余处为零。有人因内部密度可写成常数乘常数而认定独立。计算边缘密度、协方差并指出错误。\par\smallskip {\small 来源：编订补充。}
\end{exercise}

\begin{solve}\phantomsection\label{a:B17.16}积分支撑得 $f_X(x)=2(1-x)$、$f_Y(y)=2y$，各自在 $(0,1)$ 上。联合密度不等于边缘乘积，并且三角形外仍有两个边缘密度都正而联合密度为零的区域。计算得 $EX=1/3,EY=2/3$，\[EXY=\int_0^1\int_0^y2xy\,dx\,dy=\frac14.\]所以协方差为 $1/4-2/9=1/36>0$。检查独立性需要把支撑指示函数一起考虑，只看内部代数表达式能否拆成两个因子是不充分的。\end{solve}

\begin{exercise}[蒙特卡洛积分与标准误]\label{q:B17.17}
用独立 $U_i\sim U(0,1)$ 估计积分 $I=\int_0^1x^2dx$，取 $\widehat I=n^{-1}\sum_iU_i^2$。求偏差、方差、近似置信区间与切比雪夫误差界；比较误差随计算量的变化。\par\smallskip {\small 来源：编订补充。}
\end{exercise}

\begin{solve}\phantomsection\label{a:B17.17}$EU_i^2=1/3$，$EU_i^4=1/5$，所以 $\operatorname{Var}(U_i^2)=1/5-1/9=4/45$。因此 $\widehat I$ 无偏且方差为 $4/(45n)$。中心极限定理给出近似区间\[\widehat I\pm z_{1-\alpha/2}\sqrt{4/(45n)}.\]切比雪夫不等式给出有限样本保证 $P(|\widehat I-I|\ge\varepsilon)\le4/(45n\varepsilon^2)$。标准误按 $n^{-1/2}$ 下降，将标准误减半通常要将独立样本数增至四倍。数值积分真值在此已知是为了核查估计方法，实际复杂积分中通常需用样本方差估计未知方差。\end{solve}

\begin{exercise}[反向变量的方差缩减]\label{q:B17.18}
仍估计 $\int_0^1x^2dx$。对一个均匀变量 $U$，用 $A=[U^2+(1-U)^2]/2$ 作为一次配对估计。求其期望和方差，与使用两个独立均匀变量的平均比较。\par\smallskip {\small 来源：编订补充。}
\end{exercise}

\begin{solve}\phantomsection\label{a:B17.18}由于 $U$ 与 $1-U$ 的边缘都均匀，$EA=1/3$。化简为 $A=U^2-U+1/2$，平方积分得\[EA^2=\int_0^1(U^2-U+1/2)^2\,dU=\frac7{60}.\]所以 $\operatorname{Var}(A)=7/60-1/9=1/180$。两个独立样本平方的平均方差为 $(4/45)/2=2/45$，是配对法的八倍。原因是两项负相关，反向取样抵消波动。反向变量法并非对所有函数都有改善，必须检查变换后函数的协方差方向。\end{solve}

\begin{exercise}[切比雪夫界与正态近似的角色]\label{q:B17.19}
独立同分布样本均值为 $\mu$，方差为 $4$，容量 $100$。估计 $P(|\overline X-\mu|<0.5)$。分别给出仅用方差的下界、中心极限定理近似，以及总体正态时的精确值。\par\smallskip {\small 来源：编订补充。}
\end{exercise}

\begin{solve}\phantomsection\label{a:B17.19}样本均值方差为 $0.04$。切比雪夫不等式给出\[P(|\overline X-\mu|\ge0.5)\le0.04/0.25=0.16,\]所以所求概率至少为 $0.84$。标准化阈值为 $0.5/0.2=2.5$，正态近似给出 $2\Phi(2.5)-1\approx0.98758$。如果总体本来正态，该数值就是精确概率；一般总体下只能视为大样本近似，容量 $100$ 不能单独保证误差很小。前者保守但不依赖分布形状，后者通常更精细但依赖额外近似依据。\end{solve}

\begin{exercise}[学生化与斯卢茨基定理]\label{q:B17.20}
设独立同分布样本具有有限正方差 $\sigma^2$，但不要求正态。证明 $\sqrt n(\overline X-\mu)/S\xrightarrow dN(0,1)$，并解释它与有限样本 $t_{n-1}$ 分布的区别。\par\smallskip {\small 来源：编订补充。}
\end{exercise}

\begin{solve}\phantomsection\label{a:B17.20}中心极限定理给出 $Z_n=\sqrt n(\overline X-\mu)/\sigma\xrightarrow dN(0,1)$。有限二阶矩保证样本一阶、二阶矩的强大数定律适用，故 $S^2\to\sigma^2$，于是 $\sigma/S\to1$ 依概率成立。斯卢茨基定理得到\[Z_n\frac\sigma S\xrightarrow dN(0,1).\]这里不要求 $Z_n$ 与 $S$ 独立。若总体正态，则有任意有限 $n$ 下的精确 $t_{n-1}$ 分布；一般总体中没有这一精确结论。小样本强偏态时，形式上使用 $t$ 分位数并不能自动保证所标注的置信水平。\end{solve}

\begin{exercise}[三种收敛之间的辨析]\label{q:B17.21}
证明均方收敛蕴含依概率收敛。给出一个依概率收敛到零但期望不趋于零的非负随机变量序列，并解释这对交换极限与期望的警示。\par\smallskip {\small 来源：编订补充。}
\end{exercise}

\begin{solve}\phantomsection\label{a:B17.21}若 $E|X_n-X|^2\to0$，马尔可夫不等式给出\[P(|X_n-X|>\varepsilon)\le\frac{E|X_n-X|^2}{\varepsilon^2}\to0.\]取 $U\sim U(0,1)$，令 $X_n=n\mathbf1_{\{U\le1/n\}}$。对任意固定 $\varepsilon>0$ 和足够大的 $n$，$P(X_n>\varepsilon)=1/n\to0$，但 $EX_n=1$ 始终不变。罕见但越来越大的取值保留了期望。依概率收敛本身不足以交换极限和期望；常见充分条件包括支配收敛条件或一致可积性。\end{solve}

\begin{exercise}[参数空间约束下的正态估计]\label{q:B17.22}
样本来自 $N(\mu,1)$，已知 $\mu\ge0$。求约束极大似然估计，并比较它与无约束均值估计在平方损失下的误差。它是否仍然无偏？\par\smallskip {\small 来源：编订补充。}
\end{exercise}

\begin{solve}\phantomsection\label{a:B17.22}最大化似然等价于最小化 $n(\mu-\overline X)^2$，同时满足 $\mu\ge0$，因此\[\widehat\mu=\max(0,\overline X).\]对每个真实 $\mu\ge0$ 和每个观测值，将负数投影到零都不会增大到 $\mu$ 的距离，所以逐点有 $(\widehat\mu-\mu)^2\le(\overline X-\mu)^2$，进而均方误差不大于 $1/n$。但在 $\mu=0$ 时，$E\widehat\mu=1/\sqrt{2\pi n}>0$，故有偏。这个例子说明参数约束可以使有偏估计在允许的参数空间内具有更小的平方风险。\end{solve}

\begin{exercise}[遗漏变量导致斜率变化]\label{q:B17.23}
总体满足 $Y=\beta_0+\beta_1X+\beta_2Z+\varepsilon$，其中 $E\varepsilon=0$，误差与 $X$ 不相关，二阶矩有限且 $\operatorname{Var}(X)>0$。如果只用 $X$ 对 $Y$ 做含截距的总体最小平方线性预测，斜率是多少？何时等于 $\beta_1$？\par\smallskip {\small 来源：编订补充。}
\end{exercise}

\begin{solve}\phantomsection\label{a:B17.23}总体线性预测的斜率为 $\operatorname{Cov}(X,Y)/\operatorname{Var}(X)$。利用协方差线性性，\[a=\beta_1+\beta_2\frac{\operatorname{Cov}(X,Z)}{\operatorname{Var}(X)}.\]当 $\beta_2=0$ 或 $\operatorname{Cov}(X,Z)=0$ 时额外项为零；也可能由其他特定退化关系使其为零。回归系数因此依赖模型中包含哪些变量，不能只根据两变量相关就赋予因果解释。这里推导的是总体最小平方投影系数，不是有限样本斜率的逐个数值等式。\end{solve}

\begin{exercise}[加权最小二乘]\label{q:B17.24}
模型为 $Y=X\beta+\varepsilon$，$X$ 满列秩，$E\varepsilon=0$，协方差为已知正定矩阵 $\Sigma$。推导广义最小二乘估计及其协方差。若误差独立且方差不同，解释各观测的权重。\par\smallskip {\small 来源：编订补充。}
\end{exercise}

\begin{solve}\phantomsection\label{a:B17.24}用白化矩阵 $\Sigma^{-1/2}$ 变换数据，得到协方差为单位阵的误差。最小化\[(Y-X\beta)^\mathsf T\Sigma^{-1}(Y-X\beta),\]正规方程给出\[\widehat\beta=(X^\mathsf T\Sigma^{-1}X)^{-1}X^\mathsf T\Sigma^{-1}Y.\]其均值为 $\beta$，协方差为 $(X^\mathsf T\Sigma^{-1}X)^{-1}$。独立异方差时，$\Sigma$ 为对角阵，目标函数是 $\sum_i(Y_i-x_i^\mathsf T\beta)^2/\sigma_i^2$，精度高的观测权重大。若这些方差需从数据估计，则以上公式的有限样本推断不能直接视为已知协方差情形。\end{solve}

\clearpage
\part{解题方法与公式速查}
\section*{方法一：先定义试验，再计算概率}
\addcontentsline{toc}{section}{方法一：先定义试验，再计算概率}
概率空间由样本空间、事件集合和概率规则组成。本科计算题常把这三者写在自然语言中，第一步应把它们翻译为明确的数学对象。例如“任取两件产品”必须区分有放回和无放回，“随机取两点”必须区分每点均匀、彼此独立以及相对于长度、面积还是体积均匀。若这些条件尚未确定，仅凭“随机”二字不能选定分布。

古典概率中的分子和分母必须在同一种等可能基本结果上计数。抽取有先后次序时，可以把每个有序序列视为一个基本结果；也可以转成无序子集，但分子与分母应同时转化。若某些无序结果对应的序列数不同，无序结果并不等可能。例如两次掷骰子的和有十一种可能值，却不能将每个和的概率都记为十一分之一。

组合数 $\binom nk$ 适合无序选择，排列数 $n!/(n-k)!$ 适合无放回的有序选择。有放回的 $k$ 次独立选择通常产生 $n^k$ 个有序结果。处理“恰好”时特别注意不要把“至少”计数带入；手套题中选中一副完整手套后，剩余两只不能再来自同一副。

遇到“至少一个”时先试补事件。若 $A_i$ 表示第 $i$ 个对象满足要求，通常 $P(\bigcup_iA_i)$ 比 $P(\bigcap_iA_i^c)$ 难计算。独立性成立时后者是乘积；独立性不成立时，仍可使用条件概率连乘，但不能直接相乘边缘概率。容斥公式适合许多对称事件的交集容易计数的情形，如帽子、信封与占位问题。
\begin{equation*}
P\left(\bigcup_{i=1}^nA_i\right)
=\sum_{\varnothing\ne J\subset\{1,\ldots,n\}}(-1)^{|J|-1}P\left(\bigcap_{j\in J}A_j\right).
\end{equation*}
计算后至少检查三个边界：概率是否在 $[0,1]$；样本容量为零或一时是否合乎题意；参数趋近零或一时结果是否具有合理极限。这些检查不能替代推导，却常能发现漏项、错位和重复计数。对应训练为习题~\ref{q:Y1.9}、\ref{q:Y2.3}、\ref{q:Y2.5}。

\section*{方法二：条件概率与信息产生机制}
\addcontentsline{toc}{section}{方法二：条件概率与信息产生机制}
条件概率（conditional probability）在 $P(B)>0$ 时定义为 $P(A\mid B)=P(A\cap B)/P(B)$。分母改变了参考总体：我们只在已知 $B$ 发生的情形中比较 $A$。因此 $P(A\mid B)$ 与 $P(B\mid A)$ 通常不同。检测、证人和分类问题中，题目常给后者却要求前者，此时需要贝叶斯公式。

若 $H_1,\ldots,H_k$ 构成互斥完备假设，先验为 $\pi_j$，观察数据 $D$ 的似然为 $L_j=P(D\mid H_j)$，则后验为
\[
P(H_j\mid D)=\frac{\pi_jL_j}{\sum_i\pi_iL_i}.
\]
实际解题可先计算未归一化权重 $w_j=\pi_jL_j$，最后统一除以权重和。这样避免对每个假设重复计算全概率分母，也便于发现哪些假设的似然为零。两假设情形可使用赔率：后验赔率等于先验赔率乘以似然比。

连续观测没有非零单点概率，但有密度似然。给定观测 $x$ 后，$L(\theta;x)=f_\theta(x)$ 是参数的函数，不要求对参数积分等于一。贝叶斯后验则还需乘先验密度，再对参数归一化。将似然本身视为参数的概率分布，只有在明确采用相应先验且归一化有意义的特定情形下才可能成立。

连续检测需要检查的是给定真实状态后的条件独立，而不是检测结果的无条件独立。若两个结果由同一个未知状态共同影响，即使它们条件独立，边际上也往往相关。三门问题进一步说明：观察到“某门被打开且里面没有奖品”，其似然取决于主持人如何选门。信息并不只在看到的结果里，也在产生结果的规则里。对应训练为习题~\ref{q:Y2.23}、\ref{q:M-P7-1}、\ref{q:M-P9-2}。

\section*{方法三：分布函数比密度更基本}
\addcontentsline{toc}{section}{方法三：分布函数比密度更基本}
随机变量的分布函数（cumulative distribution function）为 $F(x)=P(X\le x)$，对所有实值随机变量都存在。它单调不减、右连续，在正负无穷处的极限分别为一和零。单点概率等于跳跃量 $P(X=x)=F(x)-F(x-)$。若分布绝对连续，存在非负可积函数 $f$ 使 $F(x)=\int_{-\infty}^xf(t)dt$；此时单点概率均为零，$f$ 只是几乎处处定义的密度。

验证候选密度，先查非负性和支撑，再查积分为一。积分上下限由支撑决定，而不是看到二元密度便自动写成正方形。对 $0<x<y<1$ 的三角区域，固定 $y$ 时 $x$ 从零到 $y$；固定 $x$ 时 $y$ 从 $x$ 到一。换积分次序之前画出不等式对应的区域，能避免把条件约束漏掉。

若求 $Y=g(X)$ 的分布，最稳妥的方法是先写
\[
F_Y(y)=P(g(X)\le y).
\]
例如平方变换通常把事件变成 $-\sqrt y\le X\le\sqrt y$，自然揭示两条分支。严格单调可微变换可直接使用密度公式，但需要逆函数和导数绝对值。分段单调变换必须对所有逆像求和：
\[
f_Y(y)=\sum_{x:g(x)=y}\frac{f_X(x)}{|g'(x)|},
\]
在常规可微、导数非零的分支上成立，并按几乎处处意义理解。若变换把一段区间压成一个点，结果可能出现原子，此时单纯求导可能漏掉点质量。

独立变量的最大值和最小值常用分布函数最方便：最大值不超过 $t$ 等价于所有变量不超过 $t$；最小值大于 $t$ 等价于所有变量大于 $t$。这些事件关系对离散分布同样适用，但求单点质量时要使用分布函数差。对应训练为习题~\ref{q:Y3.20}、\ref{q:Y3.25}、\ref{q:B8.7}。

\section*{方法四：联合分布、支撑与独立性}
\addcontentsline{toc}{section}{方法四：联合分布、支撑与独立性}
联合密度描述的是随机向量在区域内的概率。求边缘分布应对另一个坐标积分，且积分区间可能依赖保留的坐标。求和变换 $Z=X+Y$ 时，若引入 $W=X$，原支撑 $x,y>0$ 变成 $z>0,0<w<z$。联合密度里即使不再显含 $w$，求边缘仍要积分区间长度，不能直接把它当成 $Z$ 的密度。

独立性（independence）的定义是联合概率对可测集合分解。密度存在时，等价于 $f_{X,Y}(x,y)=f_X(x)f_Y(y)$ 几乎处处成立。这里包括支撑的指示函数。内部表达式能够分离，但支撑不是合适的直积时，不能草率认定独立。相反，在整个直积支撑上分解为两个已经归一化的边缘密度，可以直接推出独立。

二维变换 $u=g(x,y),v=h(x,y)$ 的常规公式要求在所考虑区域内一一可逆、可微且雅可比非零。若逆变换为 $(x,y)=T^{-1}(u,v)$，则
\[
f_{U,V}(u,v)=f_{X,Y}(T^{-1}(u,v))\left|\det D T^{-1}(u,v)\right|.
\]
题目若有多个逆像，需将各分支贡献相加。极坐标中的面积元素是 $r\,dr\,d\theta$，因此从极坐标密度回到平面时除以 $r$。把这个方向搞反常使密度无法归一化。

零协方差只是二阶线性关系的信息。联合正态情形中零相关推出独立；两个边缘分别正态并不足够。判断独立时可用确定性函数关系找反例，例如 $Y^2=X^2$ 会使相应事件重复，破坏概率乘积。对应训练为习题~\ref{q:A4.23}、\ref{q:Y4.25}、\ref{q:B17.15}、\ref{q:B17.16}。

\section*{方法五：先用结构求期望，再求分布}
\addcontentsline{toc}{section}{方法五：先用结构求期望，再求分布}
期望的线性性不需要独立。若目标量能写成指示变量之和，通常可以不求整个分布就得到均值。例如成功数量 $T=\sum_iI_i$ 的期望为 $\sum_iP(I_i=1)$；即使事件相关，这个公式仍成立。求方差时才出现交叉协方差：
\[
\operatorname{Var}\left(\sum_iX_i\right)=\sum_i\operatorname{Var}(X_i)+2\sum_{i<j}\operatorname{Cov}(X_i,X_j).
\]
独立且二阶矩有限可使协方差为零，反过来零协方差一般不意味着独立。

非负变量可使用尾积分公式 $EX=\int_0^\infty P(X>t)dt$；非负整数变量对应尾和 $EX=\sum_{k\ge1}P(X\ge k)$。证明来自将变量写成指示函数积分并使用非负积分交换定理，因此即使期望为无穷也成立。这个方法特别适合最大值、最小值、停时和剩余寿命。

条件期望（conditional expectation）允许按一个容易控制的状态分层：先固定随机样本数、类别或强度，计算条件均值，再取外层期望。全方差公式为
\[
\operatorname{Var}(Y)=E[\operatorname{Var}(Y\mid X)]+\operatorname{Var}(E[Y\mid X]).
\]
前项是固定状态后的内部波动，后项是状态变化造成的均值波动。随机和中常见错误是只保留前项，漏掉随机个数本身导致的方差。

计算结果应保留量纲：均值的单位与变量相同，方差的单位是变量单位平方，标准差恢复原单位，相关系数没有单位。对协方差或方差出现负数，应检查减法和混合模型分解；但协方差本身可以为负。对应训练为习题~\ref{q:A5.9}、\ref{q:A5.25}、\ref{q:A6.4}、\ref{q:B17.5}。

\section*{方法六：生成函数与极限定理}
\addcontentsline{toc}{section}{方法六：生成函数与极限定理}
对非负整数变量，概率生成函数（probability generating function）为 $G_X(s)=Es^X$。在 $|s|\le1$ 范围通常有良好定义，系数给出概率。独立和的生成函数相乘，随机和则产生复合函数。在相应矩有限或以单侧极限解释时，$G'_X(1)=EX$，$G''_X(1)=E[X(X-1)]$，注意二阶导数不是直接给出 $EX^2$。

特征函数（characteristic function）为 $\varphi_X(t)=Ee^{itX}$，对任意实值随机变量都存在。独立性给出 $\varphi_{X+Y}=\varphi_X\varphi_Y$，仿射变换给出 $\varphi_{aX+b}(t)=e^{itb}\varphi_X(at)$。莱维连续性定理将特征函数的逐点极限与依分布收敛联系起来；使用时应确认极限在零点连续并确实对应一个概率分布。

大数定律（law of large numbers）关注平均值稳定到总体均值；中心极限定理（central limit theorem）关注偏差乘以 $\sqrt n$ 后的波动形状。有限正方差的独立同分布版本是
\[
\sqrt n\frac{\overline X-\mu}{\sigma}\xrightarrow dN(0,1).
\]
“样本均值近似正态”并不等于每个原始观测近似正态，也不等于未乘 $\sqrt n$ 的误差服从标准正态。两组样本均值之差的方差要分别除以各自容量。

依分布收敛的分布函数判据通常只对固定连续点成立。若阈值随 $n$ 移动，应使用夹逼或紧性论证。加入一个依概率趋于零的扰动不改变分布极限，这是斯卢茨基思想的常见形式，不要求扰动与原量独立。对应训练为习题~\ref{q:A6.29}、\ref{q:A6.30}、\ref{q:A7.21}、\ref{q:A7.25}。

\section*{方法七：抽样分布是统计推断的桥梁}
\addcontentsline{toc}{section}{方法七：抽样分布是统计推断的桥梁}
统计量（statistic）是计算公式不含未知参数的样本函数；它的分布可以依赖未知参数。抽样分布是把样本看成随机变量后，统计量在重复抽样中的分布。观测到具体数据后，统计量变成数值，但用于评价误差的参考分布来自观测前的模型。

正态样本中三条结构最重要：样本均值为正态；中心化残差平方和为卡方；均值与残差向量独立。令 $n\ge2$，则
\[
\overline X\sim N(\mu,\sigma^2/n),\qquad\frac{(n-1)S^2}{\sigma^2}\sim\chi^2_{n-1}.
\]
两者独立使学生化比值服从 $t_{n-1}$。两组独立正态样本的标准化方差之比服从 $F$ 分布。这里“独立”和“正态”各有作用，不能只因为公式里出现样本方差就自动使用 $t$ 或 $F$。

自由度是可独立变化方向的数目。用样本均值估计总体均值后，残差满足和为零的线性约束，故平方和只剩 $n-1$ 个正交正态方向。在回归中估计 $p$ 个满秩线性系数，残差自由度是 $n-p$；在分类计数中，总数和参数拟合也会消耗自由度，但适用条件需要另外检查。

顺序统计量从另一条途径提供抽样分布。最大值的分布来自“所有观测都不超过阈值”，最小值的分布来自“所有观测都超过阈值”。端点参数影响支撑时，极值往往比均值包含更多相关信息。对应训练为习题~\ref{q:B8.3}、\ref{q:B8.4}、\ref{q:B8.7}、\ref{q:B8.8}。

\section*{方法八：区分估计方法与估计性质}
\addcontentsline{toc}{section}{方法八：区分估计方法与估计性质}
矩估计（method of moments）通过总体矩等于样本矩构造参数估计。极大似然估计（maximum likelihood estimation）固定观察数据，在参数空间中最大化联合密度或质量函数。这两者是构造方法；无偏、一致、效率和均方误差则是评价性质。一个矩估计可以有偏，一个极大似然估计也可能没有有限期望，但仍在其他意义下有用。

均方误差分解为方差与偏差平方：
\[
E_\theta(T-\theta)^2=\operatorname{Var}_\theta(T)+[E_\theta T-\theta]^2.
\]
比较无偏估计时可以只比较方差；比较有偏与无偏估计时必须同时考虑偏差。参数空间受约束时，把估计投影回允许集合可能引入偏差，却减小每个允许真值下的平方误差。

证明一致性常先对样本矩用大数定律，再用连续映射定理。对比值要检查分母极限非零；对数与倒数要处理有限样本中未定义的事件。渐近正态性常由中心极限定理加德尔塔方法得到：若 $\sqrt n(T_n-\theta)\xrightarrow dN(0,v)$ 且 $g$ 在 $\theta$ 处可微，则
\[
\sqrt n[g(T_n)-g(\theta)]\xrightarrow dN(0,[g'(\theta)]^2v).
\]
若一阶导数为零，不能套用这个式子声称得到非退化极限，应继续分析更高阶项。

极大似然解题应写清参数空间和支撑指标，随后才求导。驻点还需判断是最大值，并检查边界和极大值是否达到。均匀分布端点估计正是最大值位于约束边界的典型例子。对应训练为习题~\ref{q:B9.5}、\ref{q:B10.2}、\ref{q:B10.4}、\ref{q:B17.22}。

\section*{方法九：充分性、完全性与信息下界}
\addcontentsline{toc}{section}{方法九：充分性、完全性与信息下界}
充分统计量（sufficient statistic）压缩了模型中关于参数的全部样本信息。因子分解定理要求联合密度写成参数只通过统计量进入的部分与不含参数部分之积。检查支撑时也应把指示函数包括进去。伯努利样本的成功总数、泊松样本的总计数、正态样本的两个平方和型统计量都是常见例子。

拉奥—布莱克化以 $E(U\mid T)$ 替代原无偏估计 $U$，保留期望并使方差不增大。条件期望首先需要作为样本函数可计算；充分性使有关条件分布不依赖未知参数。如果统计量进一步完全，即其函数在所有参数下期望为零只能几乎处处为零，则莱曼—谢费定理给出唯一的均匀最小方差无偏估计。

费希尔信息量（Fisher information）衡量似然对参数的局部敏感度。在正规可微模型中，得分期望为零，信息可写成得分平方的期望，也可写成负的对数似然二阶导数期望。克拉美—拉奥不等式给出无偏估计 $g(\theta)$ 的方差下界
\[
\operatorname{Var}_\theta(U)\ge\frac{[g'(\theta)]^2}{I_n(\theta)}.
\]
使用它需核对支撑不随参数产生未处理的边界项、微分积分交换等条件。它只是一个下界，不保证存在达到下界的估计。

充分性、完全性与信息量回答不同问题：前者关乎数据压缩，中者关乎无偏估计唯一性，后者关乎局部精度限制。它们之间不能仅凭名称相似互相替代。对应训练为习题~\ref{q:B11.3}、\ref{q:B11.4}、\ref{q:B11.6}、\ref{q:B11.8}。

\section*{方法十：区间估计从枢轴反演}
\addcontentsline{toc}{section}{方法十：区间估计从枢轴反演}
枢轴量（pivotal quantity）是样本与未知参数的函数，但其分布不含未知参数。构造区间的一般路线是先找到枢轴，再在其已知分布上选择概率为 $1-\alpha$ 的区域，最后解不等式把参数放在中间。例如正态均值的标准化量、正态方差的卡方量、指数样本和的伽马量都可以这样处理。

反演时注意单调方向。若 $a\le C/\theta\le b$ 且所有量正，则 $C/b\le\theta\le C/a$，分位数顺序应交换。很多方差区间错误正是把较小的卡方分位数放在下端点分母，从而得到端点颠倒。比值参数还需检查分子和分母的自由度。

置信水平（confidence level）是重复抽样中随机区间覆盖固定参数的概率。它不自动等于给定数据以后参数位于已实现区间的后验概率。贝叶斯可信区间（credible interval）则基于后验分布积分。两者数值上有时相同，但逻辑解释不同；一个区间可以有 $95\%$ 的频率覆盖，而在某个特定先验与数据下的后验概率远小于 $95\%$。

样本量规划通常解误差半径不等式。若标准差已知，半径为 $z\sigma/\sqrt n$，所以将半径减半需四倍样本量。未知比例可以用方差上界 $1/4$ 规划，但这仍通常基于正态近似，不能直接声称对所有有限样本精确覆盖。对应训练为习题~\ref{q:M-C22-1}、\ref{q:B13.1}、\ref{q:M-P10-5}、\ref{q:B17.11}。

\section*{方法十一：检验的四个概率方向}
\addcontentsline{toc}{section}{方法十一：检验的四个概率方向}
假设检验（hypothesis testing）由零假设、备择、统计量和拒绝域构成。第一类错误是零假设真实时拒绝，第二类错误是某个备择真实时不拒绝。功效（power）等于备择下拒绝概率，是参数的函数而非通常意义上的一个固定数。对于复合零假设，检验水平要求其中所有参数的误拒概率都不超过 $\alpha$。

$p$ 值是在零假设参考分布下，出现至少同样极端结果的概率。它不是零假设为真的概率，也不是给定这一次拒绝后判断错误的概率。后一类逆向概率需要零假设与备择的先验权重和相应似然。显著性阈值应预先规定；看到数据后更改单侧、双侧方向，会改变原定错误控制。

构造拒绝域后，再在备择分布下积分得到功效。不能把零假设临界值重新以备择均值为中心设定，否则就改变了检验规则。中央区间通常是双侧检验的接受区域，其补集才是拒绝区域。离散检验中拒绝域大小常无法恰好等于 $\alpha$，除非允许边界随机化。

不拒绝零假设不等于证明其为真。若要证明实际差异足够小，应预先规定等效界并使用等效检验；若要比较多个假设，还要控制多重检验产生的至少一次误报概率。置信区间能同时显示差异大小与不确定性，通常比仅报“显著／不显著”更有解释力。对应训练为习题~\ref{q:M-P8-3}、\ref{q:M-P9-1}、\ref{q:B14.3}、\ref{q:B14.4}。

\section*{方法十二：分类计数检验的建模检查}
\addcontentsline{toc}{section}{方法十二：分类计数检验的建模检查}
皮尔逊拟合优度统计量为 $Q=\sum_j(O_j-E_j)^2/E_j$。它衡量观察计数与零假设期望计数的偏离。多项独立抽样、正类别概率和足够大的样本保证常用卡方近似有依据。若零假设完全指定 $k$ 个概率，自由度为 $k-1$；若从同一分组多项似然适当估计 $r$ 个参数，常规正则条件下自由度为 $k-1-r$。

“每个期望计数至少五”是常见经验检查，并不是覆盖所有情形的精确数学边界。极稀疏类别可事先合理合并，也可用精确条件法或模拟零分布。合并类别会改变自由度，且不应在看见哪个类别偏离最大后任意调整规则并仍沿用未经调整的显著性水平。

列联表独立性检验中，期望计数为行和乘列和除以总数。对于 $r\times c$ 表，自由度通常为 $(r-1)(c-1)$。两组类别分布同质性与联合表独立性在数值形式上相同，但抽样设计和零假设叙述不同。固定边际和的二乘二表，可用超几何分布作费希尔精确检验。

实际解释应区分统计偏离与产生偏离的原因。拟合优度拒绝只能说明指定模型与数据不相容，不能单凭它断言欺骗、因果效应或某种机制成立。包装内相关、重复测量和选择性收集都会使简单独立计数模型失效。对应训练为习题~\ref{q:M-E2-II5}、\ref{q:B15.2}、\ref{q:B15.4}、\ref{q:M-F-8}。

\section*{方法十三：回归中的线性与随机性}
\addcontentsline{toc}{section}{方法十三：回归中的线性与随机性}
线性回归（linear regression）中的“线性”指对未知系数线性。模型可以包含 $x^2,x^3,\log x,|x-c|$ 等已知基函数。最小二乘（least squares）本身是一个确定性的最优化问题，不需概率模型就能求出拟合系数；要解释其无偏性、方差或检验分布，才需要对误差建立随机模型。

满列秩设计矩阵保证最小二乘解唯一。正规方程 $X^\mathsf T(Y-X\widehat\beta)=0$ 表明残差与设计空间正交。包含截距时，残差和为零，拟合线经过均值点。高斯—马尔可夫定理只需零均值、同方差和不相关误差，给出在线性无偏类中的最优性；精确 $t$、$F$ 分布则进一步使用正态误差。

对均值响应的区间与新个体预测区间不同。前者衡量回归曲线估计误差，后者还包含新观测的不可约噪声。两者的中心可相同，半径却差一个额外的误差方差项。观测横坐标越远离样本中心，预测标准误通常越大；在样本范围外还存在模型形式不确定性，这不是单纯增加公式中的标准误就能完全解决的。

变换响应会改变拟合目标。最小化对数残差平方与最小化原尺度残差平方是不同的问题；对数正态模型中把拟合的对数均值直接取指数得到条件中位数，均值还需相应校正。回归相关不能自动解释为因果，遗漏变量可能改变斜率。对应训练为习题~\ref{q:M-C26-1c}、\ref{q:B16.2}、\ref{q:B16.4}、\ref{q:B17.23}。

\section*{方法十四：重抽样与模拟误差}
\addcontentsline{toc}{section}{方法十四：重抽样与模拟误差}
非参数自助法（nonparametric bootstrap）从经验分布有放回重抽，参数自助法（parametric bootstrap）从拟合的参数模型重抽。每次重抽样容量通常与原样本相同，从而模拟同一统计量在同样样本规模下的波动。增加重抽样次数 $B$ 降低的是模拟误差，不会把原始的 $n$ 个观测变成 $nB$ 个独立总体观测。

百分位区间直接取自助统计量的两个经验分位数；基本区间先估计统计量相对原估计的误差分布，再作反演。学生化、偏差校正等方法又有不同构造。报告区间时应说明方法、重抽样方式、重抽样次数和分位数约定。对相关数据，普通逐点重抽会破坏依赖结构，需要与抽样设计相适应的方法。

自助法并非“完全无假设”。它依赖经验分布能在所关心的统计量层面近似真实抽样分布。边界参数、不光滑统计量、极值、重尾或小样本，都可能使普通方法失效。尤其参数自助法继承拟合模型的结构；若模型本身错误，模拟得再多也只是精确地模拟了错误模型。

蒙特卡洛积分把积分写成期望，再用样本平均估计。独立样本通常使标准误按平方根速率下降。方差缩减可利用负相关配对、控制变量或重要性抽样，但必须重新核对无偏性和方差存在性。不能仅凭模拟曲线“看起来符合”就把随机证据当作严格证明。对应训练为习题~\ref{q:B13.6}、\ref{q:M-PF-29}、\ref{q:B17.17}、\ref{q:B17.18}。

\section*{方法十五：解答的自检与复习路线}
\addcontentsline{toc}{section}{方法十五：解答的自检与复习路线}
概率题完成后，先核查事件或积分区域与文字是否一致，再检查归一化、支持范围与边界。当密度不积分为一时，不应先把常数“修到正确”，而应回查是不是漏掉雅可比、混淆圆周与圆内均匀、或者在联合变换后漏了边缘积分。条件概率若大于一，多半是分母或条件事件写错；所谓独立性若由一句“显然”带过，应回到定义或联合密度分解核实。

统计题完成后，核查总体模型、抽样独立性、已知与未知参数、统计量分布、自由度和尾方向。标准差与方差单位不同，标准误还要考虑样本容量。极大似然估计的数值并不能直接代替未知参数出现在精确枢轴中；一旦这样替代，常常需要说明近似依据。所有小数先由精确式计算，在最终结果处统一舍入，避免多步后验更新中的舍入传播。

建议第一轮先做事件、条件概率、分布与期望的基本计算；第二轮做联合变换、极限定理和抽样分布；第三轮做估计与检验，并把每个区间都追溯到所用枢轴；最后做综合题，练习从自然语言挑选合适模型。证明题与计算题交替进行，能避免只会套用符号而不知道条件何时失效。

错误记录宜包括“原判断、出错原因、修正后的条件或公式、一个反例”。例如把零相关当作独立，应补记联合正态的额外条件和 $Y=SX$ 的反例；把不拒绝当作证明相等，应补记功效和区间宽度；把均匀分布的极大似然用内点求导处理，应补记参数相关支撑。这样形成的复习记录比只抄最终答案更能迁移到新题。


\clearpage
\section*{常用分布：定义、矩与适用情形}
\addcontentsline{toc}{section}{常用分布：定义、矩与适用情形}
\subsection*{伯努利分布与二项分布}
伯努利变量 $X\sim\operatorname{Bern}(p)$ 取零或一，$P(X=1)=p$。它表示一次二分类试验，不要求实际结果的文字名称为“成功”或“失败”。
\[
EX=p,\quad\operatorname{Var}(X)=p(1-p),\quad G_X(s)=1-p+ps.
\]
若 $n$ 次试验独立且成功概率相同，则成功数 $S$ 服从二项分布：
\[
P(S=k)=\binom nkp^k(1-p)^{n-k},\quad k=0,\ldots,n,
\]
\[
ES=np,\quad\operatorname{Var}(S)=np(1-p),\quad G_S(s)=(1-p+ps)^n.
\]
二项模型有三个独立的检查点：次数固定、试验独立、单次概率相同。不同成功概率的独立和仍可用生成函数乘积，但一般不是单参数二项分布；无放回抽样通常具有负相关，常应使用超几何分布。
\subsection*{超几何分布}
总体规模 $N$，其中 $M$ 个特征对象，无放回抽取 $n$ 个，特征数量 $X$ 的质量函数为
\[
P(X=k)=\frac{\binom Mk\binom{N-M}{n-k}}{\binom Nn},
\]
\[
\max(0,n+M-N)\le k\le\min(n,M).
\]
令 $p=M/N$，当 $N>1$ 时，
\[
EX=np,\qquad\operatorname{Var}(X)=np(1-p)\frac{N-n}{N-1}.
\]
支撑范围来自不能抽到超过总体已有数量的对象。若 $n=N$，方差为零，体现普查没有抽样不确定性。若总体本身的数值被当作随机生成，另一个层次的总体随机性不包括在这个条件抽样方差中。
\subsection*{几何分布与负二项分布}
本文几何分布的基本约定是记录首次成功的试验编号：
\[
P(T=k)=p(1-p)^{k-1},\quad k\ge1,
\]
\[
ET=1/p,\qquad\operatorname{Var}(T)=(1-p)/p^2.
\]
若记录成功前的失败次数 $F=T-1$，则均值变为 $(1-p)/p$，方差不变。记录第 $r$ 次成功前的失败次数 $K$ 时，
\[
P(K=k)=\binom{k+r-1}{r-1}p^r(1-p)^k,\quad k\ge0,
\]
\[
EK=\frac{r(1-p)}p,\qquad\operatorname{Var}(K)=\frac{r(1-p)}{p^2}.
\]
记录总次数时应再加 $r$。分布名称不足以唯一确定参数约定；解题时应先写出质量函数或明确计数对象。
\subsection*{泊松分布}
\[
P(X=k)=e^{-\lambda}\frac{\lambda^k}{k!},\quad k=0,1,\ldots,
\]
\[
EX=\operatorname{Var}(X)=\lambda,\qquad G_X(s)=e^{\lambda(s-1)}.
\]
独立泊松变量之和仍为泊松，参数相加。泊松总数经过相互独立的分类后，各类别计数仍是相互独立泊松变量。这两条性质分别对应事件来源叠加和事件流分裂。均值等于方差是泊松模型的性质，不是所有计数变量的性质；混合强度通常使方差大于均值。

\subsection*{均匀分布与指数分布}
$X\sim U(a,b)$，$a<b$ 时，
\[
f(x)=\frac{\mathbf1_{(a,b)}(x)}{b-a},\qquad EX=\frac{a+b}2,
\]
\[
\operatorname{Var}(X)=\frac{(b-a)^2}{12}.
\]
端点是否包含不影响连续分布，但在参数相关支撑的极大似然问题中，应明确采用的密度版本以讨论最大值是否达到。

$T\sim\operatorname{Exp}(\lambda)$，$\lambda>0$ 时，
\[
f(t)=\lambda e^{-\lambda t}\mathbf1_{(0,\infty)}(t),\quad P(T>t)=e^{-\lambda t}\quad(t\ge0),
\]
\[
ET=1/\lambda,\quad\operatorname{Var}(T)=1/\lambda^2.
\]
指数分布满足无记忆性 $P(T>s+t\mid T>s)=P(T>t)$，这里 $s,t\ge0$。恒定风险率模型可导出指数寿命，但不能从“等待时间”这个词自动推断指数分布。

\subsection*{伽马分布与贝塔分布}
本文伽马采用形状 $a>0$、率 $b>0$：
\[
f(x)=\frac{b^a}{\Gamma(a)}x^{a-1}e^{-bx}\mathbf1_{(0,\infty)}(x),
\]
\[
EX=a/b,\qquad\operatorname{Var}(X)=a/b^2.
\]
相同率的独立伽马变量相加，形状相加；不同率时一般不保持伽马形式。$n$ 个同率指数变量之和是形状 $n$ 的伽马变量。

贝塔分布采用两个正形状参数：
\[
f(x)=\frac{x^{a-1}(1-x)^{b-1}}{B(a,b)},\quad0<x<1,
\]
\[
B(a,b)=\frac{\Gamma(a)\Gamma(b)}{\Gamma(a+b)},\quad EX=\frac a{a+b},
\]
\[
\operatorname{Var}(X)=\frac{ab}{(a+b)^2(a+b+1)}.
\]
它既用于比例的先验，也出现在均匀顺序统计量和同率伽马变量份额的分布中。分布参数为 $a,b$ 时，密度幂次为 $a-1,b-1$，更新共轭参数时要避免差一。

\subsection*{正态分布与三种抽样分布}
$X\sim N(\mu,\sigma^2)$ 时，标准化 $Z=(X-\mu)/\sigma\sim N(0,1)$。正态密度为
\[
f(x)=\frac1{\sqrt{2\pi}\sigma}\exp\left[-\frac{(x-\mu)^2}{2\sigma^2}\right].
\]
独立正态变量线性组合仍为正态；联合正态向量的任意线性组合也为正态。仅有各边缘正态不够。

若 $Z_1,\ldots,Z_\nu$ 独立标准正态，则 $V=\sum_iZ_i^2\sim\chi^2_\nu$，其均值为 $\nu$，方差为 $2\nu$。若 $Z\sim N(0,1)$ 与 $V\sim\chi^2_\nu$ 独立，则 $T=Z/\sqrt{V/\nu}\sim t_\nu$。若两个独立卡方变量自由度为 $a,b$，则标准化比值 $(V_1/a)/(V_2/b)\sim F_{a,b}$。比值中每个卡方变量都必须先除以自己的自由度。

\clearpage
\section*{估计与检验公式的条件清单}
\addcontentsline{toc}{section}{估计与检验公式的条件清单}
\subsection*{一个正态总体的均值}
方差已知时使用 $Z=\sqrt n(\overline X-\mu_0)/\sigma$。双侧水平 $\alpha$ 拒绝域为 $|Z|>z_{1-\alpha/2}$，相应均值区间为
\[
\overline X\pm z_{1-\alpha/2}\frac\sigma{\sqrt n}.
\]
方差未知时，正态总体下用 $T=\sqrt n(\overline X-\mu_0)/S\sim t_{n-1}$，相应区间将 $z,\sigma$ 替为 $t_{n-1;1-\alpha/2},S$。在非正态有限方差总体中，学生化均值通常只是渐近标准正态，不能直接声称精确 $t$ 分布。
\subsection*{一个正态总体的方差}
未知均值时用 $Q=(n-1)S^2/\sigma_0^2\sim\chi^2_{n-1}$。右侧方差检验拒绝大 $Q$，左侧检验拒绝小 $Q$。等尾双侧区间为
\[
\left[\frac{(n-1)S^2}{\chi^2_{n-1;1-\alpha/2}},\frac{(n-1)S^2}{\chi^2_{n-1;\alpha/2}}\right].
\]
若均值为已知常数，改用相对该常数的平方和，自由度为 $n$。不能先用样本均值中心化又声称保留了 $n$ 个自由度。
\subsection*{两组独立正态均值}
共同方差未知时使用合并方差
\[
S_p^2=\frac{(n-1)S_1^2+(m-1)S_2^2}{n+m-2}.
\]
均值差的标准误为 $S_p\sqrt{1/n+1/m}$，自由度为 $n+m-2$。若不假设共同方差，使用两组各自估计方差构成标准误，并用韦尔奇近似自由度。配对数据先构造差值样本，再作单样本分析；两独立样本公式不适用于任意前后测量。
\subsection*{贝塔—二项与伽马—泊松更新}
$\theta\sim\operatorname{Beta}(a,b)$，观测 $n$ 次伯努利中的 $s$ 次成功后，后验为 $\operatorname{Beta}(a+s,b+n-s)$，下一次成功概率为 $(a+s)/(a+b+n)$。

$\lambda\sim\operatorname{Gamma}(a,b)$，观测独立的 $X_i\sim\operatorname{Pois}(t_i\lambda)$ 后，后验为
\[
\operatorname{Gamma}\left(a+\sum_iX_i,\ b+\sum_it_i\right).
\]
指数寿命 $X_i\sim\operatorname{Exp}(\lambda)$ 配伽马先验时，形状增加样本容量，率增加总寿命：$\operatorname{Gamma}(a+n,b+\sum_iX_i)$。泊松计数与指数等待对参数的作用不同，不要把两条更新公式交换。
\subsection*{一元回归的常用量}
令 $S_{xx}=\sum_i(x_i-\overline x)^2$、$S_{xy}=\sum_i(x_i-\overline x)(y_i-\overline y)$。含截距且 $S_{xx}>0$ 时，
\[
\widehat a=S_{xy}/S_{xx},\qquad\widehat b=\overline y-\widehat a\overline x.
\]
独立零均值同方差误差下，斜率方差为 $\sigma^2/S_{xx}$，截距方差为 $\sigma^2(1/n+\overline x^2/S_{xx})$。正态误差且 $n>2$ 时，残差方差估计为 $s^2=SSE/(n-2)$，再用 $t_{n-2}$ 分位数构造精确区间。误差异方差或相关时，上述方差表达式需要修改。

\clearpage
\section*{常用分位数表与计算约定}
\addcontentsline{toc}{section}{常用分位数表与计算约定}
以下所有分位数都采用下侧概率约定 $P(X\le q_p)=p$，保留四位小数。表格用于手工练习，正式计算可直接按分布函数求逆。表中舍入可能使临界附近的结果略有差异，应以未舍入表达式为准。


\subsection*{标准正态分位数}

\begin{center}\begin{tabular}{cc}\toprule 下侧概率 $p$ & $z_p$\\\midrule

0.800 & 0.8416\\

0.850 & 1.0364\\

0.900 & 1.2816\\

0.950 & 1.6449\\

0.960 & 1.7507\\

0.975 & 1.9600\\

0.990 & 2.3263\\

0.995 & 2.5758\\

0.999 & 3.0902\\

\bottomrule\end{tabular}\end{center}

\subsection*{学生 t 分位数}

\begin{center}\begin{tabular}{rrrr}\toprule $\nu$ & $t_{\nu;0.95}$ & $t_{\nu;0.975}$ & $t_{\nu;0.99}$\\\midrule

1 & 6.3138 & 12.7062 & 31.8205\\

2 & 2.9200 & 4.3027 & 6.9646\\

3 & 2.3534 & 3.1824 & 4.5407\\

4 & 2.1318 & 2.7764 & 3.7469\\

5 & 2.0150 & 2.5706 & 3.3649\\

6 & 1.9432 & 2.4469 & 3.1427\\

7 & 1.8946 & 2.3646 & 2.9980\\

8 & 1.8595 & 2.3060 & 2.8965\\

9 & 1.8331 & 2.2622 & 2.8214\\

10 & 1.8125 & 2.2281 & 2.7638\\

12 & 1.7823 & 2.1788 & 2.6810\\

15 & 1.7531 & 2.1314 & 2.6025\\

19 & 1.7291 & 2.0930 & 2.5395\\

20 & 1.7247 & 2.0860 & 2.5280\\

25 & 1.7081 & 2.0595 & 2.4851\\

30 & 1.6973 & 2.0423 & 2.4573\\

40 & 1.6839 & 2.0211 & 2.4233\\

60 & 1.6706 & 2.0003 & 2.3901\\

100 & 1.6602 & 1.9840 & 2.3642\\

\bottomrule\end{tabular}\end{center}

\subsection*{卡方下侧分位数}

\begin{center}\begin{tabular}{rrrrr}\toprule $\nu$ & $0.025$ & $0.05$ & $0.95$ & $0.975$\\\midrule

1 & 0.0010 & 0.0039 & 3.8415 & 5.0239\\

2 & 0.0506 & 0.1026 & 5.9915 & 7.3778\\

3 & 0.2158 & 0.3518 & 7.8147 & 9.3484\\

4 & 0.4844 & 0.7107 & 9.4877 & 11.1433\\

5 & 0.8312 & 1.1455 & 11.0705 & 12.8325\\

6 & 1.2373 & 1.6354 & 12.5916 & 14.4494\\

7 & 1.6899 & 2.1673 & 14.0671 & 16.0128\\

8 & 2.1797 & 2.7326 & 15.5073 & 17.5345\\

9 & 2.7004 & 3.3251 & 16.9190 & 19.0228\\

10 & 3.2470 & 3.9403 & 18.3070 & 20.4832\\

11 & 3.8157 & 4.5748 & 19.6751 & 21.9200\\

12 & 4.4038 & 5.2260 & 21.0261 & 23.3367\\

15 & 6.2621 & 7.2609 & 24.9958 & 27.4884\\

19 & 8.9065 & 10.1170 & 30.1435 & 32.8523\\

20 & 9.5908 & 10.8508 & 31.4104 & 34.1696\\

30 & 16.7908 & 18.4927 & 43.7730 & 46.9792\\

40 & 24.4330 & 26.5093 & 55.7585 & 59.3417\\

60 & 40.4817 & 43.1880 & 79.0819 & 83.2977\\

\bottomrule\end{tabular}\end{center}

\clearpage\subsection*{F 分布的 $0.95$ 下侧分位数}

行表示分子自由度，列表示分母自由度。

\begin{center}\begin{tabular}{rrrrr}\toprule $a\backslash b$ & $5$ & $10$ & $20$ & $60$\\\midrule

1 & 6.6079 & 4.9646 & 4.3512 & 4.0012\\

2 & 5.7861 & 4.1028 & 3.4928 & 3.1504\\

3 & 5.4095 & 3.7083 & 3.0984 & 2.7581\\

4 & 5.1922 & 3.4780 & 2.8661 & 2.5252\\

5 & 5.0503 & 3.3258 & 2.7109 & 2.3683\\

10 & 4.7351 & 2.9782 & 2.3479 & 1.9926\\

20 & 4.5581 & 2.7740 & 2.1242 & 1.7480\\

\bottomrule\end{tabular}\end{center}

\clearpage\section*{八组图解：把公式与图形对应起来}\addcontentsline{toc}{section}{八组图解：把公式与图形对应起来}

\subsection*{图解一：圆周与圆内的均匀取点}

\begin{center}\begin{tikzpicture}\begin{axis}[width=0.87\linewidth,height=0.69\linewidth,axis lines=left,tick label style={font=\small},label style={font=\small},legend style={font=\footnotesize,draw=none},xlabel={$x$},ylabel={横坐标密度},ymin=0,ymax=1.8,legend pos=north west]
\addplot[black,thick] coordinates {(-0.97000,1.309352) (-0.95916,1.125337) (-0.94832,1.003171) (-0.93749,0.914623) (-0.92665,0.846726) (-0.91581,0.792581) (-0.90497,0.748135) (-0.89413,0.710832) (-0.88330,0.678966) (-0.87246,0.651355) (-0.86162,0.627145) (-0.85078,0.605706) (-0.83994,0.586559) (-0.82911,0.569334) (-0.81827,0.553740) (-0.80743,0.539542) (-0.79659,0.526552) (-0.78575,0.514616) (-0.77492,0.503604) (-0.76408,0.493409) (-0.75324,0.483941) (-0.74240,0.475122) (-0.73156,0.466886) (-0.72073,0.459176) (-0.70989,0.451943) (-0.69905,0.445144) (-0.68821,0.438740) (-0.67737,0.432699) (-0.66654,0.426991) (-0.65570,0.421590) (-0.64486,0.416472) (-0.63402,0.411618) (-0.62318,0.407007) (-0.61235,0.402623) (-0.60151,0.398452) (-0.59067,0.394479) (-0.57983,0.390691) (-0.56899,0.387078) (-0.55816,0.383628) (-0.54732,0.380333) (-0.53648,0.377183) (-0.52564,0.374172) (-0.51480,0.371290) (-0.50397,0.368532) (-0.49313,0.365892) (-0.48229,0.363363) (-0.47145,0.360940) (-0.46061,0.358619) (-0.44978,0.356394) (-0.43894,0.354261) (-0.42810,0.352217) (-0.41726,0.350258) (-0.40642,0.348381) (-0.39559,0.346581) (-0.38475,0.344856) (-0.37391,0.343204) (-0.36307,0.341622) (-0.35223,0.340107) (-0.34140,0.338657) (-0.33056,0.337269) (-0.31972,0.335943) (-0.30888,0.334675) (-0.29804,0.333465) (-0.28721,0.332311) (-0.27637,0.331210) (-0.26553,0.330162) (-0.25469,0.329165) (-0.24385,0.328218) (-0.23302,0.327320) (-0.22218,0.326470) (-0.21134,0.325666) (-0.20050,0.324908) (-0.18966,0.324194) (-0.17883,0.323525) (-0.16799,0.322899) (-0.15715,0.322315) (-0.14631,0.321773) (-0.13547,0.321272) (-0.12464,0.320811) (-0.11380,0.320391) (-0.10296,0.320011) (-0.09212,0.319669) (-0.08128,0.319367) (-0.07045,0.319103) (-0.05961,0.318877) (-0.04877,0.318689) (-0.03793,0.318539) (-0.02709,0.318427) (-0.01626,0.318352) (-0.00542,0.318315) (0.00542,0.318315) (0.01626,0.318352) (0.02709,0.318427) (0.03793,0.318539) (0.04877,0.318689) (0.05961,0.318877) (0.07045,0.319103) (0.08128,0.319367) (0.09212,0.319669) (0.10296,0.320011) (0.11380,0.320391) (0.12464,0.320811) (0.13547,0.321272) (0.14631,0.321773) (0.15715,0.322315) (0.16799,0.322899) (0.17883,0.323525) (0.18966,0.324194) (0.20050,0.324908) (0.21134,0.325666) (0.22218,0.326470) (0.23302,0.327320) (0.24385,0.328218) (0.25469,0.329165) (0.26553,0.330162) (0.27637,0.331210) (0.28721,0.332311) (0.29804,0.333465) (0.30888,0.334675) (0.31972,0.335943) (0.33056,0.337269) (0.34140,0.338657) (0.35223,0.340107) (0.36307,0.341622) (0.37391,0.343204) (0.38475,0.344856) (0.39559,0.346581) (0.40642,0.348381) (0.41726,0.350258) (0.42810,0.352217) (0.43894,0.354261) (0.44978,0.356394) (0.46061,0.358619) (0.47145,0.360940) (0.48229,0.363363) (0.49313,0.365892) (0.50397,0.368532) (0.51480,0.371290) (0.52564,0.374172) (0.53648,0.377183) (0.54732,0.380333) (0.55816,0.383628) (0.56899,0.387078) (0.57983,0.390691) (0.59067,0.394479) (0.60151,0.398452) (0.61235,0.402623) (0.62318,0.407007) (0.63402,0.411618) (0.64486,0.416472) (0.65570,0.421590) (0.66654,0.426991) (0.67737,0.432699) (0.68821,0.438740) (0.69905,0.445144) (0.70989,0.451943) (0.72073,0.459176) (0.73156,0.466886) (0.74240,0.475122) (0.75324,0.483941) (0.76408,0.493409) (0.77492,0.503604) (0.78575,0.514616) (0.79659,0.526552) (0.80743,0.539542) (0.81827,0.553740) (0.82911,0.569334) (0.83994,0.586559) (0.85078,0.605706) (0.86162,0.627145) (0.87246,0.651355) (0.88330,0.678966) (0.89413,0.710832) (0.90497,0.748135) (0.91581,0.792581) (0.92665,0.846726) (0.93749,0.914623) (0.94832,1.003171) (0.95916,1.125337) (0.97000,1.309352)};
\addlegendentry{圆周均匀}
\addplot[black,dashed,thick] coordinates {(-0.97000,0.154765) (-0.95916,0.180073) (-0.94832,0.202002) (-0.93749,0.221558) (-0.92665,0.239325) (-0.91581,0.255674) (-0.90497,0.270863) (-0.89413,0.285078) (-0.88330,0.298457) (-0.87246,0.311109) (-0.86162,0.323119) (-0.85078,0.334556) (-0.83994,0.345476) (-0.82911,0.355929) (-0.81827,0.365952) (-0.80743,0.375582) (-0.79659,0.384848) (-0.78575,0.393774) (-0.77492,0.402384) (-0.76408,0.410698) (-0.75324,0.418734) (-0.74240,0.426506) (-0.73156,0.434030) (-0.72073,0.441317) (-0.70989,0.448380) (-0.69905,0.455229) (-0.68821,0.461874) (-0.67737,0.468322) (-0.66654,0.474583) (-0.65570,0.480663) (-0.64486,0.486569) (-0.63402,0.492307) (-0.62318,0.497884) (-0.61235,0.503305) (-0.60151,0.508574) (-0.59067,0.513697) (-0.57983,0.518677) (-0.56899,0.523519) (-0.55816,0.528226) (-0.54732,0.532803) (-0.53648,0.537252) (-0.52564,0.541576) (-0.51480,0.545779) (-0.50397,0.549863) (-0.49313,0.553832) (-0.48229,0.557686) (-0.47145,0.561430) (-0.46061,0.565064) (-0.44978,0.568591) (-0.43894,0.572014) (-0.42810,0.575333) (-0.41726,0.578551) (-0.40642,0.581670) (-0.39559,0.584690) (-0.38475,0.587614) (-0.37391,0.590442) (-0.36307,0.593177) (-0.35223,0.595820) (-0.34140,0.598371) (-0.33056,0.600832) (-0.31972,0.603205) (-0.30888,0.605489) (-0.29804,0.607687) (-0.28721,0.609798) (-0.27637,0.611825) (-0.26553,0.613767) (-0.25469,0.615625) (-0.24385,0.617401) (-0.23302,0.619095) (-0.22218,0.620708) (-0.21134,0.622240) (-0.20050,0.623692) (-0.18966,0.625064) (-0.17883,0.626358) (-0.16799,0.627573) (-0.15715,0.628710) (-0.14631,0.629769) (-0.13547,0.630751) (-0.12464,0.631656) (-0.11380,0.632484) (-0.10296,0.633236) (-0.09212,0.633913) (-0.08128,0.634513) (-0.07045,0.635038) (-0.05961,0.635488) (-0.04877,0.635862) (-0.03793,0.636162) (-0.02709,0.636386) (-0.01626,0.636536) (-0.00542,0.636610) (0.00542,0.636610) (0.01626,0.636536) (0.02709,0.636386) (0.03793,0.636162) (0.04877,0.635862) (0.05961,0.635488) (0.07045,0.635038) (0.08128,0.634513) (0.09212,0.633913) (0.10296,0.633236) (0.11380,0.632484) (0.12464,0.631656) (0.13547,0.630751) (0.14631,0.629769) (0.15715,0.628710) (0.16799,0.627573) (0.17883,0.626358) (0.18966,0.625064) (0.20050,0.623692) (0.21134,0.622240) (0.22218,0.620708) (0.23302,0.619095) (0.24385,0.617401) (0.25469,0.615625) (0.26553,0.613767) (0.27637,0.611825) (0.28721,0.609798) (0.29804,0.607687) (0.30888,0.605489) (0.31972,0.603205) (0.33056,0.600832) (0.34140,0.598371) (0.35223,0.595820) (0.36307,0.593177) (0.37391,0.590442) (0.38475,0.587614) (0.39559,0.584690) (0.40642,0.581670) (0.41726,0.578551) (0.42810,0.575333) (0.43894,0.572014) (0.44978,0.568591) (0.46061,0.565064) (0.47145,0.561430) (0.48229,0.557686) (0.49313,0.553832) (0.50397,0.549863) (0.51480,0.545779) (0.52564,0.541576) (0.53648,0.537252) (0.54732,0.532803) (0.55816,0.528226) (0.56899,0.523519) (0.57983,0.518677) (0.59067,0.513697) (0.60151,0.508574) (0.61235,0.503305) (0.62318,0.497884) (0.63402,0.492307) (0.64486,0.486569) (0.65570,0.480663) (0.66654,0.474583) (0.67737,0.468322) (0.68821,0.461874) (0.69905,0.455229) (0.70989,0.448380) (0.72073,0.441317) (0.73156,0.434030) (0.74240,0.426506) (0.75324,0.418734) (0.76408,0.410698) (0.77492,0.402384) (0.78575,0.393774) (0.79659,0.384848) (0.80743,0.375582) (0.81827,0.365952) (0.82911,0.355929) (0.83994,0.345476) (0.85078,0.334556) (0.86162,0.323119) (0.87246,0.311109) (0.88330,0.298457) (0.89413,0.285078) (0.90497,0.270863) (0.91581,0.255674) (0.92665,0.239325) (0.93749,0.221558) (0.94832,0.202002) (0.95916,0.180073) (0.97000,0.154765)};
\addlegendentry{圆内均匀}
\end{axis}\end{tikzpicture}\end{center}

取单位圆。圆周均匀指角度按弧长均匀；圆内均匀指平面面积均匀。前者横坐标密度在两端增大，后者在中心最大。两者均值都为零，但方差分别为 $1/2$ 和 $1/4$。图中为显示内部形状，仅绘制 $|x|\le0.97$；圆周密度在端点的可积奇性没有被改成有限值。

用截条面积可理解圆内密度：横坐标附近宽 $dx$ 的竖条长度约为 $2\sqrt{1-x^2}$，除以圆面积 $\pi$ 得到密度。圆周投影则需要对角度变换，两条对称弧贡献相加。参见习题~\ref{q:Y3.26}、\ref{q:A5.35}。

\clearpage\subsection*{图解二：二项计数与连续性修正}

\begin{center}\begin{tikzpicture}\begin{axis}[width=0.87\linewidth,height=0.69\linewidth,axis lines=left,tick label style={font=\small},label style={font=\small},legend style={font=\footnotesize,draw=none},xlabel={成功次数},ylabel={单位宽度概率或密度},xmin=3,xmax=28,ymin=0,legend pos=north east]
\addplot[ybar,bar width=4pt,fill=gray!25,draw=gray] coordinates {(3.00000,0.000028) (4.00000,0.000140) (5.00000,0.000551) (6.00000,0.001771) (7.00000,0.004770) (8.00000,0.010989) (9.00000,0.021978) (10.00000,0.038619) (11.00000,0.060185) (12.00000,0.083830) (13.00000,0.105017) (14.00000,0.118948) (15.00000,0.122347) (16.00000,0.114700) (17.00000,0.098314) (18.00000,0.077247) (19.00000,0.055757) (20.00000,0.037039) (21.00000,0.022677) (22.00000,0.012811) (23.00000,0.006684) (24.00000,0.003223) (25.00000,0.001436) (26.00000,0.000592) (27.00000,0.000225) (28.00000,0.000079)};
\addlegendentry{二项概率}
\addplot[black,thick] coordinates {(3.00000,0.000130) (3.12563,0.000149) (3.25126,0.000172) (3.37688,0.000198) (3.50251,0.000227) (3.62814,0.000261) (3.75377,0.000298) (3.87940,0.000341) (4.00503,0.000389) (4.13065,0.000444) (4.25628,0.000505) (4.38191,0.000574) (4.50754,0.000651) (4.63317,0.000737) (4.75879,0.000834) (4.88442,0.000942) (5.01005,0.001063) (5.13568,0.001197) (5.26131,0.001346) (5.38693,0.001511) (5.51256,0.001694) (5.63819,0.001896) (5.76382,0.002119) (5.88945,0.002365) (6.01508,0.002635) (6.14070,0.002932) (6.26633,0.003257) (6.39196,0.003613) (6.51759,0.004002) (6.64322,0.004427) (6.76884,0.004888) (6.89447,0.005390) (7.02010,0.005935) (7.14573,0.006524) (7.27136,0.007162) (7.39698,0.007850) (7.52261,0.008591) (7.64824,0.009388) (7.77387,0.010243) (7.89950,0.011160) (8.02513,0.012140) (8.15075,0.013187) (8.27638,0.014302) (8.40201,0.015489) (8.52764,0.016748) (8.65327,0.018083) (8.77889,0.019495) (8.90452,0.020986) (9.03015,0.022557) (9.15578,0.024209) (9.28141,0.025942) (9.40704,0.027759) (9.53266,0.029657) (9.65829,0.031638) (9.78392,0.033701) (9.90955,0.035845) (10.03518,0.038067) (10.16080,0.040366) (10.28643,0.042740) (10.41206,0.045186) (10.53769,0.047700) (10.66332,0.050278) (10.78894,0.052916) (10.91457,0.055608) (11.04020,0.058350) (11.16583,0.061135) (11.29146,0.063957) (11.41709,0.066808) (11.54271,0.069682) (11.66834,0.072571) (11.79397,0.075465) (11.91960,0.078357) (12.04523,0.081238) (12.17085,0.084098) (12.29648,0.086928) (12.42211,0.089718) (12.54774,0.092459) (12.67337,0.095141) (12.79899,0.097753) (12.92462,0.100286) (13.05025,0.102730) (13.17588,0.105076) (13.30151,0.107313) (13.42714,0.109434) (13.55276,0.111429) (13.67839,0.113291) (13.80402,0.115010) (13.92965,0.116580) (14.05528,0.117993) (14.18090,0.119245) (14.30653,0.120329) (14.43216,0.121240) (14.55779,0.121975) (14.68342,0.122530) (14.80905,0.122903) (14.93467,0.123091) (15.06030,0.123095) (15.18593,0.122914) (15.31156,0.122548) (15.43719,0.122001) (15.56281,0.121273) (15.68844,0.120369) (15.81407,0.119292) (15.93970,0.118047) (16.06533,0.116639) (16.19095,0.115075) (16.31658,0.113362) (16.44221,0.111506) (16.56784,0.109517) (16.69347,0.107401) (16.81910,0.105167) (16.94472,0.102826) (17.07035,0.100385) (17.19598,0.097856) (17.32161,0.095247) (17.44724,0.092568) (17.57286,0.089829) (17.69849,0.087041) (17.82412,0.084212) (17.94975,0.081353) (18.07538,0.078473) (18.20101,0.075581) (18.32663,0.072686) (18.45226,0.069798) (18.57789,0.066923) (18.70352,0.064071) (18.82915,0.061247) (18.95477,0.058461) (19.08040,0.055717) (19.20603,0.053022) (19.33166,0.050382) (19.45729,0.047802) (19.58291,0.045285) (19.70854,0.042837) (19.83417,0.040460) (19.95980,0.038157) (20.08543,0.035932) (20.21106,0.033785) (20.33668,0.031719) (20.46231,0.029735) (20.58794,0.027833) (20.71357,0.026013) (20.83920,0.024276) (20.96482,0.022621) (21.09045,0.021047) (21.21608,0.019553) (21.34171,0.018138) (21.46734,0.016800) (21.59296,0.015538) (21.71859,0.014348) (21.84422,0.013230) (21.96985,0.012181) (22.09548,0.011198) (22.22111,0.010279) (22.34673,0.009421) (22.47236,0.008622) (22.59799,0.007878) (22.72362,0.007188) (22.84925,0.006549) (22.97487,0.005957) (23.10050,0.005411) (23.22613,0.004908) (23.35176,0.004444) (23.47739,0.004019) (23.60302,0.003628) (23.72864,0.003271) (23.85427,0.002944) (23.97990,0.002646) (24.10553,0.002375) (24.23116,0.002128) (24.35678,0.001904) (24.48241,0.001701) (24.60804,0.001518) (24.73367,0.001352) (24.85930,0.001202) (24.98492,0.001068) (25.11055,0.000947) (25.23618,0.000838) (25.36181,0.000741) (25.48744,0.000654) (25.61307,0.000577) (25.73869,0.000507) (25.86432,0.000446) (25.98995,0.000391) (26.11558,0.000343) (26.24121,0.000300) (26.36683,0.000262) (26.49246,0.000228) (26.61809,0.000199) (26.74372,0.000173) (26.86935,0.000150) (26.99497,0.000130) (27.12060,0.000113) (27.24623,0.000097) (27.37186,0.000084) (27.49749,0.000072) (27.62312,0.000062) (27.74874,0.000054) (27.87437,0.000046) (28.00000,0.000039)};
\addlegendentry{正态近似}
\end{axis}\end{tikzpicture}\end{center}

本图采用 $n=50,p=0.3$。柱子表示各整数的概率，可理解为宽为一的直方图；曲线表示均值 $np$、方差 $np(1-p)$ 的正态密度。直接把一般离散单点概率与任意宽度的密度高度比较没有意义，这里以单位宽度建立可比尺度。

近似 $P(X\ge k)$ 时，将连续阈值放在 $k-1/2$，使连续区域覆盖完整整数柱。近似 $P(a\le X\le b)$ 时使用区间 $[a-1/2,b+1/2]$。修正改善许多中等容量计算，但不提供统一的误差保证。参见习题~\ref{q:M-P8-1ab}。

\clearpage\subsection*{图解三：零相关与依赖形状}

\begin{center}\begin{tikzpicture}\begin{axis}[width=0.87\linewidth,height=0.69\linewidth,axis lines=left,tick label style={font=\small},label style={font=\small},legend style={font=\footnotesize,draw=none},xlabel={$X$},ylabel={$Y$},xmin=-1.2,xmax=1.2,ymin=-1.2,ymax=1.2,axis equal image]
\addplot[only marks,mark=*,mark size=.65pt,black] coordinates {(1.00000,0.000000) (0.99939,0.034899) (0.99756,0.069756) (0.99452,0.104528) (0.99027,0.139173) (0.98481,0.173648) (0.97815,0.207912) (0.97030,0.241922) (0.96126,0.275637) (0.95106,0.309017) (0.93969,0.342020) (0.92718,0.374607) (0.91355,0.406737) (0.89879,0.438371) (0.88295,0.469472) (0.86603,0.500000) (0.84805,0.529919) (0.82904,0.559193) (0.80902,0.587785) (0.78801,0.615661) (0.76604,0.642788) (0.74314,0.669131) (0.71934,0.694658) (0.69466,0.719340) (0.66913,0.743145) (0.64279,0.766044) (0.61566,0.788011) (0.58779,0.809017) (0.55919,0.829038) (0.52992,0.848048) (0.50000,0.866025) (0.46947,0.882948) (0.43837,0.898794) (0.40674,0.913545) (0.37461,0.927184) (0.34202,0.939693) (0.30902,0.951057) (0.27564,0.961262) (0.24192,0.970296) (0.20791,0.978148) (0.17365,0.984808) (0.13917,0.990268) (0.10453,0.994522) (0.06976,0.997564) (0.03490,0.999391) (0.00000,1.000000) (-0.03490,0.999391) (-0.06976,0.997564) (-0.10453,0.994522) (-0.13917,0.990268) (-0.17365,0.984808) (-0.20791,0.978148) (-0.24192,0.970296) (-0.27564,0.961262) (-0.30902,0.951057) (-0.34202,0.939693) (-0.37461,0.927184) (-0.40674,0.913545) (-0.43837,0.898794) (-0.46947,0.882948) (-0.50000,0.866025) (-0.52992,0.848048) (-0.55919,0.829038) (-0.58779,0.809017) (-0.61566,0.788011) (-0.64279,0.766044) (-0.66913,0.743145) (-0.69466,0.719340) (-0.71934,0.694658) (-0.74314,0.669131) (-0.76604,0.642788) (-0.78801,0.615661) (-0.80902,0.587785) (-0.82904,0.559193) (-0.84805,0.529919) (-0.86603,0.500000) (-0.88295,0.469472) (-0.89879,0.438371) (-0.91355,0.406737) (-0.92718,0.374607) (-0.93969,0.342020) (-0.95106,0.309017) (-0.96126,0.275637) (-0.97030,0.241922) (-0.97815,0.207912) (-0.98481,0.173648) (-0.99027,0.139173) (-0.99452,0.104528) (-0.99756,0.069756) (-0.99939,0.034899) (-1.00000,0.000000) (-0.99939,-0.034899) (-0.99756,-0.069756) (-0.99452,-0.104528) (-0.99027,-0.139173) (-0.98481,-0.173648) (-0.97815,-0.207912) (-0.97030,-0.241922) (-0.96126,-0.275637) (-0.95106,-0.309017) (-0.93969,-0.342020) (-0.92718,-0.374607) (-0.91355,-0.406737) (-0.89879,-0.438371) (-0.88295,-0.469472) (-0.86603,-0.500000) (-0.84805,-0.529919) (-0.82904,-0.559193) (-0.80902,-0.587785) (-0.78801,-0.615661) (-0.76604,-0.642788) (-0.74314,-0.669131) (-0.71934,-0.694658) (-0.69466,-0.719340) (-0.66913,-0.743145) (-0.64279,-0.766044) (-0.61566,-0.788011) (-0.58779,-0.809017) (-0.55919,-0.829038) (-0.52992,-0.848048) (-0.50000,-0.866025) (-0.46947,-0.882948) (-0.43837,-0.898794) (-0.40674,-0.913545) (-0.37461,-0.927184) (-0.34202,-0.939693) (-0.30902,-0.951057) (-0.27564,-0.961262) (-0.24192,-0.970296) (-0.20791,-0.978148) (-0.17365,-0.984808) (-0.13917,-0.990268) (-0.10453,-0.994522) (-0.06976,-0.997564) (-0.03490,-0.999391) (-0.00000,-1.000000) (0.03490,-0.999391) (0.06976,-0.997564) (0.10453,-0.994522) (0.13917,-0.990268) (0.17365,-0.984808) (0.20791,-0.978148) (0.24192,-0.970296) (0.27564,-0.961262) (0.30902,-0.951057) (0.34202,-0.939693) (0.37461,-0.927184) (0.40674,-0.913545) (0.43837,-0.898794) (0.46947,-0.882948) (0.50000,-0.866025) (0.52992,-0.848048) (0.55919,-0.829038) (0.58779,-0.809017) (0.61566,-0.788011) (0.64279,-0.766044) (0.66913,-0.743145) (0.69466,-0.719340) (0.71934,-0.694658) (0.74314,-0.669131) (0.76604,-0.642788) (0.78801,-0.615661) (0.80902,-0.587785) (0.82904,-0.559193) (0.84805,-0.529919) (0.86603,-0.500000) (0.88295,-0.469472) (0.89879,-0.438371) (0.91355,-0.406737) (0.92718,-0.374607) (0.93969,-0.342020) (0.95106,-0.309017) (0.96126,-0.275637) (0.97030,-0.241922) (0.97815,-0.207912) (0.98481,-0.173648) (0.99027,-0.139173) (0.99452,-0.104528) (0.99756,-0.069756) (0.99939,-0.034899)};
\end{axis}\end{tikzpicture}\end{center}

图示单位圆周上的等角度点，用来展示支撑形状；它不是一份用于估计相关系数的独立随机样本。圆周均匀模型中两个坐标的协方差为零，但满足 $X^2+Y^2=1$ 的确定关系，因此不独立。给定横坐标的绝对值，纵坐标的绝对值已经确定。

相关系数只概括线性共变，不能识别所有非线性依赖。把散点图与模型支撑一起看，比单报一个接近零的相关系数更有信息。联合正态是零相关推出独立的特殊结构，圆周分布不具有该结构。参见习题~\ref{q:A5.35}、\ref{q:B17.15}。

\clearpage\subsection*{图解四：指数样本平均的正态化}

\begin{center}\begin{tikzpicture}\begin{axis}[width=0.87\linewidth,height=0.69\linewidth,axis lines=left,tick label style={font=\small},label style={font=\small},legend style={font=\footnotesize,draw=none},xlabel={标准化均值},ylabel={密度},ymin=0,ymax=1.1,legend pos=north east]
\addplot[black,dotted,thick] coordinates {(-3.50000,0.000000) (-3.46575,0.000000) (-3.43151,0.000000) (-3.39726,0.000000) (-3.36301,0.000000) (-3.32877,0.000000) (-3.29452,0.000000) (-3.26027,0.000000) (-3.22603,0.000000) (-3.19178,0.000000) (-3.15753,0.000000) (-3.12329,0.000000) (-3.08904,0.000000) (-3.05479,0.000000) (-3.02055,0.000000) (-2.98630,0.000000) (-2.95205,0.000000) (-2.91781,0.000000) (-2.88356,0.000000) (-2.84932,0.000000) (-2.81507,0.000000) (-2.78082,0.000000) (-2.74658,0.000000) (-2.71233,0.000000) (-2.67808,0.000000) (-2.64384,0.000000) (-2.60959,0.000000) (-2.57534,0.000000) (-2.54110,0.000000) (-2.50685,0.000000) (-2.47260,0.000000) (-2.43836,0.000000) (-2.40411,0.000000) (-2.36986,0.000000) (-2.33562,0.000000) (-2.30137,0.000000) (-2.26712,0.000000) (-2.23288,0.000000) (-2.19863,0.000000) (-2.16438,0.000000) (-2.13014,0.000000) (-2.09589,0.000000) (-2.06164,0.000000) (-2.02740,0.000000) (-1.99315,0.000000) (-1.95890,0.000000) (-1.92466,0.000000) (-1.89041,0.000000) (-1.85616,0.000000) (-1.82192,0.000000) (-1.78767,0.000000) (-1.75342,0.000000) (-1.71918,0.000000) (-1.68493,0.000000) (-1.65068,0.000000) (-1.61644,0.000000) (-1.58219,0.000000) (-1.54795,0.000000) (-1.51370,0.000000) (-1.47945,0.000000) (-1.44521,0.000000) (-1.41096,0.000000) (-1.37671,0.000000) (-1.34247,0.000000) (-1.30822,0.000000) (-1.27397,0.000000) (-1.23973,0.000000) (-1.20548,0.000000) (-1.17123,0.000000) (-1.13699,0.000000) (-1.10274,0.000000) (-1.06849,0.000000) (-1.03425,0.000000) (-1.00000,0.000000) (-0.96575,0.966333) (-0.93151,0.933800) (-0.89726,0.902362) (-0.86301,0.871982) (-0.82877,0.842625) (-0.79452,0.814257) (-0.76027,0.786843) (-0.72603,0.760353) (-0.69178,0.734754) (-0.65753,0.710017) (-0.62329,0.686113) (-0.58904,0.663014) (-0.55479,0.640693) (-0.52055,0.619123) (-0.48630,0.598279) (-0.45205,0.578137) (-0.41781,0.558673) (-0.38356,0.539864) (-0.34932,0.521688) (-0.31507,0.504125) (-0.28082,0.487152) (-0.24658,0.470752) (-0.21233,0.454903) (-0.17808,0.439588) (-0.14384,0.424788) (-0.10959,0.410487) (-0.07534,0.396667) (-0.04110,0.383313) (-0.00685,0.370408) (0.02740,0.357937) (0.06164,0.345887) (0.09589,0.334242) (0.13014,0.322989) (0.16438,0.312115) (0.19863,0.301607) (0.23288,0.291453) (0.26712,0.281641) (0.30137,0.272159) (0.33562,0.262996) (0.36986,0.254142) (0.40411,0.245586) (0.43836,0.237318) (0.47260,0.229328) (0.50685,0.221607) (0.54110,0.214146) (0.57534,0.206937) (0.60959,0.199970) (0.64384,0.193237) (0.67808,0.186732) (0.71233,0.180445) (0.74658,0.174370) (0.78082,0.168500) (0.81507,0.162827) (0.84932,0.157345) (0.88356,0.152048) (0.91781,0.146929) (0.95205,0.141982) (0.98630,0.137202) (1.02055,0.132583) (1.05479,0.128119) (1.08904,0.123806) (1.12329,0.119638) (1.15753,0.115610) (1.19178,0.111718) (1.22603,0.107956) (1.26027,0.104322) (1.29452,0.100810) (1.32877,0.097416) (1.36301,0.094136) (1.39726,0.090967) (1.43151,0.087904) (1.46575,0.084945) (1.50000,0.082085) (1.53425,0.079321) (1.56849,0.076651) (1.60274,0.074070) (1.63699,0.071577) (1.67123,0.069167) (1.70548,0.066838) (1.73973,0.064588) (1.77397,0.062414) (1.80822,0.060312) (1.84247,0.058282) (1.87671,0.056320) (1.91096,0.054424) (1.94521,0.052591) (1.97945,0.050821) (2.01370,0.049110) (2.04795,0.047456) (2.08219,0.045859) (2.11644,0.044315) (2.15068,0.042823) (2.18493,0.041381) (2.21918,0.039988) (2.25342,0.038642) (2.28767,0.037341) (2.32192,0.036084) (2.35616,0.034869) (2.39041,0.033695) (2.42466,0.032560) (2.45890,0.031464) (2.49315,0.030405) (2.52740,0.029381) (2.56164,0.028392) (2.59589,0.027436) (2.63014,0.026513) (2.66438,0.025620) (2.69863,0.024757) (2.73288,0.023924) (2.76712,0.023118) (2.80137,0.022340) (2.83562,0.021588) (2.86986,0.020861) (2.90411,0.020159) (2.93836,0.019480) (2.97260,0.018824) (3.00685,0.018191) (3.04110,0.017578) (3.07534,0.016986) (3.10959,0.016415) (3.14384,0.015862) (3.17808,0.015328) (3.21233,0.014812) (3.24658,0.014313) (3.28082,0.013831) (3.31507,0.013366) (3.34932,0.012916) (3.38356,0.012481) (3.41781,0.012061) (3.45205,0.011655) (3.48630,0.011262) (3.52055,0.010883) (3.55479,0.010517) (3.58904,0.010163) (3.62329,0.009820) (3.65753,0.009490) (3.69178,0.009170) (3.72603,0.008862) (3.76027,0.008563) (3.79452,0.008275) (3.82877,0.007996) (3.86301,0.007727) (3.89726,0.007467) (3.93151,0.007216) (3.96575,0.006973) (4.00000,0.006738)};
\addlegendentry{$n=1$}
\addplot[black,dashed,thick] coordinates {(-3.50000,0.000000) (-3.46575,0.000000) (-3.43151,0.000000) (-3.39726,0.000000) (-3.36301,0.000000) (-3.32877,0.000000) (-3.29452,0.000000) (-3.26027,0.000000) (-3.22603,0.000000) (-3.19178,0.000000) (-3.15753,0.000000) (-3.12329,0.000000) (-3.08904,0.000000) (-3.05479,0.000000) (-3.02055,0.000000) (-2.98630,0.000000) (-2.95205,0.000000) (-2.91781,0.000000) (-2.88356,0.000000) (-2.84932,0.000000) (-2.81507,0.000000) (-2.78082,0.000000) (-2.74658,0.000000) (-2.71233,0.000000) (-2.67808,0.000000) (-2.64384,0.000000) (-2.60959,0.000000) (-2.57534,0.000000) (-2.54110,0.000000) (-2.50685,0.000000) (-2.47260,0.000000) (-2.43836,0.000000) (-2.40411,0.000000) (-2.36986,0.000000) (-2.33562,0.000000) (-2.30137,0.000000) (-2.26712,0.000000) (-2.23288,0.000000) (-2.19863,0.000004) (-2.16438,0.000052) (-2.13014,0.000231) (-2.09589,0.000657) (-2.06164,0.001460) (-2.02740,0.002770) (-1.99315,0.004711) (-1.95890,0.007396) (-1.92466,0.010918) (-1.89041,0.015351) (-1.85616,0.020748) (-1.82192,0.027143) (-1.78767,0.034548) (-1.75342,0.042955) (-1.71918,0.052342) (-1.68493,0.062666) (-1.65068,0.073876) (-1.61644,0.085904) (-1.58219,0.098675) (-1.54795,0.112107) (-1.51370,0.126110) (-1.47945,0.140591) (-1.44521,0.155454) (-1.41096,0.170603) (-1.37671,0.185941) (-1.34247,0.201374) (-1.30822,0.216809) (-1.27397,0.232158) (-1.23973,0.247336) (-1.20548,0.262264) (-1.17123,0.276866) (-1.13699,0.291073) (-1.10274,0.304823) (-1.06849,0.318058) (-1.03425,0.330728) (-1.00000,0.342786) (-0.96575,0.354194) (-0.93151,0.364919) (-0.89726,0.374934) (-0.86301,0.384216) (-0.82877,0.392749) (-0.79452,0.400521) (-0.76027,0.407526) (-0.72603,0.413761) (-0.69178,0.419228) (-0.65753,0.423932) (-0.62329,0.427882) (-0.58904,0.431091) (-0.55479,0.433573) (-0.52055,0.435346) (-0.48630,0.436430) (-0.45205,0.436847) (-0.41781,0.436620) (-0.38356,0.435774) (-0.34932,0.434336) (-0.31507,0.432332) (-0.28082,0.429790) (-0.24658,0.426740) (-0.21233,0.423209) (-0.17808,0.419226) (-0.14384,0.414822) (-0.10959,0.410024) (-0.07534,0.404861) (-0.04110,0.399361) (-0.00685,0.393553) (0.02740,0.387464) (0.06164,0.381119) (0.09589,0.374545) (0.13014,0.367767) (0.16438,0.360810) (0.19863,0.353696) (0.23288,0.346448) (0.26712,0.339087) (0.30137,0.331634) (0.33562,0.324109) (0.36986,0.316530) (0.40411,0.308914) (0.43836,0.301279) (0.47260,0.293641) (0.50685,0.286012) (0.54110,0.278409) (0.57534,0.270842) (0.60959,0.263325) (0.64384,0.255868) (0.67808,0.248482) (0.71233,0.241176) (0.74658,0.233958) (0.78082,0.226837) (0.81507,0.219819) (0.84932,0.212911) (0.88356,0.206118) (0.91781,0.199446) (0.95205,0.192899) (0.98630,0.186481) (1.02055,0.180194) (1.05479,0.174043) (1.08904,0.168029) (1.12329,0.162154) (1.15753,0.156419) (1.19178,0.150826) (1.22603,0.145374) (1.26027,0.140065) (1.29452,0.134898) (1.32877,0.129873) (1.36301,0.124988) (1.39726,0.120244) (1.43151,0.115639) (1.46575,0.111171) (1.50000,0.106840) (1.53425,0.102642) (1.56849,0.098577) (1.60274,0.094643) (1.63699,0.090836) (1.67123,0.087155) (1.70548,0.083598) (1.73973,0.080162) (1.77397,0.076844) (1.80822,0.073642) (1.84247,0.070553) (1.87671,0.067575) (1.91096,0.064704) (1.94521,0.061939) (1.97945,0.059276) (2.01370,0.056712) (2.04795,0.054245) (2.08219,0.051872) (2.11644,0.049591) (2.15068,0.047398) (2.18493,0.045291) (2.21918,0.043267) (2.25342,0.041324) (2.28767,0.039459) (2.32192,0.037670) (2.35616,0.035953) (2.39041,0.034307) (2.42466,0.032730) (2.45890,0.031218) (2.49315,0.029769) (2.52740,0.028382) (2.56164,0.027054) (2.59589,0.025783) (2.63014,0.024567) (2.66438,0.023403) (2.69863,0.022290) (2.73288,0.021226) (2.76712,0.020209) (2.80137,0.019237) (2.83562,0.018308) (2.86986,0.017421) (2.90411,0.016574) (2.93836,0.015766) (2.97260,0.014994) (3.00685,0.014258) (3.04110,0.013555) (3.07534,0.012885) (3.10959,0.012246) (3.14384,0.011637) (3.17808,0.011056) (3.21233,0.010502) (3.24658,0.009975) (3.28082,0.009473) (3.31507,0.008994) (3.34932,0.008539) (3.38356,0.008105) (3.41781,0.007692) (3.45205,0.007299) (3.48630,0.006926) (3.52055,0.006570) (3.55479,0.006232) (3.58904,0.005910) (3.62329,0.005604) (3.65753,0.005314) (3.69178,0.005037) (3.72603,0.004775) (3.76027,0.004525) (3.79452,0.004288) (3.82877,0.004063) (3.86301,0.003849) (3.89726,0.003646) (3.93151,0.003454) (3.96575,0.003271) (4.00000,0.003097)};
\addlegendentry{$n=5$}
\addplot[black,thick] coordinates {(-3.50000,0.000012) (-3.46575,0.000017) (-3.43151,0.000023) (-3.39726,0.000031) (-3.36301,0.000041) (-3.32877,0.000054) (-3.29452,0.000071) (-3.26027,0.000092) (-3.22603,0.000119) (-3.19178,0.000153) (-3.15753,0.000195) (-3.12329,0.000247) (-3.08904,0.000311) (-3.05479,0.000390) (-3.02055,0.000486) (-2.98630,0.000602) (-2.95205,0.000741) (-2.91781,0.000908) (-2.88356,0.001107) (-2.84932,0.001342) (-2.81507,0.001620) (-2.78082,0.001945) (-2.74658,0.002326) (-2.71233,0.002767) (-2.67808,0.003278) (-2.64384,0.003866) (-2.60959,0.004541) (-2.57534,0.005311) (-2.54110,0.006187) (-2.50685,0.007179) (-2.47260,0.008299) (-2.43836,0.009556) (-2.40411,0.010964) (-2.36986,0.012534) (-2.33562,0.014278) (-2.30137,0.016208) (-2.26712,0.018339) (-2.23288,0.020680) (-2.19863,0.023245) (-2.16438,0.026046) (-2.13014,0.029095) (-2.09589,0.032401) (-2.06164,0.035976) (-2.02740,0.039829) (-1.99315,0.043969) (-1.95890,0.048404) (-1.92466,0.053139) (-1.89041,0.058181) (-1.85616,0.063533) (-1.82192,0.069197) (-1.78767,0.075175) (-1.75342,0.081466) (-1.71918,0.088066) (-1.68493,0.094973) (-1.65068,0.102179) (-1.61644,0.109677) (-1.58219,0.117456) (-1.54795,0.125506) (-1.51370,0.133813) (-1.47945,0.142360) (-1.44521,0.151132) (-1.41096,0.160108) (-1.37671,0.169270) (-1.34247,0.178594) (-1.30822,0.188057) (-1.27397,0.197634) (-1.23973,0.207299) (-1.20548,0.217026) (-1.17123,0.226786) (-1.13699,0.236551) (-1.10274,0.246291) (-1.06849,0.255976) (-1.03425,0.265577) (-1.00000,0.275065) (-0.96575,0.284408) (-0.93151,0.293577) (-0.89726,0.302544) (-0.86301,0.311279) (-0.82877,0.319755) (-0.79452,0.327946) (-0.76027,0.335824) (-0.72603,0.343367) (-0.69178,0.350551) (-0.65753,0.357354) (-0.62329,0.363756) (-0.58904,0.369738) (-0.55479,0.375285) (-0.52055,0.380380) (-0.48630,0.385011) (-0.45205,0.389166) (-0.41781,0.392837) (-0.38356,0.396016) (-0.34932,0.398698) (-0.31507,0.400879) (-0.28082,0.402558) (-0.24658,0.403735) (-0.21233,0.404412) (-0.17808,0.404594) (-0.14384,0.404286) (-0.10959,0.403495) (-0.07534,0.402231) (-0.04110,0.400503) (-0.00685,0.398324) (0.02740,0.395708) (0.06164,0.392667) (0.09589,0.389218) (0.13014,0.385377) (0.16438,0.381161) (0.19863,0.376589) (0.23288,0.371679) (0.26712,0.366450) (0.30137,0.360923) (0.33562,0.355118) (0.36986,0.349054) (0.40411,0.342752) (0.43836,0.336234) (0.47260,0.329519) (0.50685,0.322628) (0.54110,0.315581) (0.57534,0.308399) (0.60959,0.301100) (0.64384,0.293704) (0.67808,0.286230) (0.71233,0.278696) (0.74658,0.271119) (0.78082,0.263517) (0.81507,0.255906) (0.84932,0.248301) (0.88356,0.240717) (0.91781,0.233169) (0.95205,0.225670) (0.98630,0.218232) (1.02055,0.210868) (1.05479,0.203589) (1.08904,0.196403) (1.12329,0.189323) (1.15753,0.182354) (1.19178,0.175507) (1.22603,0.168788) (1.26027,0.162203) (1.29452,0.155758) (1.32877,0.149459) (1.36301,0.143309) (1.39726,0.137312) (1.43151,0.131472) (1.46575,0.125790) (1.50000,0.120269) (1.53425,0.114910) (1.56849,0.109713) (1.60274,0.104680) (1.63699,0.099810) (1.67123,0.095103) (1.70548,0.090558) (1.73973,0.086173) (1.77397,0.081946) (1.80822,0.077877) (1.84247,0.073962) (1.87671,0.070200) (1.91096,0.066587) (1.94521,0.063120) (1.97945,0.059797) (2.01370,0.056615) (2.04795,0.053569) (2.08219,0.050657) (2.11644,0.047874) (2.15068,0.045218) (2.18493,0.042684) (2.21918,0.040269) (2.25342,0.037969) (2.28767,0.035779) (2.32192,0.033697) (2.35616,0.031718) (2.39041,0.029839) (2.42466,0.028056) (2.45890,0.026365) (2.49315,0.024763) (2.52740,0.023245) (2.56164,0.021809) (2.59589,0.020451) (2.63014,0.019168) (2.66438,0.017955) (2.69863,0.016811) (2.73288,0.015732) (2.76712,0.014714) (2.80137,0.013756) (2.83562,0.012853) (2.86986,0.012004) (2.90411,0.011206) (2.93836,0.010455) (2.97260,0.009750) (3.00685,0.009089) (3.04110,0.008468) (3.07534,0.007886) (3.10959,0.007340) (3.14384,0.006829) (3.17808,0.006351) (3.21233,0.005904) (3.24658,0.005485) (3.28082,0.005094) (3.31507,0.004729) (3.34932,0.004388) (3.38356,0.004070) (3.41781,0.003773) (3.45205,0.003496) (3.48630,0.003239) (3.52055,0.002999) (3.55479,0.002775) (3.58904,0.002568) (3.62329,0.002374) (3.65753,0.002195) (3.69178,0.002028) (3.72603,0.001873) (3.76027,0.001729) (3.79452,0.001596) (3.82877,0.001472) (3.86301,0.001358) (3.89726,0.001251) (3.93151,0.001153) (3.96575,0.001062) (4.00000,0.000978)};
\addlegendentry{$n=30$}
\addplot[gray,very thick] coordinates {(-3.50000,0.000873) (-3.46575,0.000983) (-3.43151,0.001106) (-3.39726,0.001244) (-3.36301,0.001396) (-3.32877,0.001566) (-3.29452,0.001754) (-3.26027,0.001962) (-3.22603,0.002193) (-3.19178,0.002448) (-3.15753,0.002729) (-3.12329,0.003038) (-3.08904,0.003380) (-3.05479,0.003754) (-3.02055,0.004166) (-2.98630,0.004617) (-2.95205,0.005112) (-2.91781,0.005652) (-2.88356,0.006242) (-2.84932,0.006886) (-2.81507,0.007588) (-2.78082,0.008351) (-2.74658,0.009180) (-2.71233,0.010079) (-2.67808,0.011054) (-2.64384,0.012108) (-2.60959,0.013248) (-2.57534,0.014478) (-2.54110,0.015804) (-2.50685,0.017230) (-2.47260,0.018764) (-2.43836,0.020410) (-2.40411,0.022175) (-2.36986,0.024063) (-2.33562,0.026083) (-2.30137,0.028238) (-2.26712,0.030536) (-2.23288,0.032982) (-2.19863,0.035582) (-2.16438,0.038342) (-2.13014,0.041267) (-2.09589,0.044364) (-2.06164,0.047638) (-2.02740,0.051093) (-1.99315,0.054734) (-1.95890,0.058567) (-1.92466,0.062594) (-1.89041,0.066819) (-1.85616,0.071246) (-1.82192,0.075878) (-1.78767,0.080716) (-1.75342,0.085761) (-1.71918,0.091016) (-1.68493,0.096478) (-1.65068,0.102149) (-1.61644,0.108027) (-1.58219,0.114109) (-1.54795,0.120392) (-1.51370,0.126872) (-1.47945,0.133544) (-1.44521,0.140402) (-1.41096,0.147439) (-1.37671,0.154647) (-1.34247,0.162018) (-1.30822,0.169542) (-1.27397,0.177206) (-1.23973,0.185000) (-1.20548,0.192911) (-1.17123,0.200923) (-1.13699,0.209024) (-1.10274,0.217196) (-1.06849,0.225423) (-1.03425,0.233687) (-1.00000,0.241971) (-0.96575,0.250254) (-0.93151,0.258518) (-0.89726,0.266741) (-0.86301,0.274904) (-0.82877,0.282984) (-0.79452,0.290960) (-0.76027,0.298810) (-0.72603,0.306513) (-0.69178,0.314045) (-0.65753,0.321385) (-0.62329,0.328512) (-0.58904,0.335403) (-0.55479,0.342037) (-0.52055,0.348393) (-0.48630,0.354452) (-0.45205,0.360193) (-0.41781,0.365598) (-0.38356,0.370650) (-0.34932,0.375330) (-0.31507,0.379625) (-0.28082,0.383518) (-0.24658,0.386997) (-0.21233,0.390050) (-0.17808,0.392666) (-0.14384,0.394837) (-0.10959,0.396554) (-0.07534,0.397812) (-0.04110,0.398606) (-0.00685,0.398933) (0.02740,0.398793) (0.06164,0.398185) (0.09589,0.397112) (0.13014,0.395578) (0.16438,0.393588) (0.19863,0.391149) (0.23288,0.388270) (0.26712,0.384960) (0.30137,0.381231) (0.33562,0.377095) (0.36986,0.372567) (0.40411,0.367662) (0.43836,0.362396) (0.47260,0.356787) (0.50685,0.350853) (0.54110,0.344614) (0.57534,0.338088) (0.60959,0.331298) (0.64384,0.324263) (0.67808,0.317005) (0.71233,0.309547) (0.74658,0.301910) (0.78082,0.294116) (0.81507,0.286188) (0.84932,0.278147) (0.88356,0.270015) (0.91781,0.261813) (0.95205,0.253563) (0.98630,0.245285) (1.02055,0.236999) (1.05479,0.228725) (1.08904,0.220481) (1.12329,0.212285) (1.15753,0.204154) (1.19178,0.196104) (1.22603,0.188151) (1.26027,0.180309) (1.29452,0.172591) (1.32877,0.165010) (1.36301,0.157577) (1.39726,0.150302) (1.43151,0.143195) (1.46575,0.136265) (1.50000,0.129518) (1.53425,0.122960) (1.56849,0.116598) (1.60274,0.110435) (1.63699,0.104476) (1.67123,0.098722) (1.70548,0.093176) (1.73973,0.087838) (1.77397,0.082709) (1.80822,0.077788) (1.84247,0.073074) (1.87671,0.068566) (1.91096,0.064260) (1.94521,0.060154) (1.97945,0.056244) (2.01370,0.052527) (2.04795,0.048998) (2.08219,0.045652) (2.11644,0.042485) (2.15068,0.039492) (2.18493,0.036666) (2.21918,0.034003) (2.25342,0.031496) (2.28767,0.029140) (2.32192,0.026928) (2.35616,0.024855) (2.39041,0.022915) (2.42466,0.021101) (2.45890,0.019409) (2.49315,0.017831) (2.52740,0.016362) (2.56164,0.014996) (2.59589,0.013729) (2.63014,0.012554) (2.66438,0.011466) (2.69863,0.010460) (2.73288,0.009531) (2.76712,0.008674) (2.80137,0.007885) (2.83562,0.007160) (2.86986,0.006493) (2.90411,0.005882) (2.93836,0.005322) (2.97260,0.004810) (3.00685,0.004342) (3.04110,0.003914) (3.07534,0.003525) (3.10959,0.003171) (3.14384,0.002849) (3.17808,0.002557) (3.21233,0.002292) (3.24658,0.002052) (3.28082,0.001835) (3.31507,0.001639) (3.34932,0.001462) (3.38356,0.001303) (3.41781,0.001160) (3.45205,0.001031) (3.48630,0.000915) (3.52055,0.000812) (3.55479,0.000719) (3.58904,0.000636) (3.62329,0.000563) (3.65753,0.000497) (3.69178,0.000438) (3.72603,0.000386) (3.76027,0.000339) (3.79452,0.000298) (3.82877,0.000262) (3.86301,0.000229) (3.89726,0.000201) (3.93151,0.000176) (3.96575,0.000153) (4.00000,0.000134)};
\addlegendentry{标准正态}
\end{axis}\end{tikzpicture}\end{center}

令原始观测服从率为一的指数分布，图中画的是 $\sqrt n(\overline X-1)$ 的精确伽马变换密度，及标准正态密度。样本容量增大后偏斜减弱，但有限样本的分布仍有其自身支撑和尾部特征。

横轴已经中心化并乘以 $\sqrt n$。若画未缩放的样本均值，曲线会同时向均值收缩，无法直接看清中心极限定理描述的波动形状。图中 $n=1$ 的左端跳跃来自指数变量非负，而非绘图误差。参见习题~\ref{q:A6.30}、\ref{q:B17.20}。

\clearpage\subsection*{图解五：方差统计量的偏斜}

\begin{center}\begin{tikzpicture}\begin{axis}[width=0.87\linewidth,height=0.69\linewidth,axis lines=left,tick label style={font=\small},label style={font=\small},legend style={font=\footnotesize,draw=none},xlabel={$Q/\nu$},ylabel={密度},ymin=0,legend pos=north east]
\addplot[black,dotted,thick] coordinates {(0.00100,0.065455) (0.01698,0.263309) (0.03295,0.358163) (0.04893,0.426099) (0.06491,0.479136) (0.08089,0.522199) (0.09686,0.557918) (0.11284,0.587916) (0.12882,0.613285) (0.14479,0.634809) (0.16077,0.653075) (0.17675,0.668542) (0.19273,0.681573) (0.20870,0.692466) (0.22468,0.701469) (0.24066,0.708790) (0.25663,0.714607) (0.27261,0.719074) (0.28859,0.722325) (0.30457,0.724479) (0.32054,0.725638) (0.33652,0.725896) (0.35250,0.725335) (0.36847,0.724029) (0.38445,0.722046) (0.40043,0.719447) (0.41641,0.716285) (0.43238,0.712613) (0.44836,0.708476) (0.46434,0.703915) (0.48032,0.698969) (0.49629,0.693674) (0.51227,0.688062) (0.52825,0.682163) (0.54422,0.676006) (0.56020,0.669616) (0.57618,0.663016) (0.59216,0.656229) (0.60813,0.649274) (0.62411,0.642172) (0.64009,0.634939) (0.65606,0.627592) (0.67204,0.620146) (0.68802,0.612616) (0.70400,0.605013) (0.71997,0.597351) (0.73595,0.589641) (0.75193,0.581893) (0.76790,0.574117) (0.78388,0.566323) (0.79986,0.558518) (0.81584,0.550711) (0.83181,0.542909) (0.84779,0.535119) (0.86377,0.527347) (0.87974,0.519599) (0.89572,0.511880) (0.91170,0.504196) (0.92768,0.496550) (0.94365,0.488949) (0.95963,0.481394) (0.97561,0.473891) (0.99158,0.466442) (1.00756,0.459050) (1.02354,0.451719) (1.03952,0.444451) (1.05549,0.437248) (1.07147,0.430112) (1.08745,0.423046) (1.10342,0.416051) (1.11940,0.409129) (1.13538,0.402281) (1.15136,0.395508) (1.16733,0.388812) (1.18331,0.382194) (1.19929,0.375654) (1.21526,0.369193) (1.23124,0.362812) (1.24722,0.356511) (1.26320,0.350291) (1.27917,0.344152) (1.29515,0.338094) (1.31113,0.332117) (1.32711,0.326222) (1.34308,0.320408) (1.35906,0.314676) (1.37504,0.309024) (1.39101,0.303454) (1.40699,0.297965) (1.42297,0.292556) (1.43895,0.287227) (1.45492,0.281978) (1.47090,0.276808) (1.48688,0.271716) (1.50285,0.266703) (1.51883,0.261768) (1.53481,0.256910) (1.55079,0.252128) (1.56676,0.247422) (1.58274,0.242792) (1.59872,0.238236) (1.61469,0.233753) (1.63067,0.229344) (1.64665,0.225007) (1.66263,0.220742) (1.67860,0.216548) (1.69458,0.212424) (1.71056,0.208369) (1.72653,0.204382) (1.74251,0.200463) (1.75849,0.196611) (1.77447,0.192826) (1.79044,0.189105) (1.80642,0.185449) (1.82240,0.181856) (1.83837,0.178326) (1.85435,0.174858) (1.87033,0.171451) (1.88631,0.168105) (1.90228,0.164817) (1.91826,0.161589) (1.93424,0.158418) (1.95021,0.155304) (1.96619,0.152246) (1.98217,0.149243) (1.99815,0.146295) (2.01412,0.143401) (2.03010,0.140559) (2.04608,0.137769) (2.06205,0.135031) (2.07803,0.132343) (2.09401,0.129705) (2.10999,0.127116) (2.12596,0.124574) (2.14194,0.122080) (2.15792,0.119633) (2.17389,0.117232) (2.18987,0.114875) (2.20585,0.112564) (2.22183,0.110295) (2.23780,0.108070) (2.25378,0.105887) (2.26976,0.103745) (2.28574,0.101644) (2.30171,0.099583) (2.31769,0.097562) (2.33367,0.095579) (2.34964,0.093635) (2.36562,0.091728) (2.38160,0.089857) (2.39758,0.088023) (2.41355,0.086225) (2.42953,0.084461) (2.44551,0.082732) (2.46148,0.081036) (2.47746,0.079373) (2.49344,0.077743) (2.50942,0.076145) (2.52539,0.074578) (2.54137,0.073042) (2.55735,0.071536) (2.57332,0.070060) (2.58930,0.068613) (2.60528,0.067194) (2.62126,0.065804) (2.63723,0.064441) (2.65321,0.063105) (2.66919,0.061796) (2.68516,0.060513) (2.70114,0.059256) (2.71712,0.058023) (2.73310,0.056815) (2.74907,0.055632) (2.76505,0.054472) (2.78103,0.053336) (2.79700,0.052222) (2.81298,0.051131) (2.82896,0.050061) (2.84494,0.049014) (2.86091,0.047987) (2.87689,0.046982) (2.89287,0.045996) (2.90884,0.045031) (2.92482,0.044085) (2.94080,0.043158) (2.95678,0.042251) (2.97275,0.041361) (2.98873,0.040490) (3.00471,0.039637) (3.02068,0.038801) (3.03666,0.037982) (3.05264,0.037180) (3.06862,0.036395) (3.08459,0.035625) (3.10057,0.034872) (3.11655,0.034133) (3.13253,0.033410) (3.14850,0.032702) (3.16448,0.032009) (3.18046,0.031330) (3.19643,0.030664) (3.21241,0.030013) (3.22839,0.029375) (3.24437,0.028750) (3.26034,0.028138) (3.27632,0.027539) (3.29230,0.026953) (3.30827,0.026378) (3.32425,0.025816) (3.34023,0.025265) (3.35621,0.024725) (3.37218,0.024197) (3.38816,0.023680) (3.40414,0.023174) (3.42011,0.022678) (3.43609,0.022193) (3.45207,0.021718) (3.46805,0.021252) (3.48402,0.020797) (3.50000,0.020351)};
\addlegendentry{$\nu=3$}
\addplot[black,dashed,thick] coordinates {(0.00100,0.000000) (0.01698,0.000010) (0.03295,0.000130) (0.04893,0.000584) (0.06491,0.001671) (0.08089,0.003720) (0.09686,0.007062) (0.11284,0.012008) (0.12882,0.018828) (0.14479,0.027748) (0.16077,0.038937) (0.17675,0.052512) (0.19273,0.068533) (0.20870,0.087009) (0.22468,0.107898) (0.24066,0.131116) (0.25663,0.156540) (0.27261,0.184013) (0.28859,0.213351) (0.30457,0.244348) (0.32054,0.276780) (0.33652,0.310413) (0.35250,0.345005) (0.36847,0.380310) (0.38445,0.416084) (0.40043,0.452086) (0.41641,0.488083) (0.43238,0.523849) (0.44836,0.559173) (0.46434,0.593852) (0.48032,0.627701) (0.49629,0.660547) (0.51227,0.692235) (0.52825,0.722624) (0.54422,0.751592) (0.56020,0.779028) (0.57618,0.804843) (0.59216,0.828958) (0.60813,0.851312) (0.62411,0.871857) (0.64009,0.890560) (0.65606,0.907400) (0.67204,0.922366) (0.68802,0.935462) (0.70400,0.946698) (0.71997,0.956097) (0.73595,0.963689) (0.75193,0.969511) (0.76790,0.973608) (0.78388,0.976030) (0.79986,0.976834) (0.81584,0.976079) (0.83181,0.973829) (0.84779,0.970151) (0.86377,0.965115) (0.87974,0.958791) (0.89572,0.951253) (0.91170,0.942574) (0.92768,0.932827) (0.94365,0.922086) (0.95963,0.910424) (0.97561,0.897914) (0.99158,0.884625) (1.00756,0.870629) (1.02354,0.855992) (1.03952,0.840782) (1.05549,0.825061) (1.07147,0.808893) (1.08745,0.792337) (1.10342,0.775449) (1.11940,0.758285) (1.13538,0.740897) (1.15136,0.723335) (1.16733,0.705645) (1.18331,0.687872) (1.19929,0.670058) (1.21526,0.652242) (1.23124,0.634460) (1.24722,0.616748) (1.26320,0.599137) (1.27917,0.581657) (1.29515,0.564334) (1.31113,0.547194) (1.32711,0.530259) (1.34308,0.513551) (1.35906,0.497087) (1.37504,0.480886) (1.39101,0.464961) (1.40699,0.449327) (1.42297,0.433994) (1.43895,0.418973) (1.45492,0.404273) (1.47090,0.389900) (1.48688,0.375861) (1.50285,0.362160) (1.51883,0.348800) (1.53481,0.335785) (1.55079,0.323116) (1.56676,0.310792) (1.58274,0.298814) (1.59872,0.287181) (1.61469,0.275890) (1.63067,0.264940) (1.64665,0.254326) (1.66263,0.244046) (1.67860,0.234094) (1.69458,0.224468) (1.71056,0.215160) (1.72653,0.206166) (1.74251,0.197481) (1.75849,0.189098) (1.77447,0.181011) (1.79044,0.173214) (1.80642,0.165700) (1.82240,0.158462) (1.83837,0.151494) (1.85435,0.144788) (1.87033,0.138339) (1.88631,0.132137) (1.90228,0.126178) (1.91826,0.120454) (1.93424,0.114957) (1.95021,0.109681) (1.96619,0.104619) (1.98217,0.099765) (1.99815,0.095111) (2.01412,0.090650) (2.03010,0.086378) (2.04608,0.082286) (2.06205,0.078369) (2.07803,0.074621) (2.09401,0.071035) (2.10999,0.067605) (2.12596,0.064327) (2.14194,0.061193) (2.15792,0.058200) (2.17389,0.055340) (2.18987,0.052610) (2.20585,0.050004) (2.22183,0.047516) (2.23780,0.045144) (2.25378,0.042881) (2.26976,0.040723) (2.28574,0.038666) (2.30171,0.036706) (2.31769,0.034839) (2.33367,0.033060) (2.34964,0.031366) (2.36562,0.029754) (2.38160,0.028219) (2.39758,0.026758) (2.41355,0.025369) (2.42953,0.024048) (2.44551,0.022791) (2.46148,0.021596) (2.47746,0.020461) (2.49344,0.019382) (2.50942,0.018357) (2.52539,0.017383) (2.54137,0.016459) (2.55735,0.015581) (2.57332,0.014747) (2.58930,0.013956) (2.60528,0.013206) (2.62126,0.012494) (2.63723,0.011818) (2.65321,0.011178) (2.66919,0.010570) (2.68516,0.009995) (2.70114,0.009449) (2.71712,0.008932) (2.73310,0.008441) (2.74907,0.007977) (2.76505,0.007537) (2.78103,0.007121) (2.79700,0.006727) (2.81298,0.006353) (2.82896,0.006000) (2.84494,0.005665) (2.86091,0.005349) (2.87689,0.005049) (2.89287,0.004766) (2.90884,0.004498) (2.92482,0.004245) (2.94080,0.004005) (2.95678,0.003779) (2.97275,0.003565) (2.98873,0.003362) (3.00471,0.003171) (3.02068,0.002990) (3.03666,0.002820) (3.05264,0.002658) (3.06862,0.002506) (3.08459,0.002362) (3.10057,0.002226) (3.11655,0.002098) (3.13253,0.001977) (3.14850,0.001863) (3.16448,0.001755) (3.18046,0.001653) (3.19643,0.001557) (3.21241,0.001467) (3.22839,0.001381) (3.24437,0.001300) (3.26034,0.001224) (3.27632,0.001153) (3.29230,0.001085) (3.30827,0.001021) (3.32425,0.000961) (3.34023,0.000905) (3.35621,0.000851) (3.37218,0.000801) (3.38816,0.000754) (3.40414,0.000709) (3.42011,0.000667) (3.43609,0.000627) (3.45207,0.000590) (3.46805,0.000555) (3.48402,0.000522) (3.50000,0.000491)};
\addlegendentry{$\nu=10$}
\addplot[black,thick] coordinates {(0.00100,0.000000) (0.01698,0.000000) (0.03295,0.000000) (0.04893,0.000000) (0.06491,0.000000) (0.08089,0.000000) (0.09686,0.000000) (0.11284,0.000000) (0.12882,0.000000) (0.14479,0.000001) (0.16077,0.000003) (0.17675,0.000010) (0.19273,0.000027) (0.20870,0.000065) (0.22468,0.000144) (0.24066,0.000297) (0.25663,0.000575) (0.27261,0.001054) (0.28859,0.001840) (0.30457,0.003079) (0.32054,0.004957) (0.33652,0.007707) (0.35250,0.011610) (0.36847,0.016993) (0.38445,0.024225) (0.40043,0.033710) (0.41641,0.045873) (0.43238,0.061151) (0.44836,0.079972) (0.46434,0.102742) (0.48032,0.129824) (0.49629,0.161522) (0.51227,0.198062) (0.52825,0.239583) (0.54422,0.286118) (0.56020,0.337591) (0.57618,0.393813) (0.59216,0.454477) (0.60813,0.519166) (0.62411,0.587357) (0.64009,0.658434) (0.65606,0.731699) (0.67204,0.806392) (0.68802,0.881702) (0.70400,0.956794) (0.71997,1.030818) (0.73595,1.102938) (0.75193,1.172340) (0.76790,1.238257) (0.78388,1.299975) (0.79986,1.356853) (0.81584,1.408329) (0.83181,1.453929) (0.84779,1.493271) (0.86377,1.526070) (0.87974,1.552135) (0.89572,1.571372) (0.91170,1.583775) (0.92768,1.589427) (0.94365,1.588487) (0.95963,1.581189) (0.97561,1.567827) (0.99158,1.548751) (1.00756,1.524356) (1.02354,1.495072) (1.03952,1.461356) (1.05549,1.423684) (1.07147,1.382540) (1.08745,1.338413) (1.10342,1.291783) (1.11940,1.243123) (1.13538,1.192889) (1.15136,1.141515) (1.16733,1.089413) (1.18331,1.036965) (1.19929,0.984527) (1.21526,0.932420) (1.23124,0.880936) (1.24722,0.830335) (1.26320,0.780844) (1.27917,0.732661) (1.29515,0.685951) (1.31113,0.640851) (1.32711,0.597474) (1.34308,0.555903) (1.35906,0.516201) (1.37504,0.478408) (1.39101,0.442544) (1.40699,0.408614) (1.42297,0.376604) (1.43895,0.346490) (1.45492,0.318234) (1.47090,0.291789) (1.48688,0.267100) (1.50285,0.244106) (1.51883,0.222738) (1.53481,0.202926) (1.55079,0.184596) (1.56676,0.167672) (1.58274,0.152078) (1.59872,0.137738) (1.61469,0.124576) (1.63067,0.112517) (1.64665,0.101489) (1.66263,0.091421) (1.67860,0.082246) (1.69458,0.073897) (1.71056,0.066314) (1.72653,0.059436) (1.74251,0.053208) (1.75849,0.047576) (1.77447,0.042491) (1.79044,0.037907) (1.80642,0.033779) (1.82240,0.030068) (1.83837,0.026736) (1.85435,0.023748) (1.87033,0.021072) (1.88631,0.018679) (1.90228,0.016540) (1.91826,0.014632) (1.93424,0.012932) (1.95021,0.011418) (1.96619,0.010072) (1.98217,0.008877) (1.99815,0.007816) (2.01412,0.006876) (2.03010,0.006043) (2.04608,0.005307) (2.06205,0.004657) (2.07803,0.004082) (2.09401,0.003576) (2.10999,0.003130) (2.12596,0.002737) (2.14194,0.002392) (2.15792,0.002089) (2.17389,0.001822) (2.18987,0.001589) (2.20585,0.001384) (2.22183,0.001205) (2.23780,0.001048) (2.25378,0.000911) (2.26976,0.000792) (2.28574,0.000687) (2.30171,0.000596) (2.31769,0.000517) (2.33367,0.000448) (2.34964,0.000388) (2.36562,0.000335) (2.38160,0.000290) (2.39758,0.000251) (2.41355,0.000216) (2.42953,0.000187) (2.44551,0.000161) (2.46148,0.000139) (2.47746,0.000120) (2.49344,0.000103) (2.50942,0.000089) (2.52539,0.000076) (2.54137,0.000066) (2.55735,0.000056) (2.57332,0.000048) (2.58930,0.000041) (2.60528,0.000036) (2.62126,0.000030) (2.63723,0.000026) (2.65321,0.000022) (2.66919,0.000019) (2.68516,0.000016) (2.70114,0.000014) (2.71712,0.000012) (2.73310,0.000010) (2.74907,0.000009) (2.76505,0.000007) (2.78103,0.000006) (2.79700,0.000005) (2.81298,0.000005) (2.82896,0.000004) (2.84494,0.000003) (2.86091,0.000003) (2.87689,0.000002) (2.89287,0.000002) (2.90884,0.000002) (2.92482,0.000001) (2.94080,0.000001) (2.95678,0.000001) (2.97275,0.000001) (2.98873,0.000001) (3.00471,0.000001) (3.02068,0.000001) (3.03666,0.000000) (3.05264,0.000000) (3.06862,0.000000) (3.08459,0.000000) (3.10057,0.000000) (3.11655,0.000000) (3.13253,0.000000) (3.14850,0.000000) (3.16448,0.000000) (3.18046,0.000000) (3.19643,0.000000) (3.21241,0.000000) (3.22839,0.000000) (3.24437,0.000000) (3.26034,0.000000) (3.27632,0.000000) (3.29230,0.000000) (3.30827,0.000000) (3.32425,0.000000) (3.34023,0.000000) (3.35621,0.000000) (3.37218,0.000000) (3.38816,0.000000) (3.40414,0.000000) (3.42011,0.000000) (3.43609,0.000000) (3.45207,0.000000) (3.46805,0.000000) (3.48402,0.000000) (3.50000,0.000000)};
\addlegendentry{$\nu=30$}
\end{axis}\end{tikzpicture}\end{center}

若 $Q\sim\chi^2_\nu$，则 $Q/\nu$ 的均值为一，方差为 $2/\nu$。用相同均值尺度比较，可见自由度较小时明显右偏，自由度增大后集中程度提高。正态样本的 $S^2/\sigma^2$ 正具有这种分布，自由度为 $n-1$。

方差区间通常不对称，并非因为计算不精确，而是参考分布偏斜和倒数反演共同造成。对方差区间端点取平方根可得到标准差区间，但不能只对半径取平方根后保留原中心。参见习题~\ref{q:M-C22-1}、\ref{q:M-P10-3}。

\clearpage\subsection*{图解六：置信水平的重复抽样含义}

\begin{center}\begin{tikzpicture}\begin{axis}[width=0.87\linewidth,height=0.69\linewidth,axis lines=left,tick label style={font=\small},label style={font=\small},legend style={font=\footnotesize,draw=none},xlabel={重复试验编号},ylabel={以标准误为单位的均值区间},xmin=0,xmax=41]
\addplot[black,thick] coordinates {(1.00000,-2.827854) (1.00000,1.092074)};
\addplot[black,thick] coordinates {(2.00000,-3.473417) (2.00000,0.446511)};
\addplot[black,thick] coordinates {(3.00000,-1.673673) (3.00000,2.246255)};
\addplot[black,thick] coordinates {(4.00000,-2.426365) (4.00000,1.493563)};
\addplot[black,thick] coordinates {(5.00000,-0.682072) (5.00000,3.237856)};
\addplot[black,thick] coordinates {(6.00000,-1.216195) (6.00000,2.703733)};
\addplot[black,thick] coordinates {(7.00000,-3.777209) (7.00000,0.142719)};
\addplot[black,thick] coordinates {(8.00000,-2.377894) (8.00000,1.542034)};
\addplot[black,thick] coordinates {(9.00000,-1.772638) (9.00000,2.147290)};
\addplot[black,thick] coordinates {(10.00000,-2.660165) (10.00000,1.259763)};
\addplot[black,thick] coordinates {(11.00000,-3.422815) (11.00000,0.497113)};
\addplot[black,thick] coordinates {(12.00000,-1.601206) (12.00000,2.318722)};
\addplot[black,thick] coordinates {(13.00000,-1.766438) (13.00000,2.153490)};
\addplot[black,thick] coordinates {(14.00000,-2.707712) (14.00000,1.212216)};
\addplot[black,thick] coordinates {(15.00000,-3.196906) (15.00000,0.723022)};
\addplot[black,thick] coordinates {(16.00000,-0.057448) (16.00000,3.862480)};
\addplot[black,thick] coordinates {(17.00000,-2.927236) (17.00000,0.992692)};
\addplot[black,thick] coordinates {(18.00000,-2.192602) (18.00000,1.727326)};
\addplot[black,thick] coordinates {(19.00000,-0.855217) (19.00000,3.064711)};
\addplot[black,thick] coordinates {(20.00000,-3.760312) (20.00000,0.159616)};
\addplot[black,thick] coordinates {(21.00000,-0.760053) (21.00000,3.159875)};
\addplot[black,thick] coordinates {(22.00000,-1.705289) (22.00000,2.214639)};
\addplot[black,thick] coordinates {(23.00000,-2.861649) (23.00000,1.058278)};
\addplot[black,thick] coordinates {(24.00000,-1.339421) (24.00000,2.580507)};
\addplot[black,thick] coordinates {(25.00000,-1.516632) (25.00000,2.403296)};
\addplot[black,thick] coordinates {(26.00000,-0.821718) (26.00000,3.098210)};
\addplot[black,thick] coordinates {(27.00000,-2.156287) (27.00000,1.763641)};
\addplot[black,thick] coordinates {(28.00000,-0.514062) (28.00000,3.405866)};
\addplot[black,thick] coordinates {(29.00000,-2.637694) (29.00000,1.282234)};
\addplot[black,thick] coordinates {(30.00000,-0.879699) (30.00000,3.040229)};
\addplot[black,thick] coordinates {(31.00000,-2.377251) (31.00000,1.542677)};
\addplot[black,thick] coordinates {(32.00000,-1.464345) (32.00000,2.455583)};
\addplot[black,thick] coordinates {(33.00000,-2.332668) (33.00000,1.587260)};
\addplot[black,thick] coordinates {(34.00000,-1.730235) (34.00000,2.189693)};
\addplot[black,thick] coordinates {(35.00000,-1.699929) (35.00000,2.219999)};
\addplot[black,thick] coordinates {(36.00000,-0.849247) (36.00000,3.070681)};
\addplot[black,thick] coordinates {(37.00000,-2.769283) (37.00000,1.150645)};
\addplot[black,thick] coordinates {(38.00000,-2.431050) (38.00000,1.488878)};
\addplot[black,thick] coordinates {(39.00000,-0.980838) (39.00000,2.939090)};
\addplot[black,thick] coordinates {(40.00000,-2.125817) (40.00000,1.794111)};
\addplot[gray] coordinates {(0.00000,0.000000) (41.00000,0.000000)};
\end{axis}\end{tikzpicture}\end{center}

固定真实均值为零，独立重复模拟 40 次已知标准误的正态均值区间。实线覆盖零，粗虚线未覆盖。本次固定随机种子的模拟中有 0 个区间未覆盖。模拟比例无需恰好等于 $5\%$，因为 40 次有限重复本身也有随机波动。

每条竖线是不同样本产生的随机区间，水平零线是同一个固定参数。$95\%$ 描述的是长期构造程序的覆盖率。给定其中一条已经算出的区间，它是否包含零已是确定事实；若要讨论参数给定数据后的概率，需要另建贝叶斯模型。参见习题~\ref{q:M-P10-5}。

\clearpage\subsection*{图解七：样本容量与检验功效}

\begin{center}\begin{tikzpicture}\begin{axis}[width=0.87\linewidth,height=0.69\linewidth,axis lines=left,tick label style={font=\small},label style={font=\small},legend style={font=\footnotesize,draw=none},xlabel={真实均值 $\mu$},ylabel={拒绝概率},ymin=0,ymax=1.02,legend pos=south east]
\addplot[black,dotted,thick] coordinates {(-1.00000,0.885379) (-0.99000,0.879139) (-0.98000,0.872663) (-0.97000,0.865950) (-0.96000,0.858997) (-0.95000,0.851804) (-0.94000,0.844369) (-0.93000,0.836693) (-0.92000,0.828775) (-0.91000,0.820615) (-0.90000,0.812215) (-0.89000,0.803577) (-0.88000,0.794702) (-0.87000,0.785593) (-0.86000,0.776253) (-0.85000,0.766686) (-0.84000,0.756897) (-0.83000,0.746890) (-0.82000,0.736670) (-0.81000,0.726243) (-0.80000,0.715617) (-0.79000,0.704797) (-0.78000,0.693791) (-0.77000,0.682608) (-0.76000,0.671256) (-0.75000,0.659744) (-0.74000,0.648081) (-0.73000,0.636277) (-0.72000,0.624342) (-0.71000,0.612288) (-0.70000,0.600124) (-0.69000,0.587863) (-0.68000,0.575516) (-0.67000,0.563094) (-0.66000,0.550611) (-0.65000,0.538077) (-0.64000,0.525506) (-0.63000,0.512911) (-0.62000,0.500303) (-0.61000,0.487695) (-0.60000,0.475101) (-0.59000,0.462532) (-0.58000,0.450002) (-0.57000,0.437523) (-0.56000,0.425107) (-0.55000,0.412767) (-0.54000,0.400514) (-0.53000,0.388360) (-0.52000,0.376317) (-0.51000,0.364396) (-0.50000,0.352608) (-0.49000,0.340963) (-0.48000,0.329472) (-0.47000,0.318143) (-0.46000,0.306988) (-0.45000,0.296013) (-0.44000,0.285229) (-0.43000,0.274642) (-0.42000,0.264261) (-0.41000,0.254092) (-0.40000,0.244141) (-0.39000,0.234416) (-0.38000,0.224921) (-0.37000,0.215661) (-0.36000,0.206641) (-0.35000,0.197865) (-0.34000,0.189337) (-0.33000,0.181060) (-0.32000,0.173036) (-0.31000,0.165267) (-0.30000,0.157756) (-0.29000,0.150504) (-0.28000,0.143511) (-0.27000,0.136779) (-0.26000,0.130308) (-0.25000,0.124097) (-0.24000,0.118147) (-0.23000,0.112457) (-0.22000,0.107025) (-0.21000,0.101852) (-0.20000,0.096935) (-0.19000,0.092274) (-0.18000,0.087867) (-0.17000,0.083711) (-0.16000,0.079806) (-0.15000,0.076150) (-0.14000,0.072740) (-0.13000,0.069575) (-0.12000,0.066654) (-0.11000,0.063973) (-0.10000,0.061533) (-0.09000,0.059330) (-0.08000,0.057363) (-0.07000,0.055632) (-0.06000,0.054134) (-0.05000,0.052869) (-0.04000,0.051835) (-0.03000,0.051032) (-0.02000,0.050458) (-0.01000,0.050115) (0.00000,0.050000) (0.01000,0.050115) (0.02000,0.050458) (0.03000,0.051032) (0.04000,0.051835) (0.05000,0.052869) (0.06000,0.054134) (0.07000,0.055632) (0.08000,0.057363) (0.09000,0.059330) (0.10000,0.061533) (0.11000,0.063973) (0.12000,0.066654) (0.13000,0.069575) (0.14000,0.072740) (0.15000,0.076150) (0.16000,0.079806) (0.17000,0.083711) (0.18000,0.087867) (0.19000,0.092274) (0.20000,0.096935) (0.21000,0.101852) (0.22000,0.107025) (0.23000,0.112457) (0.24000,0.118147) (0.25000,0.124097) (0.26000,0.130308) (0.27000,0.136779) (0.28000,0.143511) (0.29000,0.150504) (0.30000,0.157756) (0.31000,0.165267) (0.32000,0.173036) (0.33000,0.181060) (0.34000,0.189337) (0.35000,0.197865) (0.36000,0.206641) (0.37000,0.215661) (0.38000,0.224921) (0.39000,0.234416) (0.40000,0.244141) (0.41000,0.254092) (0.42000,0.264261) (0.43000,0.274642) (0.44000,0.285229) (0.45000,0.296013) (0.46000,0.306988) (0.47000,0.318143) (0.48000,0.329472) (0.49000,0.340963) (0.50000,0.352608) (0.51000,0.364396) (0.52000,0.376317) (0.53000,0.388360) (0.54000,0.400514) (0.55000,0.412767) (0.56000,0.425107) (0.57000,0.437523) (0.58000,0.450002) (0.59000,0.462532) (0.60000,0.475101) (0.61000,0.487695) (0.62000,0.500303) (0.63000,0.512911) (0.64000,0.525506) (0.65000,0.538077) (0.66000,0.550611) (0.67000,0.563094) (0.68000,0.575516) (0.69000,0.587863) (0.70000,0.600124) (0.71000,0.612288) (0.72000,0.624342) (0.73000,0.636277) (0.74000,0.648081) (0.75000,0.659744) (0.76000,0.671256) (0.77000,0.682608) (0.78000,0.693791) (0.79000,0.704797) (0.80000,0.715617) (0.81000,0.726243) (0.82000,0.736670) (0.83000,0.746890) (0.84000,0.756897) (0.85000,0.766686) (0.86000,0.776253) (0.87000,0.785593) (0.88000,0.794702) (0.89000,0.803577) (0.90000,0.812215) (0.91000,0.820615) (0.92000,0.828775) (0.93000,0.836693) (0.94000,0.844369) (0.95000,0.851804) (0.96000,0.858997) (0.97000,0.865950) (0.98000,0.872663) (0.99000,0.879139) (1.00000,0.885379)};
\addlegendentry{$n=10$}
\addplot[black,dashed,thick] coordinates {(-1.00000,0.999782) (-0.99000,0.999732) (-0.98000,0.999672) (-0.97000,0.999600) (-0.96000,0.999513) (-0.95000,0.999409) (-0.94000,0.999285) (-0.93000,0.999137) (-0.92000,0.998962) (-0.91000,0.998754) (-0.90000,0.998509) (-0.89000,0.998220) (-0.88000,0.997882) (-0.87000,0.997486) (-0.86000,0.997024) (-0.85000,0.996488) (-0.84000,0.995866) (-0.83000,0.995147) (-0.82000,0.994319) (-0.81000,0.993368) (-0.80000,0.992278) (-0.79000,0.991035) (-0.78000,0.989619) (-0.77000,0.988012) (-0.76000,0.986193) (-0.75000,0.984141) (-0.74000,0.981834) (-0.73000,0.979246) (-0.72000,0.976352) (-0.71000,0.973126) (-0.70000,0.969541) (-0.69000,0.965569) (-0.68000,0.961180) (-0.67000,0.956346) (-0.66000,0.951038) (-0.65000,0.945227) (-0.64000,0.938883) (-0.63000,0.931978) (-0.62000,0.924487) (-0.61000,0.916383) (-0.60000,0.907642) (-0.59000,0.898242) (-0.58000,0.888165) (-0.57000,0.877394) (-0.56000,0.865914) (-0.55000,0.853718) (-0.54000,0.840797) (-0.53000,0.827151) (-0.52000,0.812782) (-0.51000,0.797697) (-0.50000,0.781908) (-0.49000,0.765431) (-0.48000,0.748288) (-0.47000,0.730505) (-0.46000,0.712114) (-0.45000,0.693151) (-0.44000,0.673656) (-0.43000,0.653676) (-0.42000,0.633259) (-0.41000,0.612458) (-0.40000,0.591331) (-0.39000,0.569935) (-0.38000,0.548332) (-0.37000,0.526587) (-0.36000,0.504764) (-0.35000,0.482930) (-0.34000,0.461148) (-0.33000,0.439487) (-0.32000,0.418009) (-0.31000,0.396779) (-0.30000,0.375856) (-0.29000,0.355301) (-0.28000,0.335168) (-0.27000,0.315510) (-0.26000,0.296376) (-0.25000,0.277810) (-0.24000,0.259854) (-0.23000,0.242542) (-0.22000,0.225908) (-0.21000,0.209978) (-0.20000,0.194775) (-0.19000,0.180319) (-0.18000,0.166623) (-0.17000,0.153698) (-0.16000,0.141551) (-0.15000,0.130186) (-0.14000,0.119603) (-0.13000,0.109799) (-0.12000,0.100771) (-0.11000,0.092513) (-0.10000,0.085016) (-0.09000,0.078271) (-0.08000,0.072271) (-0.07000,0.067004) (-0.06000,0.062462) (-0.05000,0.058635) (-0.04000,0.055517) (-0.03000,0.053099) (-0.02000,0.051376) (-0.01000,0.050344) (0.00000,0.050000) (0.01000,0.050344) (0.02000,0.051376) (0.03000,0.053099) (0.04000,0.055517) (0.05000,0.058635) (0.06000,0.062462) (0.07000,0.067004) (0.08000,0.072271) (0.09000,0.078271) (0.10000,0.085016) (0.11000,0.092513) (0.12000,0.100771) (0.13000,0.109799) (0.14000,0.119603) (0.15000,0.130186) (0.16000,0.141551) (0.17000,0.153698) (0.18000,0.166623) (0.19000,0.180319) (0.20000,0.194775) (0.21000,0.209978) (0.22000,0.225908) (0.23000,0.242542) (0.24000,0.259854) (0.25000,0.277810) (0.26000,0.296376) (0.27000,0.315510) (0.28000,0.335168) (0.29000,0.355301) (0.30000,0.375856) (0.31000,0.396779) (0.32000,0.418009) (0.33000,0.439487) (0.34000,0.461148) (0.35000,0.482930) (0.36000,0.504764) (0.37000,0.526587) (0.38000,0.548332) (0.39000,0.569935) (0.40000,0.591331) (0.41000,0.612458) (0.42000,0.633259) (0.43000,0.653676) (0.44000,0.673656) (0.45000,0.693151) (0.46000,0.712114) (0.47000,0.730505) (0.48000,0.748288) (0.49000,0.765431) (0.50000,0.781908) (0.51000,0.797697) (0.52000,0.812782) (0.53000,0.827151) (0.54000,0.840797) (0.55000,0.853718) (0.56000,0.865914) (0.57000,0.877394) (0.58000,0.888165) (0.59000,0.898242) (0.60000,0.907642) (0.61000,0.916383) (0.62000,0.924487) (0.63000,0.931978) (0.64000,0.938883) (0.65000,0.945227) (0.66000,0.951038) (0.67000,0.956346) (0.68000,0.961180) (0.69000,0.965569) (0.70000,0.969541) (0.71000,0.973126) (0.72000,0.976352) (0.73000,0.979246) (0.74000,0.981834) (0.75000,0.984141) (0.76000,0.986193) (0.77000,0.988012) (0.78000,0.989619) (0.79000,0.991035) (0.80000,0.992278) (0.81000,0.993368) (0.82000,0.994319) (0.83000,0.995147) (0.84000,0.995866) (0.85000,0.996488) (0.86000,0.997024) (0.87000,0.997486) (0.88000,0.997882) (0.89000,0.998220) (0.90000,0.998509) (0.91000,0.998754) (0.92000,0.998962) (0.93000,0.999137) (0.94000,0.999285) (0.95000,0.999409) (0.96000,0.999513) (0.97000,0.999600) (0.98000,0.999672) (0.99000,0.999732) (1.00000,0.999782)};
\addlegendentry{$n=30$}
\addplot[black,thick] coordinates {(-1.00000,1.000000) (-0.99000,1.000000) (-0.98000,1.000000) (-0.97000,1.000000) (-0.96000,1.000000) (-0.95000,1.000000) (-0.94000,1.000000) (-0.93000,1.000000) (-0.92000,1.000000) (-0.91000,1.000000) (-0.90000,1.000000) (-0.89000,1.000000) (-0.88000,1.000000) (-0.87000,1.000000) (-0.86000,1.000000) (-0.85000,1.000000) (-0.84000,1.000000) (-0.83000,1.000000) (-0.82000,1.000000) (-0.81000,1.000000) (-0.80000,1.000000) (-0.79000,1.000000) (-0.78000,1.000000) (-0.77000,1.000000) (-0.76000,1.000000) (-0.75000,1.000000) (-0.74000,1.000000) (-0.73000,1.000000) (-0.72000,1.000000) (-0.71000,1.000000) (-0.70000,1.000000) (-0.69000,1.000000) (-0.68000,0.999999) (-0.67000,0.999999) (-0.66000,0.999998) (-0.65000,0.999997) (-0.64000,0.999996) (-0.63000,0.999993) (-0.62000,0.999989) (-0.61000,0.999983) (-0.60000,0.999973) (-0.59000,0.999959) (-0.58000,0.999938) (-0.57000,0.999908) (-0.56000,0.999864) (-0.55000,0.999800) (-0.54000,0.999709) (-0.53000,0.999581) (-0.52000,0.999402) (-0.51000,0.999155) (-0.50000,0.998817) (-0.49000,0.998359) (-0.48000,0.997745) (-0.47000,0.996928) (-0.46000,0.995855) (-0.45000,0.994458) (-0.44000,0.992657) (-0.43000,0.990359) (-0.42000,0.987456) (-0.41000,0.983824) (-0.40000,0.979327) (-0.39000,0.973812) (-0.38000,0.967119) (-0.37000,0.959074) (-0.36000,0.949501) (-0.35000,0.938224) (-0.34000,0.925071) (-0.33000,0.909883) (-0.32000,0.892519) (-0.31000,0.872865) (-0.30000,0.850839) (-0.29000,0.826401) (-0.28000,0.799557) (-0.27000,0.770363) (-0.26000,0.738928) (-0.25000,0.705418) (-0.24000,0.670051) (-0.23000,0.633096) (-0.22000,0.594865) (-0.21000,0.555709) (-0.20000,0.516005) (-0.19000,0.476149) (-0.18000,0.436540) (-0.17000,0.397572) (-0.16000,0.359622) (-0.15000,0.323041) (-0.14000,0.288142) (-0.13000,0.255196) (-0.12000,0.224427) (-0.11000,0.196011) (-0.10000,0.170075) (-0.09000,0.146699) (-0.08000,0.125922) (-0.07000,0.107749) (-0.06000,0.092155) (-0.05000,0.079098) (-0.04000,0.068523) (-0.03000,0.060373) (-0.02000,0.054595) (-0.01000,0.051146) (0.00000,0.050000) (0.01000,0.051146) (0.02000,0.054595) (0.03000,0.060373) (0.04000,0.068523) (0.05000,0.079098) (0.06000,0.092155) (0.07000,0.107749) (0.08000,0.125922) (0.09000,0.146699) (0.10000,0.170075) (0.11000,0.196011) (0.12000,0.224427) (0.13000,0.255196) (0.14000,0.288142) (0.15000,0.323041) (0.16000,0.359622) (0.17000,0.397572) (0.18000,0.436540) (0.19000,0.476149) (0.20000,0.516005) (0.21000,0.555709) (0.22000,0.594865) (0.23000,0.633096) (0.24000,0.670051) (0.25000,0.705418) (0.26000,0.738928) (0.27000,0.770363) (0.28000,0.799557) (0.29000,0.826401) (0.30000,0.850839) (0.31000,0.872865) (0.32000,0.892519) (0.33000,0.909883) (0.34000,0.925071) (0.35000,0.938224) (0.36000,0.949501) (0.37000,0.959074) (0.38000,0.967119) (0.39000,0.973812) (0.40000,0.979327) (0.41000,0.983824) (0.42000,0.987456) (0.43000,0.990359) (0.44000,0.992657) (0.45000,0.994458) (0.46000,0.995855) (0.47000,0.996928) (0.48000,0.997745) (0.49000,0.998359) (0.50000,0.998817) (0.51000,0.999155) (0.52000,0.999402) (0.53000,0.999581) (0.54000,0.999709) (0.55000,0.999800) (0.56000,0.999864) (0.57000,0.999908) (0.58000,0.999938) (0.59000,0.999959) (0.60000,0.999973) (0.61000,0.999983) (0.62000,0.999989) (0.63000,0.999993) (0.64000,0.999996) (0.65000,0.999997) (0.66000,0.999998) (0.67000,0.999999) (0.68000,0.999999) (0.69000,1.000000) (0.70000,1.000000) (0.71000,1.000000) (0.72000,1.000000) (0.73000,1.000000) (0.74000,1.000000) (0.75000,1.000000) (0.76000,1.000000) (0.77000,1.000000) (0.78000,1.000000) (0.79000,1.000000) (0.80000,1.000000) (0.81000,1.000000) (0.82000,1.000000) (0.83000,1.000000) (0.84000,1.000000) (0.85000,1.000000) (0.86000,1.000000) (0.87000,1.000000) (0.88000,1.000000) (0.89000,1.000000) (0.90000,1.000000) (0.91000,1.000000) (0.92000,1.000000) (0.93000,1.000000) (0.94000,1.000000) (0.95000,1.000000) (0.96000,1.000000) (0.97000,1.000000) (0.98000,1.000000) (0.99000,1.000000) (1.00000,1.000000)};
\addlegendentry{$n=100$}
\end{axis}\end{tikzpicture}\end{center}

检验 $H_0:\mu=0$，独立正态观测方差为一，双侧水平为 $0.05$。在 $\mu=0$，三条曲线都等于第一类错误概率 $0.05$；离开零假设后，容量较大的检验更快获得高功效。

功效取决于真实差异、噪声、容量和检验规则。没有显著结果可能只是样本精度不足；反过来，大样本可以使很小的差异达到显著，所以实际意义仍应结合效应大小判断。图中的双侧曲线关于零对称，是该正态模型的对称性所致。参见习题~\ref{q:B14.1}、\ref{q:M-F-11}。

\clearpage\subsection*{图解八：拟合直线与两种不确定性带}

\begin{center}\begin{tikzpicture}\begin{axis}[width=0.87\linewidth,height=0.69\linewidth,axis lines=left,tick label style={font=\small},label style={font=\small},legend style={font=\footnotesize,draw=none},xlabel={$x$},ylabel={$y$},legend pos=north west]
\addplot[only marks,mark=*,black] coordinates {(-2.00000,-0.500000) (-1.00000,-0.200000) (0.00000,1.200000) (1.00000,1.500000) (2.00000,2.900000) (3.00000,2.800000)};
\addlegendentry{观测}
\addplot[black,thick] coordinates {(-2.50000,-0.953810) (-2.43939,-0.908615) (-2.37879,-0.863420) (-2.31818,-0.818225) (-2.25758,-0.773030) (-2.19697,-0.727835) (-2.13636,-0.682641) (-2.07576,-0.637446) (-2.01515,-0.592251) (-1.95455,-0.547056) (-1.89394,-0.501861) (-1.83333,-0.456667) (-1.77273,-0.411472) (-1.71212,-0.366277) (-1.65152,-0.321082) (-1.59091,-0.275887) (-1.53030,-0.230693) (-1.46970,-0.185498) (-1.40909,-0.140303) (-1.34848,-0.095108) (-1.28788,-0.049913) (-1.22727,-0.004719) (-1.16667,0.040476) (-1.10606,0.085671) (-1.04545,0.130866) (-0.98485,0.176061) (-0.92424,0.221255) (-0.86364,0.266450) (-0.80303,0.311645) (-0.74242,0.356840) (-0.68182,0.402035) (-0.62121,0.447229) (-0.56061,0.492424) (-0.50000,0.537619) (-0.43939,0.582814) (-0.37879,0.628009) (-0.31818,0.673203) (-0.25758,0.718398) (-0.19697,0.763593) (-0.13636,0.808788) (-0.07576,0.853983) (-0.01515,0.899177) (0.04545,0.944372) (0.10606,0.989567) (0.16667,1.034762) (0.22727,1.079957) (0.28788,1.125152) (0.34848,1.170346) (0.40909,1.215541) (0.46970,1.260736) (0.53030,1.305931) (0.59091,1.351126) (0.65152,1.396320) (0.71212,1.441515) (0.77273,1.486710) (0.83333,1.531905) (0.89394,1.577100) (0.95455,1.622294) (1.01515,1.667489) (1.07576,1.712684) (1.13636,1.757879) (1.19697,1.803074) (1.25758,1.848268) (1.31818,1.893463) (1.37879,1.938658) (1.43939,1.983853) (1.50000,2.029048) (1.56061,2.074242) (1.62121,2.119437) (1.68182,2.164632) (1.74242,2.209827) (1.80303,2.255022) (1.86364,2.300216) (1.92424,2.345411) (1.98485,2.390606) (2.04545,2.435801) (2.10606,2.480996) (2.16667,2.526190) (2.22727,2.571385) (2.28788,2.616580) (2.34848,2.661775) (2.40909,2.706970) (2.46970,2.752165) (2.53030,2.797359) (2.59091,2.842554) (2.65152,2.887749) (2.71212,2.932944) (2.77273,2.978139) (2.83333,3.023333) (2.89394,3.068528) (2.95455,3.113723) (3.01515,3.158918) (3.07576,3.204113) (3.13636,3.249307) (3.19697,3.294502) (3.25758,3.339697) (3.31818,3.384892) (3.37879,3.430087) (3.43939,3.475281) (3.50000,3.520476)};
\addlegendentry{拟合均值}
\addplot[black,dashed] coordinates {(-2.50000,-1.853465) (-2.43939,-1.794578) (-2.37879,-1.735762) (-2.31818,-1.677021) (-2.25758,-1.618358) (-2.19697,-1.559777) (-2.13636,-1.501281) (-2.07576,-1.442876) (-2.01515,-1.384565) (-1.95455,-1.326353) (-1.89394,-1.268246) (-1.83333,-1.210249) (-1.77273,-1.152368) (-1.71212,-1.094607) (-1.65152,-1.036975) (-1.59091,-0.979478) (-1.53030,-0.922122) (-1.46970,-0.864916) (-1.40909,-0.807868) (-1.34848,-0.750986) (-1.28788,-0.694279) (-1.22727,-0.637757) (-1.16667,-0.581429) (-1.10606,-0.525308) (-1.04545,-0.469403) (-0.98485,-0.413726) (-0.92424,-0.358290) (-0.86364,-0.303109) (-0.80303,-0.248195) (-0.74242,-0.193562) (-0.68182,-0.139226) (-0.62121,-0.085202) (-0.56061,-0.031505) (-0.50000,0.021848) (-0.43939,0.074841) (-0.37879,0.127457) (-0.31818,0.179678) (-0.25758,0.231489) (-0.19697,0.282871) (-0.13636,0.333809) (-0.07576,0.384286) (-0.01515,0.434285) (0.04545,0.483793) (0.10606,0.532795) (0.16667,0.581279) (0.22727,0.629233) (0.28788,0.676647) (0.34848,0.723514) (0.40909,0.769827) (0.46970,0.815582) (0.53030,0.860776) (0.59091,0.905411) (0.65152,0.949488) (0.71212,0.993011) (0.77273,1.035986) (0.83333,1.078422) (0.89394,1.120328) (0.95455,1.161715) (1.01515,1.202597) (1.07576,1.242987) (1.13636,1.282900) (1.19697,1.322352) (1.25758,1.361359) (1.31818,1.399938) (1.37879,1.438106) (1.43939,1.475880) (1.50000,1.513276) (1.56061,1.550313) (1.62121,1.587006) (1.68182,1.623371) (1.74242,1.659425) (1.80303,1.695182) (1.86364,1.730657) (1.92424,1.765865) (1.98485,1.800819) (2.04545,1.835532) (2.10606,1.870017) (2.16667,1.904285) (2.22727,1.938347) (2.28788,1.972215) (2.34848,2.005897) (2.40909,2.039405) (2.46970,2.072746) (2.53030,2.105930) (2.59091,2.138964) (2.65152,2.171856) (2.71212,2.204613) (2.77273,2.237243) (2.83333,2.269751) (2.89394,2.302143) (2.95455,2.334426) (3.01515,2.366604) (3.07576,2.398683) (3.13636,2.430667) (3.19697,2.462561) (3.25758,2.494369) (3.31818,2.526096) (3.37879,2.557744) (3.43939,2.589318) (3.50000,2.620821)};
\addplot[black,dashed] coordinates {(-2.50000,-0.054154) (-2.43939,-0.022651) (-2.37879,0.008923) (-2.31818,0.040571) (-2.25758,0.072297) (-2.19697,0.104106) (-2.13636,0.136000) (-2.07576,0.167984) (-2.01515,0.200063) (-1.95455,0.232241) (-1.89394,0.264523) (-1.83333,0.296916) (-1.77273,0.329424) (-1.71212,0.362053) (-1.65152,0.394811) (-1.59091,0.427703) (-1.53030,0.460737) (-1.46970,0.493921) (-1.40909,0.527262) (-1.34848,0.560770) (-1.28788,0.594452) (-1.22727,0.628320) (-1.16667,0.662382) (-1.10606,0.696650) (-1.04545,0.731134) (-0.98485,0.765847) (-0.92424,0.800801) (-0.86364,0.836009) (-0.80303,0.871485) (-0.74242,0.907242) (-0.68182,0.943296) (-0.62121,0.979661) (-0.56061,1.016354) (-0.50000,1.053390) (-0.43939,1.090787) (-0.37879,1.128561) (-0.31818,1.166729) (-0.25758,1.205308) (-0.19697,1.244315) (-0.13636,1.283767) (-0.07576,1.323680) (-0.01515,1.364070) (0.04545,1.404951) (0.10606,1.446339) (0.16667,1.488245) (0.22727,1.530680) (0.28788,1.573656) (0.34848,1.617179) (0.40909,1.661256) (0.46970,1.705890) (0.53030,1.751085) (0.59091,1.796840) (0.65152,1.843153) (0.71212,1.890019) (0.77273,1.937434) (0.83333,1.985388) (0.89394,2.033871) (0.95455,2.082873) (1.01515,2.132381) (1.07576,2.182381) (1.13636,2.232858) (1.19697,2.283795) (1.25758,2.335178) (1.31818,2.386988) (1.37879,2.439210) (1.43939,2.491826) (1.50000,2.544819) (1.56061,2.598172) (1.62121,2.651869) (1.68182,2.705893) (1.74242,2.760229) (1.80303,2.814861) (1.86364,2.869776) (1.92424,2.924957) (1.98485,2.980393) (2.04545,3.036069) (2.10606,3.091974) (2.16667,3.148096) (2.22727,3.204423) (2.28788,3.260946) (2.34848,3.317653) (2.40909,3.374535) (2.46970,3.431583) (2.53030,3.488789) (2.59091,3.546145) (2.65152,3.603642) (2.71212,3.661274) (2.77273,3.719034) (2.83333,3.776916) (2.89394,3.834913) (2.95455,3.893020) (3.01515,3.951231) (3.07576,4.009542) (3.13636,4.067948) (3.19697,4.126443) (3.25758,4.185025) (3.31818,4.243688) (3.37879,4.302429) (3.43939,4.361245) (3.50000,4.420132)};
\addlegendentry{均值区间边界}
\addplot[gray,dotted,thick] coordinates {(-2.50000,-2.367309) (-2.43939,-2.313439) (-2.37879,-2.259694) (-2.31818,-2.206076) (-2.25758,-2.152588) (-2.19697,-2.099231) (-2.13636,-2.046009) (-2.07576,-1.992923) (-2.01515,-1.939975) (-1.95455,-1.887170) (-1.89394,-1.834508) (-1.83333,-1.781992) (-1.77273,-1.729625) (-1.71212,-1.677408) (-1.65152,-1.625346) (-1.59091,-1.573439) (-1.53030,-1.521690) (-1.46970,-1.470103) (-1.40909,-1.418678) (-1.34848,-1.367419) (-1.28788,-1.316328) (-1.22727,-1.265408) (-1.16667,-1.214660) (-1.10606,-1.164087) (-1.04545,-1.113691) (-0.98485,-1.063475) (-0.92424,-1.013440) (-0.86364,-0.963589) (-0.80303,-0.913924) (-0.74242,-0.864447) (-0.68182,-0.815160) (-0.62121,-0.766065) (-0.56061,-0.717163) (-0.50000,-0.668457) (-0.43939,-0.619948) (-0.37879,-0.571638) (-0.31818,-0.523528) (-0.25758,-0.475620) (-0.19697,-0.427916) (-0.13636,-0.380416) (-0.07576,-0.333121) (-0.01515,-0.286033) (0.04545,-0.239153) (0.10606,-0.192482) (0.16667,-0.146020) (0.22727,-0.099769) (0.28788,-0.053728) (0.34848,-0.007898) (0.40909,0.037721) (0.46970,0.083127) (0.53030,0.128322) (0.59091,0.173305) (0.65152,0.218076) (0.71212,0.262636) (0.77273,0.306985) (0.83333,0.351123) (0.89394,0.395050) (0.95455,0.438769) (1.01515,0.482278) (1.07576,0.525580) (1.13636,0.568675) (1.19697,0.611565) (1.25758,0.654250) (1.31818,0.696731) (1.37879,0.739011) (1.43939,0.781091) (1.50000,0.822971) (1.56061,0.864655) (1.62121,0.906143) (1.68182,0.947437) (1.74242,0.988540) (1.80303,1.029452) (1.86364,1.070177) (1.92424,1.110716) (1.98485,1.151071) (2.04545,1.191244) (2.10606,1.231238) (2.16667,1.271054) (2.22727,1.310696) (2.28788,1.350165) (2.34848,1.389464) (2.40909,1.428594) (2.46970,1.467560) (2.53030,1.506362) (2.59091,1.545003) (2.65152,1.583486) (2.71212,1.621812) (2.77273,1.659986) (2.83333,1.698008) (2.89394,1.735882) (2.95455,1.773610) (3.01515,1.811193) (3.07576,1.848636) (3.13636,1.885939) (3.19697,1.923107) (3.25758,1.960139) (3.31818,1.997041) (3.37879,2.033812) (3.43939,2.070457) (3.50000,2.106977)};
\addplot[gray,dotted,thick] coordinates {(-2.50000,0.459690) (-2.43939,0.496210) (-2.37879,0.532854) (-2.31818,0.569626) (-2.25758,0.606527) (-2.19697,0.643560) (-2.13636,0.680727) (-2.07576,0.718031) (-2.01515,0.755473) (-1.95455,0.793057) (-1.89394,0.830785) (-1.83333,0.868658) (-1.77273,0.906681) (-1.71212,0.944854) (-1.65152,0.983181) (-1.59091,1.021664) (-1.53030,1.060305) (-1.46970,1.099107) (-1.40909,1.138072) (-1.34848,1.177203) (-1.28788,1.216502) (-1.22727,1.255971) (-1.16667,1.295612) (-1.10606,1.335429) (-1.04545,1.375423) (-0.98485,1.415596) (-0.92424,1.455951) (-0.86364,1.496490) (-0.80303,1.537215) (-0.74242,1.578127) (-0.68182,1.619230) (-0.62121,1.660524) (-0.56061,1.702012) (-0.50000,1.743695) (-0.43939,1.785576) (-0.37879,1.827655) (-0.31818,1.869935) (-0.25758,1.912417) (-0.19697,1.955102) (-0.13636,1.997991) (-0.07576,2.041086) (-0.01515,2.084388) (0.04545,2.127898) (0.10606,2.171616) (0.16667,2.215544) (0.22727,2.259682) (0.28788,2.304031) (0.34848,2.348590) (0.40909,2.393362) (0.46970,2.438345) (0.53030,2.483539) (0.59091,2.528946) (0.65152,2.574564) (0.71212,2.620394) (0.77273,2.666435) (0.83333,2.712687) (0.89394,2.759149) (0.95455,2.805820) (1.01515,2.852700) (1.07576,2.899788) (1.13636,2.947082) (1.19697,2.994582) (1.25758,3.042287) (1.31818,3.090195) (1.37879,3.138305) (1.43939,3.186615) (1.50000,3.235124) (1.56061,3.283830) (1.62121,3.332732) (1.68182,3.381827) (1.74242,3.431114) (1.80303,3.480591) (1.86364,3.530256) (1.92424,3.580107) (1.98485,3.630142) (2.04545,3.680358) (2.10606,3.730754) (2.16667,3.781327) (2.22727,3.832075) (2.28788,3.882995) (2.34848,3.934086) (2.40909,3.985345) (2.46970,4.036769) (2.53030,4.088357) (2.59091,4.140105) (2.65152,4.192012) (2.71212,4.244075) (2.77273,4.296291) (2.83333,4.348658) (2.89394,4.401174) (2.95455,4.453836) (3.01515,4.506642) (3.07576,4.559589) (3.13636,4.612675) (3.19697,4.665898) (3.25758,4.719254) (3.31818,4.772743) (3.37879,4.826361) (3.43939,4.880106) (3.50000,4.933975)};
\addlegendentry{预测区间边界}
\end{axis}\end{tikzpicture}\end{center}

图中使用六个编订示例点，在独立正态同方差模型下绘出逐点 $95\%$ 均值响应区间和新个体预测区间。点线预测带较宽，因为新个体还有自身误差；离样本横坐标中心越远，两类区间都变宽。

这里是每个固定 $x$ 各自的逐点区间，不是整条曲线同时覆盖率为 $95\%$ 的置信带。若同时查看许多位置，不能把逐点覆盖解释为整体覆盖。参见习题~\ref{q:B16.4}。

\clearpage
\section*{题源、许可与编订记录}
\addcontentsline{toc}{section}{题源、许可与编订记录}
\subsection*{中文公开题源}
\textbf{Y：闫一博（yyb100001），《何书元〈概率引论〉答案》。}本册选取其公开 \texttt{main.tex} 中的数学问题与解答，保留原题号作为定位键，重排题面与解答并校订。项目署名为“Solutions-of-Introduction-to-Probability © 2024 by yyb”，公开声明采用 CC BY 4.0。其问题背景来源于何书元《概率引论》；本册未冒称这些教材题目为公开解答作者或编订者的原创。项目链接：\href{https://github.com/yyb100001/Solutions-of-Introduction-to-Probability}{公开仓库及许可说明}。许可链接：\href{https://creativecommons.org/licenses/by/4.0/}{知识共享署名 4.0 国际许可}。

\textbf{A：ayhe123，《概率论讲义及习题解答》。}本册选取公开 \texttt{lecturenote.tex} 的习题与证明，保留原题号，删除依赖未收入讲义的外部交叉引用，并补齐相关推导。项目公开说明其笔记与仿真采用 CC BY 4.0。该讲义基于何书元《概率引论》及相应课程。项目链接：\href{https://github.com/ayhe123/probabilityTheory-lecturenote}{公开仓库及许可说明}；许可同上。中文题源均为学习者编写的公开解答，不把它们等同于教材官方答案。

\subsection*{MIT 开放课程题源}
\textbf{M：Jeremy Orloff、Jennifer French Kamrin，MIT OpenCourseWare 18.05，2022 年春季。}课程原名为 \emph{Introduction to Probability and Statistics}。本册将所选题目与答案翻译、重排为中文，补充解释并校订数值与条件；删除了考试分数、交卷说明、软件选择题的非必要细节及依赖未收录图形的题目。涉及原资料的情境或数据，仅作为该课教学材料，不表述为当前现实统计事实。

课程主页：\href{https://ocw.mit.edu/courses/18-05-introduction-to-probability-and-statistics-spring-2022/}{MIT OCW 18.05}。本册所用课程材料在当前使用条款下采用 \href{https://creativecommons.org/licenses/by-nc-sa/4.0/}{CC BY-NC-SA 4.0（署名—非商业性使用—相同方式共享）}；\href{https://ocw.mit.edu/pages/privacy-and-terms-of-use/}{原站使用条款}提供许可和署名要求。本册不是 MIT 的官方译本，MIT 与原作者不对本册校订和新增解答提供背书。

\textbf{B：编订补充。}为衔接本科数理统计课程而新增的标准模型题、证明题和综合题，均在逐题来源处标为“编订补充”。标注“编订延伸”的多问问题，其原始部分来自列明的课程题，新增部分由编订者给出。方法说明、图解和数值表也是本册的新增整理，不声称在原资料中原样出现。

本册 MIT 材料的中文改编与新增编订内容按 CC BY-NC-SA 4.0 提供；中文公开项目的原有 CC BY 4.0 署名和许可信息保留。各部分原有权利归其相应权利人。再使用时应保留来源、许可链接与修改说明，不据本册新增编排改变原材料自身的许可。检索与版本核对日期：2026 年 10 月 9 日。

\subsection*{课程文件定位表}
下列缩写与逐题索引配合使用。标题中的“题”对应原文件的题号；本册重排题号只用于内部学习顺序。


\noindent\textbf{C10}：\href{https://ocw.mit.edu/courses/18-05-introduction-to-probability-and-statistics-spring-2022/mit18_05_s22_class10_pset_sol.pdf}{第 10 讲课堂题及解答}。\par

\noindent\textbf{C22}：\href{https://ocw.mit.edu/courses/18-05-introduction-to-probability-and-statistics-spring-2022/mit18_05_s22_class22_pset_sol.pdf}{第 22 讲课堂题及解答}。\par

\noindent\textbf{C26}：\href{https://ocw.mit.edu/courses/18-05-introduction-to-probability-and-statistics-spring-2022/mit18_05_s22_class26_pset_sol.pdf}{第 26 讲课堂题及解答}。\par

\noindent\textbf{P7}：\href{https://ocw.mit.edu/courses/18-05-introduction-to-probability-and-statistics-spring-2022/mit18_05_s22_pset07_sol.pdf}{习题集 7及解答}。\par

\noindent\textbf{P8}：\href{https://ocw.mit.edu/courses/18-05-introduction-to-probability-and-statistics-spring-2022/mit18_05_s22_pset08_sol.pdf}{习题集 8及解答}。\par

\noindent\textbf{P9}：\href{https://ocw.mit.edu/courses/18-05-introduction-to-probability-and-statistics-spring-2022/mit18_05_s22_pset09_sol.pdf}{习题集 9及解答}。\par

\noindent\textbf{P10}：\href{https://ocw.mit.edu/courses/18-05-introduction-to-probability-and-statistics-spring-2022/mit18_05_s22_pset10_sol.pdf}{习题集 10及解答}。\par

\noindent\textbf{P11}：\href{https://ocw.mit.edu/courses/18-05-introduction-to-probability-and-statistics-spring-2022/mit18_05_s22_pset11_sol.pdf}{习题集 11及解答}。\par

\noindent\textbf{E2}：\href{https://ocw.mit.edu/courses/18-05-introduction-to-probability-and-statistics-spring-2022/mit18_05_s22_exam02_sol.pdf}{第二次考试及解答}。\par

\noindent\textbf{F}：\href{https://ocw.mit.edu/courses/18-05-introduction-to-probability-and-statistics-spring-2022/mit18_05_s22_exam_final_sol.pdf}{期末考试及解答}。\par

\noindent\textbf{PF}：\href{https://ocw.mit.edu/courses/18-05-introduction-to-probability-and-statistics-spring-2022/mit18_05_s22_prac_exam_final_sol.pdf}{模拟期末考试及解答}。\par

\subsection*{主要勘误与条件补全}
以下是影响结论或推导的代表性修正。索引中的原题号用于返回公开题源核对；本册的译写和校订责任由编订承担。
\begin{enumerate}
\item Y1.9：恰好一副手套后，另外两只应选自两副；修正组合数的文字解释。
\item Y1.13：圆心须远离格线至少半径；补充方格边长不超过硬币直径的情形。
\item Y2.23、Y2.25：补明确给定真实状态后的条件独立，以未舍入先验直接进行多次更新。
\item Y3.6：计算 $P(X\ge3)$ 的补事件必须包含零次成功项。
\item Y3.20、Y3.26、Y3.28、Y3.29、Y4.25、Y4.36 等：用分布函数或密度变换替代把连续单点概率等同于密度的写法。
\item Y3.21：参数 $n=19,p=0.7$ 时，二项众数为 $13$ 与 $14$ 两个。
\item A4.19：混合密度的求导结果必须保留有限求和。
\item A4.23：变换后的支撑为 $0<w<z$，边缘积分贡献额外的 $z$，正确密度为 $z^2e^{-z}/2$。
\item A5.14：按原解的积分模型校订联合密度为 $3/(2x^3y^2)$，原题面幂次在该区域不能归一化。
\item A5.35：圆周均匀与圆内均匀的坐标方差分别为 $R^2/2$ 与 $R^2/4$。
\item A6.15：补全伽马—泊松条件分布的解答，并将不恰当的连续金额泊松表述改为计数模型。
\item A6.17：第三小指数寿命密度中应为 $e^{-5\lambda y}$，不是 $e^{-2\lambda y}$。
\item A6.20、A6.24：校正密度与生成函数证明的变量和下标。
\item A7.5：切比雪夫给出的是下界，不能将最终不等号写成等号。
\item A7.20、A7.21、A7.22、A7.25：补全原资料缺失或未结束的极限定理论证。
\item C22、P10：统一下侧分位数约定，复算方差、标准差和均值区间，校正舍入差异。
\item F11：中央区间是不拒绝区域，外侧才是拒绝域。
\item PF30：正态统计量为 $10$ 时，双侧尾概率极小但严格为正，不写作精确的零。
\end{enumerate}

\clearpage
\section*{逐题主题与来源索引}
\addcontentsline{toc}{section}{逐题主题与来源索引}
各题的题目与详解连续排列。右栏“题／解”依次给出题目和详解的起始页码，均可点击跳转；两者可能位于同一页，也可能因内容较长而跨页。Y、A 后的数字为公开中文资料原题号；M 后接上表文件缩写与原题位置；B 表示编订补充。题源为选编与校订，不保证与原资料逐字相同。
\begin{longtable}{@{}p{13mm}p{72mm}r@{}}
\toprule 本册题号 & 主题与来源键 & 题／解\\\midrule\endfirsthead
\toprule 本册题号 & 主题与来源键 & 题／解\\\midrule\endhead


\ref{q:Y1.1} & 事件分解与德摩根律\newline {\small Y1.1} & \pageref{q:Y1.1}/\pageref{a:Y1.1}\\[2pt]

\ref{q:Y1.2} & 用集合语言表示事件\newline {\small Y1.2} & \pageref{q:Y1.2}/\pageref{a:Y1.2}\\[2pt]

\ref{q:Y1.3} & 有放回抽球与样本空间\newline {\small Y1.3} & \pageref{q:Y1.3}/\pageref{a:Y1.3}\\[2pt]

\ref{q:Y1.4} & 将事件列化为互斥事件列\newline {\small Y1.4} & \pageref{q:Y1.4}/\pageref{a:Y1.4}\\[2pt]

\ref{q:Y1.5} & 两阶段抽样的对称性\newline {\small Y1.5} & \pageref{q:Y1.5}/\pageref{a:Y1.5}\\[2pt]

\ref{q:Y1.7} & 无放回试钥匙\newline {\small Y1.7} & \pageref{q:Y1.7}/\pageref{a:Y1.7}\\[2pt]

\ref{q:Y1.8} & 生日相同的概率\newline {\small Y1.8} & \pageref{q:Y1.8}/\pageref{a:Y1.8}\\[2pt]

\ref{q:Y1.9} & 恰好得到一副手套\newline {\small Y1.9} & \pageref{q:Y1.9}/\pageref{a:Y1.9}\\[2pt]

\ref{q:Y1.10} & 两次抽样的交集\newline {\small Y1.10} & \pageref{q:Y1.10}/\pageref{a:Y1.10}\\[2pt]

\ref{q:Y1.11} & 单位正方形中的圆域\newline {\small Y1.11} & \pageref{q:Y1.11}/\pageref{a:Y1.11}\\[2pt]

\ref{q:Y1.12} & 到达时差与等待\newline {\small Y1.12} & \pageref{q:Y1.12}/\pageref{a:Y1.12}\\[2pt]

\ref{q:Y1.13} & 硬币落格与边界区域\newline {\small Y1.13} & \pageref{q:Y1.13}/\pageref{a:Y1.13}\\[2pt]

\ref{q:Y1.14} & 两次服务是否相遇\newline {\small Y1.14} & \pageref{q:Y1.14}/\pageref{a:Y1.14}\\[2pt]

\ref{q:Y1.15} & 立方体内的球域\newline {\small Y1.15} & \pageref{q:Y1.15}/\pageref{a:Y1.15}\\[2pt]

\ref{q:Y1.17} & 概率的次可加性\newline {\small Y1.17} & \pageref{q:Y1.17}/\pageref{a:Y1.17}\\[2pt]

\ref{q:Y1.18} & 可数个概率一事件\newline {\small Y1.18} & \pageref{q:Y1.18}/\pageref{a:Y1.18}\\[2pt]

\ref{q:Y1.19} & 四个互斥区域的分割\newline {\small Y1.19} & \pageref{q:Y1.19}/\pageref{a:Y1.19}\\[2pt]

\ref{q:Y1.20} & 递增事件列的极限\newline {\small Y1.20} & \pageref{q:Y1.20}/\pageref{a:Y1.20}\\[2pt]

\ref{q:Y2.1} & 至少一次求职成功\newline {\small Y2.1} & \pageref{q:Y2.1}/\pageref{a:Y2.1}\\[2pt]

\ref{q:Y2.2} & 自习室占用问题\newline {\small Y2.2} & \pageref{q:Y2.2}/\pageref{a:Y2.2}\\[2pt]

\ref{q:Y2.3} & 无人拿到自己帽子\newline {\small Y2.3} & \pageref{q:Y2.3}/\pageref{a:Y2.3}\\[2pt]

\ref{q:Y2.4} & 有放回点名的覆盖概率\newline {\small Y2.4} & \pageref{q:Y2.4}/\pageref{a:Y2.4}\\[2pt]

\ref{q:Y2.5} & 信封匹配与容斥\newline {\small Y2.5} & \pageref{q:Y2.5}/\pageref{a:Y2.5}\\[2pt]

\ref{q:Y2.6} & 独立任务的共同成功\newline {\small Y2.6} & \pageref{q:Y2.6}/\pageref{a:Y2.6}\\[2pt]

\ref{q:Y2.7} & 发现信号的试验次数\newline {\small Y2.7} & \pageref{q:Y2.7}/\pageref{a:Y2.7}\\[2pt]

\ref{q:Y2.9} & 保证成功概率的人员数\newline {\small Y2.9} & \pageref{q:Y2.9}/\pageref{a:Y2.9}\\[2pt]

\ref{q:Y2.10} & 德梅雷的骰子问题\newline {\small Y2.10} & \pageref{q:Y2.10}/\pageref{a:Y2.10}\\[2pt]

\ref{q:Y2.11} & 条件成绩概率\newline {\small Y2.11} & \pageref{q:Y2.11}/\pageref{a:Y2.11}\\[2pt]

\ref{q:Y2.12} & 连续受损的生存概率\newline {\small Y2.12} & \pageref{q:Y2.12}/\pageref{a:Y2.12}\\[2pt]

\ref{q:Y2.13} & 已知一枚骰子的点数\newline {\small Y2.13} & \pageref{q:Y2.13}/\pageref{a:Y2.13}\\[2pt]

\ref{q:Y2.15} & 两孩家庭中的条件事件\newline {\small Y2.15} & \pageref{q:Y2.15}/\pageref{a:Y2.15}\\[2pt]

\ref{q:Y2.16} & 通知渠道的全概率\newline {\small Y2.16} & \pageref{q:Y2.16}/\pageref{a:Y2.16}\\[2pt]

\ref{q:Y2.17} & 隐私调查的随机化回答\newline {\small Y2.17} & \pageref{q:Y2.17}/\pageref{a:Y2.17}\\[2pt]

\ref{q:Y2.19} & 比分差达到二的比赛\newline {\small Y2.19} & \pageref{q:Y2.19}/\pageref{a:Y2.19}\\[2pt]

\ref{q:Y2.20} & 加强抽样的波利亚球袋\newline {\small Y2.20} & \pageref{q:Y2.20}/\pageref{a:Y2.20}\\[2pt]

\ref{q:Y2.21} & 已知正点反推交通方式\newline {\small Y2.21} & \pageref{q:Y2.21}/\pageref{a:Y2.21}\\[2pt]

\ref{q:Y2.22} & 观察合格品后的状态\newline {\small Y2.22} & \pageref{q:Y2.22}/\pageref{a:Y2.22}\\[2pt]

\ref{q:Y2.23} & 低先验概率下的连续检测\newline {\small Y2.23} & \pageref{q:Y2.23}/\pageref{a:Y2.23}\\[2pt]

\ref{q:Y2.25} & 目击证据与先验比例\newline {\small Y2.25} & \pageref{q:Y2.25}/\pageref{a:Y2.25}\\[2pt]

\ref{q:Y2.26} & 无穷多次成功的条件\newline {\small Y2.26} & \pageref{q:Y2.26}/\pageref{a:Y2.26}\\[2pt]

\ref{q:Y2.27} & 最终不再失败的条件\newline {\small Y2.27} & \pageref{q:Y2.27}/\pageref{a:Y2.27}\\[2pt]

\ref{q:Y2.28} & 加红球后的黑球出现次数\newline {\small Y2.28} & \pageref{q:Y2.28}/\pageref{a:Y2.28}\\[2pt]

\ref{q:Y2.29} & 黑红同时增长的抽球\newline {\small Y2.29} & \pageref{q:Y2.29}/\pageref{a:Y2.29}\\[2pt]

\ref{q:Y2.30} & 全概率公式的证明\newline {\small Y2.30} & \pageref{q:Y2.30}/\pageref{a:Y2.30}\\[2pt]

\ref{q:Y2.31} & 贝叶斯公式的证明\newline {\small Y2.31} & \pageref{q:Y2.31}/\pageref{a:Y2.31}\\[2pt]

\ref{q:Y2.32} & 串联与并联可靠性\newline {\small Y2.32} & \pageref{q:Y2.32}/\pageref{a:Y2.32}\\[2pt]

\ref{q:Y2.33} & 比赛局数与胜率\newline {\small Y2.33} & \pageref{q:Y2.33}/\pageref{a:Y2.33}\\[2pt]

\ref{q:Y2.34} & 集体迟到的提前量\newline {\small Y2.34} & \pageref{q:Y2.34}/\pageref{a:Y2.34}\\[2pt]

\ref{q:Y3.1} & 分布函数的无穷远极限\newline {\small Y3.1} & \pageref{q:Y3.1}/\pageref{a:Y3.1}\\[2pt]

\ref{q:Y3.2} & 指数型密度的归一化\newline {\small Y3.2} & \pageref{q:Y3.2}/\pageref{a:Y3.2}\\[2pt]

\ref{q:Y3.3} & 二次指数可积的条件\newline {\small Y3.3} & \pageref{q:Y3.3}/\pageref{a:Y3.3}\\[2pt]

\ref{q:Y3.4} & 最大值与最小值恒等式\newline {\small Y3.4} & \pageref{q:Y3.4}/\pageref{a:Y3.4}\\[2pt]

\ref{q:Y3.5} & 分段密度的归一化\newline {\small Y3.5} & \pageref{q:Y3.5}/\pageref{a:Y3.5}\\[2pt]

\ref{q:Y3.6} & 二项计数的尾概率\newline {\small Y3.6} & \pageref{q:Y3.6}/\pageref{a:Y3.6}\\[2pt]

\ref{q:Y3.7} & 两组二项计数的比较\newline {\small Y3.7} & \pageref{q:Y3.7}/\pageref{a:Y3.7}\\[2pt]

\ref{q:Y3.8} & 首次与第多次成功\newline {\small Y3.8} & \pageref{q:Y3.8}/\pageref{a:Y3.8}\\[2pt]

\ref{q:Y3.9} & 泊松分裂与独立性\newline {\small Y3.9} & \pageref{q:Y3.9}/\pageref{a:Y3.9}\\[2pt]

\ref{q:Y3.13} & 周期信号的均匀等待\newline {\small Y3.13} & \pageref{q:Y3.13}/\pageref{a:Y3.13}\\[2pt]

\ref{q:Y3.14} & 恒定风险率与指数寿命\newline {\small Y3.14} & \pageref{q:Y3.14}/\pageref{a:Y3.14}\\[2pt]

\ref{q:Y3.15} & 正态尺寸的合格概率\newline {\small Y3.15} & \pageref{q:Y3.15}/\pageref{a:Y3.15}\\[2pt]

\ref{q:Y3.16} & 泊松变量的平方根\newline {\small Y3.16} & \pageref{q:Y3.16}/\pageref{a:Y3.16}\\[2pt]

\ref{q:Y3.17} & 均匀电流的平方变换\newline {\small Y3.17} & \pageref{q:Y3.17}/\pageref{a:Y3.17}\\[2pt]

\ref{q:Y3.18} & 指数变量的仿射变换\newline {\small Y3.18} & \pageref{q:Y3.18}/\pageref{a:Y3.18}\\[2pt]

\ref{q:Y3.19} & 半柯西变量的对数\newline {\small Y3.19} & \pageref{q:Y3.19}/\pageref{a:Y3.19}\\[2pt]

\ref{q:Y3.20} & 倒数、绝对值与正切变换\newline {\small Y3.20} & \pageref{q:Y3.20}/\pageref{a:Y3.20}\\[2pt]

\ref{q:Y3.21} & 二项众数与参数估计\newline {\small Y3.21} & \pageref{q:Y3.21}/\pageref{a:Y3.21}\\[2pt]

\ref{q:Y3.22} & 迟到人数的众数和估计\newline {\small Y3.22} & \pageref{q:Y3.22}/\pageref{a:Y3.22}\\[2pt]

\ref{q:Y3.23} & 泊松众数与似然\newline {\small Y3.23} & \pageref{q:Y3.23}/\pageref{a:Y3.23}\\[2pt]

\ref{q:Y3.24} & 计数结果与参数的区别\newline {\small Y3.24} & \pageref{q:Y3.24}/\pageref{a:Y3.24}\\[2pt]

\ref{q:Y3.25} & 骰子样本的最大最小值\newline {\small Y3.25} & \pageref{q:Y3.25}/\pageref{a:Y3.25}\\[2pt]

\ref{q:Y3.26} & 圆周投影的反正弦密度\newline {\small Y3.26} & \pageref{q:Y3.26}/\pageref{a:Y3.26}\\[2pt]

\ref{q:Y3.28} & 正弦变换的两条分支\newline {\small Y3.28} & \pageref{q:Y3.28}/\pageref{a:Y3.28}\\[2pt]

\ref{q:Y3.29} & 逆分布函数生成正态数\newline {\small Y3.29} & \pageref{q:Y3.29}/\pageref{a:Y3.29}\\[2pt]

\ref{q:Y3.31} & 分段单调变换\newline {\small Y3.31} & \pageref{q:Y3.31}/\pageref{a:Y3.31}\\[2pt]

\ref{q:Y4.3} & 联合概率表与边缘分布\newline {\small Y4.3} & \pageref{q:Y4.3}/\pageref{a:Y4.3}\\[2pt]

\ref{q:Y4.6} & 同参数二项变量之和\newline {\small Y4.6} & \pageref{q:Y4.6}/\pageref{a:Y4.6}\\[2pt]

\ref{q:Y4.8} & 函数依赖的离散向量\newline {\small Y4.8} & \pageref{q:Y4.8}/\pageref{a:Y4.8}\\[2pt]

\ref{q:Y4.9} & 最小值的尾概率\newline {\small Y4.9} & \pageref{q:Y4.9}/\pageref{a:Y4.9}\\[2pt]

\ref{q:Y4.15} & 三角支撑上的边缘密度\newline {\small Y4.15} & \pageref{q:Y4.15}/\pageref{a:Y4.15}\\[2pt]

\ref{q:Y4.16} & 随机家庭规模与男孩数\newline {\small Y4.16} & \pageref{q:Y4.16}/\pageref{a:Y4.16}\\[2pt]

\ref{q:Y4.22} & 连续变量的离散随机缩放\newline {\small Y4.22} & \pageref{q:Y4.22}/\pageref{a:Y4.22}\\[2pt]

\ref{q:Y4.25} & 博克斯—穆勒正态生成\newline {\small Y4.25} & \pageref{q:Y4.25}/\pageref{a:Y4.25}\\[2pt]

\ref{q:Y4.30} & 指数变量的极值与比值\newline {\small Y4.30} & \pageref{q:Y4.30}/\pageref{a:Y4.30}\\[2pt]

\ref{q:Y4.34} & 平方和与平方比的变换\newline {\small Y4.34} & \pageref{q:Y4.34}/\pageref{a:Y4.34}\\[2pt]

\ref{q:Y4.36} & 线性变换的密度公式\newline {\small Y4.36} & \pageref{q:Y4.36}/\pageref{a:Y4.36}\\[2pt]

\ref{q:A4.4} & 常数与任意变量独立\newline {\small A4.4} & \pageref{q:A4.4}/\pageref{a:A4.4}\\[2pt]

\ref{q:A4.5} & 联合分布函数的右连续性\newline {\small A4.5} & \pageref{q:A4.5}/\pageref{a:A4.5}\\[2pt]

\ref{q:A4.11} & 曲边区域上的联合密度\newline {\small A4.11} & \pageref{q:A4.11}/\pageref{a:A4.11}\\[2pt]

\ref{q:A4.14} & 指数寿命的比较\newline {\small A4.14} & \pageref{q:A4.14}/\pageref{a:A4.14}\\[2pt]

\ref{q:A4.17} & 不独立变量的平方独立\newline {\small A4.17} & \pageref{q:A4.17}/\pageref{a:A4.17}\\[2pt]

\ref{q:A4.18} & 曲线下均匀点的投影\newline {\small A4.18} & \pageref{q:A4.18}/\pageref{a:A4.18}\\[2pt]

\ref{q:A4.19} & 离散与连续变量之差\newline {\small A4.19} & \pageref{q:A4.19}/\pageref{a:A4.19}\\[2pt]

\ref{q:A4.23} & 求和变换后的边缘积分\newline {\small A4.23} & \pageref{q:A4.23}/\pageref{a:A4.23}\\[2pt]

\ref{q:A4.27} & 独立正态的线性组合\newline {\small A4.27} & \pageref{q:A4.27}/\pageref{a:A4.27}\\[2pt]

\ref{q:A4.28} & 二维正态的条件分布\newline {\small A4.28} & \pageref{q:A4.28}/\pageref{a:A4.28}\\[2pt]

\ref{q:A4.32} & 尾分布幂次与竞争概率\newline {\small A4.32} & \pageref{q:A4.32}/\pageref{a:A4.32}\\[2pt]

\ref{q:A5.5} & 广告信息之间的等待\newline {\small A5.5} & \pageref{q:A5.5}/\pageref{a:A5.5}\\[2pt]

\ref{q:A5.8} & 标准正态半径的期望\newline {\small A5.8} & \pageref{q:A5.8}/\pageref{a:A5.8}\\[2pt]

\ref{q:A5.9} & 交换对称与比例期望\newline {\small A5.9} & \pageref{q:A5.9}/\pageref{a:A5.9}\\[2pt]

\ref{q:A5.12} & 离散极值的尾和公式\newline {\small A5.12} & \pageref{q:A5.12}/\pageref{a:A5.12}\\[2pt]

\ref{q:A5.14} & 曲边密度的归一化勘误\newline {\small A5.14} & \pageref{q:A5.14}/\pageref{a:A5.14}\\[2pt]

\ref{q:A5.17} & 全部指数寿命结束的等待\newline {\small A5.17} & \pageref{q:A5.17}/\pageref{a:A5.17}\\[2pt]

\ref{q:A5.20} & 样本均值与方差的无偏性\newline {\small A5.20} & \pageref{q:A5.20}/\pageref{a:A5.20}\\[2pt]

\ref{q:A5.22} & 最优逆方差加权\newline {\small A5.22} & \pageref{q:A5.22}/\pageref{a:A5.22}\\[2pt]

\ref{q:A5.25} & 非负变量的尾积分\newline {\small A5.25} & \pageref{q:A5.25}/\pageref{a:A5.25}\\[2pt]

\ref{q:A5.26} & 和与差的相关系数\newline {\small A5.26} & \pageref{q:A5.26}/\pageref{a:A5.26}\\[2pt]

\ref{q:A5.28} & 正方形密度的协方差\newline {\small A5.28} & \pageref{q:A5.28}/\pageref{a:A5.28}\\[2pt]

\ref{q:A5.30} & 正态尺寸差与装配\newline {\small A5.30} & \pageref{q:A5.30}/\pageref{a:A5.30}\\[2pt]

\ref{q:A5.32} & 不相关但不独立的反例\newline {\small A5.32} & \pageref{q:A5.32}/\pageref{a:A5.32}\\[2pt]

\ref{q:A5.33} & 中心正态的高阶矩\newline {\small A5.33} & \pageref{q:A5.33}/\pageref{a:A5.33}\\[2pt]

\ref{q:A5.35} & 圆周与圆内均匀模型\newline {\small A5.35} & \pageref{q:A5.35}/\pageref{a:A5.35}\\[2pt]

\ref{q:A6.1} & 随机组数下的条件二项\newline {\small A6.1} & \pageref{q:A6.1}/\pageref{a:A6.1}\\[2pt]

\ref{q:A6.3} & 几何个指数等待之和\newline {\small A6.3} & \pageref{q:A6.3}/\pageref{a:A6.3}\\[2pt]

\ref{q:A6.4} & 复合泊松等待的均值方差\newline {\small A6.4} & \pageref{q:A6.4}/\pageref{a:A6.4}\\[2pt]

\ref{q:A6.9} & 条件指数与条件伽马\newline {\small A6.9} & \pageref{q:A6.9}/\pageref{a:A6.9}\\[2pt]

\ref{q:A6.14} & 威布尔寿命的剩余等待\newline {\small A6.14} & \pageref{q:A6.14}/\pageref{a:A6.14}\\[2pt]

\ref{q:A6.15} & 泊松观察后的伽马后验\newline {\small A6.15} & \pageref{q:A6.15}/\pageref{a:A6.15}\\[2pt]

\ref{q:A6.17} & 七个接口的第三次失效\newline {\small A6.17} & \pageref{q:A6.17}/\pageref{a:A6.17}\\[2pt]

\ref{q:A6.20} & 指数顺序统计量的间隔\newline {\small A6.20} & \pageref{q:A6.20}/\pageref{a:A6.20}\\[2pt]

\ref{q:A6.21} & 仿射变换的概率母函数\newline {\small A6.21} & \pageref{q:A6.21}/\pageref{a:A6.21}\\[2pt]

\ref{q:A6.24} & 两组硬币成功数之差\newline {\small A6.24} & \pageref{q:A6.24}/\pageref{a:A6.24}\\[2pt]

\ref{q:A6.29} & 均匀和的特征函数极限\newline {\small A6.29} & \pageref{q:A6.29}/\pageref{a:A6.29}\\[2pt]

\ref{q:A6.30} & 指数和的特征函数极限\newline {\small A6.30} & \pageref{q:A6.30}/\pageref{a:A6.30}\\[2pt]

\ref{q:A6.34} & 二项分布的泊松极限\newline {\small A6.34} & \pageref{q:A6.34}/\pageref{a:A6.34}\\[2pt]

\ref{q:A6.35} & 给定总和后的部分和\newline {\small A6.35} & \pageref{q:A6.35}/\pageref{a:A6.35}\\[2pt]

\ref{q:A7.2} & 截尾期望的概率上界\newline {\small A7.2} & \pageref{q:A7.2}/\pageref{a:A7.2}\\[2pt]

\ref{q:A7.3} & 柯西不等式与尾概率\newline {\small A7.3} & \pageref{q:A7.3}/\pageref{a:A7.3}\\[2pt]

\ref{q:A7.4} & 最小值的概率界\newline {\small A7.4} & \pageref{q:A7.4}/\pageref{a:A7.4}\\[2pt]

\ref{q:A7.5} & 伽马矩与切比雪夫界\newline {\small A7.5} & \pageref{q:A7.5}/\pageref{a:A7.5}\\[2pt]

\ref{q:A7.10} & 三角函数样本均值之比\newline {\small A7.10} & \pageref{q:A7.10}/\pageref{a:A7.10}\\[2pt]

\ref{q:A7.14} & 比例调查的样本量\newline {\small A7.14} & \pageref{q:A7.14}/\pageref{a:A7.14}\\[2pt]

\ref{q:A7.15} & 两组样本均值差的近似\newline {\small A7.15} & \pageref{q:A7.15}/\pageref{a:A7.15}\\[2pt]

\ref{q:A7.20} & 泊松概率和的二分极限\newline {\small A7.20} & \pageref{q:A7.20}/\pageref{a:A7.20}\\[2pt]

\ref{q:A7.21} & 缩放收敛与紧性\newline {\small A7.21} & \pageref{q:A7.21}/\pageref{a:A7.21}\\[2pt]

\ref{q:A7.22} & 不同参数二项和的极限\newline {\small A7.22} & \pageref{q:A7.22}/\pageref{a:A7.22}\\[2pt]

\ref{q:A7.25} & 小扰动不改变分布极限\newline {\small A7.25} & \pageref{q:A7.25}/\pageref{a:A7.25}\\[2pt]

\ref{q:M-F-C4} & 统计量与未知参数\newline {\small M-F-C4} & \pageref{q:M-F-C4}/\pageref{a:M-F-C4}\\[2pt]

\ref{q:B8.1} & 样本均值的均值与方差\newline {\small B8.1} & \pageref{q:B8.1}/\pageref{a:B8.1}\\[2pt]

\ref{q:B8.2} & 样本方差的无偏性\newline {\small B8.2} & \pageref{q:B8.2}/\pageref{a:B8.2}\\[2pt]

\ref{q:B8.3} & 正态样本均值和残差的独立性\newline {\small B8.3} & \pageref{q:B8.3}/\pageref{a:B8.3}\\[2pt]

\ref{q:B8.4} & 卡方分布及自由度\newline {\small B8.4} & \pageref{q:B8.4}/\pageref{a:B8.4}\\[2pt]

\ref{q:B8.5} & 学生 t 统计量\newline {\small B8.5} & \pageref{q:B8.5}/\pageref{a:B8.5}\\[2pt]

\ref{q:B8.6} & 方差比与 F 分布\newline {\small B8.6} & \pageref{q:B8.6}/\pageref{a:B8.6}\\[2pt]

\ref{q:B8.7} & 顺序统计量的分布\newline {\small B8.7} & \pageref{q:B8.7}/\pageref{a:B8.7}\\[2pt]

\ref{q:B8.8} & 均匀样本最大值\newline {\small B8.8} & \pageref{q:B8.8}/\pageref{a:B8.8}\\[2pt]

\ref{q:B8.9} & 指数样本的最小值与和\newline {\small B8.9} & \pageref{q:B8.9}/\pageref{a:B8.9}\\[2pt]

\ref{q:M-F-C13} & 中心极限定理的正确标准化\newline {\small M-F-C13} & \pageref{q:M-F-C13}/\pageref{a:M-F-C13}\\[2pt]

\ref{q:B9.1} & 伯努利矩估计\newline {\small B9.1} & \pageref{q:B9.1}/\pageref{a:B9.1}\\[2pt]

\ref{q:B9.2} & 均匀分布的矩估计与极值估计\newline {\small B9.2} & \pageref{q:B9.2}/\pageref{a:B9.2}\\[2pt]

\ref{q:B9.3} & 伽马分布的两个矩参数\newline {\small B9.3} & \pageref{q:B9.3}/\pageref{a:B9.3}\\[2pt]

\ref{q:B9.4} & 正态分布的矩估计\newline {\small B9.4} & \pageref{q:B9.4}/\pageref{a:B9.4}\\[2pt]

\ref{q:B9.5} & 偏差与均方误差分解\newline {\small B9.5} & \pageref{q:B9.5}/\pageref{a:B9.5}\\[2pt]

\ref{q:B9.6} & 样本方差的最小均方误差缩放\newline {\small B9.6} & \pageref{q:B9.6}/\pageref{a:B9.6}\\[2pt]

\ref{q:B9.7} & 一致性与渐近正态性\newline {\small B9.7} & \pageref{q:B9.7}/\pageref{a:B9.7}\\[2pt]

\ref{q:B9.8} & 样本分位数的直观比较\newline {\small B9.8} & \pageref{q:B9.8}/\pageref{a:B9.8}\\[2pt]

\ref{q:M-C10-1} & 离散参数空间中的硬币似然\newline {\small M-C10-1} & \pageref{q:M-C10-1}/\pageref{a:M-C10-1}\\[2pt]

\ref{q:M-C10-2} & 指数寿命的极大似然估计\newline {\small M-C10-2} & \pageref{q:M-C10-2}/\pageref{a:M-C10-2}\\[2pt]

\ref{q:M-E2-II1} & 已知形状的伽马率估计\newline {\small M-E2-II1} & \pageref{q:M-E2-II1}/\pageref{a:M-E2-II1}\\[2pt]

\ref{q:B10.1} & 正态模型的联合极大似然\newline {\small B10.1} & \pageref{q:B10.1}/\pageref{a:B10.1}\\[2pt]

\ref{q:B10.2} & 均匀端点与支撑集\newline {\small B10.2} & \pageref{q:B10.2}/\pageref{a:B10.2}\\[2pt]

\ref{q:B10.3} & 几何分布的估计与边界\newline {\small B10.3} & \pageref{q:B10.3}/\pageref{a:B10.3}\\[2pt]

\ref{q:B10.4} & 带删失的指数寿命\newline {\small B10.4} & \pageref{q:B10.4}/\pageref{a:B10.4}\\[2pt]

\ref{q:B10.5} & 极大似然的不变性\newline {\small B10.5} & \pageref{q:B10.5}/\pageref{a:B10.5}\\[2pt]

\ref{q:B10.6} & 指数率估计的有限样本偏差\newline {\small B10.6} & \pageref{q:B10.6}/\pageref{a:B10.6}\\[2pt]

\ref{q:B11.1} & 伯努利样本的充分统计量\newline {\small B11.1} & \pageref{q:B11.1}/\pageref{a:B11.1}\\[2pt]

\ref{q:B11.2} & 泊松样本的充分性与条件二项分布\newline {\small B11.2} & \pageref{q:B11.2}/\pageref{a:B11.2}\\[2pt]

\ref{q:B11.3} & 拉奥—布莱克改进\newline {\small B11.3} & \pageref{q:B11.3}/\pageref{a:B11.3}\\[2pt]

\ref{q:B11.4} & 二项族的完全性\newline {\small B11.4} & \pageref{q:B11.4}/\pageref{a:B11.4}\\[2pt]

\ref{q:B11.5} & 估计平方概率\newline {\small B11.5} & \pageref{q:B11.5}/\pageref{a:B11.5}\\[2pt]

\ref{q:B11.6} & 伯努利模型的信息量与下界\newline {\small B11.6} & \pageref{q:B11.6}/\pageref{a:B11.6}\\[2pt]

\ref{q:B11.7} & 信息下界并非总能达到\newline {\small B11.7} & \pageref{q:B11.7}/\pageref{a:B11.7}\\[2pt]

\ref{q:B11.8} & 非正则模型中的错误下界\newline {\small B11.8} & \pageref{q:B11.8}/\pageref{a:B11.8}\\[2pt]

\ref{q:M-E2-II2} & 三类别观察的贝叶斯更新\newline {\small M-E2-II2} & \pageref{q:M-E2-II2}/\pageref{a:M-E2-II2}\\[2pt]

\ref{q:M-E2-II3} & 先验赔率与后验赔率\newline {\small M-E2-II3} & \pageref{q:M-E2-II3}/\pageref{a:M-E2-II3}\\[2pt]

\ref{q:M-F-C14} & 贝塔与二项共轭\newline {\small M-F-C14} & \pageref{q:M-F-C14}/\pageref{a:M-F-C14}\\[2pt]

\ref{q:B12.1} & 平方损失下的贝叶斯估计\newline {\small B12.1} & \pageref{q:B12.1}/\pageref{a:B12.1}\\[2pt]

\ref{q:B12.2} & 伽马与泊松共轭的多次观察\newline {\small B12.2} & \pageref{q:B12.2}/\pageref{a:B12.2}\\[2pt]

\ref{q:B12.3} & 正态均值的共轭更新\newline {\small B12.3} & \pageref{q:B12.3}/\pageref{a:B12.3}\\[2pt]

\ref{q:M-P7-1} & 主持人规则改变三门问题\newline {\small M-P7-1} & \pageref{q:M-P7-1}/\pageref{a:M-P7-1}\\[2pt]

\ref{q:M-P7-2} & 未知面数骰子的预测\newline {\small M-P7-2} & \pageref{q:M-P7-2}/\pageref{a:M-P7-2}\\[2pt]

\ref{q:M-P8-1fg} & 显著性与后验概率的区别\newline {\small M-P8-1fg} & \pageref{q:M-P8-1fg}/\pageref{a:M-P8-1fg}\\[2pt]

\ref{q:M-C22-1} & 小样本正态总体的三类区间\newline {\small M-C22-1} & \pageref{q:M-C22-1}/\pageref{a:M-C22-1}\\[2pt]

\ref{q:M-P10-1} & 已知标准差的置信水平比较\newline {\small M-P10-1} & \pageref{q:M-P10-1}/\pageref{a:M-P10-1}\\[2pt]

\ref{q:M-P10-2} & 灌装量的正态均值区间\newline {\small M-P10-2} & \pageref{q:M-P10-2}/\pageref{a:M-P10-2}\\[2pt]

\ref{q:M-P10-3} & 正态方差的区间\newline {\small M-P10-3} & \pageref{q:M-P10-3}/\pageref{a:M-P10-3}\\[2pt]

\ref{q:M-P11-1} & 身高区间与宽度要求\newline {\small M-P11-1} & \pageref{q:M-P11-1}/\pageref{a:M-P11-1}\\[2pt]

\ref{q:M-P11-2} & 顺序统计量构造端点区间\newline {\small M-P11-2} & \pageref{q:M-P11-2}/\pageref{a:M-P11-2}\\[2pt]

\ref{q:M-C22-3} & 比例误差与样本量\newline {\small M-C22-3} & \pageref{q:M-C22-3}/\pageref{a:M-C22-3}\\[2pt]

\ref{q:B13.1} & 指数率的精确区间\newline {\small B13.1} & \pageref{q:B13.1}/\pageref{a:B13.1}\\[2pt]

\ref{q:B13.2} & 均匀端点的单侧区间\newline {\small B13.2} & \pageref{q:B13.2}/\pageref{a:B13.2}\\[2pt]

\ref{q:B13.3} & 两个均值差的合并方差区间\newline {\small B13.3} & \pageref{q:B13.3}/\pageref{a:B13.3}\\[2pt]

\ref{q:B13.4} & 韦尔奇近似区间\newline {\small B13.4} & \pageref{q:B13.4}/\pageref{a:B13.4}\\[2pt]

\ref{q:B13.5} & 配对差的区间\newline {\small B13.5} & \pageref{q:B13.5}/\pageref{a:B13.5}\\[2pt]

\ref{q:M-P10-5} & 置信区间与后验区间不能互换\newline {\small M-P10-5} & \pageref{q:M-P10-5}/\pageref{a:M-P10-5}\\[2pt]

\ref{q:B13.6} & 自助法百分位区间\newline {\small B13.6} & \pageref{q:B13.6}/\pageref{a:B13.6}\\[2pt]

\ref{q:M-P8-1ab} & 硬币的精确双侧检验\newline {\small M-P8-1ab} & \pageref{q:M-P8-1ab}/\pageref{a:M-P8-1ab}\\[2pt]

\ref{q:M-P8-1cd} & 离散拒绝域与功效\newline {\small M-P8-1cd} & \pageref{q:M-P8-1cd}/\pageref{a:M-P8-1cd}\\[2pt]

\ref{q:M-P8-2} & 第一类与第二类错误\newline {\small M-P8-2} & \pageref{q:M-P8-2}/\pageref{a:M-P8-2}\\[2pt]

\ref{q:M-P8-3} & 复合零假设与功效规划\newline {\small M-P8-3} & \pageref{q:M-P8-3}/\pageref{a:M-P8-3}\\[2pt]

\ref{q:M-P8-4} & 直接计算两类错误\newline {\small M-P8-4} & \pageref{q:M-P8-4}/\pageref{a:M-P8-4}\\[2pt]

\ref{q:M-P9-1} & 单侧与双侧 t 检验\newline {\small M-P9-1} & \pageref{q:M-P9-1}/\pageref{a:M-P9-1}\\[2pt]

\ref{q:M-P9-2} & 停止规则改变抽样分布\newline {\small M-P9-2} & \pageref{q:M-P9-2}/\pageref{a:M-P9-2}\\[2pt]

\ref{q:M-E2-II4} & 正态右尾检验的分位数\newline {\small M-E2-II4} & \pageref{q:M-E2-II4}/\pageref{a:M-E2-II4}\\[2pt]

\ref{q:M-F-C10} & 检验与置信区间对偶\newline {\small M-F-C10} & \pageref{q:M-F-C10}/\pageref{a:M-F-C10}\\[2pt]

\ref{q:B14.1} & 奈曼—皮尔逊检验\newline {\small B14.1} & \pageref{q:B14.1}/\pageref{a:B14.1}\\[2pt]

\ref{q:B14.2} & 正态均值的似然比检验\newline {\small B14.2} & \pageref{q:B14.2}/\pageref{a:B14.2}\\[2pt]

\ref{q:B14.3} & 多重检验与邦费罗尼控制\newline {\small B14.3} & \pageref{q:B14.3}/\pageref{a:B14.3}\\[2pt]

\ref{q:B14.4} & 不拒绝不等于等效\newline {\small B14.4} & \pageref{q:B14.4}/\pageref{a:B14.4}\\[2pt]

\ref{q:M-P9-3} & 单个正态总体的方差检验\newline {\small M-P9-3} & \pageref{q:M-P9-3}/\pageref{a:M-P9-3}\\[2pt]

\ref{q:M-E2-II5} & 两个总体的类别比例比较\newline {\small M-E2-II5} & \pageref{q:M-E2-II5}/\pageref{a:M-E2-II5}\\[2pt]

\ref{q:B15.1} & 多项分类的拟合优度\newline {\small B15.1} & \pageref{q:B15.1}/\pageref{a:B15.1}\\[2pt]

\ref{q:B15.2} & 估计参数后的自由度\newline {\small B15.2} & \pageref{q:B15.2}/\pageref{a:B15.2}\\[2pt]

\ref{q:B15.3} & 列联表独立性检验\newline {\small B15.3} & \pageref{q:B15.3}/\pageref{a:B15.3}\\[2pt]

\ref{q:B15.4} & 小样本的费希尔精确检验\newline {\small B15.4} & \pageref{q:B15.4}/\pageref{a:B15.4}\\[2pt]

\ref{q:M-P9-5} & 两个正态方差的 F 检验\newline {\small M-P9-5} & \pageref{q:M-P9-5}/\pageref{a:M-P9-5}\\[2pt]

\ref{q:M-C26-1a} & 三个点的最小二乘直线\newline {\small M-C26-1a} & \pageref{q:M-C26-1a}/\pageref{a:M-C26-1a}\\[2pt]

\ref{q:M-C26-1b} & 二次多项式仍是线性回归\newline {\small M-C26-1b} & \pageref{q:M-C26-1b}/\pageref{a:M-C26-1b}\\[2pt]

\ref{q:M-C26-1c} & 对数尺度的指数拟合\newline {\small M-C26-1c} & \pageref{q:M-C26-1c}/\pageref{a:M-C26-1c}\\[2pt]

\ref{q:M-C26-2d} & 正态误差下的最小二乘与似然\newline {\small M-C26-2d} & \pageref{q:M-C26-2d}/\pageref{a:M-C26-2d}\\[2pt]

\ref{q:B16.1} & 正规方程与残差正交\newline {\small B16.1} & \pageref{q:B16.1}/\pageref{a:B16.1}\\[2pt]

\ref{q:B16.2} & 高斯—马尔可夫最优性\newline {\small B16.2} & \pageref{q:B16.2}/\pageref{a:B16.2}\\[2pt]

\ref{q:B16.3} & 斜率估计的方差与区间\newline {\small B16.3} & \pageref{q:B16.3}/\pageref{a:B16.3}\\[2pt]

\ref{q:B16.4} & 均值响应区间与预测区间\newline {\small B16.4} & \pageref{q:B16.4}/\pageref{a:B16.4}\\[2pt]

\ref{q:B16.5} & 单因素方差分析\newline {\small B16.5} & \pageref{q:B16.5}/\pageref{a:B16.5}\\[2pt]

\ref{q:B16.6} & 相关系数与回归解释度\newline {\small B16.6} & \pageref{q:B16.6}/\pageref{a:B16.6}\\[2pt]

\ref{q:M-F-5} & 帕累托模型从估计到区间\newline {\small M-F-5} & \pageref{q:M-F-5}/\pageref{a:M-F-5}\\[2pt]

\ref{q:M-PF-26} & 瑞利参数与变换统计量\newline {\small M-PF-26} & \pageref{q:M-PF-26}/\pageref{a:M-PF-26}\\[2pt]

\ref{q:M-F-7} & 定位中的连续更新与预测\newline {\small M-F-7} & \pageref{q:M-F-7}/\pageref{a:M-F-7}\\[2pt]

\ref{q:M-F-8} & 六类别概率的综合检验\newline {\small M-F-8} & \pageref{q:M-F-8}/\pageref{a:M-F-8}\\[2pt]

\ref{q:M-F-9} & 两样本检验的完整流程\newline {\small M-F-9} & \pageref{q:M-F-9}/\pageref{a:M-F-9}\\[2pt]

\ref{q:M-F-11} & 双侧检验在指定备择下的功效\newline {\small M-F-11} & \pageref{q:M-F-11}/\pageref{a:M-F-11}\\[2pt]

\ref{q:M-F-12} & 折线基函数的最小二乘\newline {\small M-F-12} & \pageref{q:M-F-12}/\pageref{a:M-F-12}\\[2pt]

\ref{q:M-F-13} & 自助分位数的读法\newline {\small M-F-13} & \pageref{q:M-F-13}/\pageref{a:M-F-13}\\[2pt]

\ref{q:M-PF-18} & 抽球机制与逆向条件概率\newline {\small M-PF-18} & \pageref{q:M-PF-18}/\pageref{a:M-PF-18}\\[2pt]

\ref{q:M-PF-19} & 停时分布的边界计算\newline {\small M-PF-19} & \pageref{q:M-PF-19}/\pageref{a:M-PF-19}\\[2pt]

\ref{q:M-PF-20} & 分段密度与分位数\newline {\small M-PF-20} & \pageref{q:M-PF-20}/\pageref{a:M-PF-20}\\[2pt]

\ref{q:M-PF-21} & 从分布函数回到密度与期望\newline {\small M-PF-21} & \pageref{q:M-PF-21}/\pageref{a:M-PF-21}\\[2pt]

\ref{q:M-PF-29} & 参数自助法与基本区间\newline {\small M-PF-29} & \pageref{q:M-PF-29}/\pageref{a:M-PF-29}\\[2pt]

\ref{q:M-PF-30} & 极小尾概率与数值稳定性\newline {\small M-PF-30} & \pageref{q:M-PF-30}/\pageref{a:M-PF-30}\\[2pt]

\ref{q:M-PF-31} & 正态测量的后验精度规划\newline {\small M-PF-31} & \pageref{q:M-PF-31}/\pageref{a:M-PF-31}\\[2pt]

\ref{q:B17.1} & 泊松分裂与参数推断\newline {\small B17.1} & \pageref{q:B17.1}/\pageref{a:B17.1}\\[2pt]

\ref{q:B17.2} & 有限总体抽样与修正因子\newline {\small B17.2} & \pageref{q:B17.2}/\pageref{a:B17.2}\\[2pt]

\ref{q:B17.3} & 随机化回答的估计误差\newline {\small B17.3} & \pageref{q:B17.3}/\pageref{a:B17.3}\\[2pt]

\ref{q:B17.4} & 尾和公式与收集全套\newline {\small B17.4} & \pageref{q:B17.4}/\pageref{a:B17.4}\\[2pt]

\ref{q:B17.5} & 随机和的复合泊松结构\newline {\small B17.5} & \pageref{q:B17.5}/\pageref{a:B17.5}\\[2pt]

\ref{q:B17.6} & 正态变量平方和与比值\newline {\small B17.6} & \pageref{q:B17.6}/\pageref{a:B17.6}\\[2pt]

\ref{q:B17.7} & 指数竞争与获胜标记\newline {\small B17.7} & \pageref{q:B17.7}/\pageref{a:B17.7}\\[2pt]

\ref{q:B17.8} & 条件期望作为最小平方预测\newline {\small B17.8} & \pageref{q:B17.8}/\pageref{a:B17.8}\\[2pt]

\ref{q:B17.9} & 混合正态的方差分解\newline {\small B17.9} & \pageref{q:B17.9}/\pageref{a:B17.9}\\[2pt]

\ref{q:B17.10} & 德尔塔方法估计风险函数\newline {\small B17.10} & \pageref{q:B17.10}/\pageref{a:B17.10}\\[2pt]

\ref{q:B17.11} & 威尔逊比例区间的反演\newline {\small B17.11} & \pageref{q:B17.11}/\pageref{a:B17.11}\\[2pt]

\ref{q:B17.12} & 正态模型的充分统计量\newline {\small B17.12} & \pageref{q:B17.12}/\pageref{a:B17.12}\\[2pt]

\ref{q:B17.13} & 最小值模型的非对称估计\newline {\small B17.13} & \pageref{q:B17.13}/\pageref{a:B17.13}\\[2pt]

\ref{q:B17.14} & 贝塔—二项预测的额外变异\newline {\small B17.14} & \pageref{q:B17.14}/\pageref{a:B17.14}\\[2pt]

\ref{q:B17.15} & 正态性、独立性与零相关\newline {\small B17.15} & \pageref{q:B17.15}/\pageref{a:B17.15}\\[2pt]

\ref{q:B17.16} & 支撑集破坏独立性\newline {\small B17.16} & \pageref{q:B17.16}/\pageref{a:B17.16}\\[2pt]

\ref{q:B17.17} & 蒙特卡洛积分与标准误\newline {\small B17.17} & \pageref{q:B17.17}/\pageref{a:B17.17}\\[2pt]

\ref{q:B17.18} & 反向变量的方差缩减\newline {\small B17.18} & \pageref{q:B17.18}/\pageref{a:B17.18}\\[2pt]

\ref{q:B17.19} & 切比雪夫界与正态近似的角色\newline {\small B17.19} & \pageref{q:B17.19}/\pageref{a:B17.19}\\[2pt]

\ref{q:B17.20} & 学生化与斯卢茨基定理\newline {\small B17.20} & \pageref{q:B17.20}/\pageref{a:B17.20}\\[2pt]

\ref{q:B17.21} & 三种收敛之间的辨析\newline {\small B17.21} & \pageref{q:B17.21}/\pageref{a:B17.21}\\[2pt]

\ref{q:B17.22} & 参数空间约束下的正态估计\newline {\small B17.22} & \pageref{q:B17.22}/\pageref{a:B17.22}\\[2pt]

\ref{q:B17.23} & 遗漏变量导致斜率变化\newline {\small B17.23} & \pageref{q:B17.23}/\pageref{a:B17.23}\\[2pt]

\ref{q:B17.24} & 加权最小二乘\newline {\small B17.24} & \pageref{q:B17.24}/\pageref{a:B17.24}\\[2pt]

\bottomrule\end{longtable}

\noindent 全册题目、详解、方法说明与索引组成完整学习资料。引用原题时，应同时给出原项目或课程文件及原题号；本册题号用于检索，不替代原作者署名。

\end{document}
